\documentclass[11pt]{article}
\usepackage[margin=1.1in]{geometry}
\usepackage{amsmath,amssymb,amsthm}
\usepackage{aliascnt}
\usepackage{enumitem}
\usepackage{booktabs}
\usepackage{xcolor}
\definecolor{linkblue}{RGB}{25,60,120}
\usepackage[colorlinks=true,linkcolor=linkblue,citecolor=linkblue,urlcolor=linkblue]{hyperref}
\emergencystretch=2em

\newtheorem{theorem}{Theorem}[section]
\newaliascnt{lemma}{theorem}
\newtheorem{lemma}[lemma]{Lemma}
\aliascntresetthe{lemma}
\newaliascnt{proposition}{theorem}
\newtheorem{proposition}[proposition]{Proposition}
\aliascntresetthe{proposition}
\newaliascnt{corollary}{theorem}
\newtheorem{corollary}[corollary]{Corollary}
\aliascntresetthe{corollary}
\newaliascnt{fact}{theorem}
\newtheorem{fact}[fact]{Fact}
\aliascntresetthe{fact}
\theoremstyle{definition}
\newaliascnt{definition}{theorem}
\newtheorem{definition}[definition]{Definition}
\aliascntresetthe{definition}
\newaliascnt{problem}{theorem}
\newtheorem{problem}[problem]{Problem}
\aliascntresetthe{problem}
\newaliascnt{conjecture}{theorem}
\newtheorem{conjecture}[conjecture]{Conjecture}
\aliascntresetthe{conjecture}
\newaliascnt{residualinput}{theorem}
\newtheorem{residualinput}[residualinput]{Residual Input}
\aliascntresetthe{residualinput}
\newaliascnt{researchquestion}{theorem}
\newtheorem{researchquestion}[researchquestion]{Question}
\aliascntresetthe{researchquestion}
\theoremstyle{remark}
\newaliascnt{remark}{theorem}
\newtheorem{remark}[remark]{Remark}
\aliascntresetthe{remark}
\usepackage[capitalize]{cleveref}
\crefname{problem}{Problem}{Problems}
\Crefname{problem}{Problem}{Problems}
\crefname{conjecture}{Conjecture}{Conjectures}
\Crefname{conjecture}{Conjecture}{Conjectures}
\crefname{residualinput}{Residual Input}{Residual Inputs}
\Crefname{residualinput}{Residual Input}{Residual Inputs}
\crefname{researchquestion}{Question}{Questions}
\Crefname{researchquestion}{Question}{Questions}
\crefname{fact}{Fact}{Facts}
\Crefname{fact}{Fact}{Facts}
\crefname{lemma}{Lemma}{Lemmas}
\Crefname{lemma}{Lemma}{Lemmas}
\crefname{proposition}{Proposition}{Propositions}
\Crefname{proposition}{Proposition}{Propositions}
\crefname{corollary}{Corollary}{Corollaries}
\Crefname{corollary}{Corollary}{Corollaries}
\crefname{definition}{Definition}{Definitions}
\Crefname{definition}{Definition}{Definitions}
\crefname{remark}{Remark}{Remarks}
\Crefname{remark}{Remark}{Remarks}

\newcommand{\F}{\mathbb{F}}
\newcommand{\B}{B}
\newcommand{\RS}{\mathrm{RS}}
\newcommand{\eca}{\varepsilon_{\mathrm{ca}}}
\newcommand{\emca}{\varepsilon_{\mathrm{mca}}}
\newcommand{\dmin}{\delta_{\min}}
\newcommand{\dist}{\Delta}
\newcommand{\distS}{\Delta_S}
\newcommand{\Lst}{\Lambda}
\newcommand{\Syn}{\operatorname{Syn}}
\newcommand{\Hb}{\mathrm{H}_2}
\newcommand{\eone}{e_1}
\newcommand{\UU}{\mathbb{U}}

\title{Universal Field-Size Caps and a Two-Sided Sandwich\\ for Mutual Correlated Agreement on Smooth Reed--Solomon Domains}
\author{Przemek Chojecki\\ \small ulam.ai}
\date{July 2, 2026}

\begin{document}
\maketitle

\begin{abstract}
We study the mutual correlated agreement (MCA), correlated agreement (CA), and list challenges of Arnon--Boneh--Fenzi \cite{ABF26} for Reed--Solomon codes $C=\RS[\F,D,k]$ on smooth evaluation domains.  The goal is to locate the threshold $\delta_C^*(\varepsilon)$ at cryptographic targets such as $2^{-128}$.

On the unsafe side we prove a field-size-universal cap.  For every finite field with $n<|\F|<2^{256}$ and every smooth multiplicative row of rate $\rho\in\{\frac12,\frac14,\frac18,\frac1{16}\}$ with $k\le2^{40}$,
\[
\emca(C,\delta)\;>\;\frac1{2k}\Bigl(1-\frac n{|\F|}\Bigr)\;\ge\;2^{-86}\;\gg\;2^{-128}
\]
throughout the printed near-capacity band, with the same bound for $\eca$ on the integer-admissible subinterval.  The proof is elementary: quotient-locator pigeonhole floors are composed with a self-contained deep-point conversion, with no value-set counting, exponential-sum estimate, or equidistribution input.  The same mechanism extends to map-smooth domains, Chebyshev line rounds, bivariate circle codes, and rational scale maps relevant to genus-one rows.

On the safe side we prove two self-contained reductions.  First, for every linear code, MCA has error at most $(\lfloor\delta n\rfloor+1)/q$ below one third of the minimum distance.  Second, below one half of the minimum distance, MCA is bounded by ordinary CA up to an explicit tangent term.  Combining this with the classical unique-decoding proximity gap of Ben-Sasson--Carmon--Ishai--Kopparty--Saraf \cite{BCIKS20} gives a conditional half-distance safe frontier; a self-contained half-Johnson CA bound gives an import-free alternative.

For the deployed KoalaBear-sextic and Mersenne-$31$ circle rows, these results give explicit two-sided threshold intervals at the natural target of each row ($2^{-128}$ and $2^{-100}$ respectively).  The remaining gap is a three-part frontier package: quotient/prefix flatness, a base-field-normalized split-pencil census, and primitive shift-pair control in the exact second-moment ledger.  The paper also records explicit witness words, explicit certifying pairs in residue-line normal form, an optimized failure profile, and exact integer certificates for the printed deployed-row verdicts; the sole optional non-elementary input is the imported half-distance CA theorem, isolated explicitly.
\end{abstract}

\medskip
\noindent\textbf{Keywords:} Reed--Solomon codes, proximity gaps, correlated agreement, list decoding, FRI, STARKs.\\
\textbf{MSC 2020:} 94B35 (primary); 11T71, 68Q25 (secondary).

{\small\setcounter{tocdepth}{2}\tableofcontents}

\section{Introduction}\label{sec:intro}

Proximity testing for Reed--Solomon codes underlies the soundness of modern SNARKs, and the sharpest available soundness accounting is phrased through \emph{mutual correlated agreement} (MCA): a random linear combination $f_1+\gamma f_2$ being $\delta$-close to the code should force $f_1,f_2$ to be simultaneously explained by codewords on a \emph{common} agreement set, except with small probability $\emca(C,\delta)$ over $\gamma$. The survey of Arnon, Boneh and Fenzi \cite{ABF26} crystallizes the state of the art into a \emph{grand MCA challenge}: for $C=\RS[\F,L,k]$ over a smooth domain $L$, with rate $\rho\in\{1/2,1/4,1/8,1/16\}$, degree $k\le2^{40}$, and field size $|\F|<2^{256}$, determine the largest proximity parameter $\delta^*$ at which $\emca(C,\delta)\le2^{-128}$ can hold.

Positive results reach the Johnson bound $1-\sqrt\rho$. On the negative side, \cite{Cho26b} showed $\delta^*_C(2^{-128})<1-\rho-1/64$ for all prime fields of size up to roughly $2^{150}$ (per-rate reaches between $2^{150.6}$ and $2^{161.8}$), by exhibiting quotient-locator witnesses whose bad-slope sets are value sets of restricted symmetric sums over small subgroups; it further proved a \emph{conditional} universal cap over every field of size at most $2^{256}$ --- error $2^{-42}$ at gap $2^{-11}$, assuming $|\F|\ge2n$ --- by composing quotient-core lists with the same conversion isolated and proved directly here (its universal-cap theorem). Above the value-set reach, what remained open was the \emph{error-one} question, identified there with the \emph{per-fiber collision problem}: controlling, at a \emph{fixed} large prime, the collisions of characteristic-zero locator values under reduction. Extension fields, the setting actually deployed (\cite[\S6.3]{ABF26} works over a sextic extension of KoalaBear), posed an additional obstruction: by the subfield-confinement theorem of \cite{Cho26b}, refined in \cref{sec:subfield}, every explicit witness of \cite{Cho26a,Cho26b} is $\B$-rational and contributes only $p^{-5}<2^{-154}$ over the deployed extension.

This paper sharpens the universal cap --- gap $2^{-9}$ (resp.\ $2^{-10}$) in place of $2^{-11}$, error $2^{-86}$ with no condition relating $n$ and $|\F|$, and extension fields handled directly --- for every field below $2^{256}$ at once, by changing what is counted.


\begin{theorem}[universal unsafe cap, informal; see \cref{thm:main,cor:grand}]\label{thm:informal}
Let $\F$ be any finite field with $|\F|<2^{256}$, let $D\subseteq\F^\times$ be a smooth multiplicative evaluation domain of order $n$, and let $C=\RS[\F,D,k]$ have $k=\rho n\le2^{40}$ with $\rho\in\{1/2,1/4,1/8,1/16\}$. Assume $2^{10}\mid n$ for $\rho\in\{1/2,1/4,1/8\}$, and $2^{11}\mid n$ for $\rho=1/16$. Then
\[
\emca(C,\delta)>2^{-86}\gg2^{-128}
\]
for every $\delta\in[1-\rho-2^{-9},1-\rho)$ at rates $1/2,1/4,1/8$, and for every $\delta\in[1-\rho-2^{-10},1-\rho)$ at rate $1/16$. The same lower bound holds for $\eca(C,\delta)$ on the integer-admissible subinterval ending at $1-\rho-1/n$. If $|\F|\ge2n$, the lower bound is $\ge2^{-42}$. Hence $\delta^*_C(2^{-128})\le1-\rho-2^{-9}$ for $\rho\in\{1/2,1/4,1/8\}$ and $\delta^*_C(2^{-128})\le1-\rho-2^{-10}$ for $\rho=1/16$.
\end{theorem}

Combining the unsafe caps, the safe-side theorems, and the certificate refinements gives the following two-sided summary.

\begin{theorem}[two-sided deployed intervals, informal; see \cref{thm:sandwich2}, \cref{tab:scanner2}, \cref{cor:circle-anyfield}]\label{thm:informal-sandwich}
For every deployed row --- the KoalaBear-sextic multiplicative row and the Mersenne-$31$ circle rows, at their stated challenge targets --- the grand MCA threshold is pinned two-sidedly: at rate $\rho=\tfrac12$,
\begin{align*}
\delta^*_C\ &\in\ \bigl[\tfrac{1-\rho}3,\ \tfrac{15331}{32768}\bigr] &&\text{unconditionally, and}\\
\delta^*_C\ &\in\ \bigl[\tfrac{1-\rho}2,\ \tfrac{15331}{32768}\bigr] &&\text{modulo one classical import \cite{BCIKS20},}
\end{align*}
so the threshold is known within a factor smaller than two \textup(the upper edge is $\tfrac{15331}{32768}\approx0.46787$\textup). The same two-sided sandwich holds for genus-one \textup(ECFFT\textup) rows and for bivariate circle rows over every challenge field, whether or not it contains $i$.
\end{theorem}

The proof composes two ingredients. The first (\cref{lem:fiber}) is a pigeonhole lower bound on list sizes at a single received word: the quotient locators of \cite{Cho26a} attach to each $\ell$-subset $A$ of the order-$N$ quotient coset $Q=D^{a}$ a codeword $c_A$ at distance $\le1-\rho-t/N$ from the word $u_{z_A}=x^{k+ta}+z_Ax^{k+(t-1)a}$, where $z_A=-\eone(A)$; pigeonholing the $\binom N\ell$ subsets over their $\le|\B|$ slope values produces one word $u_z$ carrying a list of size $\binom N\ell/|\B|$. Two refinements of \cite[Main Thm.\,(d)]{Cho26a} matter here: we run the construction at \emph{slack two} ($t=2$), so that the listed codewords have degree $\le k$ and live in $\RS[\F,D,k+1]$ without any divisibility demand on $k+1$ (which is odd in all challenge instantiations, while $a$ is a $2$-power); and we pigeonhole over the \emph{subfield} $\B$ containing $D$ rather than over $\F$, which is what makes the extension-field case go through. The second ingredient is the direct deep-point conversion of \cref{thm:A}: in the integer-error range $f\le n-k-1$, if $\eca(C,\delta)$ is below $\eta(\frac1k-\frac n{k|\F|})$, then \emph{every} list of $\RS[\F,D,k+1]$ at radius $\delta$ has size at most $\lceil|\F|\,\eca/(1-\eta)\rceil$. Taking $\eta=\frac12$, the heavy fiber of \cref{lem:fiber} violates this whenever $\binom N\ell\ge|\B|(|\F|/k+1)$ --- an inequality with hundreds of bits of slack throughout the challenge envelope --- and the contrapositive yields $\eca(C,\delta)>\frac1{2k}(1-n/|\F|)$ throughout the integer-admissible range. The MCA cap then persists up to capacity by support-wise MCA monotonicity.

No value-set counting, exponential-sum estimate, or equidistribution hypothesis enters: pigeonhole fibers exist over every field unconditionally, and \cref{thm:A} converts list mass into correlated-agreement error directly. The cost of this route, relative to the error-one results of \cite{Cho26a,Cho26b} below $2^{150}$, is an error of order $1/k$ rather than $1-o(1)$, and a gap $2^{-9}$ at the first three prize rates (uniformly $2^{-10}$ across all four), rather than the fixed gap $1/64$; \cref{sec:discussion} quantifies the trade-off and what remains to be settled.

Beyond the universal cap (\cref{cor:grand}, with a two-tier summary combining it with \cite{Cho26b}), we prove: a deployed-parameter corollary at the KoalaBear-sextic instantiation of \cite[\S6.3]{ABF26}, with $\emca>2^{-22}$ at gap $2^{-7}$ and the same $\eca$ lower bound on the integer-admissible subinterval (\cref{cor:deployed}), extended to every interleaved row of \cite[Tables~2--3]{ABF26} via an interleaving-transfer lemma (\cref{lem:inter,cor:rows}); an independent slacked variant through \cite[Thm.~1.9]{BCHKS25} for robustness (\cref{prop:slacked}); and a subfield-confinement lemma (\cref{lem:confine}) showing that for lines with $\B$-valued words over an extension $\F/\B$, every CA- or MCA-bad slope lies in $\B$. The latter has a structural consequence (\cref{cor:Fvalued}): at deployed parameters the lines certifying \cref{cor:deployed} are necessarily \emph{not} $\B$-rational, so the failure proved here is fiber-borne and currently non-constructive; exhibiting explicit certifying lines is posed as the explicit-witness question \cref{prob:explicit}.

A further axis is that the counting side of the composition is not tied to the multiplicative structure of $D$. \Cref{sec:isogeny} proves the fiber lemma for arbitrary polynomial folding maps with complete uniform fibers (\cref{def:map-smooth,lem:phi-fiber}), removes the divisibility hypothesis $a\mid k$, and instantiates the resulting universal cap (\cref{thm:phi-cap}) on the deployed isogeny-smooth family: twin-coset $x$-coordinate domains of the circle group over $p\equiv3\pmod4$ carry exact Chebyshev fibers at every dyadic scale, and over challenge fields containing $i$ the bivariate circle code is diagonally equivalent to a Reed--Solomon code on a torus twin coset. The Mersenne-$31$/QM31 circle family of \cite{HLP24} thereby acquires the same negative bands as the multiplicative rows --- error $>2^{-23}$ at gap $2^{-7}$ at the deployed sizes, gaps $2^{-9}$--$2^{-11}$ universally (\cref{cor:circle-grand,cor:circle-deployed}) --- answering the circle case of the isogeny-smooth problem with the elliptic (ECFFT) case reduced to a precise lacunarity obstruction (\cref{rem:ecfft-obstruction}, the genus-one/circle transport question \cref{prob:isogeny}).

The paper has three main components.  First, it gives a field-size-universal unsafe construction: quotient-fiber and graded-prefix floors produce large lists at one received word, and a direct deep-point conversion turns that list mass into CA and MCA failure near capacity.  Second, it supplies the safe side needed to interpret these failures as a threshold statement: a self-contained deep-regime theorem proves MCA soundness below one third of the distance, while a pincer theorem reduces MCA to ordinary CA up to half the distance, with the half-distance endpoint obtained from a single imported proximity-gap theorem.  Third, it packages the deployed consequences as finite certificates: the scanner outputs the safe edge, unsafe edge, and remaining open band for each row, and every printed inequality is checkable by exact integer arithmetic.

After these components, the remaining mathematical question has a sharper form.  The unsafe side reaches the entropy--subfield envelope, and the safe side reduces to three explicit residual inputs: quotient/prefix flatness, a base-field-normalized split-pencil census, and primitive shift-pair control.  Proving these with polynomial loss gives the frontier with reserve; proving them with printed constants gives the finite adjacent pairs.

\paragraph{Roadmap.} \Cref{sec:prelim} fixes definitions, aligned with \cite[Defs.~4.1, 4.3]{ABF26}, proves the direct deep-point conversion used for the main cap, and records the separate BCHKS/ABF slacked fallback. \Cref{sec:fiber} proves the fiber and interleaving lemmas and the conservative quotient-support upper ledger, \cref{sec:main} the main theorem, \cref{sec:grand} the challenge-level and deployed corollaries, and \cref{sec:subfield} subfield confinement. \Cref{sec:isogeny} extends the cap beyond multiplicative smoothness: it proves the map-smooth, divisibility-free fiber lemma and universal cap, the exact Chebyshev-fiber and torus twin-coset geometry of circle domains, and the resulting circle-code caps at the deployed Mersenne-$31$/QM31 parameters. \Cref{sec:prize-program} formulates the integer staircase and the certificate interface used to compare safe and unsafe radii.  \Cref{sec:discharge} develops the structural ledgers, extension-line lift, residual-band language, scanner, and finite certificate core.  \Cref{sec:pincer} proves the safe-side pincer below half the distance and the scale-optimized unsafe floors, then assembles the narrowed two-sided sandwich.  \Cref{sec:answers} treats rational scales, circle-code transport over every challenge field, explicit witnesses, optimized failure profiles, and the deployed-row certificate table.  \Cref{sec:discussion} discusses scope, limitations, and the remaining frontier inputs.  Throughout, companion results of \cite{Cho26a,Cho26b} are cited by their printed theorem names rather than numbers, since the numbering of those manuscripts is not yet frozen.

\section{Preliminaries}\label{sec:prelim}

\paragraph{Codes and distances.} Throughout, $\B\subseteq\F$ are finite fields, $q:=|\F|$, and $D=\alpha H\subseteq\B^\times$ is a multiplicative coset of a cyclic subgroup $H$ of order $n$; we write $D$ for the evaluation domain (denoted $L$ in \cite{ABF26}). The subgroup case is $\alpha=1$, and taking $\B=\F$ is always allowed. For $k\le n$,
\[
\RS[\F,D,k]\;:=\;\bigl\{(p(x))_{x\in D}\;:\;p\in\F[X],\ \deg p<k\bigr\}\subseteq\F^n,
\qquad \rho:=k/n .
\]
For $u,v\in\F^n$, $\dist(u,v)$ is the relative Hamming distance, $\dist(u,C):=\min_{c\in C}\dist(u,c)$, and the relative minimum distance of $C=\RS[\F,D,k]$ is $\dmin(C)=1-\rho+1/n$. For $S\subseteq D$, $\distS(u,v)$ and $\distS(u,C)$ denote the corresponding restricted relative distances on $S$. Lists are denoted
\[
\Lst(C,\delta,u):=\{c\in C:\dist(u,c)\le\delta\},\qquad \Lst(C,\delta):=\max_{u\in\F^n}|\Lst(C,\delta,u)| .
\]
For $s\ge1$ the $s$-interleaved code is $C^{\equiv s}:=\{(c^{(1)},\dots,c^{(s)}):c^{(i)}\in C\}\subseteq(\F^n)^s$ with the column distance $\dist_s(F,G):=\frac1n\,\bigl|\{j:\exists i,\ f^{(i)}_j\ne g^{(i)}_j\}\bigr|$, and $\dist_s(F,C^{\equiv s}):=\min_{\mathbf c\in C^{\equiv s}}\dist_s(F,\mathbf c)$.

\begin{definition}[correlated agreement error; {\cite[Def.~4.1]{ABF26}} up to notation]\label{def:ca}
For $0\le\delta_{\mathrm{fld}}\le\delta_{\mathrm{int}}<1$,
\[
\eca(C,\delta_{\mathrm{fld}},\delta_{\mathrm{int}})\;:=\;\max_{f_1,f_2\in\F^n}\ \Pr_{\gamma\leftarrow\F}\Bigl[\ \dist(f_1+\gamma f_2,\,C)\le\delta_{\mathrm{fld}}\ \wedge\ \dist_2\bigl((f_1,f_2),\,C^{\equiv2}\bigr)>\delta_{\mathrm{int}}\ \Bigr],
\]
and the no-proximity-loss variant is $\eca(C,\delta):=\eca(C,\delta,\delta)$. A slope $\gamma$ realizing the event for a fixed pair $(f_1,f_2)$ is called \emph{CA-bad} for that pair.
\end{definition}

\begin{definition}[mutual correlated agreement error, support-wise; {\cite[Def.~4.3]{ABF26}} up to notation]\label{def:mca}
For $\delta\in(0,1)$,
\[
\emca(C,\delta):=\max_{f_1,f_2\in\F^n}\Pr_{\gamma\leftarrow\F}\left[
\begin{array}{l}
\exists S\subseteq D,\ |S|\ge(1-\delta)n,\ \distS(f_1+\gamma f_2,C)=0,\\[2pt]
\text{and }\distS\bigl((f_1,f_2),C^{\equiv2}\bigr)>0
\end{array}
\right].
\]
Equivalently, $\gamma$ is MCA-bad for $(f_1,f_2)$ if some large agreement set $S$ explains the point $f_1+\gamma f_2$, but the same set $S$ does not simultaneously explain $f_1$ and $f_2$ by two codewords of $C$.
\end{definition}

\begin{definition}[challenge threshold]\label{def:dstar}
Following the grand MCA challenge of \cite[\S1]{ABF26}, for $\varepsilon^*\in(0,1)$ set
$\delta^*_C(\varepsilon^*):=\sup\{\delta\in(0,1-\rho):\emca(C,\delta)\le\varepsilon^*\}$, with $\sup\emptyset:=0$.
\end{definition}

All displayed threshold intervals use the integer-ball convention: distances are integral, so a radius $\delta$ corresponds to the agreement requirement $\lceil(1-\delta)n\rceil$ and to $\lfloor\delta n\rfloor$ permitted errors. When a theorem certifies that every $\delta\in[\delta_0,1-\rho)$ is unsafe at target $\varepsilon^*$, it follows that $\delta^*_C(\varepsilon^*)\le\delta_0$ under the supremum convention; when it certifies safety at a single radius $\delta_1<1-\rho$, it follows that $\delta^*_C(\varepsilon^*)\ge\delta_1$.

\begin{fact}[{cf.\ \cite[Fact~4.5]{ABF26}}]\label{fact:chain}
For every $\delta\in(0,1)$: $\eca(C,\delta)\le\emca(C,\delta)$.
\end{fact}

\begin{proof}
Fix $(f_1,f_2)$ and let $\gamma$ be CA-bad: $\dist(f_1+\gamma f_2,C)\le\delta$ and $\dist_2((f_1,f_2),C^{\equiv2})>\delta$. Choose $S\subseteq D$ of size at least $(1-\delta)n$ on which $f_1+\gamma f_2$ agrees with a codeword of $C$. Since $(f_1,f_2)$ is not $\delta$-close to $C^{\equiv2}$, there is no set of size at least $(1-\delta)n$ on which $f_1$ and $f_2$ are simultaneously explained by codewords of $C$; in particular this fails on the chosen $S$. Thus $\gamma$ is MCA-bad for the same pair. Taking maxima over pairs gives the claim.
\end{proof}

\begin{lemma}[support-wise MCA monotonicity]\label{lem:mca-monotone}
For every linear code $C\subseteq\F^D$, the function
\[
        \delta\longmapsto \emca(C,\delta)
\]
is nondecreasing.
\end{lemma}

\begin{proof}
Fix a pair $(f_1,f_2)$ and suppose that $\gamma$ is MCA-bad at radius $\delta$. Then there is a set $S\subseteq D$ with
\[
        |S|\ge (1-\delta)n
\]
such that $f_1+\gamma f_2$ is explained by $C$ on $S$, but $(f_1,f_2)$ is not simultaneously explained by two codewords of $C$ on the same $S$. If $\delta'\ge\delta$, then
\[
        |S|\ge (1-\delta)n\ge (1-\delta')n,
\]
so the same set $S$ witnesses that $\gamma$ is MCA-bad at radius $\delta'$. Taking the maximum over pairs gives the claim.
\end{proof}

\paragraph{Conversion theorem.} The main composition below uses the deep-point list-to-CA conversion proved here directly. This is the same consequence previously imported from Crites--Stewart, with the integer-radius admissibility condition made explicit; see \cref{rem:import}. We keep the slacked BCHKS/ABF fallback as a separate imported route.

\begin{theorem}[deep-point list-to-CA conversion]
\label{thm:A}
Let $C=\RS[\F,D,k]$, let $C^+:=\RS[\F,D,k+1]$, let
$q:=|\F|$ and $n:=|D|$, and assume $q>n$.  Let $\delta\in(0,1)$ and set
$f_\delta:=\lfloor \delta n\rfloor$.  Assume
\[
        f_\delta \le n-k-1 .
\]
For any $\eta\in[0,1)$, if
\[
        \eca(C,\delta)
        \le
        \eta\Bigl(\frac1k-\frac n{kq}\Bigr),
\]
then
\[
        \Lst(C^+,\delta)
        \le
        \Bigl\lceil
        \frac{q\,\eca(C,\delta)}{1-\eta}
        \Bigr\rceil .
\]
\end{theorem}

\begin{proof}
Put $\epsilon:=\eca(C,\delta)$ and $f=f_\delta$.
Let $U:D\to\F$ be any received word and suppose that
$P_1,\ldots,P_L\in\F[X]_{<k+1}$ are distinct polynomials whose
restrictions to $D$ all lie in the Hamming ball of relative radius
$\delta$ around $U$.  Since distances are integral, each $P_i$ agrees
with $U$ on at least
\[
        a:=n-f
\]
points of $D$.  The hypothesis $f\le n-k-1$ gives $a\ge k+1$, in particular
$a>k$.

Let $\Omega:=\F\setminus D$, so $|\Omega|=q-n$.  For each
$\alpha\in\Omega$ define the simple-pole line
\[
        f_\alpha(x)=\frac{U(x)}{x-\alpha},
        \qquad
        g_\alpha(x)=-\frac1{x-\alpha}
        \qquad (x\in D).
\]
If $P_i$ agrees with $U$ on a set $S_i$ of size at least $a$, then for
$z_i(\alpha):=P_i(\alpha)$ and every $x\in S_i$,
\[
        f_\alpha(x)+z_i(\alpha)g_\alpha(x)
        =\frac{P_i(x)-P_i(\alpha)}{x-\alpha}.
\]
The right-hand side is the restriction of a polynomial of degree $<k$.
Thus each distinct value among the $P_i(\alpha)$ gives a slope for
which the line point is $\delta$-close to $C$.

We next check the correlated-agreement far condition.  If a polynomial
$G\in\F[X]_{<k}$ agreed with $g_\alpha$ on more than $k$ points of $D$,
then
\[
        (X-\alpha)G(X)+1
\]
would be a polynomial of degree at most $k$ with more than $k$ roots,
but its value at $X=\alpha$ is $1$, a contradiction.  Hence
$(f_\alpha,g_\alpha)$ is not $\delta$-close to $C^{\equiv2}$, because
any common explanation on a support of size at least $a>k$ would in
particular explain $g_\alpha$ there.  Therefore the distinct values of
$P_i(\alpha)$ are CA-bad slopes.

It remains to lower-bound the number of distinct values for some
$\alpha$.  For $i\ne j$, the polynomial $P_i-P_j$ is nonzero of degree
at most $k$, so it has at most $k$ roots in $\Omega$.  Let $M(\alpha)$
be the number of distinct values among $P_1(\alpha),\ldots,P_L(\alpha)$.
Over all $\alpha\in\Omega$, the total number of colliding unordered pairs
$(i,j)$ with $P_i(\alpha)=P_j(\alpha)$ is at most $k\binom L2$.
Equivalently, if the fiber sizes of the map $i\mapsto P_i(\alpha)$ are
$m_v(\alpha)$, then
\[
        \sum_{\alpha\in\Omega}\sum_v \binom{m_v(\alpha)}2
        \le k\binom L2 .
\]
Choose $\alpha\in\Omega$ for which the inner collision count is at most
$k\binom L2/(q-n)$.  If the fiber sizes of the map $i\mapsto P_i(\alpha)$ are
$m_1,\ldots,m_{M(\alpha)}$, then
\[
        \sum_r m_r=L,
        \qquad
        \sum_r m_r^2=L+2\sum_r\binom{m_r}{2}
        \le L+\frac{kL(L-1)}{q-n}.
\]
Cauchy--Schwarz therefore gives
\[
        L^2\le M(\alpha)\sum_r m_r^2,
        \qquad
        M(\alpha)\ge
        \frac{L(q-n)}{q-n+k(L-1)}
        \ge
        \frac{L(q-n)}{q-n+kL} .
\]
Consequently
\[
        \epsilon
        \ge
        \frac1q\cdot \frac{L(q-n)}{q-n+kL}.
\]
Solving this inequality for $L$ gives
\[
        L
        \le
        \frac{\epsilon q(q-n)}{q-n-k\epsilon q},
\]
provided $\epsilon<(q-n)/(kq)$.  Under the stated hypothesis,
$\epsilon\le\eta(q-n)/(kq)$ with $\eta<1$, so the denominator is at
least $(1-\eta)(q-n)$ and hence
\[
        L\le \frac{q\epsilon}{1-\eta}.
\]
Since $U$ and the list were arbitrary, the displayed ceiling bounds the
maximum list size of $C^+$ at radius $\delta$.
\end{proof}

\begin{remark}[relation to Crites--Stewart]
The displayed conversion is the only consequence of Crites--Stewart
needed by the universal-cap argument.  The proof above is included to make the main cap
self-contained.  The separate slacked fallback through BCHKS/ABF remains
an independent imported route; it is not used in the final deployed certificate
grammar and is retained as a robustness check.
\end{remark}

\begin{proposition}[quantitative deep-list floor]
\label{prop:quantitative-deep-list-floor}
Let $C=\RS[\F,D,k]$, let $C^+:=\RS[\F,D,k+1]$, put
$q=|\F|$ and $n=|D|$, and assume $q>n$.  Let
$\delta_0\in[0,1-k/n)$, and suppose that for some received word
$U:D\to\F$ there are $L\ge1$ distinct codewords of $C^+$ in the
closed ball of radius $\delta_0$ around $U$.  Then, for every
$\delta\in[\delta_0,1-k/n)$,
\[
        \eca(C,\delta)
        \ge
        \frac{L(q-n)}{q(q-n+kL)} .
\]
For every such $\delta>0$,
\[
        \emca(C,\delta)
        \ge
        \frac{L(q-n)}{q(q-n+kL)}
\]
as well.  More precisely, some received line has at least
\[
        \left\lceil\frac{L(q-n)}{q-n+kL}\right\rceil
\]
CA-bad finite slopes at every such radius; for $\delta>0$, the same slopes are
MCA-bad.
\end{proposition}

\begin{proof}
Fix $\delta\in[\delta_0,1-k/n)$ and choose distinct
$P_1,\ldots,P_L\in\F[X]_{<k+1}$ in the $\delta_0$-ball around $U$.
They are also in the $\delta$-ball.  Since $\delta<1-k/n$, if
$f=\lfloor\delta n\rfloor$ then $f\le n-k-1$, so every codeword in this
$\delta$-ball agrees with $U$ on at least $a=n-f>k$ points.

Repeat the simple-pole construction in the proof of \cref{thm:A}.  For
$\alpha\in\Omega:=\F\setminus D$ set
\[
        f_\alpha(x)=\frac{U(x)}{x-\alpha},
        \qquad
        g_\alpha(x)=-\frac1{x-\alpha}
        \qquad (x\in D).
\]
For every distinct value of $P_i(\alpha)$, the word
$f_\alpha+P_i(\alpha)g_\alpha$ is $\delta$-close to $C$.  The pair
$(f_\alpha,g_\alpha)$ is not $\delta$-close to $C^{\equiv2}$, because a
common explanation on more than $k$ points would in particular explain
$g_\alpha$ there, and then $(X-\alpha)G(X)+1$ would be a polynomial of degree at most $k$ with more than $k$ roots and value $1$ at $\alpha$.
Thus the distinct values of the map $i\mapsto P_i(\alpha)$ are CA-bad
slopes.

For $i\ne j$, the polynomial $P_i-P_j$ has degree at most $k$ and is
nonzero, hence it has at most $k$ roots in $\Omega$.  Averaging over
$|\Omega|=q-n$ gives an $\alpha$ for which the number of colliding
unordered pairs among the $L$ values is at most $k\binom L2/(q-n)$.  If
$m_1,\ldots,m_M$ are the value multiplicities for this $\alpha$, then
\[
        \sum_r m_r=L,
        \qquad
        \sum_r m_r^2
        \le L+\frac{kL(L-1)}{q-n}.
\]
By Cauchy--Schwarz,
\[
        M\ge
        \frac{L^2}{L+kL(L-1)/(q-n)}
        =
        \frac{L(q-n)}{q-n+k(L-1)}
        \ge
        \frac{L(q-n)}{q-n+kL}.
\]
Since $M$ is an integer, $M\ge\lceil L(q-n)/(q-n+kL)\rceil$.  Dividing
by the $q$ possible finite slopes gives the CA lower bound.  For $\delta>0$, the MCA
bound follows from \cref{fact:chain}.
\end{proof}

\begin{corollary}[deep-list trigger and ceiling]
\label{cor:deep-list-trigger-ceiling}
In the setting of \cref{prop:quantitative-deep-list-floor}, define
\[
        \mathcal E_{q,k}(L)
        :=\frac{L(q-n)}{q(q-n+kL)},
        \qquad
        \mathcal N_{q,k}(L)
        :=\left\lceil\frac{L(q-n)}{q-n+kL}\right\rceil .
\]
Then $\mathcal E_{q,k}(L)$ is increasing in $L$, and
\[
        \mathcal E_{q,k}(L)>
        \frac1{2k}\left(1-\frac nq\right)
        \quad\Longleftrightarrow\quad
        L>\frac{q-n}{k}.
\]
Moreover
\[
        \mathcal E_{q,k}(L)<\frac1k\left(1-\frac nq\right)
\]
for every finite $L$, and the right-hand side is the limiting value as
$L\to\infty$.
\end{corollary}

\begin{proof}
The derivative of $L/(q-n+kL)$ is positive.  The trigger inequality is
\[
        \frac{L}{q-n+kL}>\frac1{2k},
\]
which is equivalent to $kL>q-n$.  The upper bound follows from
$q-n+kL>kL$.
\end{proof}

\begin{remark}[why this is stronger than the contrapositive]
\label{rem:quantitative-floor-vs-contrapositive}
\Cref{thm:A} is optimized for contradiction: once $L>(q-n)/k$, taking
$\eta=1/2$ already forces the printed $(2k)^{-1}(1-n/q)$ lower bound.
\Cref{prop:quantitative-deep-list-floor} keeps the whole interpolation count.
It is the form a staircase scanner should use when a quotient fiber is too small
to cross the $1/(2k)$ trigger but still contributes a nonzero explicit bad-slope
numerator, and also when a very large fiber improves the constant toward
$(1/k)(1-n/q)$.
\end{remark}

\begin{proposition}[punctured simple-pole deep-list floor]
\label{prop:punctured-simple-pole-floor}
Let $C=\RS[\F,D,k]$, let $C^+:=\RS[\F,D,k+1]$, put
$q=|\F|$ and $n=|D|$.  Let
\[
        \Omega_0\subseteq \F\setminus D,
        \qquad m:=|\Omega_0|>0 .
\]
Let $\delta_0\in[0,1-k/n)$, and suppose that for some received word
$U:D\to\F$ there are $L\ge1$ distinct codewords of $C^+$ in the closed ball of
radius $\delta_0$ around $U$.  Then, for every
$\delta\in[\delta_0,1-k/n)$,
\[
        \eca(C,\delta)
        \ge
        \frac{Lm}{q(m+kL)} .
\]
For every such $\delta>0$,
\[
        \emca(C,\delta)
        \ge
        \frac{Lm}{q(m+kL)} .
\]
More precisely, there is some $\alpha\in\Omega_0$ such that the simple-pole line
\[
        f_\alpha(x)=\frac{U(x)}{x-\alpha},
        \qquad
        g_\alpha(x)=-\frac1{x-\alpha}
        \qquad (x\in D)
\]
has at least
\[
        \left\lceil\frac{Lm}{m+kL}\right\rceil
\]
CA-bad finite slopes at every such radius; for $\delta>0$, the same slopes are
MCA-bad.
\end{proposition}

\begin{proof}
Repeat the proof of \cref{prop:quantitative-deep-list-floor}, but average only
over $\Omega_0$ instead of over all of $\F\setminus D$.  The simple-pole
argument works for every $\alpha\notin D$, so every distinct value among
$P_1(\alpha),\ldots,P_L(\alpha)$ gives a CA-bad slope.  For $i\ne j$, the
nonzero polynomial $P_i-P_j$ has degree at most $k$, hence has at most $k$ roots
inside the smaller set $\Omega_0$.  Therefore the total number of colliding
unordered pairs over all $\alpha\in\Omega_0$ is at most $k\binom L2$.  Choose
$\alpha\in\Omega_0$ for which the collision count is at most
$k\binom L2/m$.  If $M$ is the number of distinct values at this $\alpha$, the
same Cauchy--Schwarz calculation gives
\[
        M\ge
        \frac{Lm}{m+k(L-1)}
        \ge
        \frac{Lm}{m+kL} .
\]
Since $M$ is integral, this gives the displayed bad-slope numerator.  Dividing
by the $q$ finite slopes gives the CA lower bound, and the MCA statement for
$\delta>0$ follows from \cref{fact:chain}.
\end{proof}

\begin{corollary}[extension-pole deep-list floor]
\label{cor:extension-pole-deep-list-floor}
Assume in addition that $\B\subsetneq\F$, $D\subseteq\B$, and $|D|\ge2$.  Put
$b:=|\B|$.  In the setting of
\cref{prop:punctured-simple-pole-floor}, take
\[
        \Omega_0:=\F\setminus\B,
        \qquad m=q-b .
\]
Then for every $\delta\in[\delta_0,1-k/n)$,
\[
        \eca(C,\delta)
        \ge
        \mathcal E^{\rm ext}_{q,b,k}(L)
        :=
        \frac{L(q-b)}{q(q-b+kL)} .
\]
For $\delta>0$ the same bound holds for $\emca(C,\delta)$.  Moreover the
witnessing line may be chosen with pole $\alpha\in\F\setminus\B$, and its pole
vector $((x-\alpha)^{-1})_{x\in D}$ is not a scalar multiple over $\F$ of any
$\B$-valued vector.  In particular the witness line is genuinely
extension-valued, not merely a $\B$-rational line viewed over $\F$.
\end{corollary}

\begin{proof}
The lower bound is \cref{prop:punctured-simple-pole-floor} with
$\Omega_0=\F\setminus\B$.  It remains to prove genuine extension-valuedness.
Suppose that, for some $\lambda\in\F^\times$, all numbers
$\lambda/(x-\alpha)$ with $x\in D$ lay in $\B$.  Choose two distinct points
$x,y\in D$.  Then
\[
        \frac{y-\alpha}{x-\alpha}
        =
        \frac{\lambda/(x-\alpha)}{\lambda/(y-\alpha)}
        \in \B .
\]
Writing this ratio as $r\in\B$, we cannot have $r=1$ because $x\ne y$.  Hence
\[
        \alpha=\frac{rx-y}{r-1}\in\B,
\]
contradicting $\alpha\notin\B$.  Thus the pole vector is not projectively
$\B$-valued.  In particular the displayed affine parametrization has no
$\B$-valued direction vector after rescaling.
\end{proof}

\begin{corollary}[extension-pole quotient-remainder floor]
\label{cor:extension-pole-quotient-remainder-floor}
Assume $\B\subsetneq\F$, $D\subseteq\B^\times$ is a multiplicative coset of order
$n\ge2$, put $b:=|\B|$, and let $C=\RS[\F,D,k]$.
Fix a quotient-remainder profile as in
\cref{lem:heaviest-prefix-locator-floor} with $K=k+1$ and agreement $A_0>k$.
Let
\[
        L:=H_{c,m,s}^{k+1}
\]
be the certified heaviest prefix-fiber size, or replace $L$ by either fallback
\[
        \left\lceil \frac{M_{c,m,s}}{I_{c,m,s}^{k+1}}\right\rceil,
        \qquad
        \left\lceil\frac{M_{c,m,s}}{|\B|^{w_c(s,A_0-k-1)}}\right\rceil .
\]
Then, for every
\[
        \delta\in[1-A_0/n,1-k/n),
\]
there is a genuinely extension-valued simple-pole line with at least
\[
        \mathcal N^{\rm ext}_{q,b,k}(L)
        :=
        \left\lceil\frac{L(q-b)}{q-b+kL}\right\rceil
\]
CA-bad finite slopes.  For $\delta>0$, these slopes are MCA-bad.  Equivalently,
\[
        \eca(C,\delta)
        \ge
        \frac{L(q-b)}{q(q-b+kL)},
        \qquad
        \emca(C,\delta)
        \ge
        \frac{L(q-b)}{q(q-b+kL)}
        \quad(\delta>0).
\]
If $L>(q-b)/k$, this gives the clean trigger
\[
        \eca(C,\delta)>
        \frac1{2k}\left(1-\frac bq\right),
        \qquad
        \emca(C,\delta)>
        \frac1{2k}\left(1-\frac bq\right)
        \quad(\delta>0).
\]
\end{corollary}

\begin{proof}
\Cref{lem:heaviest-prefix-locator-floor} gives a $\B$-valued received word with
at least $L$ codewords of $C^+$ in the closed ball of radius $1-A_0/n$.  Apply
\cref{cor:extension-pole-deep-list-floor}.  The trigger inequality is
\[
        \frac{L(q-b)}{q(q-b+kL)}>
        \frac{q-b}{2kq},
\]
which is equivalent to $kL>q-b$.
\end{proof}

\begin{definition}[extension-pole certificate]
\label{def:extension-pole-certificate}
An extension-pole certificate for a list $P_1,\ldots,P_L\in\F[X]_{<k+1}$ and a
received word $U$ consists of:
\begin{enumerate}[label=\textup{(\roman*)},leftmargin=2.2em]
\item a declared proper subfield $\B\subsetneq\F$ containing $D$ and a pole
$\alpha\in\F\setminus\B$;
\item the agreement supports showing that each $P_i$ lies in the declared ball
around $U$;
\item the bucket table for the values $P_i(\alpha)$, with duplicate values
coalesced;
\item the simple-pole line $(f_\alpha,g_\alpha)$ and the proof that
$g_\alpha$ is not projectively $\B$-valued;
\item the final bad-slope numerator, namely the number of distinct buckets in
\textup{(iii)}.
\end{enumerate}
The verifier checks $\alpha\notin\B$, the support agreements, the value buckets,
and the simple-pole far condition.  For quotient-remainder floors, the list in
\textup{(ii)} may itself be supplied by a prefix-fiber certificate from
\cref{def:prefix-fiber-certificate}.
\end{definition}

\begin{remark}[what the extension-pole construction changes]
\label{rem:extension-pole-status}
The subfield-confinement theorem says that $\B$-rational lines are inert over a
proper extension.  \Cref{cor:extension-pole-deep-list-floor} gives the matching
existence theorem on the failure side: any large list can be converted using
poles restricted to $\F\setminus\B$, so the resulting bad line is genuinely
extension-valued.  Thus the extension-line bucket in a prize scanner is no
longer only an abstract warning; it has an explicit simple-pole source and a
stable numerator with $q-b$ replacing $q-n$.  What remains to be settled is not existence
of extension-valued witnesses, but their sharp classification inside the
aperiodic/residue-line ledgers and, for a fully explicit counterexample, printing
a concrete pole $\alpha$ with its value-bucket table.
\end{remark}

\begin{theorem}[{\cite[Thm.~1.9]{BCHKS25}}, as stated in {\cite[Thm.~5.2]{ABF26}}]\label{thm:B}
Let $C=\RS[\F,D,k]$, let $\rho=k/n$, and let $\delta\in(0,1-\rho)$ satisfy $\delta+2/n\le 1-\rho-1/n$. If
\[
\eca\Bigl(C,\ \delta+\tfrac2n,\ 1-\rho-\tfrac1n\Bigr)\;<\;\frac1{2n},
\qquad\text{then}\qquad
\Lst(C,\delta)\;<\;|\F| .
\]
\end{theorem}

\begin{remark}[on imports and radius conventions]\label{rem:import}
\Cref{thm:A} is the only Crites--Stewart-type consequence used by the main cap, and it is now proved above by the simple-pole deep-point argument.  The important admissibility condition is the integer-radius condition
\[
        \lfloor \delta n\rfloor \le n-k-1 .
\]
Consequently, the no-loss CA lower bound proved below should be read only in this admissible range. The Proximity-Prize MCA cap is unaffected: we prove the lower bound at the endpoint
\[
        \delta_N:=1-\rho-\frac2N
\]
and then use support-wise MCA monotonicity to extend the MCA failure to all larger sub-capacity radii.

\Cref{thm:B} remains a separate slacked fallback imported through the BCHKS/ABF route. It is not needed for the main field-size-universal MCA cap.
\end{remark}

\section{Locator fibers and interleaving transfer}\label{sec:fiber}

Fix $N\mid n$, set $a:=n/N$, and let $Q:=D^{a}:=\{x^a:x\in D\}\subseteq\B^\times$, a multiplicative coset of order $N$. The $a$-th power map $D\to Q$ is surjective and its fibers $S_b:=\{x\in D:x^a=b\}$ have exactly $a$ elements each. Since $a\mid n\mid|\B|-1$, the characteristic of $\F$ does not divide $a$, so $X^a-b$ is separable for $b\ne0$ and its full root set inside $D$ is $S_b$ whenever $b\in Q$. For a set $A$, $e_j(A)$ denotes the $j$-th elementary symmetric function of its elements.

\begin{lemma}[locator fibers are lists]\label{lem:fiber}
Let $\B\subseteq\F$, let $D\subseteq\B^\times$ be a multiplicative coset of order $n$, and let $k$ satisfy $a\mid k$ and $\rho=k/n$. Write $\ell_1:=\rho N+1$ and $\ell_2:=\rho N+2$.
\begin{enumerate}[label=\textup{(\roman*)}]
\item If $\ell_1\le N$, there exists $z\in\B$ such that, with $u_z:=\bigl(x^{k+a}+z\,x^{k}\bigr)_{x\in D}\in\F^n$,
\[
\bigl|\Lst\bigl(\RS[\F,D,k],\ 1-\rho-\tfrac1N,\ u_z\bigr)\bigr|\;\ge\;\binom{N}{\ell_1}\Big/|\B| .
\]
\item If $\ell_2\le N$, there exists $z\in\B$ such that, with $u_z:=\bigl(x^{k+2a}+z\,x^{k+a}\bigr)_{x\in D}\in\F^n$,
\[
\bigl|\Lst\bigl(\RS[\F,D,k+1],\ 1-\rho-\tfrac2N,\ u_z\bigr)\bigr|\;\ge\;\binom{N}{\ell_2}\Big/|\B| .
\]
\end{enumerate}
In both cases the words $u_z$ are $\B$-valued and all listed codewords lie in $\RS[\B,D,k]$ \textup(resp.\ $\RS[\B,D,k+1]$\textup).
\end{lemma}

\begin{proof}
We prove (ii); part (i) is the identical computation one degree rung lower and is \cite[Main Thm.\,(d)]{Cho26a} sharpened by pigeonholing over $\B$ in place of $\F$.

Let $A\subseteq Q$ with $|A|=\ell_2$, and put
\[
L_A(X)\;:=\;\prod_{b\in A}\bigl(X^a-b\bigr)\;=\;\sum_{j=0}^{\ell_2}(-1)^j e_j(A)\,X^{a(\ell_2-j)}\;=\;X^{k+2a}-e_1(A)\,X^{k+a}+R_A(X),
\]
where $a\ell_2=k+2a$, $a(\ell_2-1)=k+a$, and $R_A$ collects the terms with $j\ge2$, so $\deg R_A\le a(\ell_2-2)=k$. The polynomial $L_A$ is squarefree with root set $S_A:=\bigcup_{b\in A}S_b\subseteq D$ of size $a\ell_2=k+2a$.

Set $z_A:=-e_1(A)\in\B$ and $c_A:=\bigl(-R_A(x)\bigr)_{x\in D}$. Since $\deg R_A\le k$, we have $c_A\in\RS[\B,D,k+1]\subseteq\RS[\F,D,k+1]$. For every $x\in S_A$,
\[
0=L_A(x)=x^{k+2a}+z_A\,x^{k+a}+R_A(x),\qquad\text{i.e.}\qquad u_{z_A}(x)=c_A(x),
\]
so $u_{z_A}$ and $c_A$ agree on at least $k+2a$ points of $D$ and $\dist(u_{z_A},c_A)\le1-\rho-2/N$.

The map $A\mapsto c_A$ is injective on the $\ell_2$-subsets of $Q$ sharing a common slope value: if $z_A=z_{A'}$ and $c_A=c_{A'}$, then $R_A$ and $R_{A'}$ are polynomials of degree $\le k<n$ agreeing on all $n$ points of $D$, hence equal as polynomials; then $L_A=L_{A'}$, so $S_A=S_{A'}$, and $A=\{x^a:x\in S_A\}=A'$.

Finally, $z_A=-e_1(A)$ takes at most $|\B|$ values as $A$ ranges over the $\binom{N}{\ell_2}$ subsets; some $z\in\B$ is attained by at least $\binom{N}{\ell_2}/|\B|$ of them, and the corresponding codewords $c_A$ are pairwise distinct elements of $\Lst(\RS[\F,D,k+1],\,1-\rho-2/N,\,u_z)$.
\end{proof}

\begin{remark}[why slack two]\label{rem:slacktwo}
Part (i) applied to $C^+=\RS[\F,D,k+1]$ directly would require $a\mid k+1$; in every challenge instantiation $a$ is a power of two and $k+1$ is odd, so this fails. At slack two the listed codewords have degree $\le k$, i.e.\ lie in $C^+$, while the divisibility demand stays on $k$. This is exactly the interface needed by \cref{thm:A}, whose list conclusion concerns $C^+$ while its hypothesis concerns $C$.
\end{remark}

\begin{remark}[the subfield pigeonhole elsewhere]\label{rem:subpigeon}
The same one-line strengthening --- pigeonholing locator data over the field of definition $\B$ rather than over $\F$ --- also upgrades the coefficient-pigeonhole bound of \cite{Cho26b} to the generated field: the prefix map $\Phi_\sigma$ there has image in $\B^\sigma$ whenever $D\subseteq\B$, since the elementary symmetric prefixes of subsets of $D$ are $\B$-valued. This is the form consumed by the ``necessary'' half of the list-profile conjecture of \cite{Cho26b}, which states the entropy floor at the generated field $\tau^{*}(\rho,q_D)$; \cref{lem:fiber} supplies the missing subfield step.
\end{remark}

\begin{remark}[lists, not lines]\label{rem:lines}
The words $u_z$ of \cref{lem:fiber} lie on the $\B$-rational line $f+zg$ with $f=(x^{k+ta})_{x\in D}$, $g=(x^{k+(t-1)a})_{x\in D}$, but the lemma is purely a statement about lists at a received word and no line structure is used in \cref{sec:main}. This matters: by \cref{cor:Fvalued} below, over a proper extension $\F\supsetneq\B$ the lines that end up certifying the CA failure cannot be these (or any) $\B$-rational lines.
\end{remark}

\subsection{Quotient-profile floors with remainder supports}
\label{sec:quotient-remainder}

The fiber lemma above uses supports that are unions of complete quotient fibers.
For threshold work one also needs agreements that are not multiples of the fiber
size.  The following locator-prefix pigeonhole gives a finite lower ledger for
such remainder supports.  It is a lower-floor theorem only; the matching upper
ledger over all divisor scales is not settled here; it is the upper-ledger side of the residual band problem.

Let $c\mid n$ and put $Q_c:=D^c$, so the map
\[
        \pi_c:D\to Q_c,\qquad x\mapsto x^c
\]
has fibers of size $c$ and $|Q_c|=N=n/c$.  For $0\le s<c$ and an integer
$\sigma\ge0$, define the prefix weight
\[
        w_c(s,\sigma)
        :=\#\{h:1\le h\le\sigma,
              \ h\bmod c\in\{0,1,\ldots,s\}\},
\]
where residues are taken in $\{0,1,\ldots,c-1\}$.  Thus
$w_c(0,\sigma)=\lfloor \sigma/c\rfloor$ for full-fiber supports.

\begin{lemma}[quotient-remainder locator-prefix lower floor]
\label{lem:quotient-remainder-prefix}
Let $\B\subseteq\F$, let $D\subseteq\B^\times$ be a multiplicative coset of
order $n$, and let $K<n$.  Fix $c\mid n$, write $N=n/c$, and let
\[
        A_0=mc+s,
        \qquad 0\le s<c,
        \qquad 0\le m\le N,
        \qquad A_0\ge K .
\]
If $s>0$, assume $mc+s\le n$.  Put $\sigma=A_0-K$ and
\[
        M_{c,m,s}:=\binom Nm\binom{n-mc}{s}.
\]
Then there is a $\B$-valued received word $U:D\to\B$ such that
\[
        \bigl|\Lst(\RS[\F,D,K],\,1-A_0/n,\,U)\bigr|
        \ge
        \left\lceil \frac{M_{c,m,s}}{|\B|^{w_c(s,\sigma)}}\right\rceil .
\]
The listed codewords may be chosen in $\RS[\B,D,K]$.
\end{lemma}

\begin{proof}
For $M\subseteq Q_c$ with $|M|=m$, write
\[
        P_M(X)=\prod_{b\in M}(X^c-b).
\]
For $R\subseteq D\setminus \pi_c^{-1}(M)$ with $|R|=s$, put
\[
        E_R(X)=\prod_{r\in R}(X-r),
        \qquad
        L_{M,R}(X)=P_M(X)E_R(X).
\]
Then $L_{M,R}$ is a squarefree locator of degree $A_0$ whose roots are exactly
\(\pi_c^{-1}(M)\cup R\).  Write
\[
        L_{M,R}(X)=X^{A_0}+\sum_{h=1}^{A_0}\lambda_h(M,R)X^{A_0-h}.
\]
The coefficients lie in $\B$.  Moreover, because $P_M$ has monomials only in
degrees congruent to $0$ modulo $c$ and $E_R$ has degree $s$, the coefficient
$\lambda_h(M,R)$ is zero unless $h\bmod c\in\{0,1,\ldots,s\}$.
Consequently the high-degree prefix
\[
        \bigl(\lambda_h(M,R):1\le h\le\sigma,
              \ h\bmod c\in\{0,1,\ldots,s\}\bigr)
\]
takes at most $|\B|^{w_c(s,\sigma)}$ values.

Pigeonhole over these prefix values.  There is a fixed prefix vector $\xi$ and a
subfamily $\mathcal T_\xi$ of at least
$M_{c,m,s}/|\B|^{w_c(s,\sigma)}$ pairs $(M,R)$ with this prefix.  Let
\[
        U_\xi(X)=X^{A_0}+\sum_{1\le h\le\sigma}\lambda_h(\xi)X^{A_0-h}
\]
be the corresponding high-degree part, restricted to $D$.  For every
$(M,R)\in\mathcal T_\xi$, the difference
\[
        R_{M,R}(X):=L_{M,R}(X)-U_\xi(X)
\]
has degree $<K$.  Hence the codeword $-R_{M,R}$ agrees with $U_\xi$ on all
roots of $L_{M,R}$, a set of size $A_0$.

It remains only to check that the listed codewords are distinct.  If two pairs
$(M,R)$ and $(M',R')$ in the same prefix class gave the same codeword on $D$,
then the two remainders of degree $<K<n$ would be equal as polynomials.  Since
the high prefixes are also equal, $L_{M,R}=L_{M',R'}$.  Their root sets in $D$
are therefore equal.  A full fiber has size $c$ while $|R|=s<c$, so the full
fibers contained in this root set are exactly $M$, and then the residual set is
forced as well.  Thus $(M,R)=(M',R')$.
\end{proof}

\begin{definition}[prefix fibers and prefix-image certificates]
\label{def:prefix-fiber-certificate}
In the notation of \cref{lem:quotient-remainder-prefix}, let
$\Phi_{c,m,s}^{K}$ be the map from pairs $(M,R)$ to the high-degree prefix
\[
        \bigl(\lambda_h(M,R):1\le h\le A_0-K,
              \ h\bmod c\in\{0,1,\ldots,s\}\bigr)
        \in \B^{w_c(s,A_0-K)} .
\]
If $A_0=K$ this is the unique empty prefix.  For a prefix value $\xi$, put
\[
        \mu_{c,m,s}^{K}(\xi)
        :=\bigl| (\Phi_{c,m,s}^{K})^{-1}(\xi)\bigr|.
\]
Define the image size and the heaviest prefix fiber by
\[
        I_{c,m,s}^{K}:=\bigl|\operatorname{im}\Phi_{c,m,s}^{K}\bigr|,
        \qquad
        H_{c,m,s}^{K}:=\max_{\xi}\mu_{c,m,s}^{K}(\xi).
\]
Thus
\[
        H_{c,m,s}^{K}
        \ge
        \left\lceil\frac{M_{c,m,s}}{I_{c,m,s}^{K}}\right\rceil
        \ge
        \left\lceil\frac{M_{c,m,s}}{|\B|^{w_c(s,A_0-K)}}\right\rceil .
\]
A prefix-fiber certificate is an orbit decomposition, hash table, or other
verifiable bucket table for the finite map $\Phi_{c,m,s}^{K}$, including the
image size and the maximum fiber size.  If only $I_{c,m,s}^{K}$ is certified,
the middle lower bound is still valid; if no image certificate is supplied, the
ambient-box lower bound is valid.
\end{definition}

\begin{lemma}[heaviest-prefix locator floor]
\label{lem:heaviest-prefix-locator-floor}
Under the hypotheses of \cref{lem:quotient-remainder-prefix}, there is a
$\B$-valued received word $U:D\to\B$ such that
\[
        \bigl|\Lst(\RS[\F,D,K],\,1-A_0/n,\,U)\bigr|
        \ge H_{c,m,s}^{K}.
\]
Consequently
\[
        \bigl|\Lst(\RS[\F,D,K],\,1-A_0/n,\,U)\bigr|
        \ge
        \left\lceil \frac{M_{c,m,s}}{I_{c,m,s}^{K}}\right\rceil
        \ge
        \left\lceil \frac{M_{c,m,s}}{|\B|^{w_c(s,A_0-K)}}\right\rceil
\]
when only the coarser certificates are available.  The listed codewords may be
chosen in $\RS[\B,D,K]$.
\end{lemma}

\begin{proof}
Choose a prefix value $\xi$ with
$\mu_{c,m,s}^{K}(\xi)=H_{c,m,s}^{K}$ and form the high-degree word
\[
        U_\xi(X)=X^{A_0}+\sum_{1\le h\le A_0-K}
        \lambda_h(\xi)X^{A_0-h},
\]
where $\lambda_h(\xi)=0$ for the high coefficients outside the allowed
residue classes.  For every pair $(M,R)$ in this heaviest prefix fiber, the difference
$L_{M,R}(X)-U_\xi(X)$ has degree $<K$, so the negative of this remainder is a
codeword of $\RS[\B,D,K]$ agreeing with $U_\xi$ on the $A_0$ roots of
$L_{M,R}$.  If two pairs in the same prefix fiber gave the same codeword on
$D$, then their degree-$<K<n$ remainders would be equal as polynomials; together
with the common prefix this gives $L_{M,R}=L_{M',R'}$.  Their root sets in $D$
are equal.  A full quotient fiber has size $c$ while the residual set has size
$s<c$, so the full fibers contained in the root set determine $M$, and then the
residual set determines $R$.  Thus the codewords are pairwise distinct.  The two
coarser displayed bounds follow from \cref{def:prefix-fiber-certificate}.
\end{proof}

For $A\in\{k+1,
\ldots,n\}$ and $c\mid n$, write
\[
        A=m_c(A)c+s_c(A),
        \qquad 0\le s_c(A)<c,
\]
and set
\[
\begin{aligned}
        M_c(A)&:=
        \binom{n/c}{m_c(A)}
        \binom{n-cm_c(A)}{s_c(A)},\\
        L_{c,{\rm heavy}}^+(A)&:=
        H_{c,m_c(A),s_c(A)}^{k+1},\\
        L_{c,{\rm image}}^+(A)&:=
        \left\lceil\frac{M_c(A)}{I_{c,m_c(A),s_c(A)}^{k+1}}\right\rceil,\\
        L_{c,{\rm box}}^+(A)&:=
        \left\lceil
        \frac{M_c(A)}{|\B|^{w_c(s_c(A),A-k-1)}}
        \right\rceil .
\end{aligned}
\]
The three levels satisfy
\[
        L_{c,{\rm heavy}}^+(A)
        \ge L_{c,{\rm image}}^+(A)
        \ge L_{c,{\rm box}}^+(A).
\]
Here ``heavy'' uses a full prefix-fiber certificate, ``image'' uses only the
prefix-image size, and ``box'' is the certificate-free bound already implicit in
\cref{lem:quotient-remainder-prefix}.

\begin{theorem}[quotient-remainder deep-point lower ledger]
\label{thm:quotient-remainder-deep-floor}
Let $\B\subseteq\F$, let $q=|\F|>n$, let
$D\subseteq\B^\times$ be a multiplicative coset of order $n$, and let
$C=\RS[\F,D,k]$.  Fix an agreement value $A\in\{k+1,\ldots,n\}$ and a
divisor $c\mid n$.  For each
\[
        \star\in\{{\rm heavy},{\rm image},{\rm box}\}
\]
and every
\[
        \delta\in[1-A/n,\,1-k/n)
\]
one has
\[
        \eca(C,\delta)
        \ge \mathcal E_{q,k}\bigl(L_{c,\star}^+(A)\bigr).
\]
If $\delta>0$, then also
\[
        \emca(C,\delta)
        \ge \mathcal E_{q,k}\bigl(L_{c,\star}^+(A)\bigr).
\]
Here $\mathcal E_{q,k}$ is as in \cref{cor:deep-list-trigger-ceiling}.  More
precisely, some received line has at least
\[
        \mathcal N_{q,k}\bigl(L_{c,\star}^+(A)\bigr)
\]
CA-bad finite slopes at that radius, and for $\delta>0$ these slopes are
MCA-bad.

Consequently, for a closed agreement threshold $a\ge k+1$, define
\[
        B_{Q,\star}^-(a)
        :=\max_{\substack{c\mid n\\ A\in\{a,a+1,\ldots,n\}}}
          \mathcal N_{q,k}\bigl(L_{c,\star}^+(A)\bigr).
\]
Then $B_{Q,\star}^-(a)$ is a lower numerator for $q\,\eca(C,\delta)$ for
every $\delta\in[1-a/n,1-k/n)$, and for $q\,\emca(C,\delta)$ on the
same interval with $\delta>0$.  The strongest
certified ledger uses $\star={\rm heavy}$; the certificate-free ledger uses
$\star={\rm box}$.
\end{theorem}

\begin{proof}
Apply \cref{lem:heaviest-prefix-locator-floor} with $K=k+1$ and $A_0=A$.  This
produces a $C^+=\RS[\F,D,k+1]$ list of size at least
$L_{c,{\rm heavy}}^+(A)$ at radius $1-A/n$, with the image and ambient-box
variants giving the corresponding smaller certified list sizes.  Then apply
\cref{prop:quantitative-deep-list-floor}.  If $A\ge a$ and
$\delta\ge1-a/n$, then $\delta\ge1-A/n$, so the same converted floor applies.
The maximum over $(c,A)$ is the correct lower-ledger operation: the definitions
of $\eca$ and $\emca$ maximize over received lines, so lower witnesses from
different divisor scales or exact agreement values need not share one line and
must not be summed.
\end{proof}

\begin{corollary}[quotient-remainder trigger]
\label{cor:quotient-remainder-trigger}
In the setting of \cref{thm:quotient-remainder-deep-floor}, if
$L_{c,\star}^+(A)>(q-n)/k$ for any
$\star\in\{{\rm heavy},{\rm image},{\rm box}\}$, then, for every
$\delta\in[1-A/n,1-k/n)$,
\[
        \eca(C,\delta)>
        \frac1{2k}\left(1-\frac nq\right),
\]
and for every such $\delta>0$,
\[
        \emca(C,\delta)>
        \frac1{2k}\left(1-\frac nq\right).
\]
\end{corollary}

\begin{proof}
This is \cref{cor:deep-list-trigger-ceiling} applied inside
\cref{thm:quotient-remainder-deep-floor}.
\end{proof}

\begin{corollary}[quantitative first-grid floor]
\label{cor:quantitative-first-grid-floor}
Let $C=\RS[\F,D,k]$, let $q=|\F|>n$, and assume only that $D\subseteq\F$ has
$n$ distinct points.  Then for every
$\delta\in[1-(k+1)/n,1-k/n)$,
\[
        \eca(C,\delta)
        \ge
        \mathcal E_{q,k}\left(\binom n{k+1}\right),
\]
and for every such $\delta>0$,
\[
        \emca(C,\delta)
        \ge
        \mathcal E_{q,k}\left(\binom n{k+1}\right).
\]
In particular, if $\binom n{k+1}>(q-n)/k$, then the printed
$(2k)^{-1}(1-n/q)$ lower bound holds on this entire first grid interval without
any multiplicative smoothness assumption.
\end{corollary}

\begin{proof}
Use the ordinary locator identity for all $(k+1)$-subsets of $D$; this is the
$c=1$ argument behind \cref{lem:heaviest-prefix-locator-floor}, but it does not
use multiplicative smoothness.  The prefix is empty, so the list size is at
least
$\binom n{k+1}$.  Apply \cref{prop:quantitative-deep-list-floor} and then
\cref{cor:deep-list-trigger-ceiling}.
\end{proof}

\begin{definition}[quotient lower-ledger certificate]
\label{def:quotient-lower-ledger-certificate}
A quotient-remainder lower-ledger certificate for a row and a threshold $a$
contains, for each divisor scale $c$ and each exact agreement $A\ge a$ that is
used in the maximum:
\begin{enumerate}[label=\textup{(\roman*)},leftmargin=2.2em]
\item the decomposition $A=m_c(A)c+s_c(A)$;
\item the support count $M_c(A)$ and prefix weight $w_c(s_c(A),A-k-1)$;
\item either a full prefix-fiber bucket table certifying
$H_{c,m_c(A),s_c(A)}^{k+1}$, or an image-size certificate for
$I_{c,m_c(A),s_c(A)}^{k+1}$, or the declaration that the ambient-box denominator
is being used;
\item the resulting list floor $L_{c,\star}^+(A)$;
\item the converted slope numerator
$\mathcal N_{q,k}(L_{c,\star}^+(A))$ and error floor
$\mathcal E_{q,k}(L_{c,\star}^+(A))$.
\end{enumerate}
The verifier checks the locator-prefix formula, distinctness of codewords inside
a prefix fiber, the integer arithmetic, and the final maximum over divisor
scales and exact agreement values.
\end{definition}

\begin{remark}[quotient lower ledger versus quotient upper ledger]
\label{rem:T1-lower-ledger-status}
The quotient-profile lower side is now scanner-ready at every agreement value
$A=mc+s$: the generated-field contribution is the actual heaviest prefix fiber
$H_{c,m,s}^{k+1}$ when certified, the exact prefix-image denominator
$I_{c,m,s}^{k+1}$ when only image data are certified, and the ambient box
$|\B|^{w_c(s,A-k-1)}$ when no certificate is supplied.  The challenge-field size
$q$ enters only through the independent conversion functions
$\mathcal E_{q,k}$ and $\mathcal N_{q,k}$.  Across divisor scales the lower-side
operation is a maximum, not a union count, because the witness line may change.
The remaining upper-ledger problem is sharper and unchanged in logical status:
replace the safe support ledger and the exact declared-branch image ledger by an
exhaustive quotient-periodic classification, with duplicate slopes coalesced
across all divisor scales.
\end{remark}

\begin{corollary}[augmented slack-one quotient floor]
\label{cor:augmented-slack-one}
Let $\B\subseteq\F$ and let $D\subseteq\B^\times$ be a multiplicative coset
of order $n$.  Put $C=\RS[\F,D,k]$ and $C^+=\RS[\F,D,k+1]$.  If
$c\mid\gcd(n,k)$, then there is a $\B$-valued received word $U_c$ such that
\[
        \bigl|\Lst(C^+,\,1-(k+c)/n,\,U_c)\bigr|
        \ge
        \binom{n/c}{k/c+1} .
\]
One may take $U_c(x)=x^{k+c}$.  In particular, for $c=1$ this gives the
universal first-grid floor
\[
        \bigl|\Lst(C^+,\,1-(k+1)/n,\,x^{k+1})\bigr|
        \ge
        \binom n{k+1},
\]
with no quotient denominator.  The $c=1$ statement in fact holds for any set
$D\subseteq\B$ of $n$ distinct points, not only for multiplicative cosets.
\end{corollary}

\begin{proof}
For $c>1$, apply Lemma~\ref{lem:quotient-remainder-prefix} with $K=k+1$,
$s=0$, and $m=k/c+1$.  Then $A_0=k+c$ and $\sigma=A_0-K=c-1$, so
$w_c(0,c-1)=0$.  The high-degree prefix is only the monomial $X^{k+c}$, giving
the displayed word.  For $c=1$, the same argument is just the ordinary locator
identity for all $(k+1)$-subsets of $D$: the high-degree prefix of a monic
locator of degree $k+1$ is $X^{k+1}$, and its remainder has degree at most $k$.
\end{proof}

\begin{corollary}[first-grid deep-point cap]
\label{cor:first-grid-cap}
Let $C=\RS[\F,D,k]$, $q=|\F|$, and suppose $q>n$.  Let
$c\mid\gcd(n,k)$.  If $c>1$, assume $D\subseteq\B^\times$ is a
multiplicative coset as above; for $c=1$, no smoothness is needed.  If
\[
        \binom{n/c}{k/c+1}\ge q/k+1,
\]
then
\[
        \eca\bigl(C,1-(k+c)/n\bigr)
        >\frac1{2k}\left(1-\frac nq\right),
        \qquad
        \emca\bigl(C,1-(k+c)/n\bigr)
        >\frac1{2k}\left(1-\frac nq\right).
\]
For $c=1$ this is the first closed grid point below capacity.
\end{corollary}

\begin{proof}
Corollary~\ref{cor:augmented-slack-one} gives a $C^+$-list of size
$L\ge q/k+1$ at agreement $k+c$.  The proof of the deep-point conversion
Theorem~\ref{thm:A} applies whenever the agreement is $>k$, so it applies here
for every $c\ge1$.  If the displayed CA lower bound failed, then
\[
        \eca(C,1-(k+c)/n)\le \frac{q-n}{2kq}.
\]
Theorem~\ref{thm:A} with $\eta=1/2$ would give
\[
        \Lst(C^+,1-(k+c)/n)
        \le
        \left\lceil \frac{q-n}{k}\right\rceil
        < q/k+1,
\]
contradicting the list lower bound.  Every no-loss CA-bad slope is also
support-wise MCA-bad on the same explaining support, giving the MCA inequality.
\end{proof}

\begin{corollary}[first-grid cap for large challenge-envelope rows]
\label{cor:first-grid-grand}
Let $\rho\in\{\frac12,\frac14,\frac18,\frac1{16}\}$, let
$C=\RS[\F,D,k]$ with $|D|=n$, $k=\rho n$, $q=|\F|<2^{256}$, and
$q>n$.  Assume
\[
\begin{array}{c|cccc}
\rho & \frac12 & \frac14 & \frac18 & \frac1{16}\\ \hline
k\ge & 127 & 78 & 58 & 47
\end{array}
\]
for the corresponding rate.  Then the first closed grid point below capacity is
already CA- and MCA-unsafe:
\[
        \eca\bigl(C,1-\rho-1/n\bigr)
        >\frac1{2k}\left(1-\frac nq\right),
        \qquad
        \emca\bigl(C,1-\rho-1/n\bigr)
        >\frac1{2k}\left(1-\frac nq\right).
\]
If also $k\le2^{40}$, then this lower bound is $>2^{-86}$ and hence is far
above the prize target $2^{-128}$.
\end{corollary}

\begin{proof}
By Corollary~\ref{cor:first-grid-cap} with $c=1$, it is enough to check
\[
        \binom n{k+1}\ge q/k+1.
\]
Since $q<2^{256}$, the stronger sufficient condition is
$k(\binom n{k+1}-1)\ge2^{256}$.  The four printed lower bounds on $k$ are exactly
small finite checks of this inequality at the corresponding rate; the left side
is increasing thereafter.  For the error size, $q\ge n+1$, $n=k/\rho\le16k$,
and $k\le2^{40}$ give
\[
        \frac1{2k}\left(1-\frac nq\right)
        \ge \frac1{2k(n+1)}
        >2^{-86}.
\]
\end{proof}


\subsection{A conservative quotient-support upper ledger}
\label{sec:quotient-support-upper}

The remainder lower floor above supplies the unsafe side of the quotient-profile
ledger.  For a prize threshold one also needs a safe-side accounting rule: if a
later argument declares a bad parameter to be quotient-periodic, the declaration
must come with a divisor-lattice numerator that is counted with overlaps under
control.  The next statements give such a conservative upper ledger.  They do
not prove that every bad line parameter is quotient-periodic; they only say that
once a permitted support family has been specified, every ledger term witnessed
inside that family is bounded by the size of the family, with a factor $d$ for
finite-parameter degree-$d$ power curves.

Let $\mathcal S$ be a family of subsets of $D$, all of size at least $a$.  For a
received line $f+zg$, say that a finite slope is
\emph{$\mathcal S$-witnessed support-wise bad} if it has a support-wise
noncontained explanation on some $S\in\mathcal S$.  Define the analogous
$\mathcal S$-witnessed no-loss CA slopes by requiring that $f+zg$ be explained
on some $S\in\mathcal S$ while $(f,g)$ is not globally explained on any support
of size at least $a$.

\begin{lemma}[one support pays for at most one line parameter]
\label{lem:one-support-one-line}
Let $C=\RS[\F,D,k]$ and let $a\ge k$.  Fix a support
$S\subseteq D$ with $|S|\ge a$.  For any received line $f+zg$, at most one finite
slope $z\in\F$ can be support-wise noncontained on the fixed support $S$.
Likewise, at most one finite slope can be no-loss CA-bad while being explained on
this fixed support $S$.  Consequently, for every support family $\mathcal S$,
\[
        \#\{\mathcal S\text{-witnessed support-wise MCA-bad finite slopes}\}
        \le |\mathcal S|,
\]
and the same bound holds for no-loss CA slopes witnessed by $\mathcal S$.
\end{lemma}

\begin{proof}
Suppose two distinct finite slopes $z_1\ne z_2$ are both explained on the same
support $S$, by codewords $p_1,p_2\in C$:
\[
        f+z_i g=p_i\qquad\text{on }S,\quad i=1,2.
\]
Then on $S$,
\[
        g=\frac{p_1-p_2}{z_1-z_2},\qquad
        f=p_1-z_1g.
\]
Both right-hand sides are codewords of $C$.  Thus $(f,g)$ is simultaneously
explained on $S$.  This contradicts support-wise noncontainment on $S$.  It also
contradicts no-loss CA badness at agreement $a$, because $|S|\ge a$ gives a
global simultaneous explanation support of the forbidden size.  Hence each fixed
support contributes at most one bad finite line parameter, and summing over
supports gives the displayed bound.
\end{proof}

For a finite-parameter power curve
\[
        W_\gamma=f_0+\gamma f_1+\cdots+\gamma^d f_d,
        \qquad \gamma\in\F,
\]
say that a parameter is support-wise curve-bad on $S$ if $W_\gamma$ is explained
by a codeword on $S$, but the coefficient tuple $(f_0,\ldots,f_d)$ is not
simultaneously explained by $C^{\equiv(d+1)}$ on $S$.  The no-loss curve-CA
variant replaces this support-wise noncontainment by the corresponding global
far condition at agreement $a$.

\begin{lemma}[one support pays for at most $d$ curve parameters]
\label{lem:one-support-d-curve}
Let $C=\RS[\F,D,k]$, let $a\ge k$, and fix a support $S\subseteq D$ with
$|S|\ge a$.  For a degree-$d$ power curve $W_\gamma$, at most $d$ finite
parameters $\gamma\in\F$ can be support-wise curve-bad on $S$.  The same bound
holds for no-loss curve-CA parameters explained on $S$.  Consequently any
support family $\mathcal S$ witnesses at most $d|\mathcal S|$ bad finite curve
parameters.
\end{lemma}

\begin{proof}
If $d+1$ distinct parameters $\gamma_0,\ldots,\gamma_d$ were explained on the
same support $S$, with codewords $P_j\in C$, then Lagrange interpolation in the
parameter variable gives codewords $Q_0,\ldots,Q_d\in C$ such that
\[
        \sum_{i=0}^d \gamma_j^i Q_i=P_j
        \qquad (0\le j\le d).
\]
On every point of $S$, the same interpolation applied to the values of
$W_{\gamma_j}$ gives $f_i=Q_i$ for every coefficient $i$.  Hence
$(f_0,\ldots,f_d)$ is simultaneously explained by $C^{\equiv(d+1)}$ on $S$.
This contradicts support-wise curve-noncontainment on $S$, and also contradicts
no-loss curve-CA badness at agreement $a$.  Thus every support contributes at
most $d$ parameters.
\end{proof}

\begin{lemma}[one support pays for at most one interleaved tuple]
\label{lem:one-support-one-interleaved}
Let $C=\RS[\F,D,k]$, let $a\ge k$, and let $\mu\ge1$.  Fix a $\mu$-row received
word $U=(U_1,\ldots,U_\mu)$.  For a fixed support $S\subseteq D$ with
$|S|\ge a$, there is at most one tuple $(P_1,\ldots,P_\mu)\in C^\mu$ agreeing
with $U$ on every row over $S$.  Hence the number of interleaved list tuples whose
common agreement support belongs to a support family $\mathcal S$ is at most
$|\mathcal S|$.
\end{lemma}

\begin{proof}
For each row $i$, two degree-$<k$ polynomials agreeing with $U_i$ on the fixed
support $S$ agree with each other on at least $k$ points, so they are equal.  Thus
each row has at most one codeword compatible with $S$, and therefore the whole
interleaved tuple is unique if it exists.
\end{proof}

We now instantiate the support-family ledger on quotient-remainder supports.  For
$c\mid n$ and an integer $A$ with $0\le A\le n$, write
\[
        A=m_c(A)c+s_c(A),\qquad 0\le s_c(A)<c,
\]
and define
\[
\begin{aligned}
        \mathcal S_c(A)
        :=\{&\pi_c^{-1}(M)\cup R:
          M\subseteq Q_c,
          |M|=m_c(A),\\
        & R\subseteq D\setminus\pi_c^{-1}(M),
          |R|=s_c(A)\} .
\end{aligned}
\]
For a divisor set $\mathcal C\subseteq\{c:c\mid n\}$ put
\[
        \mathcal S_{\mathcal C}^{\ge a}
        :=\bigcup_{c\in\mathcal C}\bigcup_{A=a}^n \mathcal S_c(A)
\]
and
\[
        U_{\mathcal C}^{\rm supp}(a)
        :=\left|\mathcal S_{\mathcal C}^{\ge a}\right|.
\]
When the exact union is inconvenient, use the safe sum
\[
\label{eq:quotient-support-safe-sum}
        U_{\mathcal C}^{\rm sum}(a)
        :=\sum_{c\in\mathcal C}\sum_{A=a}^n
          \binom{n/c}{\lfloor A/c\rfloor}
          \binom{n-c\lfloor A/c\rfloor}{A-c\lfloor A/c\rfloor},
\]
so that $U_{\mathcal C}^{\rm supp}(a)\le U_{\mathcal C}^{\rm sum}(a)$.

\begin{proposition}[divisor-union quotient support ledger]
\label{prop:divisor-union-support-ledger}
Let $C=\RS[\F,D,k]$ and let $a\ge k$.  Fix a divisor set $\mathcal C$ and let
$\mathcal S_{\mathcal C}^{\ge a}$ be the quotient-remainder support family above.
Then, for every received line, the number of finite support-wise MCA-bad slopes
whose witness support lies in $\mathcal S_{\mathcal C}^{\ge a}$ is at most
\[
        U_{\mathcal C}^{\rm supp}(a)
        \le U_{\mathcal C}^{\rm sum}(a).
\]
The same bound holds for no-loss CA slopes witnessed by this family.  For a
finite-parameter degree-$d$ power curve, the corresponding support-wise curve-MCA
and no-loss curve-CA numerator is at most
\[
        d\,U_{\mathcal C}^{\rm supp}(a)
        \le d\,U_{\mathcal C}^{\rm sum}(a).
\]
For every interleaving arity $\mu\ge1$, the number of interleaved list tuples whose
common agreement support belongs to $\mathcal S_{\mathcal C}^{\ge a}$ is at most
$U_{\mathcal C}^{\rm supp}(a)$.
\end{proposition}

\begin{proof}
Apply Lemmas~\ref{lem:one-support-one-line}, \ref{lem:one-support-d-curve}, and
\ref{lem:one-support-one-interleaved} to the support family
$\mathcal S_{\mathcal C}^{\ge a}$.  The inequality
$U_{\mathcal C}^{\rm supp}(a)\le U_{\mathcal C}^{\rm sum}(a)$ is only the union
bound over the divisor and agreement-size indices.
\end{proof}

\begin{remark}[what this does and does not finish]
\label{rem:quotient-upper-status}
Proposition~\ref{prop:divisor-union-support-ledger} supplies a certified
safe-side quotient ledger at the support-family level, and it is already enough
for a staircase scanner to avoid double-counting divisor scales: one may compute
the exact union $U_{\mathcal C}^{\rm supp}$ or use the printed safe sum.  This is
not yet the sharp quotient upper theorem needed for a full prize solution.  Near
capacity the support count can be exponentially larger than the number of
distinct bad slopes, so the next refinement is a slope-collision quotient ledger:
classify when many quotient-remainder supports map to the same line or curve
parameter, and replace the support count by the corresponding distinct-parameter
union count.
\end{remark}


\begin{definition}[block support-profile polynomial]
\label{def:block-support-profile-polynomial}
Let $D$ be a multiplicative coset of order $n$, let
$\mathcal C$ be a nonempty finite set of divisors of $n$, and put
\[
        h:=\operatorname{lcm}\mathcal C,
        \qquad N_h:=n/h .
\]
For each $c\in\mathcal C$ let $H_c\le H_h$ denote the subgroup of the cyclic
kernel of order $h$ with $|H_c|=c$.  If $T\subseteq H_h$, define
\[
        \phi_c(T):=\#\{\text{cosets }uH_c\subseteq H_h:\ uH_c\subseteq T\},
\]
the number of complete $c$-subfibers contained in the block subset $T$.
The block support-profile polynomial of $\mathcal C$ is
\[
        G_{\mathcal C,h}(Y,(Z_c)_{c\in\mathcal C})
        :=\sum_{T\subseteq H_h}
          Y^{|T|}\prod_{c\in\mathcal C}Z_c^{\phi_c(T)} .
\]
For a vector $\nu=(\nu_c)_{c\in\mathcal C}$ write
$Z^\nu:=\prod_{c\in\mathcal C}Z_c^{\nu_c}$, and define
$N_{\mathcal C,h}(A;\nu)$ by
\[
        G_{\mathcal C,h}(Y,Z)^{N_h}
        =\sum_{A=0}^{n}\sum_{\nu}
          N_{\mathcal C,h}(A;\nu)Y^A Z^\nu .
\]
\end{definition}

\begin{lemma}[block polynomial counts exact quotient profiles]
\label{lem:block-polynomial-counts-profiles}
For a support $S\subseteq D$ and $c\mid n$, let
\[
        F_c(S):=\#\{\text{$c$-fibers of }D\to D^c\text{ contained in }S\}.
\]
Then $N_{\mathcal C,h}(A;\nu)$ is exactly the number of supports
$S\subseteq D$ such that
\[
        |S|=A,
        \qquad F_c(S)=\nu_c\quad(c\in\mathcal C).
\]
\end{lemma}

\begin{proof}
The $h$-fibers partition $D$ into $N_h$ blocks, and every $c$-fiber with
$c\in\mathcal C$ is contained in a unique $h$-fiber because $c\mid h$.  After
choosing an identification of each $h$-fiber with the cyclic kernel $H_h$, the
number of selected points and the number of complete internal $c$-subfibers are
recorded by the monomial
\[
        Y^{|T|}\prod_{c\in\mathcal C}Z_c^{\phi_c(T)}
\]
for the block subset $T$.  The choices in distinct $h$-fibers are independent,
and the global statistics $|S|$ and $F_c(S)$ are the sums of the block statistics.
Multiplying the block polynomial over the $N_h$ blocks and extracting the
coefficient gives the asserted count.
\end{proof}

\begin{lemma}[membership by full-fiber profile]
\label{lem:quotient-support-membership-profile}
Fix $A$ and $c\mid n$, and write
\[
        A=m_c(A)c+s_c(A),
        \qquad 0\le s_c(A)<c .
\]
For a support $S\subseteq D$ of size $A$, one has
\[
        S\in\mathcal S_c(A)
        \qquad\Longleftrightarrow\qquad
        F_c(S)=m_c(A)=\lfloor A/c\rfloor .
\]
\end{lemma}

\begin{proof}
If $S\in\mathcal S_c(A)$, then by definition it is the union of
$m_c(A)$ complete $c$-fibers and $s_c(A)<c$ residual points outside those fibers;
there is no room for any further complete $c$-fiber.  Hence
$F_c(S)=m_c(A)$.  Conversely, if $F_c(S)=m_c(A)$, take the complete $c$-fibers
contained in $S$ as the set $M$ in the definition of $\mathcal S_c(A)$; the
remaining $A-m_c(A)c=s_c(A)$ selected points form the residual set.
\end{proof}

\begin{theorem}[exact divisor-lattice support-union formula]
\label{thm:exact-divisor-block-support-ledger}
Let $\mathcal C$ be a nonempty finite set of divisors of $n$, and define
$G_{\mathcal C,h}$ and $N_{\mathcal C,h}(A;\nu)$ as in
\cref{def:block-support-profile-polynomial}.  Then the exact support-family
union from \cref{prop:divisor-union-support-ledger} is
\[
\boxed{
        U_{\mathcal C}^{\rm supp}(a)
        =\sum_{A=a}^{n}\sum_{\nu}
          N_{\mathcal C,h}(A;\nu)\,
          \mathbf 1\!
          \left[\exists c\in\mathcal C:
                 \nu_c=\left\lfloor A/c\right\rfloor\right].
      }
\]
Consequently the quotient-support branch numerators in
\cref{prop:divisor-union-support-ledger} may be computed by this coefficient
formula instead of by the safe sum
$U_{\mathcal C}^{\rm sum}(a)$.
\end{theorem}

\begin{proof}
For a fixed exact size $A$, the union
$\bigcup_{c\in\mathcal C}\mathcal S_c(A)$ consists precisely of those supports
$S$ of size $A$ for which $S\in\mathcal S_c(A)$ for at least one divisor
$c\in\mathcal C$.  By \cref{lem:quotient-support-membership-profile}, this is
exactly the condition
\[
        F_c(S)=\lfloor A/c\rfloor
\]
for at least one $c\in\mathcal C$.  By
\cref{lem:block-polynomial-counts-profiles}, the number of supports with a fixed
profile vector $\nu=(F_c(S))_{c\in\mathcal C}$ is
$N_{\mathcal C,h}(A;\nu)$.  Summing over all exact sizes $A\ge a$ and over the
profile vectors satisfying the displayed predicate gives the formula.
\end{proof}

\begin{corollary}[single-divisor sanity check]
\label{cor:single-divisor-support-ledger}
For $\mathcal C=\{c\}$ the block polynomial is
\[
        G_{\{c\},c}(Y,Z)=(1+Y)^c+(Z-1)Y^c,
\]
and \cref{thm:exact-divisor-block-support-ledger} gives
\[
        \left|\mathcal S_c(A)\right|
        =\binom{n/c}{\lfloor A/c\rfloor}
          \binom{n-c\lfloor A/c\rfloor}{A-c\lfloor A/c\rfloor},
\]
the summand used in the safe ledger.
\end{corollary}

\begin{proof}
Inside one $c$-fiber, every proper subset contributes only $Y^{|T|}$, while the
full subset contributes $Y^cZ$.  Thus the block polynomial is
$(1+Y)^c-Y^c+Y^cZ$.  Raising to $n/c$ and selecting exactly
$m=\lfloor A/c\rfloor$ full blocks and $s=A-mc$ residual points in the remaining
blocks gives the displayed count.
\end{proof}

\begin{definition}[divisor-block support-ledger certificate]
\label{def:divisor-block-support-ledger-certificate}
A divisor-block support-ledger certificate for a divisor set $\mathcal C$ and a
threshold $a$ consists of:
\begin{enumerate}[label=\textup{(\roman*)},leftmargin=2.2em]
\item the lcm block size $h=\operatorname{lcm}\mathcal C$ and the block count
$N_h=n/h$;
\item the internal coset table of every subgroup $H_c\le H_h$ for
$c\in\mathcal C$;
\item the block polynomial $G_{\mathcal C,h}$, either as a full coefficient table
or as an orbit-compressed table under translations of $H_h$;
\item the coefficient extraction from $G_{\mathcal C,h}^{N_h}$ for every exact
size $A\ge a$;
\item the predicate table
$\mathbf 1[\exists c\in\mathcal C:\nu_c=\lfloor A/c\rfloor]$ and the final
integer $U_{\mathcal C}^{\rm supp}(a)$.
\end{enumerate}
The verifier checks the block enumeration, the subgroup containment relation
$c\mid h$, coefficient arithmetic, and the final predicate sum.  This is an exact
support-union certificate; it does not by itself coalesce distinct line or curve
parameters that arise from different supports.
\end{definition}

\begin{remark}[support ledger versus slope ledger]
\label{rem:support-vs-slope-ledger}
\Cref{thm:exact-divisor-block-support-ledger} closes the divisor-lattice overlap
problem at the support-family level: intersections between quotient descriptions
at incomparable divisor scales are now counted by one coefficient extraction in
the common $h$-block.  This is still weaker than the desired distinct-slope
quotient ledger.  To finish that stronger ledger one must apply the exact
support-image maps of
\cref{prop:distinct-parameter-line-ledger,prop:distinct-parameter-curve-ledger}
to the support set counted here and then prove or certify the resulting parameter
image size after duplicate coalescing.
\end{remark}

\subsection{Distinct-parameter quotient image ledger}
\label{sec:quotient-distinct-parameter}

The support-union ledger of \cref{sec:quotient-support-upper} is deliberately
safe but often very loose: many quotient-remainder supports can certify the same
line slope or the same curve parameter.  The next refinement is to count the
image of the actual parameter map.  This subsection gives an exact image ledger
for the quotient branch.  It is still a branch ledger, not an aperiodic upper
bound: it counts only witnesses whose agreement support lies in a declared
quotient-remainder support family.

We first record the support-to-parameter map in syndrome form.  Let
$C=\RS[\F,D,k]$, $|D|=n$, and put $r=n-k$.  For each $x\in D$ set
\[
        \lambda_x:=\left(\prod_{y\in D,\ y\ne x}(x-y)\right)^{-1}.
\]
For a word $Y:D\to\F$, define its syndrome vector
\[
        \Syn(Y)_m:=\sum_{x\in D}\lambda_x x^mY(x),
        \qquad 0\le m<r .
\]
Thus $\Syn(Y)=0$ if and only if $Y\in C$.  If
$T\subseteq D$ has size $j\le r$, write
\[
        L_T(X)=\prod_{x\in T}(X-x)=\ell_0+\ell_1X+\cdots+\ell_jX^j,
        \qquad
        \boldsymbol\ell_T=(\ell_0,\ldots,\ell_j)^t,
\]
and put $t=r-j$.  For $w\in\F^r$ let $H_{t,j}(w)$ be the $t\times(j+1)$ Hankel
window
\[
        H_{t,j}(w)_{a,b}=w_{a+b},
        \qquad 0\le a<t,\quad 0\le b\le j .
\]
When $S=D\setminus T$, the integer $t$ is exactly $|S|-k$.


\begin{lemma}[support-locator syndrome recurrence]
\label{lem:support-locator-syndrome-recurrence}
Let $Y:D\to\F$ and let $T\subseteq D$ with $|T|=j\le r$.  Put
$S=D\setminus T$ and $t=r-j$.  Then $Y$ is explained by $C$ on $S$ if and only if
\[
        H_{t,j}(\Syn(Y))\boldsymbol\ell_T=0 .
\]
\end{lemma}

\begin{proof}
Let $W_T\subseteq\F^r$ be the span of the parity-check columns indexed by $T$,
\[
        W_T=\operatorname{span}_\F\{\lambda_x(1,x,\ldots,x^{r-1}):x\in T\}.
\]
The word $Y$ is explained by $C$ on $S$ if and only if $Y-c$ is supported on $T$
for some $c\in C$, equivalently $\Syn(Y)\in W_T$.  If
$w\in W_T$, write $w_m=\sum_{x\in T}a_xx^m$.  Then, for $0\le h<t$,
\[
        \bigl(H_{t,j}(w)\boldsymbol\ell_T\bigr)_h
        =\sum_{b=0}^j\ell_b w_{h+b}
        =\sum_{x\in T}a_xx^hL_T(x)=0.
\]
Conversely, because $L_T$ is monic, the displayed recurrence determines all
coordinates $w_j,\ldots,w_{r-1}$ from the first $j$ coordinates, so its solution
space has dimension at most $j$.  The $j$ parity-check columns indexed by the
distinct points of $T$ have Vandermonde rank $j$, hence $W_T$ has dimension $j$
and is exactly the recurrence space.
\end{proof}

\begin{lemma}[exact affine and projective support-image map]
\label{lem:exact-line-image-map}
Let $f,g:D\to\F$, put $u=\Syn(f)$ and $v=\Syn(g)$, and fix
$T\subseteq D$ with $|T|=j\le r$ and $S=D\setminus T$.  Define
\[
        A_T:=H_{t,j}(u)\boldsymbol\ell_T,
        \qquad
        B_T:=H_{t,j}(v)\boldsymbol\ell_T .
\]
A finite slope $z\in\F$ is explained on the fixed support $S$ and is
support-wise noncontained there if and only if
\[
        A_T+zB_T=0,
        \qquad B_T\ne0 .
\]
Consequently the fixed support contributes either no finite bad slope or the
unique value
\[
        \theta_T=-A_{T,m}/B_{T,m}
\]
for any coordinate $m$ with $B_{T,m}\ne0$, provided all nonzero coordinates give
the same value.

For the projective line $\alpha f+\beta g$, the fixed support contributes a
support-wise noncontained point $[\alpha:\beta]\in\mathbb P^1(\F)$ if and only if
$A_T$ and $B_T$ are linearly dependent and not both zero.  In that case the
contributed projective parameter is the unique point satisfying
\[
        \alpha A_T+\beta B_T=0 .
\]
\end{lemma}

\begin{proof}
A word $Y$ is explained by $C$ on $S=D\setminus T$ if and only if
$Y-c$ is supported on $T$ for some $c\in C$, which is equivalent to
$\Syn(Y)$ lying in the span of the parity-check columns indexed by $T$.  The
standard locator recurrence says that this is equivalent to
\[
        H_{t,j}(\Syn(Y))\boldsymbol\ell_T=0 .
\]
Applying this to $Y=f+zg$ gives $A_T+zB_T=0$.  Under this incidence, simultaneous
explanation of $f$ and $g$ on $S$ is equivalent to $B_T=0$, because
$A_T=(A_T+zB_T)-zB_T$.  This proves the finite-slope statement.  The projective
statement is the same calculation with $Y=\alpha f+\beta g$; if
$A_T=B_T=0$ then both coefficient words are already explained on $S$, so the
point is contained rather than bad.
\end{proof}

Let $\mathcal S_{\mathcal C}^{\ge a}$ be the quotient-remainder support family
from \cref{sec:quotient-support-upper}, and let
\[
        \mathcal T_{\mathcal C}^{\ge a}
        :=\{D\setminus S:S\in\mathcal S_{\mathcal C}^{\ge a}\}
\]
be the corresponding co-support family.  For a received line define
\[
        \Theta_{\mathcal C}^{\rm aff}(f,g;a)
        :=\{\theta_T:T\in\mathcal T_{\mathcal C}^{\ge a},
              \theta_T\text{ is defined by }\cref{lem:exact-line-image-map}\}
\]
and define $\Theta_{\mathcal C}^{\rm proj}(f,g;a)$ analogously using the
projective point of \cref{lem:exact-line-image-map}.

\begin{proposition}[distinct-parameter quotient ledger for lines]
\label{prop:distinct-parameter-line-ledger}
For every received line $f+zg$, the finite affine slopes that are support-wise
MCA-bad and have a witness support in
$\mathcal S_{\mathcal C}^{\ge a}$ are exactly the elements of
$\Theta_{\mathcal C}^{\rm aff}(f,g;a)$.  Hence the exact quotient-branch line
numerator is
\[
        |\Theta_{\mathcal C}^{\rm aff}(f,g;a)|
        \le U_{\mathcal C}^{\rm supp}(a).
\]
The corresponding projective-slope numerator is
\[
        |\Theta_{\mathcal C}^{\rm proj}(f,g;a)|
        \le U_{\mathcal C}^{\rm supp}(a).
\]
The same image ledgers upper-bound no-loss CA parameters whose explaining
support lies in the same quotient family.
\end{proposition}

\begin{proof}
This is only \cref{lem:exact-line-image-map} applied to every co-support in the
family, followed by taking a union of the resulting parameter values.  The CA
claim follows because no-loss CA-badness is stronger than having an explained
slope on the support and forbids simultaneous explanation on any support of size
at least $a$.
\end{proof}

The curve analogue is just as explicit.  For a degree-$d$ power curve
\[
        W_\gamma=f_0+\gamma f_1+\cdots+\gamma^d f_d,
        \qquad \gamma\in\F,
\]
put $u_i=\Syn(f_i)$ and, for $T\subseteq D$, define
\[
        V_{i,T}:=H_{t,j}(u_i)\boldsymbol\ell_T\in\F^t .
\]
Let
\[
        \Gamma_T^{(d)}
        :=\{\gamma\in\F:
             V_{0,T}+\gamma V_{1,T}+\cdots+\gamma^dV_{d,T}=0,
             \ (V_{0,T},\ldots,V_{d,T})\ne0\} .
\]
The nonzero tuple condition is exactly support-wise noncontainment of the
coefficient tuple on the fixed support.

\begin{lemma}[exact support-image map for finite-parameter power curves]
\label{lem:exact-curve-image-map}
For a fixed support $S=D\setminus T$, the set of support-wise noncontained
curve-bad finite parameters is exactly $\Gamma_T^{(d)}$.  In particular
$|\Gamma_T^{(d)}|\le d$.
\end{lemma}

\begin{proof}
The same syndrome-locator recurrence applied to $W_\gamma$ gives the displayed
vector equation.  If every $V_{i,T}$ is zero, then every coefficient word
$f_i$ is already explained on $S$, so the parameter is contained.  Otherwise at
least one scalar coordinate of
$V_{0,T}+\gamma V_{1,T}+\cdots+\gamma^dV_{d,T}$ is a nonzero polynomial in
$\gamma$ of degree at most $d$, so it has at most $d$ roots.
\end{proof}

Define the exact quotient curve image
\[
        \Gamma_{\mathcal C}^{(d)}(f_0,\ldots,f_d;a)
        :=\bigcup_{T\in\mathcal T_{\mathcal C}^{\ge a}}\Gamma_T^{(d)} .
\]

\begin{proposition}[distinct-parameter quotient ledger for curves]
\label{prop:distinct-parameter-curve-ledger}
For a finite-parameter degree-$d$ power curve, the quotient-branch support-wise
curve-MCA numerator is exactly
\[
        |\Gamma_{\mathcal C}^{(d)}(f_0,
        \ldots,f_d;a)| .
\]
It satisfies the safe bounds
\[
        |\Gamma_{\mathcal C}^{(d)}(f_0,\ldots,f_d;a)|
        \le d\,U_{\mathcal C}^{\rm supp}(a)
        \le d\,U_{\mathcal C}^{\rm sum}(a).
\]
The same image ledger upper-bounds no-loss curve-CA parameters whose explaining
support lies in the quotient family.
\end{proposition}

\begin{proof}
Apply \cref{lem:exact-curve-image-map} to each support and take the union of the
resulting parameter sets.  The displayed inequalities follow from
$|\Gamma_T^{(d)}|\le d$ and from the support-union bound of
\cref{prop:divisor-union-support-ledger}.
\end{proof}

\begin{definition}[fixed-support parameter gcd]
\label{def:fixed-support-parameter-gcd}
Let $C=\RS[\F,D,k]$, let $q=|\F|$, and let
\[
        W_\Gamma=f_0+\Gamma f_1+\cdots+\Gamma^d f_d,
        \qquad \Gamma\in\F,
\]
be a finite-parameter degree-$d$ power curve of received words.  Fix a
co-support $T\subseteq D$, put $S=D\setminus T$, and assume
$|S|>k$.  With the notation of \cref{lem:exact-curve-image-map}, write
\[
        P_{T,m}(\Gamma)
        :=\sum_{i=0}^d \Gamma^i (V_{i,T})_m,
        \qquad 0\le m<t,
        \qquad t=|S|-k .
\]
Define the fixed-support contribution polynomial
$g_T^{(d)}(\Gamma)\in\F[\Gamma]$ as follows.  If all coordinate polynomials
$P_{T,m}$ are identically zero, set $g_T^{(d)}=1$.  Otherwise let
$g_T^{(d)}$ be the monic greatest common divisor of the nonzero polynomials
among $P_{T,0},\ldots,P_{T,t-1}$.
\end{definition}

\begin{lemma}[fixed-support gcd equals the noncontained parameter set]
\label{lem:fixed-support-gcd-exact}
In the setting of \cref{def:fixed-support-parameter-gcd}, the finite parameters
$\gamma\in\F$ for which $W_\gamma$ is explained on $S=D\setminus T$ and the
coefficient tuple $(f_0,\ldots,f_d)$ is not simultaneously explained on $S$ are
exactly the roots in $\F$ of $g_T^{(d)}$.  In particular
\[
        \deg g_T^{(d)}\le d,
\]
with equality possible only before taking common factors across different
supports.
\end{lemma}

\begin{proof}
By \cref{lem:exact-curve-image-map}, the fixed support explains $W_\gamma$ exactly
when
\[
        P_{T,m}(\gamma)=0\qquad(0\le m<t).
\]
If every coordinate polynomial is identically zero, then every coefficient
vector $V_{i,T}$ is zero, so every coefficient word $f_i$ is already explained on
$S$; this is the contained branch and contributes no noncontained parameter.
The convention $g_T^{(d)}=1$ gives the empty root set, as required.

Otherwise the common zeros of the coordinate polynomials are exactly the roots
of their monic gcd after zero coordinates are discarded.  Since at least one
coordinate polynomial is nonzero of degree at most $d$, the gcd also has degree
at most $d$.  The nonzero tuple condition in \cref{lem:exact-curve-image-map} is
precisely the assertion that the coordinate polynomials are not all zero.
\end{proof}

\begin{definition}[quotient image lcm polynomial]
\label{def:quotient-image-lcm-polynomial}
Let $\mathcal T$ be any finite family of co-supports $T\subseteq D$ with
$|D\setminus T|>k$.  For a degree-$d$ finite-parameter power curve define
\[
        R_{\mathcal T}^{(d)}(\Gamma)
        :=\operatorname{rad}\operatorname{lcm}_{T\in\mathcal T} g_T^{(d)}(\Gamma),
\]
where factors equal to $1$ are ignored and the empty lcm is $1$.  For the quotient branch at agreement
threshold $a>k$ and divisor set $\mathcal C$, we use
\[
        \mathcal T=\mathcal T_{\mathcal C}^{\ge a}
        :=\{D\setminus S:S\in\mathcal S_{\mathcal C}^{\ge a}\}.
\]
The finite-field root-count polynomial is
\[
        R_{\mathcal T,F}^{(d)}(\Gamma)
        :=\gcd\bigl(R_{\mathcal T}^{(d)}(\Gamma),\Gamma^q-\Gamma\bigr).
\]
Because $R_{\mathcal T}^{(d)}$ is squarefree, the degree of
$R_{\mathcal T,F}^{(d)}$ is the number of $\F$-rational roots of
$R_{\mathcal T}^{(d)}$.
\end{definition}

\begin{theorem}[exact finite-parameter quotient image ledger]
\label{thm:exact-quotient-image-lcm-ledger}
Let $a>k$, let $\mathcal C$ be a finite divisor set, and let
$\mathcal T_{\mathcal C}^{\ge a}$ be the quotient co-support family.  For every
finite-parameter degree-$d$ power curve, the exact quotient-branch support-wise
curve-MCA numerator is
\[
        \boxed{
        \left|\Gamma_{\mathcal C}^{(d)}(f_0,\ldots,f_d;a)\right|
        =\deg R_{\mathcal T_{\mathcal C}^{\ge a},F}^{(d)} .
        }
\]
Consequently
\[
        \left|\Gamma_{\mathcal C}^{(d)}(f_0,\ldots,f_d;a)\right|
        \le \deg R_{\mathcal T_{\mathcal C}^{\ge a}}^{(d)}
        \le d\,U_{\mathcal C}^{\rm supp}(a),
\]
and the final support count may be evaluated by the exact divisor-block formula
of \cref{thm:exact-divisor-block-support-ledger}.  The same lcm ledger
upper-bounds no-loss curve-CA parameters whose explaining support lies in the
same quotient family.
\end{theorem}

\begin{proof}
By \cref{lem:fixed-support-gcd-exact}, for each fixed co-support $T$ the
noncontained parameters witnessed on $D\setminus T$ are exactly the $\F$-roots of
$g_T^{(d)}$.  Taking the union over
$T\in\mathcal T_{\mathcal C}^{\ge a}$ is therefore the same as taking the
$\F$-rational root set of the squarefree lcm
$R_{\mathcal T_{\mathcal C}^{\ge a}}^{(d)}$.  Over a finite field, the
$\F$-rational roots of a squarefree polynomial are counted by the degree of its
gcd with $\Gamma^q-\Gamma$, proving the displayed equality.

The degree bound follows from $\deg g_T^{(d)}\le d$ and from the fact that a
squarefree lcm has degree at most the sum of the degrees of the factors before
coalescing.  The no-loss curve-CA claim follows because a no-loss CA parameter
with an explaining support in the quotient family is, on that support,
noncontained; otherwise the coefficient tuple would have a simultaneous
explanation on a support of size at least $a$.
\end{proof}

\begin{corollary}[affine lines as the degree-one lcm ledger]
\label{cor:affine-line-lcm-ledger}
For an affine received line $f+Zg$, take $d=1$, $f_0=f$, and $f_1=g$.  Then
$R_{\mathcal T_{\mathcal C}^{\ge a}}^{(1)}(Z)$ is the squarefree product of the
linear factors $Z-z$ for the distinct finite quotient-witnessed support-wise
MCA-bad slopes.  Hence
\[
        \#\{\text{finite quotient-witnessed MCA-bad slopes}\}
        =\deg R_{\mathcal T_{\mathcal C}^{\ge a}}^{(1)} .
\]
The same polynomial upper-bounds no-loss CA finite slopes whose explaining
support lies in the declared quotient family.
\end{corollary}

\begin{proof}
For $d=1$, each nonconstant fixed-support gcd has degree at most one.  Its root,
if present, is exactly the unique finite slope supplied by
\cref{lem:exact-line-image-map}.  Taking the squarefree lcm over the declared
co-support family therefore multiplies one copy of $Z-z$ for each distinct
finite slope and no copies for supports that contribute no finite noncontained
slope.
\end{proof}

\begin{definition}[projective line lcm ledger]
\label{def:projective-line-lcm-ledger}
For a projective received line $Z_0f+Z_1g$ and a co-support $T$ with
$|D\setminus T|>k$, form the homogeneous linear coordinate forms
\[
        L_{T,m}(Z_0,Z_1):=Z_0(A_T)_m+Z_1(B_T)_m,
        \qquad 0\le m<t .
\]
If all $L_{T,m}$ are zero, set $h_T^{\rm proj}=1$.  Otherwise let
$h_T^{\rm proj}$ be the gcd of the nonzero linear forms, normalized by a fixed
monomial order.  Thus $h_T^{\rm proj}$ is either $1$ or a homogeneous linear
form.  For a co-support family $\mathcal T$ define
\[
        R_{\mathcal T}^{\rm proj}(Z_0,Z_1)
        :=\operatorname{rad}\operatorname{lcm}_{T\in\mathcal T}h_T^{\rm proj}.
\]
\end{definition}

\begin{corollary}[exact projective quotient image ledger]
\label{cor:projective-line-lcm-ledger}
For projective slopes, the exact quotient-branch support-wise MCA numerator is
\[
        \deg R_{\mathcal T_{\mathcal C}^{\ge a}}^{\rm proj}.
\]
The same degree upper-bounds no-loss projective-CA parameters whose explaining
support lies in the quotient family.
\end{corollary}

\begin{proof}
A nonzero family of homogeneous linear forms on $\mathbb P^1$ has a common zero
if and only if the forms have a common homogeneous linear factor; this factor
cuts out exactly one $\F$-rational projective point.  The all-zero case is the
contained branch and contributes no bad point, by convention.  The squarefree
lcm therefore contains one linear factor for each distinct projective parameter
and no duplicates.
\end{proof}

\begin{definition}[lcm quotient-image certificate]
\label{def:lcm-quotient-image-certificate}
An lcm quotient-image certificate for a line or finite-parameter curve consists
of:
\begin{enumerate}[label=\textup{(\roman*)},leftmargin=2.2em]
\item the divisor set $\mathcal C$, the agreement threshold $a$, and either the
exact support-union certificate of \cref{def:divisor-block-support-ledger-certificate}
or an explicit hash/orbit enumeration of the co-support family
$\mathcal T_{\mathcal C}^{\ge a}$;
\item for every represented co-support $T$, the locator vector
$\boldsymbol\ell_T$ and the syndrome-window coefficient vectors $V_{i,T}$;
\item the coordinate polynomials $P_{T,m}(\Gamma)$, the fixed-support gcd
$g_T^{(d)}$, and a proof that the all-zero case has been assigned to the
contained branch;
\item the squarefree lcm $R_{\mathcal T}^{(d)}$ and, over the finite field, the
root-count factor $R_{\mathcal T,F}^{(d)}=\gcd(R_{\mathcal T}^{(d)},\Gamma^q-\Gamma)$;
\item the declared numerator $\deg R_{\mathcal T,F}^{(d)}$ for finite curves or
$\deg R_{\mathcal T}^{\rm proj}$ for projective lines.
\end{enumerate}
The verifier checks the locator split condition, the syndrome recurrence, the
gcd and lcm arithmetic, squarefreeness, and the final finite-field root count.
This certificate replaces a support count by an exact distinct-parameter count
inside the declared quotient branch.
\end{definition}

\begin{remark}[status after the lcm image ledger]
\label{rem:lcm-image-ledger-status}
The quotient ledgers now separate three tasks that were previously entangled.
First, \cref{thm:quotient-remainder-deep-floor} supplies lower floors by
heaviest generated-field prefix fibers.  Second,
\cref{thm:exact-divisor-block-support-ledger} counts the declared quotient
support union with divisor overlaps handled exactly.  Third,
\cref{thm:exact-quotient-image-lcm-ledger,cor:affine-line-lcm-ledger,cor:projective-line-lcm-ledger}
push that support family through the fixed-locator equations and coalesce
duplicate finite, projective, and curve parameters by gcd/lcm arithmetic.  The
remaining quotient-side obstruction is not duplicate coalescing inside a declared
quotient family; it is the structural safe-side theorem that every bad parameter
in the intended regime either belongs to this quotient image, to the tangent
branch, or to the aperiodic/extension buckets bounded elsewhere.
\end{remark}

\begin{definition}[quotient image certificate]
\label{def:quotient-image-certificate}
A quotient image certificate for a line or curve ledger consists of:
\begin{enumerate}[label=\textup{(\roman*)}]
\item the divisor set $\mathcal C$ and the agreement threshold $a$;
\item a deterministic description, hash, or orbit decomposition of
$\mathcal T_{\mathcal C}^{\ge a}$;
\item for every represented co-support $T$, the locator coefficients
$\boldsymbol\ell_T$ and the syndrome-window vectors $A_T,B_T$ for a line, or
$V_{0,T},\ldots,V_{d,T}$ for a curve;
\item the declared parameter value, parameter bucket, or the declaration that no
parameter is contributed;
\item a bucket table proving the number of distinct parameters after duplicate
coalescing.
\end{enumerate}
The verifier checks the locator split condition, the Hankel equations, the
noncontainment vector, and the bucket consistency.  This certificate is the
preferred replacement for a raw quotient-support count whenever the support count
is much larger than the actual parameter image.
\end{definition}

\begin{remark}[status of the refinement]
\label{rem:distinct-ledger-status}
The propositions above are exact for the declared quotient branch and are a real
refinement over the support ledger: the numerator is an image size, not a support
count.  They still do not prove that all bad parameters are quotient-periodic.
The remaining safe-side theorem for the full prize is the aperiodic Hankel-packing
bound: after tangent and quotient image branches are removed, one must bound the
parameters coming from co-supports outside the declared quotient families.
\end{remark}


\subsection{Aperiodic Hankel chart atlas}
\label{sec:aperiodic-hankel-certificates}

The quotient image ledger of \cref{sec:quotient-distinct-parameter} counts bad
parameters whose witness support lies in a declared quotient-remainder family.
The complementary safe-side problem is to control all other supports.  This
subsection is the contributor-facing formulation of that problem.  It does not claim
that the needed aperiodic eliminants always exist.  It gives a finite chart
atlas, with precise ideals and pivots, so that a verifier can check a supplied
eliminant and convert it into a bad-parameter upper bound.

The main point is to avoid one huge saturated incidence ideal.  We split the
aperiodic branch into:
\begin{enumerate}[label=\textup{(\roman*)},leftmargin=2.2em]
\item a regular overdetermined Hankel bucket, where a single maximal minor in the
line parameter already gives an eliminant;
\item finite affine pivot charts, one for each noncontainment coordinate;
\item a separate projective-infinity chart, if projective slopes are being
counted;
\item finite-parameter curve pivot charts, one for a nonzero coefficient of one
coordinate polynomial;
\item a singular bucket, where all regular minors vanish or the projection to the
parameter line remains positive-dimensional after the previous splits.
\end{enumerate}
Only the last bucket is genuinely unresolved.  A staircase scanner should keep it
as a named obstruction rather than hiding it inside the word ``aperiodic.''

Fix an exact agreement value $A\ge k$ and put
\[
        j=j_A:=n-A,
        \qquad
        t=t_A:=A-k,
        \qquad
        R=n-k .
\]
Let
\[
        P_D(X):=\prod_{x\in D}(X-x).
\]
A monic degree-$j$ co-support locator is written
\[
        L_{\ell}(X)=X^j+\ell_{j-1}X^{j-1}+\cdots+\ell_0,
        \qquad \ell=(\ell_0,\ldots,\ell_{j-1})\in\mathbb A^j .
\]
Let $R_{D,j}(\ell;X)$ be the remainder of $P_D(X)$ modulo $L_{\ell}(X)$, and let
$I_{\rm split,j}$ be the ideal generated by the $j$ coefficients of
$R_{D,j}$.  Since $P_D$ is squarefree, the $F$-points of
$V(I_{\rm split,j})$ are exactly the squarefree $D$-split monic degree-$j$
locators.

For a word $y:D\to F$ write $\Syn(y)\in F^R$ for the parity-check syndrome used
above.  Let
\[
        \widetilde\ell=(\ell_0,\ldots,\ell_{j-1},1)^t .
\]
For an affine line $f+zg$, put $u=\Syn(f)$ and $v=\Syn(g)$ and define
\[
        A(\ell):=H_{t,j}(u)\widetilde\ell,
        \qquad
        B(\ell):=H_{t,j}(v)\widetilde\ell
        \qquad\in F^t .
\]
For a fixed split locator, the support-wise finite-slope condition is exactly
\[
        A(\ell)+ZB(\ell)=0,
        \qquad
        B(\ell)\ne0 .
\]
The inequality is the noncontainment condition.

\subsubsection{Regular overdetermined bucket}

Assume first that
\[
        t_A\ge j_A+1,
        \qquad\text{equivalently}\qquad
        2A\ge n+k+1 .
\]
Set
\[
        M_A(Z):=H_{t_A,j_A}(u)+Z H_{t_A,j_A}(v).
\]
This is a $t_A\times(j_A+1)$ matrix with entries affine-linear in $Z$.

\begin{definition}[regular Hankel minor certificate]
\label{def:hankel-regularity-certificate}
A regular Hankel minor certificate at exact agreement $A$ is a row set
$I_A\subseteq\{0,\ldots,t_A-1\}$ of size $j_A+1$ such that
\[
        \Delta_A(Z)
        :=\det M_A(Z)_{I_A,\{0,\ldots,j_A\}}
\]
is not the zero polynomial in $F[Z]$.  If no such row set exists, the line is
called \emph{Hankel-singular at $A$}.
\end{definition}

\begin{lemma}[regular exact-agreement eliminant]
\label{lem:regular-exact-agreement-eliminant}
Assume $t_A\ge j_A+1$, and suppose $\Delta_A(Z)$ is a regular Hankel minor
certificate.  Every finite support-wise MCA-bad slope with a witness support of
exact size $A$ is a root of $\Delta_A$.  Hence the number of such slopes is at
most
\[
        \deg\Delta_A\le j_A+1=n-A+1 .
\]
The same root containment holds for no-loss CA-bad slopes with exact agreement
$A$.
\end{lemma}

\begin{proof}
Let $z$ be a bad finite slope with witness support $S$ of size $A$, and put
$T=D\setminus S$.  Then $|T|=j_A$, and the locator vector
$\widetilde\ell_T$ is a nonzero kernel vector of $M_A(z)$.  Therefore the column
rank of $M_A(z)$ is less than $j_A+1$, so every $(j_A+1)\times(j_A+1)$ maximal
minor vanishes at $z$, including $\Delta_A(z)$.  The degree bound follows because
$\Delta_A$ is the determinant of a $(j_A+1)\times(j_A+1)$ matrix whose entries
are affine-linear in $Z$.
\end{proof}

\begin{theorem}[regular closed-range Hankel packing]
\label{thm:regular-closed-ball-hankel-packing}
Let $a\ge k$ and suppose that for each exact agreement
$A\in\{a,a+1,\ldots,n\}$ with $t_A\ge j_A+1$ a regular Hankel minor certificate
$\Delta_A$ is supplied.  Then all bad finite slopes whose exact witness size is
in this regular bucket lie in the root set of
\[
        \Delta_{\ge a}(Z):=\prod_{A=a}^{n}\Delta_A(Z),
\]
with factors omitted for exact agreements outside the overdetermined range.
Consequently their number is at most
\[
        \sum_A \deg\Delta_A .
\]
If exact root sets are supplied, roots already charged to the tangent or
quotient-image ledgers may be removed from this count by duplicate coalescing.
\end{theorem}

\begin{proof}
Apply \cref{lem:regular-exact-agreement-eliminant} at the exact size of the
witness support and take the union of the resulting root sets.  The degree sum is
a safe union bound.
\end{proof}

\paragraph{Canonical regular rank-drop gcds.}
The regular-minor certificate below can be made canonical.  Instead of choosing
one nonzero maximal minor, take the greatest common divisor of all nonzero
maximal minors.  This gives exactly the finite-parameter values at which the
Hankel pencil can drop rank, before the split-locator and noncontainment tests
are imposed.  Thus it is an upper ledger for the regular bucket, and often a much
smaller one than the degree of a single selected minor.

Fix an exact agreement value $A\ge k$, and keep the notation
\[
        j=j_A:=n-A,
        \qquad
        t=t_A:=A-k,
        \qquad
        r=n-k .
\]
The regular overdetermined range is
\[
        t_A>j_A,
        \qquad\text{i.e.}\qquad
        t_A\ge j_A+1,
        \qquad\text{equivalently}\qquad
        2A\ge n+k+1 .
\]

\begin{definition}[canonical affine rank-drop gcd]
\label{def:canonical-affine-rankdrop-gcd}
Fix an affine received line $f+Zg$, put
$u=\Syn(f)$ and $v=\Syn(g)$, and assume $t_A\ge j_A+1$.  Set
\[
        M_A(Z):=H_{t_A,j_A}(u)+Z H_{t_A,j_A}(v).
\]
Let $\mathfrak M_A(Z)$ be the finite set of all $(j_A+1)\times(j_A+1)$ maximal
minors of $M_A(Z)$.  If every element of $\mathfrak M_A(Z)$ is the zero
polynomial, the exact-agreement bucket is called \emph{rank-drop singular}.
Otherwise define
\[
        G_A^{\rm aff}(Z)
        :=\operatorname{rad}\gcd\{\Delta(Z):
              \Delta\in\mathfrak M_A(Z),\ \Delta\ne0\},
\]
with the gcd normalized to be monic.  Over the finite field $\F$ of size $q$ set
\[
        G_{A,F}^{\rm aff}(Z):=
        \gcd\bigl(G_A^{\rm aff}(Z),Z^q-Z\bigr).
\]
Thus $G_{A,F}^{\rm aff}$ is the squarefree polynomial whose roots are precisely
those $F$-rational parameters at which all nonzero maximal minors vanish.
\end{definition}

\begin{theorem}[canonical exact-agreement rank-drop ledger]
\label{thm:canonical-affine-rankdrop-ledger}
Assume $t_A\ge j_A+1$ and the exact-agreement bucket is not rank-drop singular.
Every finite support-wise MCA-bad slope with an exact witness support of size
$A$ is a root of $G_A^{\rm aff}$.  Consequently the number of such slopes is at
most
\[
        \deg G_{A,F}^{\rm aff}
        \le \deg G_A^{\rm aff}
        \le j_A+1=n-A+1 .
\]
The same bound holds for no-loss CA-bad slopes with exact agreement $A$.
\end{theorem}

\begin{proof}
Let $z\in F$ be support-wise bad with exact witness support $S$ of size $A$, and
put $T=D\setminus S$.  By \cref{lem:support-locator-syndrome-recurrence}, the
locator vector $\boldsymbol\ell_T$ is a nonzero kernel vector of $M_A(z)$.
Therefore $\operatorname{rank}M_A(z)<j_A+1$, so every maximal minor of $M_A(Z)$
vanishes at $Z=z$.  In particular every nonzero maximal minor vanishes at $z$,
and hence their gcd $G_A^{\rm aff}$ vanishes at $z$.

The finite-field root count is exactly $\deg G_{A,F}^{\rm aff}$ because
$G_{A,F}^{\rm aff}$ is squarefree and divides $Z^q-Z$.  Each maximal minor has
degree at most $j_A+1$, since it is the determinant of a matrix whose entries are
affine-linear in $Z$; the gcd cannot have larger degree.  The no-loss CA claim is
identical, because no-loss CA-badness with exact agreement $A$ also supplies an
explaining support of size $A$ and hence the same locator-kernel condition.
\end{proof}

\begin{remark}[why this improves the one-minor certificate]
\label{rem:gcd-rankdrop-improves-one-minor}
\Cref{lem:regular-exact-agreement-eliminant} used one certified nonzero maximal
minor and therefore counted the roots of a single determinant.  The polynomial
$G_A^{\rm aff}$ uses all maximal minors and counts their common finite-field root
set.  It is never worse than any one-minor certificate and can be much smaller;
it is also canonical, so two scanners that use the same syndrome convention print
the same regular-bucket polynomial.
\end{remark}

\begin{definition}[closed-ball rank-drop lcm]
\label{def:closed-ball-rankdrop-lcm}
For an agreement threshold $a>k$, let
\[
        \mathcal A_{\rm reg}(a):=
        \{A\in\{a,a+1,\ldots,n\}: t_A\ge j_A+1
          \text{ and the bucket }A\text{ is not rank-drop singular}\}.
\]
Define
\[
        G_{\ge a}^{\rm aff}(Z)
        :=\operatorname{rad}\operatorname{lcm}_{A\in\mathcal A_{\rm reg}(a)}
          G_A^{\rm aff}(Z),
        \qquad
        G_{\ge a,F}^{\rm aff}(Z):=
        \gcd\bigl(G_{\ge a}^{\rm aff}(Z),Z^q-Z\bigr),
\]
with the empty lcm equal to $1$.  Exact agreement values
$A\in\{a,\ldots,n\}\setminus\mathcal A_{\rm reg}(a)$ are recorded as residual
buckets: they are either underdetermined ($t_A<j_A+1$) or rank-drop singular.
\end{definition}

\begin{corollary}[canonical regular closed-ball Hankel packing]
\label{cor:canonical-regular-closed-ball-hankel-packing}
All finite support-wise MCA-bad slopes whose exact witness size lies in
$\mathcal A_{\rm reg}(a)$ are roots of $G_{\ge a}^{\rm aff}$.  Hence their number
is at most
\[
        \deg G_{\ge a,F}^{\rm aff}
        \le \sum_{A\in\mathcal A_{\rm reg}(a)}(n-A+1).
\]
The same bound holds for no-loss CA-bad slopes in the same regular buckets.
\end{corollary}

\begin{proof}
Apply \cref{thm:canonical-affine-rankdrop-ledger} at the exact agreement size of
the witness and take the union over $A$.  The squarefree lcm coalesces duplicate
roots across exact-agreement buckets; the final gcd with $Z^q-Z$ counts only
finite-field roots.  The displayed degree sum is the safe bound obtained before
coalescing.
\end{proof}

\begin{definition}[paid-root removal for the regular branch]
\label{def:paid-root-removal-regular-branch}
Suppose a squarefree polynomial $P_{\rm paid,F}(Z)$ is supplied whose roots are
exactly the finite-field parameters already charged to tangent/common-line and
quotient-image ledgers.  Define the regular residual polynomial
\[
        G_{{\rm reg}\setminus{\rm paid},F}^{\rm aff}(Z)
        :=\frac{G_{\ge a,F}^{\rm aff}(Z)}
        {\gcd(G_{\ge a,F}^{\rm aff}(Z),P_{\rm paid,F}(Z))} .
\]
Its degree is the number of regular rank-drop parameters not already paid by the
previous ledgers.
\end{definition}

\begin{corollary}[quotient/tangent-removed regular aperiodic numerator]
\label{cor:quotient-tangent-removed-regular-numerator}
If $P_{\rm paid,F}$ is exact in the sense of
\cref{def:paid-root-removal-regular-branch}, then the number of finite slopes in
the regular overdetermined branch that remain after removing the paid
quotient/tangent parameters is at most
\[
        \deg G_{{\rm reg}\setminus{\rm paid},F}^{\rm aff} .
\]
This is a safe aperiodic numerator for the regular buckets.  The underdetermined
and rank-drop singular exact-agreement buckets must still be handled by pivot
charts or by a separately named residual obstruction.
\end{corollary}

\begin{proof}
By \cref{cor:canonical-regular-closed-ball-hankel-packing}, every regular-bucket
bad slope is a root of $G_{\ge a,F}^{\rm aff}$.  Removing the roots of the exact
paid polynomial deletes precisely those already charged parameters from this
finite root set.  The remaining root set is counted by the displayed quotient
because all polynomials are squarefree divisors of $Z^q-Z$.
\end{proof}

\begin{definition}[finite-parameter curve rank-drop gcd]
\label{def:curve-rankdrop-gcd}
For a finite-parameter degree-$d$ curve
\[
        W_\Gamma=f_0+\Gamma f_1+\cdots+\Gamma^d f_d,
\]
put $u_i=\Syn(f_i)$ and, at exact agreement $A$, define
\[
        M_A^{(d)}(\Gamma):=
        \sum_{i=0}^{d}\Gamma^i H_{t_A,j_A}(u_i).
\]
When $t_A\ge j_A+1$ and not all maximal minors of $M_A^{(d)}$ vanish
identically, set
\[
        G_A^{(d)}(\Gamma)
        :=\operatorname{rad}\gcd\{\Delta(\Gamma):
              \Delta\text{ is a nonzero maximal minor of }M_A^{(d)}\},
\]
and
\[
        G_{A,F}^{(d)}(\Gamma):=
        \gcd\bigl(G_A^{(d)}(\Gamma),\Gamma^q-\Gamma\bigr).
\]
\end{definition}

\begin{theorem}[canonical curve rank-drop ledger]
\label{thm:canonical-curve-rankdrop-ledger}
In a non-singular overdetermined exact-agreement bucket, every finite
curve-MCA-bad parameter is a root of $G_A^{(d)}$.  Consequently the number of
such parameters is at most
\[
        \deg G_{A,F}^{(d)}
        \le \deg G_A^{(d)}
        \le d(j_A+1).
\]
The same bound holds for no-loss curve-CA parameters in the same bucket.  The
closed-ball form is obtained by taking the squarefree lcm of the polynomials
$G_A^{(d)}$ over the regular exact-agreement buckets and then taking the gcd with
$\Gamma^q-\Gamma$.
\end{theorem}

\begin{proof}
A curve-bad parameter with exact witness support $S$ gives, for
$T=D\setminus S$, a nonzero locator vector in the kernel of
$M_A^{(d)}(\gamma)$.  Hence every maximal minor vanishes at $\gamma$, and the
same gcd argument as in \cref{thm:canonical-affine-rankdrop-ledger} applies.
Each matrix entry has degree at most $d$ in $\Gamma$, so each maximal minor has
degree at most $d(j_A+1)$.
\end{proof}

\begin{definition}[regular rank-drop certificate]
\label{def:regular-rankdrop-certificate}
A regular rank-drop certificate for a row, a received affine line or finite
curve, and an agreement threshold $a$ consists of:
\begin{enumerate}[label=\textup{(\roman*)},leftmargin=2.2em]
\item the syndrome vectors of the coefficient words and the exact-agreement range
$A\in\{a,\ldots,n\}$;
\item for every overdetermined exact agreement $A$, either a declaration that all
maximal minors vanish identically, or the canonical gcd polynomial
$G_A^{\rm aff}$, respectively $G_A^{(d)}$, together with ideal-membership or raw
minor arithmetic proving that it is the gcd of all nonzero maximal minors;
\item the squarefree lcm over all non-singular regular buckets and the finite-root
factor obtained by gcd with $Z^q-Z$ or $\Gamma^q-\Gamma$;
\item if paid branches are removed, the exact paid-root polynomial and the finite
root-set quotient of \cref{def:paid-root-removal-regular-branch};
\item a residual-bucket table listing all underdetermined and rank-drop singular
exact agreements that are not covered by the regular gcd ledger.
\end{enumerate}
The verifier checks the syndrome convention, the Hankel matrices, the maximal
minor gcds, squarefree lcms, finite-field root factors, paid-root removal, and the
residual-bucket table.
\end{definition}

\begin{remark}[status after the rank-drop ledger]
\label{rem:rankdrop-ledger-status}
The regular overdetermined part of the aperiodic Hankel problem is now canonical:
one no longer has to choose a minor or accept a degree sum when the common root
set of all maximal minors is smaller.  This still does not solve the near-capacity
underdetermined buckets, which form the hard aperiodic residual class.
Those buckets remain assigned to the affine pivot charts, projective-infinity
charts, curve coefficient pivots, or named singular residual obstructions.
\end{remark}

\subsubsection{Finite affine pivot charts}

The regular minor test is intentionally coarse and only works directly in the
overdetermined range.  The next atlas is the one contributors should use to
construct actual aperiodic eliminants.

\begin{definition}[aperiodic locator chart]
\label{def:aperiodic-locator-chart}
An aperiodic locator chart at exact agreement $A$ is a pair
\[
        X=(E_X,\Delta_X),
\]
where $E_X$ is a finite list of polynomial equations in the locator coefficients
$\ell$, and $\Delta_X(\ell)$ is a product of declared nonzero tests.  The chart
represents the constructible set
\[
        V(I_{\rm split,j}+\langle E_X\rangle)\setminus V(\Delta_X).
\]
The equations and inequations must include, or reference, the removal of already
paid tangent/common-code-line and quotient-image branches.  A family of charts
$\mathfrak X_A$ is a valid aperiodic chart cover if the union of their
constructible sets contains every remaining split locator of exact co-support
size $j_A$ not assigned to those earlier ledgers.
\end{definition}

For a coordinate $h\in\{0,\ldots,t-1\}$ define the affine pivot consistency ideal
\[
\begin{aligned}
        J^{\rm aff}_{X,h}:={}&
        I_{\rm split,j}+\langle E_X\rangle
        +\langle A_m(\ell)B_h(\ell)-A_h(\ell)B_m(\ell):0\le m<t\rangle \\
        &+\langle ZB_h(\ell)+A_h(\ell)\rangle
        \subseteq F[\ell_0,\ldots,\ell_{j-1},Z],
\end{aligned}
\]
and saturate it by the pivot denominator:
\[
        \widehat J^{\rm aff}_{X,h}:=
        J^{\rm aff}_{X,h}:(\Delta_X B_h)^\infty .
\]
The equations $A_mB_h-A_hB_m=0$ force $A$ and $B$ to be collinear on the pivot
chart, and $ZB_h+A_h=0$ records the unique finite slope.

\begin{theorem}[pivoted affine-line eliminant certificate]
\label{thm:aperiodic-affine-eliminant}
Let $X$ be an aperiodic locator chart and let $h$ be a noncontainment pivot.  If
there is a nonzero polynomial
\[
        Q_{X,h}(Z)\in \widehat J^{\rm aff}_{X,h}\cap F[Z]
\]
of degree $D_{X,h}$, then the number of finite support-wise MCA-bad slopes in
that chart with $B_h\ne0$ is at most $\min\{|F|,D_{X,h}\}$.  The same bound holds
for no-loss CA-bad slopes in that pivot chart.
\end{theorem}

\begin{proof}
A bad slope in the chart has a split locator with $\Delta_X\ne0$ and, on this
pivot chart, $B_h\ne0$.  The Hankel incidence equation gives
$A+ZB=0$, hence the point $(\ell,Z)$ satisfies all generators of
$J^{\rm aff}_{X,h}$ and survives the saturation.  Therefore $Q_{X,h}(Z)=0$.
A nonzero degree-$D_{X,h}$ polynomial over $F$ has at most $D_{X,h}$ roots.
\end{proof}

\begin{remark}[why the pivot split is useful]
\label{rem:pivot-split-useful}
The saturation by the whole vector $B$ in the unsplit incidence ideal is replaced
by the explicit open cover $B_h\ne0$.  A scanner can try all $h$, produce small
univariate eliminants on successful pivots, and report the residual locus
$B_0=\cdots=B_{t-1}=0$ as contained for support-wise MCA.  For no-loss CA, that
residual locus also contributes no bad slope at the same support, because if
$A=B=0$ then both $f$ and $g$ are explained on the witness support.
\end{remark}

\subsubsection{Projective-slope split}

For projective parameters $[Z_0:Z_1]$ the incidence equation is
\[
        Z_0A(\ell)+Z_1B(\ell)=0 .
\]
The finite patch $Z_0\ne0$ is exactly the affine pivot atlas above with
$Z=Z_1/Z_0$.  The extra point at infinity is $[0:1]$.

\begin{definition}[projective infinity chart]
\label{def:projective-infinity-chart}
For a chart $X$, the projective-infinity chart is the constructible locus
\[
        I_{\rm split,j}+\langle E_X\rangle+
        \langle B_0,\ldots,B_{t-1}\rangle,
        \qquad
        \Delta_X\ne0,
        \quad A\ne0 .
\]
If it is nonempty, it contributes the single projective parameter $[0:1]$.
If it is empty, the infinity point contributes nothing on $X$.
\end{definition}

\begin{corollary}[projective eliminant certificate]
\label{cor:projective-eliminant-certificate}
For a chart cover $\mathfrak X_A$, the projective-slope numerator is bounded by
the union of the finite-patch root sets from
\cref{thm:aperiodic-affine-eliminant} together with at most one extra point
$[0:1]$, provided the infinity charts are certified either empty or nonempty.
Equivalently, a homogeneous nonzero eliminant
\[
        Q_X(Z_0,Z_1)
\]
of degree $D_X$ bounds the number of projective bad parameters in that chart by
$D_X$.
\end{corollary}

\begin{proof}
The finite patch is the affine theorem.  The complement of that patch is the
single point $[0:1]$, which is governed by \cref{def:projective-infinity-chart}.
For the homogeneous formulation, every bad projective parameter is a zero of the
homogeneous eliminant, and a nonzero homogeneous polynomial of degree $D_X$ on
$\mathbb P^1(F)$ has at most $D_X$ zeros.
\end{proof}

\subsubsection{Finite-parameter curve charts}

For a degree-$d$ finite power curve
\[
        W_\Gamma=f_0+\Gamma f_1+\cdots+\Gamma^d f_d,
\]
put
\[
        V_i(\ell):=H_{t,j}(\Syn(f_i))\widetilde\ell\in F^t,
        \qquad 0\le i\le d .
\]
The curve incidence equation is
\[
        V_0(\ell)+\Gamma V_1(\ell)+\cdots+\Gamma^d V_d(\ell)=0 .
\]
Curve-MCA noncontainment means that not all coefficient vectors $V_i(\ell)$ are
zero.  Thus the noncontained branch is covered by coefficient pivots
$(i,h)$ with $(V_i)_h\ne0$.

For a chart $X$ and a coefficient pivot $(i,h)$, define
\[
\begin{aligned}
        J^{\rm curve,d}_{X,i,h}:={}&
        I_{\rm split,j}+\langle E_X\rangle \\
        &+\left\langle
          \sum_{r=0}^d \Gamma^r (V_r(\ell))_m:
          0\le m<t
          \right\rangle
        \subseteq F[\ell_0,\ldots,\ell_{j-1},\Gamma],
\end{aligned}
\]
and set
\[
        \widehat J^{\rm curve,d}_{X,i,h}
        :=J^{\rm curve,d}_{X,i,h}:(\Delta_X (V_i)_h)^\infty .
\]

\begin{theorem}[pivoted curve eliminant certificate]
\label{thm:aperiodic-curve-eliminant}
If a chart $X$ and coefficient pivot $(i,h)$ have a nonzero eliminant
\[
        Q_{X,i,h}(\Gamma)\in
        \widehat J^{\rm curve,d}_{X,i,h}\cap F[\Gamma]
\]
of degree $D_{X,i,h}$, then the number of finite degree-$d$ curve-MCA bad
parameters in that chart and pivot is at most $\min\{|F|,D_{X,i,h}\}$.  The same
bound applies to no-loss curve-CA parameters in the same pivot chart.
\end{theorem}

\begin{proof}
Every bad curve parameter on the pivot chart gives a split locator satisfying all
curve-incidence equations and the nonzero pivot condition $(V_i)_h\ne0$.  It
therefore survives the saturation and is a root of the supplied eliminant.
\end{proof}

% -----------------------------------------------------------------------------
\paragraph{Residual rank-pivot charts.}
The gcd ledger above removes the regular overdetermined branch.  What remains is
not a formless exceptional set: every remaining bad parameter is still a point of
a finite collection of determinantal rank charts.  The point of the next
statements is to make this residual bucket verifier-facing.  A certificate may
now say exactly which rank chart it uses, which noncontainment coordinate is
not settled here, and which univariate or homogeneous eliminant cuts out the parameter
projection.  If no such eliminant is supplied, the chart is a named residual
obstruction rather than an unaccounted term.

Fix an exact agreement value $A\ge k$ and keep
\[
        j=j_A:=n-A,
        \qquad
        t=t_A:=A-k,
        \qquad
        r=n-k .
\]
For a projective received line $Z_0 f+Z_1 g$, put
\[
        U:=H_{t,j}(\Syn(f)),
        \qquad
        V:=H_{t,j}(\Syn(g)),
        \qquad
        M^{\rm proj}_A(Z_0,Z_1):=Z_0U+Z_1V .
\]
For a locator vector $\widetilde\ell$ write
\[
        A(\ell):=U\widetilde\ell,
        \qquad
        B(\ell):=V\widetilde\ell .
\]
Thus the projective support equation is
\[
        Z_0A(\ell)+Z_1B(\ell)=0,
\]
and projective noncontainment is $(A(\ell),B(\ell))\ne(0,0)$.

Let $X=(E_X,\Delta_X)$ be a locator chart in the sense of
\cref{def:aperiodic-locator-chart}.  For row and column sets
$I\subseteq\{0,\ldots,t-1\}$ and
$J\subseteq\{0,\ldots,j\}$ with $|I|=|J|=s$, write
\[
        \Delta_{I,J}(Z_0,Z_1):=
        \det M^{\rm proj}_A(Z_0,Z_1)_{I,J},
\]
with the convention that the unique $0\times0$ minor is $1$.  Let
$\mathcal M_{s+1}(M^{\rm proj}_A)$ denote the set of all $(s+1)\times(s+1)$
minors, empty when $s=\min\{t,j+1\}$.

For a noncontainment pivot
\[
        \pi\in\{A_h:0\le h<t\}\cup\{B_h:0\le h<t\},
\]
let $\Pi_\pi(\ell)$ denote the corresponding scalar coordinate.  Define the
saturated projective rank-pivot ideal
\[
\begin{aligned}
        \widehat J^{\rm proj}_{X,I,J,\pi}:=
        \bigl(&I_{{\rm split},j}+\langle E_X\rangle
        +\langle M^{\rm proj}_A(Z_0,Z_1)\widetilde\ell\rangle\\
        &+\langle \Delta: \Delta\in\mathcal M_{s+1}(M^{\rm proj}_A)\rangle
        \bigr):
        \bigl(\Delta_X\Delta_{I,J}\Pi_\pi\bigr)^\infty
        \subseteq F[\ell_0,\ldots,\ell_{j-1},Z_0,Z_1].
\end{aligned}
\]
The open condition $\Delta_{I,J}\ne0$ and the vanishing of all larger minors put
$M^{\rm proj}_A$ on the exact-rank-$s$ stratum on this chart.  The saturation by
$\Pi_\pi$ chooses one of the finitely many open noncontainment coordinates.

\begin{theorem}[residual rank-pivot cover]
\label{thm:residual-rank-pivot-cover}
Fix an exact agreement value $A$ and a chart cover $\mathfrak X_A$ of the
remaining locator branch after tangent/common-line and quotient-image removals.
Every projective support-wise MCA-bad parameter with exact witness size $A$ in
that remaining branch lies on one of the rank-pivot charts
\[
        \widehat J^{\rm proj}_{X,I,J,\pi}
\]
for some $X\in\mathfrak X_A$, some rank $s$, some nonzero $s\times s$ minor
$\Delta_{I,J}$ at the parameter, and some nonzero noncontainment coordinate
$\Pi_\pi$.  The same cover contains every no-loss projective CA-bad parameter in
that branch.
\end{theorem}

\begin{proof}
Let $[z_0:z_1]$ be bad with witness support $S$, co-support $T$, and locator
vector $\widetilde\ell_T$.  Since the chart family covers the remaining branch,
$\widetilde\ell_T$ lies in some $X\in\mathfrak X_A$ with $\Delta_X\ne0$.  The
support equation gives
$M^{\rm proj}_A(z_0,z_1)\widetilde\ell_T=0$.  Let
$s=\operatorname{rank}M^{\rm proj}_A(z_0,z_1)$.  Choose a nonzero
$s\times s$ minor at this parameter; then every $(s+1)\times(s+1)$ minor
vanishes.  Finally, support-wise noncontainment says that
$(A(\ell_T),B(\ell_T))\ne(0,0)$, so some coordinate pivot $\Pi_\pi$ is nonzero.
Thus the point survives the displayed saturation and belongs to the indicated
rank-pivot chart.  The no-loss CA statement uses the same explaining support and
same noncontainment conclusion.
\end{proof}

\begin{theorem}[projective residual eliminant certificate]
\label{thm:projective-residual-eliminant-certificate}
If a rank-pivot chart has a nonzero homogeneous eliminant
\[
        Q_{X,I,J,\pi}(Z_0,Z_1)
        \in
        \widehat J^{\rm proj}_{X,I,J,\pi}\cap F[Z_0,Z_1]
\]
of degree $D_{X,I,J,\pi}$, then the number of projective line parameters in that
chart is at most $D_{X,I,J,\pi}$.  More exactly, over a finite field $F$ of size
$q$, the number of finite-patch roots is
\[
        \deg\gcd\bigl(Q_{X,I,J,\pi}(1,Z),Z^q-Z\bigr),
\]
and the point at infinity contributes precisely when
$Q_{X,I,J,\pi}(0,1)=0$.  Thus the exact finite-field count is
\[
        \deg\gcd\bigl(Q_{X,I,J,\pi}(1,Z),Z^q-Z\bigr)
        +\mathbf 1_{Q_{X,I,J,\pi}(0,1)=0}.
\]
\end{theorem}

\begin{proof}
Every parameter represented by the chart gives a point of the saturated
incidence variety, hence is a zero of every element of the elimination ideal,
including $Q_{X,I,J,\pi}$.  A nonzero homogeneous polynomial of degree $D$ on
$\mathbb P^1(F)$ has at most $D$ projective zeros.  On the affine patch
$Z_0=1$ the finite parameters are exactly the roots of $Q(1,Z)$ in $F$, which
are counted by the gcd with $Z^q-Z$ after squarefree reduction.  The complement
of the finite patch is the single point $[0:1]$, giving the displayed indicator.
\end{proof}

\begin{corollary}[finite and infinity patches]
\label{cor:finite-infinity-rank-pivot-patches}
The finite affine rank-pivot certificate is obtained from
\cref{thm:projective-residual-eliminant-certificate} by setting $Z_0=1$ and
using only noncontainment pivots $B_h\ne0$.  The infinity chart is obtained by
setting $Z_0=0$, $Z_1=1$, imposing
\[
        B(\ell)=0,
        \qquad
        A_h(\ell)\ne0
\]
for some $h$, and applying the same saturated rank-pivot test to the constant
matrix $V$.  Consequently the projective sampler has no hidden extra case: every
finite or infinite bad point is represented by one of these two patch types.
\end{corollary}

\begin{proof}
On the finite patch $[1:Z]$, the support equation is
$A(\ell)+ZB(\ell)=0$.  If $B(\ell)=0$, then also $A(\ell)=0$, so the fixed
support is contained; hence every finite noncontained point has some
$B_h\ne0$.  At infinity $[0:1]$, the support equation is $B(\ell)=0$ and
noncontainment is exactly $A(\ell)\ne0$.  These alternatives exhaust
$\mathbb P^1(F)$.
\end{proof}

\paragraph{Curve coefficient-pivot rank charts.}
For a finite-parameter degree-$d$ power curve
\[
        W_\Gamma=f_0+\Gamma f_1+\cdots+\Gamma^d f_d,
\]
put $U_i:=H_{t,j}(\Syn(f_i))$ and
\[
        M_A^{(d)}(\Gamma):=
        \sum_{i=0}^{d}\Gamma^i U_i,
        \qquad
        V_i(\ell):=U_i\widetilde\ell .
\]
For row and column sets $I,J$ of size $s$, write
$\Delta_{I,J}^{(d)}(\Gamma)$ for the corresponding $s\times s$ minor of
$M_A^{(d)}(\Gamma)$, and let $\mathcal M_{s+1}(M_A^{(d)})$ be the set of all
larger minors.  For a coefficient noncontainment pivot $(i,h)$ define
\[
\begin{aligned}
        \widehat J^{\rm curve,d}_{X,I,J,i,h}:=
        \bigl(&I_{{\rm split},j}+\langle E_X\rangle
        +\langle M_A^{(d)}(\Gamma)\widetilde\ell\rangle\\
        &+\langle \Delta: \Delta\in\mathcal M_{s+1}(M_A^{(d)})\rangle
        \bigr):
        \bigl(\Delta_X\Delta_{I,J}^{(d)}(V_i)_h\bigr)^\infty .
\end{aligned}
\]

\begin{theorem}[curve residual rank-pivot eliminant]
\label{thm:curve-residual-rank-pivot-eliminant}
The curve-bad parameters in a residual exact-agreement branch are covered by the
charts $\widehat J^{\rm curve,d}_{X,I,J,i,h}$.  If such a chart has a nonzero
eliminant
\[
        Q_{X,I,J,i,h}(\Gamma)
        \in \widehat J^{\rm curve,d}_{X,I,J,i,h}\cap F[\Gamma]
\]
of degree $D_{X,I,J,i,h}$, then that chart contributes at most
\[
        \deg\gcd\bigl(Q_{X,I,J,i,h}(\Gamma),\Gamma^q-\Gamma\bigr)
        \le D_{X,I,J,i,h}
\]
finite parameters over $F$.  The same bound applies to no-loss curve-CA
parameters in the same chart.
\end{theorem}

\begin{proof}
A curve-bad parameter with locator $\widetilde\ell_T$ satisfies
$M_A^{(d)}(\gamma)\widetilde\ell_T=0$.  Taking $s$ to be the rank of
$M_A^{(d)}(\gamma)$ gives a nonzero $s\times s$ minor and vanishing of all larger
minors.  Noncontainment of the coefficient tuple means that some coordinate
$(V_i(\ell_T))_h$ is nonzero.  Thus the point lies on one of the saturated
rank-pivot charts and any univariate eliminant for that chart must vanish at
$\gamma$.  The finite-field root count is the gcd with $\Gamma^q-\Gamma$.
\end{proof}

\paragraph{Extension-line Weil restriction.}
The previous charts are intrinsic over the line field.  For extension-field
samplers one also needs a generated-field ledger: a low-dimensional set over the
base field $\B$ can be negligible after division by $|F|$, even when it is large
as a subset of $\B$.  The following is the certificate form of that accounting.

Let $F/\B$ have degree $m$, choose a $\B$-basis
$\omega_0,\ldots,\omega_{m-1}$ of $F$, and write a finite affine line parameter
as
\[
        Z=Z_0\omega_0+\cdots+Z_{m-1}\omega_{m-1},
        \qquad Z_i\in\B .
\]
Assume $D\subseteq\B$ and write every locator coefficient in the same base field
coordinates when necessary.  Replacing each $F$-valued equation in a rank-pivot
chart by its $m$ coordinate equations over $\B$ gives a Weil-restricted saturated
ideal
\[
        \widehat J^{\rm WR}_{Y}
        \subseteq
        \B[\ell\text{-coordinates},Z_0,\ldots,Z_{m-1}]
\]
for a chart $Y$; noncontainment is split by one nonzero $\B$-coordinate of the
chosen $F$-valued pivot and included in the saturation.
Let $\pi_Z$ denote projection to the parameter-coordinate space
$\mathbb A^m_{\B}$.

\begin{theorem}[extension-line dimension--degree ledger]
\label{thm:extension-line-dimension-degree-ledger}
Suppose that for a Weil-restricted residual chart $Y$ a certificate gives a
closed set
\[
        Y_Z\supseteq \pi_Z\bigl(V(\widehat J^{\rm WR}_{Y})\bigr)
        \subseteq \mathbb A^m_{\B}
\]
of degree at most $\Delta_Y$ and dimension at most $e_Y$.  Then the number of
extension-field affine parameters contributed by this chart is at most
\[
        \Delta_Y |\B|^{e_Y}.
\]
If $Y_Z$ is zero-dimensional and radical, this bound may be replaced by the
exact number of its $\B$-rational points.  After normalization by the line field,
this chart contributes at most
\[
        \frac{\Delta_Y |\B|^{e_Y}}{|F|}
        =\Delta_Y |\B|^{e_Y-m}
\]
to the MCA or CA error numerator.
\end{theorem}

\begin{proof}
The chosen basis identifies $F$ with $\B^m$, so affine line parameters over $F$
are exactly the $\B$-rational points of $\mathbb A^m_{\B}$.  Every bad parameter
represented by the chart lies in the projection of the Weil-restricted incidence
variety and hence in $Y_Z(\B)$.  The standard affine degree point-count bound gives
$|Y_Z(\B)|\le\Delta_Y|\B|^{e_Y}$.  The normalized statement divides by
$|F|=|\B|^m$.
\end{proof}

\begin{corollary}[base-rational line inertness as a chart certificate]
\label{cor:base-rational-line-inertness-chart}
If $f$ and $g$ are $\B$-valued, then every finite affine bad slope on every
rank-pivot chart lies in $\B$.  Equivalently, in the above coordinates all
extension-valued charts with $(Z_1,\ldots,Z_{m-1})\ne0$ are empty.
\end{corollary}

\begin{proof}
For a fixed support and locator over $\B$, the vectors $A(\ell)$ and $B(\ell)$
are $\B$-valued.  On a finite noncontained chart some coordinate $B_h$ is
nonzero, and the slope is forced by
$Z=-A_h/B_h\in\B$.  Hence no genuinely extension-valued finite parameter can
survive the chart equations.
\end{proof}

\paragraph{Residual aperiodic certificate theorem.}
We can now state the safe-side theorem in the form a scanner should print.  Let
$P_{\rm paid,F}$ be the exact squarefree finite-root polynomial already charged
to tangent/common-line and quotient-image branches, and let
$G_{{\rm reg}\setminus{\rm paid},F}^{\rm aff}$ be the regular residual polynomial
of \cref{def:paid-root-removal-regular-branch}.  For the remaining exact
agreement levels, let $\mathfrak R_A$ be a finite list of rank-pivot residual
charts, each equipped with either a finite/projective/curve eliminant or a
Weil-restricted extension projection certificate.

\begin{theorem}[scanner-checkable residual aperiodic ledger]
\label{thm:scanner-checkable-residual-aperiodic-ledger}
Assume that the chart families $\mathfrak R_A$, for $A\ge a$, cover every
remaining locator after tangent/common-line and quotient-image removals.  Suppose
each residual chart $Y$ is certified by one of the following data:
\begin{enumerate}[label=\textup{(\roman*)},leftmargin=2.2em]
\item an affine eliminant $Q_Y(Z)$, contributing
$\deg\gcd(Q_Y,Z^q-Z)$;
\item a homogeneous projective eliminant $Q_Y(Z_0,Z_1)$, contributing
$\deg\gcd(Q_Y(1,Z),Z^q-Z)+\mathbf 1_{Q_Y(0,1)=0}$;
\item a finite curve eliminant $Q_Y(\Gamma)$, contributing
$\deg\gcd(Q_Y,\Gamma^q-\Gamma)$;
\item an extension Weil-restricted projection certificate $(\Delta_Y,e_Y)$,
contributing $\Delta_Y|\B|^{e_Y}$ parameters.
\end{enumerate}
Then the finite affine aperiodic numerator after quotient/tangent removals is
bounded by
\[
        \deg G_{{\rm reg}\setminus{\rm paid},F}^{\rm aff}
        +\sum_{A\ge a}\sum_{Y\in\mathfrak R_A} B_Y,
\]
where $B_Y$ is the corresponding contribution above.  The projective and curve
variants are obtained by using the projective or curve contributions for all
charts.  If exact root sets or exact zero-dimensional projections are printed,
the displayed safe sum may be replaced by their union after duplicate
coalescing.
\end{theorem}

\begin{proof}
The regular nonsingular overdetermined branch is covered by
\cref{cor:quotient-tangent-removed-regular-numerator}.  Every remaining bad
parameter is covered by the rank-pivot cover theorem, or by its curve analogue,
after applying the declared chart cover.  Each chart contribution is an upper
bound by \cref{thm:projective-residual-eliminant-certificate,thm:curve-residual-rank-pivot-eliminant,thm:extension-line-dimension-degree-ledger}.  Summing the
chart bounds is a safe union bound; replacing the sum by an exact union is valid
when the certificate supplies the corresponding finite root or projection sets.
\end{proof}

\begin{definition}[singular residual certificate]
\label{def:singular-residual-certificate}
A singular or underdetermined residual aperiodic certificate consists of:
\begin{enumerate}[label=\textup{(\roman*)},leftmargin=2.2em]
\item the paid tangent/common-line and quotient-image root ledgers already
removed;
\item the regular rank-drop gcd/lcm ledger and its paid-root removal;
\item for each remaining exact agreement $A$, a chart cover of the residual
locator set after the paid branches are removed;
\item for every chart, its rank $s$, pivot minor $\Delta_{I,J}$, noncontainment
pivot, and either an affine, projective, or curve eliminant, or a
Weil-restricted projection degree/dimension certificate;
\item a residual table for every chart on which no eliminant or projection bound
is supplied, labelled as quotient, tangent/common-line, base-field exceptional,
genuinely extension-valued, or new singular obstruction.
\end{enumerate}
The verifier checks chart coverage, rank-stratum equations, saturations by the
pivot and noncontainment coordinates, eliminant membership, finite-field root
counts, extension projection bounds, and duplicate coalescing when claimed.
\end{definition}

\begin{remark}[status after the singular-bucket charts]
\label{rem:singular-bucket-chart-status}
This does not prove that the hard aperiodic residual buckets are small.  It does remove
one ambiguity: underdetermined and rank-drop singular buckets now have a precise
finite certificate language, including finite slopes, projective infinity,
finite-parameter curves, and genuinely extension-valued parameters.  The next
mathematical step is to construct low-degree eliminants, or prove structural
classification, for these rank-pivot charts in the near-capacity regimes where
$t_A=A-k$ is small.
\end{remark}

% === Combined Insert A.2: kernel-compressed residual eliminants and exact residual images ===

\paragraph{Kernel-compressed rank-pivot eliminants.}
The rank-pivot charts above are already finite certificate objects, but the
variables are still the full locator coordinates.  The next reduction removes the
linear kernel equations on a pivot chart by Cramer's rule.  It is the algebraic
form a scanner should use before attempting Gr\"obner elimination or resultants:
all remaining equations are split-locator equations, chart-removal equations,
rank-stratum equations, and noncontainment pivots in the free kernel coordinates
and the line parameter.

We write the statement first on the finite affine patch.  Fix an exact agreement
value $A$, put $j=n-A$ and $t=A-k$, and set
\[
        M_A(Z):=U+ZV,
        \qquad
        U:=H_{t,j}(\Syn(f)),\quad V:=H_{t,j}(\Syn(g)).
\]
Let $I\subseteq\{0,\ldots,t-1\}$ and $J\subseteq\{0,\ldots,j\}$ have the same
cardinality $s$, and put
\[
        \Delta_{I,J}(Z):=\det M_A(Z)_{I,J}.
\]
Let $C_J:=\{0,\ldots,j\}\setminus J$ and introduce free kernel coordinates
$Y=(Y_c)_{c\in C_J}$.  On the open set $\Delta_{I,J}\ne0$, define the polynomial
vector
\[
        \Lambda^{\#}_{I,J}(Z,Y)\in F[Z,Y]^{j+1}
\]
by
\[
        \bigl(\Lambda^{\#}_{I,J}\bigr)_{C_J}
        :=\Delta_{I,J}Y,
        \qquad
        \bigl(\Lambda^{\#}_{I,J}\bigr)_{J}
        :=-\operatorname{adj}\bigl(M_A(Z)_{I,J}\bigr)
             M_A(Z)_{I,C_J}Y .
\]
Thus
\[
        M_A(Z)_{I,*}\Lambda^{\#}_{I,J}(Z,Y)=0.
\]
If, in addition, all $(s+1)\times(s+1)$ minors of $M_A(Z)$ vanish, then every row
of $M_A(Z)$ lies in the span of the pivot rows, and hence
\[
        M_A(Z)\Lambda^{\#}_{I,J}(Z,Y)=0 .
\]

To encode splitting without dividing by the leading coefficient, use homogeneous
locator coordinates.  For
$\lambda=(\lambda_0,\ldots,\lambda_j)$ set
\[
        L_\lambda(X):=\sum_{b=0}^{j}\lambda_bX^b .
\]
Let $Q(X)=\sum_{r=0}^{n-j}q_rX^r$ be an auxiliary quotient polynomial and let
$I_{{\rm hsplit},j}$ be the homogeneous split incidence
\[
        P_D(X)=L_\lambda(X)Q(X),
        \qquad \lambda_j\ne0,
\]
viewed as the coefficient equations saturated by $\lambda_j$.  Equivalently,
its points with $\lambda_j\ne0$ are precisely the nonzero scalar multiples of the
monic locators $\prod_{x\in T}(X-x)$ with $T\subseteq D$ and $|T|=j$.
For a monic locator chart $X=(E_X,\Delta_X)$, write $E_X^{\rm hom}$ and
$\Delta_X^{\rm hom}$ for the usual homogenizations in the coordinates
$\lambda_0,\ldots,\lambda_j$.

For a finite noncontainment pivot $B_h\ne0$, set
\[
        \Pi_h^{\#}(Z,Y):=\bigl(V\Lambda^{\#}_{I,J}(Z,Y)\bigr)_h .
\]
The \emph{kernel-compressed affine rank-pivot ideal} is
\[
\begin{aligned}
        \mathcal K^{\rm aff}_{X,I,J,h}:=
        \bigl(&\langle \operatorname{coeff}_X(P_D-L_{\Lambda^{\#}_{I,J}}Q)\rangle
        +\langle E_X^{\rm hom}(\Lambda^{\#}_{I,J})\rangle \\
        &+\langle \Delta: \Delta\in\mathcal M_{s+1}(M_A)\rangle\bigr):
        \bigl(\Delta_{I,J}\,\Delta_X^{\rm hom}(\Lambda^{\#}_{I,J})\,
              (\Lambda^{\#}_{I,J})_j\,\Pi_h^{\#}\bigr)^\infty
\end{aligned}
\]
in the ring $F[Z,Y,(q_r)]$.  Here
$L_{\Lambda^{\#}_{I,J}}$ means $L_\lambda$ after substituting
$\lambda=\Lambda^{\#}_{I,J}(Z,Y)$.

\begin{theorem}[kernel-compressed affine eliminant certificate]
\label{thm:kernel-compressed-affine-eliminant}
Every finite affine bad slope represented by the rank-pivot chart
$(X,I,J,h)$ is in the projection of
$V(\mathcal K^{\rm aff}_{X,I,J,h})$ to the $Z$-line.  Consequently, if
\[
        Q_{X,I,J,h}(Z)\in
        \mathcal K^{\rm aff}_{X,I,J,h}\cap F[Z]
\]
is nonzero, then the chart contributes at most
\[
        \deg\gcd(Q_{X,I,J,h}(Z),Z^q-Z)
\]
finite slopes over $F_q$.
\end{theorem}

\begin{proof}
Let $z$ be a bad slope in the chart, and let $\lambda$ be a nonzero scalar
multiple of the corresponding locator vector.  Since the pivot minor is nonzero,
put $Y_c=\lambda_c/\Delta_{I,J}(z)$ for $c\in C_J$.  Cramer's rule and the pivot
row equations give
\[
        \Lambda^{\#}_{I,J}(z,Y)=\lambda .
\]
The locator is split, lies in the chart $X$, has nonzero leading coefficient,
and has the chosen nonzero noncontainment coordinate.  The larger minors vanish
on the exact-rank stratum.  Thus the point survives the displayed saturation and
lies in the compressed incidence.  Every element of the elimination ideal must
vanish on its $Z$-coordinate, and the finite-field count is obtained by
intersecting with $Z^q-Z$.
\end{proof}

\begin{remark}[why the compression is useful]
The original chart has $j$ locator variables, one parameter variable, and the
split equations.  The compressed chart has only $j+1-s$ free kernel coordinates
before imposing the leading-coordinate and chart opens.  In the overdetermined
rank-drop case this is often a large saving.  In the genuinely underdetermined
near-capacity case the saving may be small, but the construction still separates
three causes of difficulty: many split locators, a dominant parameter projection,
or a high-degree finite eliminant.
\end{remark}

The projective and curve variants are identical after replacing the parameter
ring.  For projective lines, replace $M_A(Z)$ by
$M_A^{\rm proj}(Z_0,Z_1)$, use homogeneous pivot minors, and take
homogeneous eliminants in $F[Z_0,Z_1]$.  On the infinity patch $[0:1]$, this is
just the constant matrix $V$ with equations $B(\ell)=0$ and an open pivot
$A_h\ne0$.  For a degree-$d$ finite power curve, replace $M_A(Z)$ by
\[
        M_A^{(d)}(\Gamma)=\sum_{i=0}^{d}\Gamma^i
        H_{t,j}(\Syn(f_i))
\]
and use coefficient pivots $(V_i)_h\ne0$.

\begin{corollary}[kernel-compressed projective and curve eliminants]
\label{cor:kernel-compressed-projective-curve-eliminants}
The projective rank-pivot chart obtained from the compressed construction has the
same parameter projection as the uncompressed rank-pivot chart on the declared
opens.  A nonzero homogeneous eliminant $Q(Z_0,Z_1)$ contributes at most
\[
        \deg\gcd(Q(1,Z),Z^q-Z)+\mathbf 1_{Q(0,1)=0}
\]
projective parameters over $F_q$.  For a finite degree-$d$ curve chart, a nonzero
eliminant $Q(\Gamma)$ contributes at most
\[
        \deg\gcd(Q(\Gamma),\Gamma^q-\Gamma)
\]
parameters.
\end{corollary}

\begin{proof}
The same Cramer parametrization applies to the projective pencil and to the
curve pencil.  The root-count statements are the finite-field root counts for a
homogeneous polynomial on $\mathbb P^1(F_q)$ and a univariate polynomial on
$\mathbb A^1(F_q)$.
\end{proof}

\paragraph{Exact residual support-image lcm ledger.}
The compressed eliminant is a constructive way to find a small parameter
polynomial.  There is also an exact, definition-level residual image ledger that
works for any finite co-support family, quotient or not.  This is the clean
object against which scanner-produced eliminants should be audited.

Let $\mathcal T$ be any family of co-supports $T\subseteq D$ of size $j=n-A$.
For a line $f+Zg$, put
\[
        A_T:=H_{t,j}(\Syn(f))\boldsymbol\ell_T,
        \qquad
        B_T:=H_{t,j}(\Syn(g))\boldsymbol\ell_T .
\]
Define the finite affine fixed-support polynomial
\[
        g_T^{\rm aff}(Z):=
        \begin{cases}
        \operatorname{monic}\gcd\{(A_T)_h+Z(B_T)_h:0\le h<t\},
             & (A_T,B_T)\ne(0,0),\\
        1,   & (A_T,B_T)=(0,0),
        \end{cases}
\]
with the convention that the gcd of constants with no common root is $1$.  Thus
$g_T^{\rm aff}$ has degree either $0$ or $1$.  The contained case
$(A_T,B_T)=(0,0)$ is assigned the polynomial $1$, because it contributes no
support-wise bad slope.
Set
\[
        R_{\mathcal T}^{\rm aff}(Z):=
        \operatorname{rad}\operatorname{lcm}_{T\in\mathcal T}g_T^{\rm aff}(Z).
\]
For projective lines, define $g_T^{\rm proj}(Z_0,Z_1)$ as the monic homogeneous
gcd of the linear forms
\[
        Z_0(A_T)_h+Z_1(B_T)_h,
        \qquad 0\le h<t,
\]
again setting it equal to $1$ in the contained case, and put
\[
        R_{\mathcal T}^{\rm proj}:=
        \operatorname{rad}\operatorname{lcm}_{T\in\mathcal T}g_T^{\rm proj}.
\]
For a finite degree-$d$ curve, use the vector polynomials
\[
        \sum_{i=0}^{d}\Gamma^i(V_{i,T})_h
\]
and define $R_{\mathcal T}^{(d)}(\Gamma)$ by the same radical lcm rule, with the
all-zero coefficient tuple treated as contained.

\begin{theorem}[exact arbitrary-residual image ledger]
\label{thm:exact-arbitrary-residual-image-ledger}
Let $\mathcal T$ be any finite co-support family.  The exact finite affine line
parameters witnessed by $\mathcal T$ are the roots in $F$ of
$R_{\mathcal T}^{\rm aff}$.  Hence their number is
\[
        \deg\gcd(R_{\mathcal T}^{\rm aff}(Z),Z^q-Z).
\]
The exact projective parameters witnessed by $\mathcal T$ are the projective
zeros of $R_{\mathcal T}^{\rm proj}$, counted over $F_q$ by
\[
        \deg\gcd(R_{\mathcal T}^{\rm proj}(1,Z),Z^q-Z)
        +\mathbf 1_{R_{\mathcal T}^{\rm proj}(0,1)=0}.
\]
For finite degree-$d$ curves, the exact finite parameter count is
\[
        \deg\gcd(R_{\mathcal T}^{(d)}(\Gamma),\Gamma^q-\Gamma).
\]
The same ledgers upper-bound no-loss CA parameters witnessed by the same
co-support family.
\end{theorem}

\begin{proof}
For a fixed co-support $T$, the syndrome-locator recurrence says that the line
point is explained on $D\setminus T$ precisely when
$A_T+ZB_T=0$.  If $A_T=B_T=0$, then both coefficient words are explained on the
same support and the fixed support is contained, so it is correctly charged as
$1$.  Otherwise the common roots of the scalar coordinates are exactly the
finite bad slopes contributed by $T$, and this common-root set is represented by
$g_T^{\rm aff}$.  Taking the radical lcm over $T\in\mathcal T$ coalesces duplicate
slopes.  The projective and curve statements are identical, using the homogeneous
or degree-$d$ coordinate polynomials.  No-loss CA badness is stronger than
existence of an explaining support with noncontainment, so the same image set is
an upper ledger.
\end{proof}


\begin{corollary}[closed-ball residual exact image after paid-branch removal]
\label{cor:residual-exact-image-after-paid-removal}
Fix a closed-ball agreement threshold $a$.  For each exact agreement
$A\in\{a,a+1,\ldots,n\}$, let $\mathcal T_{{\rm all},A}$ be the family of
co-supports of size $n-A$, let $\mathcal T_{{\rm paid},A}\subseteq
\mathcal T_{{\rm all},A}$ be a certified paid support family, and put
\[
        \mathcal T_{{\rm res},A}
        :=\mathcal T_{{\rm all},A}\setminus\mathcal T_{{\rm paid},A}.
\]
Let $R_{{\rm res},A}^{\rm aff}$ be the polynomial of
\cref{thm:exact-arbitrary-residual-image-ledger} for the fixed-size family
$\mathcal T_{{\rm res},A}$.  Define the closed-ball residual polynomial
\[
        R_{{\rm res},\ge a}^{\rm aff}(Z)
        :=\operatorname{rad}\operatorname{lcm}_{A=a}^{n}
          R_{{\rm res},A}^{\rm aff}(Z),
        \qquad
        R_{{\rm res},\ge a,F}^{\rm aff}(Z)
        :=\gcd\bigl(R_{{\rm res},\ge a}^{\rm aff}(Z),Z^q-Z\bigr).
\]
Then the exact finite affine residual numerator after paid-support removal is
\[
        \deg R_{{\rm res},\ge a,F}^{\rm aff}.
\]
The projective and finite-curve residual numerators are obtained by the same
exact-agreement radical-lcm construction, using
$R_{{\rm res},A}^{\rm proj}$ and $R_{{\rm res},A}^{(d)}$ respectively, followed by
the finite-field root count of \cref{thm:exact-arbitrary-residual-image-ledger}.
If a scanner also prints an exact paid root polynomial $P_{\rm paid,F}$ for the
paid branch, the total distinct finite affine numerator is
\[
        \deg\operatorname{rad}\operatorname{lcm}
        \bigl(P_{\rm paid,F},R_{{\rm res},\ge a,F}^{\rm aff}\bigr),
\]
not the sum of the paid and residual degrees.
\end{corollary}

\begin{proof}
Apply \cref{thm:exact-arbitrary-residual-image-ledger} separately at each exact
agreement size $A$, where the co-support size $j=n-A$ and the Hankel window
height $t=A-k$ are fixed.  The radical lcm over $A$ is exactly union and duplicate
coalescing across exact agreement levels.  Intersecting with $Z^q-Z$ counts only
$F$-rational finite slopes.  The final displayed lcm performs the same duplicate
coalescing between paid roots and residual roots.
\end{proof}

\paragraph{Elimination trichotomy for residual charts.}
The exact image lcm exists because split locators are a finite set.  A compressed
rank-pivot certificate is useful when it finds a much smaller polynomial without
enumerating all residual supports.  The following trichotomy is the scanner
contract.

\begin{proposition}[compressed-chart elimination trichotomy]
\label{prop:compressed-chart-elimination-trichotomy}
For a compressed finite affine residual chart over the algebraic closure, exactly
one of the following occurs:
\begin{enumerate}[label=\textup{(\roman*)},leftmargin=2.2em]
\item the saturated compressed ideal is the unit ideal, and the chart is empty;
\item the saturated compressed ideal is proper and the elimination ideal in $F[Z]$ contains a nonzero polynomial, and its
squarefree generator gives a finite parameter envelope containing the exact
support-image lcm of that chart;
\item the saturated compressed ideal is proper and the elimination ideal in $F[Z]$ is zero, equivalently the projection of the
relaxed chart to the $Z$-line is Zariski dense.
\end{enumerate}
If the full homogeneous split equations are present, case \textup{(iii)} can occur
only through a positive-dimensional relaxation or through a contained component
not removed by the noncontainment saturation; an exact split, noncontained chart
has finite parameter image.
\end{proposition}

\begin{proof}
The first three alternatives are the standard elimination theorem for the
projection of an affine variety to a line.  If the full split equations are
present and the leading coordinate is nonzero, the locator coordinate belongs to
the finite set of scalar multiples of the finitely many degree-$j$ locators
supported on $D$.  After projectivizing the locator scale, each fixed support
contributes at most one finite noncontained slope by
\cref{lem:one-support-one-line}.  Hence the exact noncontained split image is a
finite subset of the parameter line.  A dense projection can therefore only come
from a relaxation of the split condition or from an unremoved contained component
on which every parameter is explained but no noncontainment pivot should survive.
\end{proof}

\begin{remark}[compressed-residual status]
The singular bucket has now been reduced to two concrete certificate tasks.  The
first is computational-algebraic: use the kernel-compressed ideals to find small
univariate or homogeneous eliminants.  The second is enumerative but exact: when
compression is not yet sharp enough, print the residual support-image lcm for a
finite residual family and coalesce it with the paid quotient/tangent roots.
Neither task proves the final aperiodic packing theorem by itself, but both
replace the phrase ``singular residual bucket'' by checkable objects with exact
root-count semantics.
\end{remark}

\subsubsection{Fallback and final certificate}

The eliminant certificates above are strongest when the projection to the
parameter line is finite.  A failed eliminant search should not be discarded; it
should be split further or recorded as a singular bucket.  The following fallback
is safe but usually too weak for a final $2^{-128}$ threshold.

\begin{proposition}[dimension-degree fallback]
\label{prop:dimension-degree-fallback}
Let $Y\subseteq\mathbb A^N_F$ be an affine variety of degree at most $\Delta$ and
Krull dimension at most $e$.  Then
\[
        |Y(F)|\le \Delta |F|^e .
\]
Consequently, if a saturated Hankel incidence chart for line or curve parameters
has degree $\Delta_X$ and dimension $e_X$, then the number of bad parameters in
that chart is at most $\Delta_X|F|^{e_X}$.
\end{proposition}

\begin{proof}
This is the standard affine degree point-count bound, proved by induction on
dimension using generic affine hyperplane sections.  The parameter image is no
larger than the incidence set.
\end{proof}

For a threshold statement, use finitely many charts and agreement levels.  Let
$\mathfrak X_A$ be a chart cover of the declared aperiodic branch at exact
agreement $A$.  If each finite affine pivot chart has a univariate eliminant, set
\[
        B_{\rm ap}^{\rm aff,elim}(a)
        :=\sum_{A=a}^n\sum_{X\in\mathfrak X_A}
          \sum_{h=0}^{t_A-1}\deg Q_{X,h} .
\]
For projective slopes add the certified infinity contribution, and for degree-$d$
curves replace the inner sum by the coefficient-pivot sum
$\sum_{i=0}^d\sum_h\deg Q_{X,i,h}$.  Exact root-union tables may replace these
safe sums whenever duplicate parameters across charts are certified.

\begin{definition}[aperiodic Hankel-packing certificate]
\label{def:aperiodic-hankel-certificate}
An aperiodic Hankel-packing certificate for a row, a line or curve family, and an
agreement threshold $a$ consists of:
\begin{enumerate}[label=\textup{(\roman*)},leftmargin=2.2em]
\item the tangent/common-code-line and quotient-image ledgers that have already
been removed;
\item for each exact agreement $A\in\{a,\ldots,n\}$, either a regular minor
certificate $\Delta_A$ or a declaration that the exact-agreement bucket is
Hankel-singular;
\item for every singular or near-capacity bucket, a finite constructible locator
chart cover $\mathfrak X_A$ as in \cref{def:aperiodic-locator-chart};
\item for affine lines, a list of pivot eliminants $Q_{X,h}(Z)$, or a failed-pivot
record explaining which residual locus remains unsplit;
\item for projective slopes, the finite-patch eliminants plus the infinity-chart
empty/nonempty certificate;
\item for curves, coefficient-pivot eliminants $Q_{X,i,h}(\Gamma)$;
\item a root-union table or a safe degree-sum table producing the declared
aperiodic numerator $B_{\rm ap}(a)$;
\item if any chart uses the dimension-degree fallback, its degree, dimension, and
an explicit statement that the resulting bound is only a fallback, not a sharp
threshold certificate.
\end{enumerate}
The verifier checks locator divisibility by $P_D$, chart coverage of the declared
aperiodic complement, nonzero pivot tests, ideal membership of eliminants after
saturation, and final root-count arithmetic.
\end{definition}

\begin{remark}[contributor workflow]
\label{rem:aperiodic-contributor-workflow}
For each row and exact agreement $A$, contributors should try the following in
order.  First test the overdetermined regular minors.  If a nonzero minor is
found, record its root set and remove roots already paid by tangent or quotient
ledgers.  If all maximal minors vanish identically, move the instance to the
singular bucket.  On the singular or near-capacity bucket, build locator charts,
then split by affine pivots $B_h\ne0$; for projective samplers also test the
infinity chart $B=0$, $A\ne0$; for curves split by coefficient pivots
$(V_i)_h\ne0$.  A chart is complete only when it has either a nonzero eliminant,
a certified empty locus, or a named residual obstruction.  Residual obstructions
should be labelled as quotient-periodic, tangent/common-code-line,
subfield/extension exceptional, or candidate new obstruction.
\end{remark}

\begin{remark}[status]
\label{rem:aperiodic-certificate-status}
\Cref{thm:regular-closed-ball-hankel-packing,thm:aperiodic-affine-eliminant,cor:projective-eliminant-certificate,thm:aperiodic-curve-eliminant}
are proved certificate theorems.  They do not prove that all rows admit small
aperiodic eliminants.  They make the remaining MCA safe-side task concrete: after
tangent and quotient branches are removed, construct small eliminants for the
pivot charts, or identify the singular chart that prevents such an eliminant.
\end{remark}


\begin{lemma}[interleaving transfer]\label{lem:inter}
For every linear code $C\subseteq\F^n$, every $s\ge1$ and every $\delta\in(0,1)$:
\[
\eca\bigl(C^{\equiv s},\delta\bigr)\ \ge\ \eca(C,\delta),
\qquad
\emca\bigl(C^{\equiv s},\delta\bigr)\ \ge\ \emca(C,\delta).
\]
\end{lemma}

\begin{proof}
For $h\in\F^n$ write $h^{\Delta}:=(h,\dots,h)\in(\F^n)^s$. For any codeword tuple $\mathbf c=(c^{(1)},\dots,c^{(s)})\in C^{\equiv s}$, the columns where $h^\Delta$ and $\mathbf c$ agree are $\bigcap_i\{j:c^{(i)}_j=h_j\}\subseteq\{j:c^{(1)}_j=h_j\}$, so $\dist_s(h^\Delta,C^{\equiv s})\ge\dist(h,C)$; taking $c^{(1)}=\dots=c^{(s)}$ gives equality:
$\dist_s(h^\Delta,C^{\equiv s})=\dist(h,C)$.

Fix $(f_1,f_2)$ attaining $\eca(C,\delta)$ (resp.\ $\emca$) and consider the pair $(f_1^\Delta,f_2^\Delta)$ for $C^{\equiv s}$, with the slope $\gamma$ acting row-wise, so $f_1^\Delta+\gamma f_2^\Delta=(f_1+\gamma f_2)^\Delta$. By the displayed equality, $\dist_s\bigl((f_1+\gamma f_2)^\Delta,C^{\equiv s}\bigr)\le\delta$ iff $\dist(f_1+\gamma f_2,C)\le\delta$. For the correlated condition, the pair $(f_1^\Delta,f_2^\Delta)$ viewed as a $2s$-row interleaved word has rows $f_1$ ($s$ times) and $f_2$ ($s$ times); by the same column argument its distance to $(C^{\equiv s})^{\equiv2}=C^{\equiv2s}$ equals $\dist_2((f_1,f_2),C^{\equiv2})$. For the MCA condition, a common set $S$ explains $f_1^\Delta$ and $f_2^\Delta$ iff it explains $f_1$ and $f_2$. Hence $\gamma$ is CA-bad (resp.\ MCA-bad) for $(f_1^\Delta,f_2^\Delta)$ iff it is so for $(f_1,f_2)$, and the maxima over pairs can only grow.
\end{proof}

\section{The main theorem}\label{sec:main}

\begin{theorem}[universal cap]\label{thm:main}
Let $\B\subseteq\F$ be finite fields, let $q=|\F|$, let $D\subseteq\B^\times$ be a multiplicative coset of order $n$, and let $k$ have $\rho=k/n\in(0,1)$; set $C:=\RS[\F,D,k]$. Let $N\mid n$ satisfy $a:=n/N\mid k$ and $(1-\rho)N\ge3$. If
\begin{equation}\label{eq:hyp}
        \binom{N}{\rho N+2}
        \ge
        |\B|\Bigl(\frac qk+1\Bigr),
\end{equation}
then, writing
\[
        \delta_N:=1-\rho-\frac2N,
\]
we have
\[
\eca(C,\delta)
>
\frac1{2k}\Bigl(1-\frac nq\Bigr)
\]
for every integer-admissible radius
\[
        \delta_N\le \delta \le 1-\rho-\frac1n .
\]
In particular,
\[
        \emca(C,\delta)
        >
        \frac1{2k}\Bigl(1-\frac nq\Bigr)
\]
for every
\[
        \delta_N\le \delta < 1-\rho .
\]
Consequently,
\[
        \delta^*_C(\varepsilon^*)
        \le
        1-\rho-\frac2N
\]
for every
\[
        \varepsilon^*
        \le
        \frac1{2k}\Bigl(1-\frac nq\Bigr).
\]
\end{theorem}

\begin{proof}
Note first that $\rho N=k/a\in\mathbb Z$, $\ell_2=\rho N+2\le N-1$ by $(1-\rho)N\ge3$, and $\delta_N=1-\rho-2/N>0$ for the same reason. Also $k+1\le n$.

Fix
\[
        \delta\in[\delta_N,\,1-\rho-\tfrac1n].
\]
Then
\[
        \lfloor \delta n\rfloor \le n-k-1,
\]
so \cref{thm:A} is applicable.

Suppose toward a contradiction that
\[
        \eca(C,\delta)
        \le
        \frac1{2k}\Bigl(1-\frac nq\Bigr).
\]
Apply \cref{thm:A} with $\eta=\frac12$. Its hypothesis holds, so
\[
        \Lst(C^+,\delta)
        \le
        \Bigl\lceil 2q\cdot \eca(C,\delta)\Bigr\rceil
        \le
        \Bigl\lceil\frac{q-n}{k}\Bigr\rceil
        <
        \frac{q-n}{k}+1
        \le
        \frac qk+1 .
\]
On the other hand, by \cref{lem:fiber}(ii) and monotonicity of lists in the radius, since $\delta\ge\delta_N$,
\[
        \Lst(C^+,\delta)
        \ge
        \bigl|\Lst(C^+,\delta_N,u_z)\bigr|
        \ge
        \binom{N}{\ell_2}\Big/|\B|
        \ge
        \frac qk+1
\]
by \eqref{eq:hyp}. This is a contradiction. Hence
\[
        \eca(C,\delta)
        >
        \frac1{2k}\Bigl(1-\frac nq\Bigr)
\]
for every $\delta\in[\delta_N,1-\rho-1/n]$.

In particular, the above holds at $\delta=\delta_N$. By \cref{fact:chain},
\[
        \emca(C,\delta_N)
        \ge
        \eca(C,\delta_N)
        >
        \frac1{2k}\Bigl(1-\frac nq\Bigr).
\]
By support-wise MCA monotonicity, \cref{lem:mca-monotone}, the same lower bound for $\emca(C,\delta)$ holds for every
\[
        \delta\in[\delta_N,1-\rho).
\]
Thus no sub-capacity radius
\[
        \delta\ge 1-\rho-\frac2N
\]
is $\varepsilon^*$-admissible whenever
\[
        \varepsilon^*
        \le
        \frac1{2k}\Bigl(1-\frac nq\Bigr),
\]
which proves the claimed bound on $\delta^*_C(\varepsilon^*)$.
\end{proof}

\begin{remark}[constants]\label{rem:eta}
Choosing $\eta=1-\theta$ in \cref{thm:A} improves the conclusion to $\eca>\frac{1-\theta}k(1-\frac nq)$ at the price of strengthening \eqref{eq:hyp} to $\binom N{\rho N+2}\ge|\B|\bigl(\frac q{\theta k}+1\bigr)$; given the hundreds of bits of slack in \eqref{eq:hyp} throughout the challenge envelope (\cref{tab:numerics}), the error constant is effectively $(1-o(1))/k$. We fix $\eta=\frac12$ for cleanliness.
\end{remark}

\begin{remark}[scope]\label{rem:scope}
No smoothness is assumed beyond the existence of a single divisor $N\mid n$ with $a=n/N\mid k$ and \eqref{eq:hyp}. The theorem applies verbatim to $2$-adic, mixed-radix, or any other FFT-friendly multiplicative domain, and --- via the subfield pigeonhole --- to any extension field $\F\supseteq\B$ with $D\subseteq\B$.

The deep-point conversion is an integer-radius theorem with the condition $\lfloor\delta n\rfloor\le n-k-1$. Therefore the no-loss CA lower bound is stated for
\[
        \delta \le 1-\rho-\frac1n .
\]
The MCA threshold cap is stronger in the range needed for the Proximity Prize: once the support-wise MCA lower bound is proved at
\[
        \delta_N=1-\rho-\frac2N,
\]
monotonicity of support-wise MCA extends it to every larger $\delta<1-\rho$.
\end{remark}

\section{Consequences for the grand MCA challenge}\label{sec:grand}

\begin{corollary}[universal field-size cap for the challenge envelope]\label{cor:grand}
Let $\rho\in\{\frac12,\frac14,\frac18,\frac1{16}\}$ and set
\[
N_\rho:=1024\quad(\rho\in\{\tfrac12,\tfrac14,\tfrac18\}),\qquad N_{1/16}:=2048 .
\]
Let $\F$ be any finite field with $q<2^{256}$, let $\B\subseteq\F$ be any subfield, let $D\subseteq\B^\times$ be a multiplicative coset of order $n$ with $N_\rho\mid n$, and let $k=\rho n\le2^{40}$. Then $C=\RS[\F,D,k]$ satisfies
\[
\emca(C,\delta)
>
\frac1{2k}\Bigl(1-\frac nq\Bigr)
\ge
2^{-86}
\gg
2^{-128}
\]
for every
\[
        \delta\in
        \Bigl[1-\rho-\frac2{N_\rho},\,1-\rho\Bigr).
\]
Moreover, the same lower bound holds for $\eca(C,\delta)$ throughout the integer-admissible subinterval
\[
        \delta\in
        \Bigl[1-\rho-\frac2{N_\rho},\,1-\rho-\frac1n\Bigr].
\]
If $q\ge2n$, the lower bound improves to $\ge2^{-42}$. In particular,
\[
\delta^*_C(2^{-128})\;\le\;1-\rho-2^{-9}\quad(\rho\in\{\tfrac12,\tfrac14,\tfrac18\}),
\qquad
\delta^*_C(2^{-128})\;\le\;1-\rho-2^{-10}\quad(\rho=\tfrac1{16}).
\]
\end{corollary}

\begin{proof}
We verify the hypotheses of \cref{thm:main} with $N=N_\rho$. Divisibility: $a=n/N_\rho$ divides $k$ because $k/a=\rho N_\rho\in\mathbb Z$ (equal to $512,256,128,128$ respectively); $(1-\rho)N_\rho\ge128\ge3$. For \eqref{eq:hyp}, the entropy bound $\binom N\ell\ge2^{N\Hb(\ell/N)}/(N+1)$ together with the conservative evaluations of $\Hb$ recorded in \cref{tab:numerics} gives $\log_2\binom{N_\rho}{\rho N_\rho+2}\ge552$ in all four cases, while
\[
|\B|\Bigl(\frac qk+1\Bigr)\;\le\;q\Bigl(\frac qk+1\Bigr)\;\le\;q^2+q\;<\;2^{512}+2^{256}\;<\;2^{513}\;\le\;2^{552},
\]
so \eqref{eq:hyp} holds with at least $39$ bits to spare (and over $500$ bits at rate $\frac12$). Since $D\subseteq\F^\times$, we have $n\le q-1$; for fixed $n$, the function $1-n/q$ is minimized at $q=n+1$, hence $1-n/q\ge1/(n+1)$. Also $n=k/\rho\le16k\le2^{44}$. Therefore
\[
\frac1{2k}\Bigl(1-\frac nq\Bigr)\ge\frac1{2k(n+1)}>2^{-86}.
\]
If $q\ge2n$ then $1-n/q\ge1/2$, giving the sharper $1/(4k)\ge2^{-42}$. The gap is $2/N_\rho=2^{-9}$ for $\rho\in\{1/2,1/4,1/8\}$ and $2/N_\rho=2^{-10}$ for $\rho=1/16$. \Cref{thm:main} gives the MCA lower bound throughout the stated sub-capacity interval, and gives the CA lower bound on the integer-admissible subinterval ending at $1-\rho-1/n$.
\end{proof}

\begin{table}[ht]
\centering
\begin{tabular}{ccccccc}
\toprule
$\rho$ & $N_\rho$ & $\ell_2=\rho N_\rho+2$ & $\Hb(\ell_2/N_\rho)\ge$ & $\log_2\binom{N_\rho}{\ell_2}\ge$ & needed & gap $2/N_\rho$\\
\midrule
$1/2$  & $1024$ & $514$ & $0.99998$ & $1013$ & $513$ & $2^{-9}$\\
$1/4$  & $1024$ & $258$ & $0.8143$  & $823$  & $513$ & $2^{-9}$\\
$1/8$  & $1024$ & $130$ & $0.5490$  & $552$  & $513$ & $2^{-9}$\\
$1/16$ & $2048$ & $130$ & $0.34105$ & $687$  & $513$ & $2^{-10}$\\
\bottomrule
\end{tabular}
\caption{Numerics for \cref{cor:grand}: $\log_2\binom{N}{\ell}\ge N\Hb(\ell/N)-\log_2(N+1)$, rounded down. The ``needed'' column is the universal upper bound $\log_2\bigl(|\B|(q/k+1)\bigr)<513$ valid throughout the envelope $q<2^{256}$, $k\ge1$.}
\label{tab:numerics}
\end{table}

\begin{remark}[hypothesis $N_\rho\mid n$]
For the $2$-adic domains of the challenge this holds whenever $n\ge2^{10}$ (resp.\ $2^{11}$), i.e.\ for all but very small instances. For mixed-radix smooth $n$, replace $N_\rho$ by any divisor $N\mid n$ of comparable size with $\rho N\in\mathbb Z$; the entropy count in \cref{tab:numerics} is insensitive to the exact value of $N$ at fixed $\log_2 N$. The divisibility-free fiber lemma removes the integrality requirement $\rho N\in\mathbb Z$ altogether: \cref{lem:phi-fiber}, \cref{cor:mixed-radix} caps every mixed-radix row possessing \emph{any} divisor in the printed dyadic windows.
\end{remark}

\begin{remark}[coarse grand-challenge caps and the remaining frontier]\label{rem:twotier}
Combining \cref{cor:grand} with the grand-challenge cap of \cite[Prop.\ ``Grand-challenge cap'']{Cho26b} gives the following coarse uniform negative caps for smooth multiplicative domains:
\[
\delta^*_C(2^{-128})\;\le\;
\begin{cases}
1-\rho-\frac1{64}, & \F\ \text{prime},\ q\lesssim2^{150},\\[2pt]
1-\rho-2^{-9}, & \rho\in\{\frac12,\frac14,\frac18\},\ \text{every}\ \F,\ q<2^{256},\\[2pt]
1-\rho-2^{-10}, & \rho=\frac1{16},\ \text{every}\ \F,\ q<2^{256}.
\end{cases}
\]
These inequalities are not the final frontier studied in the remainder of this paper.  They are first-tier caps: robust, uniform, and useful for placing the grand challenge beyond the Johnson region, but far from the capacity edge at the deployed parameters.  The rest of the paper has a different purpose.  It refines these coarse caps into the identity-scale and flexible-budget frontier certificates, proves the structural ledgers that rule out the already-paid tangent, quotient, extension-line, and split top-chart strata, and isolates the remaining safe-side inputs as quotient/prefix flatness, the base-field-normalized split-pencil census, and primitive shift-pair control.

Thus the quoted caps should be read as the broad negative band available before the frontier analysis, not as a stopping point.  Together with the positive results collected in \cite[Table~1]{ABF26}, they give the initial sandwich; the later frontier package states the sharpened active sandwich.  Circle rows --- line rounds and bivariate circle codes --- have analogous coarse caps; see \cref{cor:circle-grand,rem:circle-transfer}.
\end{remark}

\begin{corollary}[deployed parameters: KoalaBear sextic]\label{cor:deployed}
Let $\B=\F_p$ with $p=2^{31}-2^{24}+1$, let $\F=\F_{p^6}$ (so $q=p^6\approx2^{185.93}$), let $D\subseteq\B^\times$ be the subgroup of order $n=2^{21}$, and let $k=2^{20}$, $\rho=\frac12$, as in \cite[\S6.3]{ABF26}. Then
\[
\emca\bigl(\RS[\F_{p^6},D,2^{20}],\,\delta\bigr)
>
\bigl(1-2^{-164}\bigr)\,2^{-21}
>
2^{-22}
\]
for every
\[
        \delta\in\Bigl[\frac12-2^{-7},\,\frac12\Bigr).
\]
Moreover, the same lower bound holds for $\eca$ throughout the integer-admissible subinterval
\[
        \delta\in
        \Bigl[\frac12-2^{-7},\,\frac12-2^{-21}\Bigr).
\]
\end{corollary}

\begin{proof}
Apply \cref{thm:main} with $N=256$, $a=2^{13}\mid k$, $(1-\rho)N=128\ge3$, gap $2/N=2^{-7}$. Hypothesis \eqref{eq:hyp}: $\log_2\binom{256}{130}\ge256\cdot\Hb(130/256)-\log_2257\ge256\cdot0.99982-8.01>247$, while $|\B|(q/k+1)<2^{31}\,(2^{166}+1)<2^{198}$. The error bound is $\frac1{2k}(1-n/q)$ with $n/q=2^{21}/p^6<2^{-164}$. \Cref{thm:main} gives the stated MCA bound on $[\frac12-2^{-7},\frac12)$ and the stated CA bound on $[\frac12-2^{-7},\frac12-2^{-21})$.
\end{proof}

\begin{corollary}[all interleaved rows of {\cite[Tables 2--3]{ABF26}}]\label{cor:rows}
In the setting of \cref{cor:deployed}, for each $s=2^j$ with $0\le j\le12$ let $C_s:=\RS[\F_{p^6},D_s,k_s]$ with $|D_s|=n_s=2^{21}/s$ and $k_s=2^{20}/s$ (rate $\frac12$, $s\,n_s=2^{21}$ as deployed). Then for every $\delta\in\bigl[\frac12-2^{-7},\,\frac12\bigr)$,
\[
\emca\bigl(C_s^{\equiv s},\delta\bigr)
>
\frac s{2^{21}}\bigl(1-2^{-164}\bigr)
\ge
2^{-21}\bigl(1-2^{-164}\bigr).
\]
Moreover, the same lower bound holds for $\eca(C_s^{\equiv s},\delta)$ on the integer-admissible subinterval
\[
        \delta\in
        \Bigl[\frac12-2^{-7},\,\frac12-\frac1{n_s}\Bigr).
\]
Every row of the deployed parameter family therefore fails the $2^{-128}$ target at gap $2^{-7}$ below capacity.
\end{corollary}

\begin{proof}
For each $s\le2^{12}$: $256\mid n_s$ (as $n_s\ge2^9$), $a_s=n_s/256=2^{13}/s$ divides $k_s=2^{20}/s$, and $(1-\rho)\cdot256=128\ge3$. Hypothesis \eqref{eq:hyp} for $C_s$: $|\B|(q/k_s+1)<2^{31}(2^{166}\,s+1)<2^{198}\,s\le2^{210}\le2^{247}\le\binom{256}{130}$. \Cref{thm:main} gives the MCA lower bound for $C_s$ throughout $[\frac12-2^{-7},\frac12)$, and the CA lower bound throughout $[\frac12-2^{-7},\frac12-1/n_s]$. \Cref{lem:inter} transfers the corresponding lower bounds to the $s$-interleaved code.
\end{proof}

\begin{proposition}[independent slacked variant via \cref{thm:B}]\label{prop:slacked}
In the setting of \cref{thm:main}, assume instead $(1-\rho)N\ge2$, $n\ge3N$, and
\[
\binom{N}{\rho N+1}\;\ge\;|\B|\cdot q .
\]
Then
\[
\eca\Bigl(C,\ 1-\rho-\frac1N+\frac2n,\ 1-\rho-\frac1n\Bigr)\;\ge\;\frac1{2n}.
\]
In particular, with $N=N_\rho$ as in \cref{cor:grand} and $n\ge3N_\rho$, every challenge-envelope instantiation satisfies this slacked CA lower bound, and $\frac1{2n}\ge2^{-45}\gg2^{-128}$.
\end{proposition}

\begin{proof}
Set $\delta_0:=1-\rho-\frac1N$. By $(1-\rho)N\ge2$, $\delta_0\in(0,1-\rho)$. By $n\ge3N$, we also have $\delta_0+2/n\le1-\rho-1/n$, so the parameters in \cref{thm:B} are admissible. By \cref{lem:fiber}(i), some $u_z$ carries at least $\binom N{\rho N+1}/|\B|\ge q$ codewords of $C$ within radius $\delta_0$, so $\Lst(C,\delta_0)\ge q$. The contrapositive of \cref{thm:B} gives
\[
\eca\Bigl(C,\delta_0+\frac2n,1-\rho-\frac1n\Bigr)\ge\frac1{2n}.
\]
For the numerical claim, the worst case of $\log_2\binom{N_\rho}{\rho N_\rho+1}$ over the four rates is at $\rho=1/8$, where the entropy bound used in \cref{tab:numerics} certifies $\log_2\binom{1024}{129}\ge1024\,\Hb(129/1024)-\log_2 1025\ge549$ (the exact value is $\approx554.7$); hence $\binom{N_\rho}{\rho N_\rho+1}\ge2^{549}>2^{512}>|\B|q$ throughout $q<2^{256}$. Finally $n=k/\rho\le16k\le2^{44}$, so $1/(2n)\ge2^{-45}$.
\end{proof}

\begin{remark}
\Cref{prop:slacked} is logically independent of \cref{thm:A} and yields a weaker error ($\frac1{2n}$ versus $\frac1{2k}$) in the \emph{maximally slacked} form of correlated agreement --- the internal radius sits at capacity minus $\frac1n$ --- which is the weakest failure mode one can certify. That even this fails by $83$ orders of magnitude (in bits) above the $2^{-128}$ target gives a separate fallback route through the imported BCHKS/ABF theorem.
\end{remark}

\section{Subfield confinement and the shape of certifying lines}\label{sec:subfield}

The deployed instantiations run the protocol over an extension $\F/\B$ while keeping the domain in the base: $D\subseteq\B$. The subfield-confinement theorem of \cite[Thm.\ ``Subfield confinement'']{Cho26b} pins down where bad slopes can live for lines whose words are base-rational: all support-wise MCA-bad slopes lie in $\B$. The following lemma refines this by tracking two radii, so that it also covers the proximity-loss form of correlated agreement (\cref{def:ca}); part (b) is the statement of \cite{Cho26b}, reproved here for completeness in the present notation.

\begin{lemma}[subfield confinement; refines {\cite{Cho26b}}]\label{lem:confine}
Let $\B\subseteq\F$ be finite fields, $D\subseteq\B^\times$, $C=\RS[\F,D,k]$, and let $f,g\in\B^n\subseteq\F^n$ be $\B$-valued words. Then:
\begin{enumerate}[label=\textup{(\alph*)}]
\item for any $0\le\delta_{\mathrm{fld}}\le\delta_{\mathrm{int}}<1$, every CA-bad slope for $(f,g)$ at radii $(\delta_{\mathrm{fld}},\delta_{\mathrm{int}})$ lies in $\B$;
\item for any $\delta\in(0,1)$, every MCA-bad slope for $(f,g)$ at radius $\delta$ lies in $\B$.
\end{enumerate}
Consequently the bad-slope set of any $\B$-valued pair has density at most $|\B|/q$ in $\F$. The same holds for pairs $(\lambda f,\lambda g)$ with $\lambda\in\F^\times$ and $f,g$ $\B$-valued, by $\F$-linearity of $C$.
\end{lemma}

\begin{proof}
Fix a $\B$-basis $1=\beta_0,\beta_1,\dots,\beta_{m-1}$ of $\F$. Every word $w\in\F^n$ decomposes uniquely as $w=\sum_i\beta_i w_i$ with $w_i\in\B^n$, and every codeword $c\in\RS[\F,D,k]$ decomposes as $c=\sum_i\beta_i c_i$ with $c_i\in\RS[\B,D,k]$: write the coefficient vector of the underlying polynomial in the basis and evaluate at $D\subseteq\B$.

Let $z\in\F\setminus\B$, say $z=\sum_iz_i\beta_i$ with $z_{i^*}\ne0$ for some $i^*\ge1$. We first record the support-wise consequence. Suppose that $c\in C$ agrees with $f+zg$ on some set $S\subseteq D$. Since $f_j,g_j\in\B$, the basis components of $(f+zg)_j$ are $f_j+z_0g_j$ (component $0$) and $z_ig_j$ (components $i\ge1$). Agreement on $S$ holds componentwise, so on $S$:
\[
f+z_0g=c_0,\qquad z_{i^*}\,g=c_{i^*}.
\]
Hence $g$ agrees on $S$ with $z_{i^*}^{-1}c_{i^*}\in C$, and then $f=(f+z_0g)-z_0g$ agrees on $S$ with $c_0-z_0z_{i^*}^{-1}c_{i^*}\in C$. Thus every support on which the line point $f+zg$ is explained by $C$ also explains the pair $(f,g)$ on that same support.

For CA, if $\dist(f+zg,C)\le\delta_{\mathrm{fld}}$, choose such an $S$ with $|S|\ge(1-\delta_{\mathrm{fld}})n$; the preceding paragraph gives $\dist_2((f,g),C^{\equiv2})\le\delta_{\mathrm{fld}}\le\delta_{\mathrm{int}}$, so $z$ is not CA-bad. For MCA, the preceding paragraph rules out every possible bad support $S$ in the support-wise definition. Hence $z$ is not MCA-bad either.
\end{proof}

\begin{corollary}[certifying lines over extensions are genuinely $\F$-valued]\label{cor:Fvalued}
At the parameters of \cref{cor:deployed}, $|\B|/q=p^{-5}<2^{-154}$. Hence any pair $(f,g)$ whose CA- or MCA-bad-slope density at some radius
\[
        \delta\in
        \Bigl[\frac12-2^{-7},\,\frac12\Bigr)
\]
exceeds $2^{-154}$ contains a word that is not $\B$-valued, even after removing a common scalar factor. In particular, any pair attaining the MCA lower bound $>2^{-22}$ of \cref{cor:deployed} is genuinely $\F$-valued. On the integer-admissible subinterval
\[
        \delta\in
        \Bigl[\frac12-2^{-7},\,\frac12-2^{-21}\Bigr),
\]
the same statement applies to pairs attaining the CA lower bound. All explicit failure witnesses currently known on these domains (the locator lines of \cite{Cho26a,Cho26b} and the fiber words of \cref{lem:fiber}) are $\B$-rational and are therefore inert here: the failure certified by \cref{cor:deployed} is fiber-borne and, at present, non-constructive.
\end{corollary}

\begin{proof}
Immediate from \cref{lem:confine} and \cref{cor:deployed}: a $\B$-valued pair, up to common scalar, has CA- and MCA-bad-slope density at most $|\B|/q=p^{-5}<2^{-154}<2^{-22}$. Thus any pair witnessing the deployed lower bound cannot be of this form.
\end{proof}

\begin{remark}
\Cref{lem:confine,cor:Fvalued} sharpen the division of labor between the two grand challenges of \cite{ABF26} over deployed extensions: the \emph{list} side does not deflate ($\RS[\B,D,k]\subseteq\RS[\F,D,k]$, so all base-field list lower bounds, including \cref{lem:fiber}, persist verbatim), while every sub-$2^{-128}$ \emph{MCA} attack by explicit base-rational witnesses is killed by confinement --- and yet MCA still fails at $2^{-22}$.  The failure can also be made constructive once a large list is supplied: \cref{cor:extension-pole-deep-list-floor,cor:extension-pole-quotient-remainder-floor} allow one to choose a simple pole in $\F\setminus\B$ and print its value-bucket table. The deflation and decoupling corollaries of \cite{Cho26b} record the same division at the no-loss radius, using the same corrected rounding $p^{-5}=2^{-154.94\ldots}<2^{-154}$ as in \cref{cor:Fvalued}. Exhibiting deployed-size bucket tables for such lines is the explicit-witness question \cref{prob:explicit}.
\end{remark}


\section{Beyond multiplicative smoothness: map-smooth domains, Chebyshev fibers, and the circle-code cap}\label{sec:isogeny}

The fiber lemma \cref{lem:fiber} uses exactly two features of the power map $x\mapsto x^a$ on a multiplicative coset: its fibers inside the domain are complete and of uniform size $a$, and the locators $\prod_{b\in A}(X^a-b)$ are \emph{lacunary} --- expanded in powers of $X^a$, their two top terms occupy degrees $a\ell$ and $a(\ell-1)$, leaving everything from degree $a(\ell-2)$ downward to the listed codeword. This section isolates exactly these two features. \Cref{lem:phi-fiber} proves the fiber lemma for an arbitrary polynomial folding map $\varphi$ whose fibers inside the domain are complete and uniform, and simultaneously removes the divisibility hypothesis $a\mid k$ by rounding the locator length; \cref{thm:phi-cap} is the corresponding universal cap, containing \cref{thm:main} as the special case $\varphi=X^a$ with $a\mid k$. We then verify both features for the deployed isogeny-smooth family, the circle-group domains of Circle STARKs over $p=2^{31}-1$ \cite{HLP24}: twin-coset $x$-coordinate domains carry exact Chebyshev fibers at every dyadic scale (\cref{lem:cheb-fibers}), and over any challenge field containing $i$ the bivariate circle code is diagonally equivalent to a Reed--Solomon code on a torus twin coset, an equivalence preserving lists and both correlated-agreement errors exactly (\cref{lem:diag-invariance,lem:circle-rs}). The resulting universal caps (\cref{cor:circle-grand}) and deployed instantiations (\cref{cor:circle-deployed}) settle the circle-group case of the isogeny-smooth problem posed below; what remains of that problem is isolated, with a precise obstruction, in \cref{rem:ecfft-obstruction} and restated as the genus-one/circle transport question \cref{prob:isogeny}.

\subsection{Map-smooth domains and a divisibility-free fiber lemma}
\label{sec:map-smooth}

\begin{definition}[map-smooth domain]\label{def:map-smooth}
Let $\B\subseteq\F$ be finite fields, let $\varphi\in\B[X]$ have degree $a\ge1$, and let $D\subseteq\B$ be a finite set of size $n$. The set $D$ is \emph{$(\varphi,a)$-smooth} if every nonempty fiber $\varphi^{-1}(w)\cap D$, $w\in\varphi(D)$, has exactly $a$ elements. In that case $Q:=\varphi(D)\subseteq\B$ has size $N:=n/a$, and $D$ is the disjoint union of the $N$ fibers $S_b:=\{x\in D:\varphi(x)=b\}$, $b\in Q$.
\end{definition}

Two remarks on scope. First, a multiplicative coset $D\subseteq\B^\times$ of order $n$ is $(X^a,a)$-smooth for every $a\mid n$: the fibers are the cosets of the order-$a$ subgroup contained in $D$'s underlying group. Second, nothing is assumed about the roots of $\varphi-w$ outside $D$: the definition constrains fibers inside $D$ only, so separability or squarefreeness of $\varphi-w$ plays no role below.

\begin{lemma}[map-smooth fiber lemma, divisibility-free]\label{lem:phi-fiber}
Let $\B\subseteq\F$ and let $D\subseteq\B$ be a $(\varphi,a)$-smooth domain of size $n$, with $Q=\varphi(D)$ and $N=n/a$ as above. Let $1\le k<n$ and $\rho:=k/n$. Put
\[
\ell_1:=\Bigl\lfloor\frac{k-1}a\Bigr\rfloor+2,\qquad
\ell_2:=\Bigl\lfloor\frac{k}a\Bigr\rfloor+2,\qquad
A_1:=a\ell_1,\qquad A_2:=a\ell_2,
\]
so that $k+a\le A_1\le k+2a-1$ and $k+a+1\le A_2\le k+2a$, with $A_1=k+a$ and $A_2=k+2a$ when $a\mid k$.
\begin{enumerate}[label=\textup{(\roman*)}]
\item If $\ell_1\le N-1$, there exists $z\in\B$ such that, with $u_z:=\bigl(\varphi(x)^{\ell_1}+z\,\varphi(x)^{\ell_1-1}\bigr)_{x\in D}\in\B^n$,
\[
\bigl|\Lst\bigl(\RS[\F,D,k],\,1-\tfrac{A_1}n,\,u_z\bigr)\bigr|\ \ge\ \binom N{\ell_1}\Big/|\B| .
\]
\item If $\ell_2\le N-1$, there exists $z\in\B$ such that, with $u_z:=\bigl(\varphi(x)^{\ell_2}+z\,\varphi(x)^{\ell_2-1}\bigr)_{x\in D}\in\B^n$,
\[
\bigl|\Lst\bigl(\RS[\F,D,k+1],\,1-\tfrac{A_2}n,\,u_z\bigr)\bigr|\ \ge\ \binom N{\ell_2}\Big/|\B| .
\]
\end{enumerate}
In both cases $u_z$ is $\B$-valued, and the listed codewords lie in $\RS[\B,D,a(\ell_1-2)+1]\subseteq\RS[\B,D,k]$ \textup(resp.\ $\RS[\B,D,a(\ell_2-2)+1]\subseteq\RS[\B,D,k+1]$\textup).
\end{lemma}

\begin{proof}
First the bracketing of $A_1,A_2$: writing $k=a\lfloor k/a\rfloor+s$ with $0\le s\le a-1$ gives $A_2=k-s+2a\in[k+a+1,k+2a]$; writing $k-1=a\lfloor(k-1)/a\rfloor+s_1$ with $0\le s_1\le a-1$ gives $A_1=k-1-s_1+2a\in[k+a,k+2a-1]$; and $a\mid k$ forces $s=0$ and $s_1=a-1$.

We prove (ii); part (i) is the identical computation one degree rung lower, with $k+1$ replaced by $k$ throughout. Let $A\subseteq Q$ with $|A|=\ell_2$, and put
\[
P_A(Y):=\prod_{b\in A}(Y-b)=Y^{\ell_2}-\eone(A)\,Y^{\ell_2-1}+R_A(Y),
\qquad
R_A(Y):=\sum_{j=2}^{\ell_2}(-1)^je_j(A)\,Y^{\ell_2-j},
\]
so $\deg_YR_A\le\ell_2-2$, and set $L_A(X):=P_A(\varphi(X))\in\B[X]$. The polynomial $L_A$ vanishes on
\[
S_A:=\varphi^{-1}(A)\cap D=\bigsqcup_{b\in A}S_b,
\qquad
|S_A|=a\ell_2=A_2,
\]
since the fibers $S_b$ are complete, pairwise disjoint, and of size exactly $a$.

Set $z_A:=-\eone(A)\in\B$ and $c_A:=\bigl(-R_A(\varphi(x))\bigr)_{x\in D}$. Since $\deg_XR_A(\varphi(X))\le a(\ell_2-2)=A_2-2a\le k$, we have $c_A\in\RS[\B,D,a(\ell_2-2)+1]\subseteq\RS[\B,D,k+1]\subseteq\RS[\F,D,k+1]$. For every $x\in S_A$,
\[
0=L_A(x)=\varphi(x)^{\ell_2}+z_A\,\varphi(x)^{\ell_2-1}+R_A(\varphi(x)),
\qquad\text{i.e.}\qquad
u_{z_A}(x)=c_A(x),
\]
so $u_{z_A}$ and $c_A$ agree on at least $A_2$ points of $D$ and $\dist(u_{z_A},c_A)\le1-A_2/n$.

The assignment $A\mapsto c_A$ is injective among subsets sharing a common value of $z_A$. Indeed, suppose $z_A=z_{A'}$ and $c_A=c_{A'}$. Then $R_A(\varphi(X))$ and $R_{A'}(\varphi(X))$ are polynomials of degree $\le k<n$ agreeing at all $n$ points of $D$, hence equal as polynomials. Composition with the nonconstant $\varphi$ is injective on $\F[Y]$: if $R:=R_A-R_{A'}\ne0$ has degree $d\ge1$, then $R(\varphi(X))$ has degree $ad\ge1$ with leading coefficient $\mathrm{lc}(R)\,\mathrm{lc}(\varphi)^d\ne0$, and if $d=0$ then $R(\varphi(X))=R$ is a nonzero constant; either way $R(\varphi(X))\ne0$. Hence $R_A=R_{A'}$, so $P_A=Y^{\ell_2}+z_AY^{\ell_2-1}+R_A(Y)=P_{A'}$, and the root multisets give $A=A'$.

Finally, $z_A=-\eone(A)$ lies in $\B$ because $A\subseteq Q\subseteq\B$, so it takes at most $|\B|$ values as $A$ ranges over the $\binom N{\ell_2}$ subsets of $Q$ of size $\ell_2$ (here $\ell_2\le N-1\le N$ is used to have this many subsets). Some value $z\in\B$ is attained by at least $\binom N{\ell_2}/|\B|$ subsets, and the corresponding codewords $c_A$ are pairwise distinct elements of $\Lst(\RS[\F,D,k+1],1-A_2/n,u_z)$.
\end{proof}

\begin{remark}[the divisibility hypothesis is removed]\label{rem:no-divisibility}
With $\varphi=X^a$, $D$ a multiplicative coset, and $a\mid k$, \cref{lem:phi-fiber} is \cref{lem:fiber} verbatim: $A_1=k+a$, $A_2=k+2a$, and $u_z$ has the printed monomial shape. The rounding $\ell_2=\lfloor k/a\rfloor+2$ removes the hypothesis $a\mid k$ at the cost of a certified agreement $A_2\in[k+a+1,k+2a]$ rather than exactly $k+2a$; the certified radius therefore satisfies
\[
1-\rho-\frac2N\ \le\ 1-\frac{A_2}n\ \le\ 1-\rho-\frac1N-\frac1n,
\]
so the listed ball always sits at least one quotient step $\approx1/N$ below capacity. This matters precisely where $a\mid k$ genuinely fails: odd-dimension circle codes (\cref{cor:circle-grand}) and mixed-radix multiplicative rows without an integral $\rho N$ (\cref{cor:mixed-radix}). \Cref{rem:slacktwo} already used slack two to move the divisibility demand from $k+1$ onto $k$; the present rounding removes it altogether.
\end{remark}

\begin{theorem}[universal cap for map-smooth domains]\label{thm:phi-cap}
Let $\B\subseteq\F$ be finite fields, $q:=|\F|$, and let $D\subseteq\B$ be a $(\varphi,a)$-smooth domain of size $n<q$. Let $1\le k<n$, $\rho:=k/n$, $C:=\RS[\F,D,k]$, and let $N=n/a$, $\ell_2=\lfloor k/a\rfloor+2$, $A_2=a\ell_2$ be as in \cref{lem:phi-fiber}. Assume $\ell_2\le N-1$ and
\begin{equation}\label{eq:hyp-phi}
\binom N{\ell_2}\ \ge\ |\B|\Bigl(\frac qk+1\Bigr).
\end{equation}
Then
\[
\eca(C,\delta)\;>\;\frac1{2k}\Bigl(1-\frac nq\Bigr)
\qquad\text{for every }\ \delta\in\Bigl[1-\frac{A_2}n,\ 1-\rho-\frac1n\Bigr],
\]
and
\[
\emca(C,\delta)\;>\;\frac1{2k}\Bigl(1-\frac nq\Bigr)
\qquad\text{for every }\ \delta\in\Bigl[1-\frac{A_2}n,\ 1-\rho\Bigr).
\]
Consequently
\[
\delta^*_C(\varepsilon^*)\ \le\ 1-\frac{A_2}n\ \le\ 1-\rho-\frac{a+1}n
\qquad\text{for every }\ \varepsilon^*\le\frac1{2k}\Bigl(1-\frac nq\Bigr).
\]
\end{theorem}

\begin{proof}
Since $\ell_2\le N-1$ we have $A_2\le n-a<n$, and since $A_2\ge k+a+1\ge k+2$ both stated intervals are nonempty, with $1-A_2/n\le1-\rho-(a+1)/n\le1-\rho-2/n$.

Fix $\delta\in[1-A_2/n,\,1-\rho-1/n]$; then $\lfloor\delta n\rfloor\le n-k-1$, so \cref{thm:A} is applicable to $C$ at radius $\delta$. Suppose toward a contradiction that $\eca(C,\delta)\le\frac1{2k}(1-\frac nq)$. Applying \cref{thm:A} with $\eta=\frac12$ gives, for $C^+:=\RS[\F,D,k+1]$,
\[
\Lst(C^+,\delta)\ \le\ \bigl\lceil2q\,\eca(C,\delta)\bigr\rceil
\ \le\ \Bigl\lceil\frac{q-n}k\Bigr\rceil
\ <\ \frac{q-n}k+1\ \le\ \frac qk+1 .
\]
On the other hand, by \cref{lem:phi-fiber}(ii) and monotonicity of lists in the radius (using $\delta\ge1-A_2/n$),
\[
\Lst(C^+,\delta)\ \ge\ \bigl|\Lst\bigl(C^+,1-\tfrac{A_2}n,u_z\bigr)\bigr|\ \ge\ \binom N{\ell_2}\Big/|\B|\ \ge\ \frac qk+1
\]
by \eqref{eq:hyp-phi}, a contradiction. Hence the $\eca$ lower bound holds throughout the stated integer-admissible interval. In particular it holds at $\delta_0:=1-A_2/n$, where \cref{fact:chain} gives the same lower bound for $\emca(C,\delta_0)$, and support-wise MCA monotonicity (\cref{lem:mca-monotone}) extends it to every $\delta\in[\delta_0,1-\rho)$. The bound on $\delta^*_C$ follows exactly as in \cref{thm:main}, using $A_2\ge k+a+1$.
\end{proof}

\begin{remark}[relation to \cref{thm:main}]\label{rem:phi-cap-scope}
With $\varphi=X^a$, $D$ a multiplicative coset, and $a\mid k$, \cref{thm:phi-cap} is \cref{thm:main} verbatim: $A_2=k+2a$, the left endpoint is $\delta_N=1-\rho-2/N$, hypothesis \eqref{eq:hyp-phi} is \eqref{eq:hyp}, and $\ell_2\le N-1$ is the printed condition $(1-\rho)N\ge3$. The generalization is strict in two independent directions --- the folding map and the removal of $a\mid k$ --- at no cost in constants. As in \cref{rem:eta}, the constant $\frac1{2k}$ can be pushed to $\frac{1-o(1)}k$ given slack in \eqref{eq:hyp-phi}.
\end{remark}

\begin{corollary}[mixed-radix multiplicative rows, divisibility-free]\label{cor:mixed-radix}
Let $\rho\in\{\frac12,\frac14,\frac18,\frac1{16}\}$, let $\F$ be any finite field with $q:=|\F|<2^{256}$, let $\B\subseteq\F$ be any subfield, let $D\subseteq\B^\times$ be a multiplicative coset of order $n$, and let $k=\rho n\le2^{40}$, $C=\RS[\F,D,k]$. Suppose $n$ has \emph{some} divisor $N$ with
\[
2^{10}\le N\le2^{11}\ \ \bigl(\rho\in\{\tfrac12,\tfrac14,\tfrac18\}\bigr),
\qquad\text{resp.}\qquad
2^{11}\le N\le2^{12}\ \ \bigl(\rho=\tfrac1{16}\bigr),
\]
with no divisibility condition relating $N$ and $k$ \textup(in particular $\rho N$ need not be an integer\textup). Then
\[
\emca(C,\delta)\;>\;\frac1{2k}\Bigl(1-\frac nq\Bigr)\;\ge\;2^{-86}\;\gg\;2^{-128}
\qquad\text{for every }\ \delta\in\Bigl[1-\rho-\frac1N,\ 1-\rho\Bigr),
\]
and the same lower bound holds for $\eca(C,\delta)$ on the integer-admissible subinterval $[1-\frac{A_2}n,\,1-\rho-\frac1n]$, with $A_2$ as in \cref{lem:phi-fiber} for $a=n/N$. In particular
\[
\delta^*_C(2^{-128})\ \le\ 1-\rho-2^{-11}\ \ \bigl(\rho\in\{\tfrac12,\tfrac14,\tfrac18\}\bigr),
\qquad
\delta^*_C(2^{-128})\ \le\ 1-\rho-2^{-12}\ \ \bigl(\rho=\tfrac1{16}\bigr).
\]
\end{corollary}

\begin{proof}
Apply \cref{thm:phi-cap} with $\varphi=X^{n/N}$ and $a=n/N$; the coset $D$ is $(X^a,a)$-smooth, and $n<q$ as in \cref{cor:grand}. Since $\lfloor k/a\rfloor\in(k/a-1,k/a]$, we have $\ell_2/N\in(\rho+\frac1N,\rho+\frac2N]$, so $\ell_2\le\rho N+2\le N/2+2\le N-1$. For \eqref{eq:hyp-phi}: the interval $(\rho+\frac1N,\rho+\frac2N]$ is contained in $[\frac18,\,0.502]$ for the first three rates and in $[\frac1{16},\,0.0645]$ for $\rho=\frac1{16}$; on each interval the concave function $\Hb$ attains its minimum at an endpoint, and $\Hb(0.502)\ge0.999\ge\Hb(\frac18)$ while $\Hb(0.0645)\ge\Hb(\frac1{16})$, so $\Hb(\ell_2/N)\ge\Hb(\frac18)\ge0.5435$, resp.\ $\Hb(\ell_2/N)\ge\Hb(\frac1{16})\ge0.3372$. Hence, using the worst $N$ in each window,
\[
\log_2\binom N{\ell_2}\ \ge\ N\,\Hb(\ell_2/N)-\log_2(N+1)\ \ge\
\begin{cases}
2^{10}\cdot0.5435-\log_2(2^{11}+1)\ \ge\ 545,\\[2pt]
2^{11}\cdot0.3372-\log_2(2^{12}+1)\ \ge\ 678,
\end{cases}
\]
while $|\B|(q/k+1)\le q(q/k+1)\le q^2+q<2^{513}$, exactly as in \cref{cor:grand}. \Cref{thm:phi-cap} gives the lower bounds on $[1-\frac{A_2}n,1-\rho-\frac1n]$ and $[1-\frac{A_2}n,1-\rho)$; since $A_2\ge k+a+1$, the MCA interval contains $[1-\rho-\frac1N,1-\rho)$, and $\frac1N\ge2^{-11}$, resp.\ $2^{-12}$. The error numerics $\frac1{2k}(1-\frac nq)\ge\frac1{2k(n+1)}>2^{-86}$ (using $n\le16k\le2^{44}$) are those of \cref{cor:grand}.
\end{proof}

For FFT-friendly mixed-radix orders $n=2^{\alpha}3^{\beta}5^{\gamma}\cdots$ the divisor set is dense on the logarithmic scale, so the window hypothesis is mild; and taking the smallest available $N$ in the window recovers the smallest gap.

\subsection{Twin cosets on the circle, Chebyshev fibers, and torus uniformization}
\label{sec:circle-geometry}

Throughout this subsection $p$ is a prime with $p\equiv3\pmod4$, so that $X^2+1$ is irreducible over $\F_p$ and $\F_{p^2}=\F_p(i)$ with $i^2=-1$. Let
\[
\UU\ :=\ \{u\in\F_{p^2}^{\times}:\ u^{p+1}=1\}
\]
be the norm-one torus, a cyclic group of order $p+1$; for the deployed Mersenne prime $p=2^{31}-1$ one has $p+1=2^{31}$. The circle $\mathcal C(\F_p):=\{(x,y)\in\F_p^2:x^2+y^2=1\}$ is in bijection with $\UU$ via
\[
\iota(x,y)\ :=\ x+iy,
\qquad
\iota^{-1}(u)\ =\ \Bigl(\frac{u+u^{-1}}2,\ \frac{u-u^{-1}}{2i}\Bigr),
\]
using $u^p=u^{-1}$ for $u\in\UU$ and $i^p=-i$. Write $\chi(u):=\frac{u+u^{-1}}2\in\F_p$ for the $x$-coordinate map; $\chi(u)=\chi(v)$ holds if and only if $v\in\{u,u^{-1}\}$. The Chebyshev polynomials of the first kind, normalized by $T_0=1$, $T_1=X$, $T_{a+1}=2XT_a-T_{a-1}$, have integer coefficients and degree $a$, reduce modulo $p$ to polynomials in $\F_p[X]$, and satisfy the semiconjugacy
\[
T_a(\chi(u))=\chi(u^a)\quad(u\in\UU,\ a\ge0),
\qquad\qquad
T_{ab}=T_a\circ T_b .
\]
Following the twin-coset geometry of circle FFTs \cite{HLP24}, for a subgroup $H\le\UU$ of order $M$ and $g\in\UU$ with $g^2\notin H$ we call
\[
\mathcal D\ :=\ gH\ \cup\ g^{-1}H\ \subseteq\ \UU
\]
a \emph{twin coset}. It has size $2M$, is closed under inversion, and contains no self-inverse element: if $u\in gH$ satisfied $u=u^{-1}$, then $u\in gH\cap(gH)^{-1}=gH\cap g^{-1}H$, forcing $g^2\in H$; likewise for $u\in g^{-1}H$. Consequently $\chi$ is exactly two-to-one on $\mathcal D$, and the $x$-coordinate domain $D:=\chi(\mathcal D)\subseteq\F_p$ has size exactly $M$. For $a\mid M$ write $H^{(a)}:=\{h^a:h\in H\}$ for the image subgroup, of order $M/a$, and $\mathcal D^{(a)}:=\{u^a:u\in\mathcal D\}$.

\begin{lemma}[power fibers on torus twin cosets]\label{lem:torus-fibers}
Let $\mathcal D=gH\cup g^{-1}H$ be a twin coset of size $2M$ and let $a\mid M$.
\begin{enumerate}[label=\textup{(\alph*)}]
\item If $g^{2a}\notin H^{(a)}$, then $\mathcal D$ is $(X^a,a)$-smooth \textup(over any field containing $\F_{p^2}$\textup), and $\mathcal D^{(a)}=g^aH^{(a)}\sqcup g^{-a}H^{(a)}$ is again a twin coset, of size $2M/a$.
\item If $g^{2a}\in H^{(a)}$, then $\mathcal D$ is $(X^a,2a)$-smooth, with $\mathcal D^{(a)}=g^aH^{(a)}$ a single coset of size $M/a$.
\end{enumerate}
Moreover, $g^{2a}\notin H^{(a)}$ holds for every $a\mid M$ simultaneously if and only if $g^{2M}\ne1$, i.e.\ $\operatorname{ord}(g)\nmid2M$.
\end{lemma}

\begin{proof}
Since $a\mid M\mid p+1$ and $\UU$ is cyclic of order $p+1$, the kernel $K_a$ of $u\mapsto u^a$ on $\UU$ has order $a$ and, being the unique subgroup of $\UU$ of that order, is contained in $H$. Hence the $a$-th power map sends the coset $gH$ onto $g^aH^{(a)}$ with every fiber a coset of $K_a$, i.e.\ exactly $a$-to-one, and likewise sends $g^{-1}H$ onto $g^{-a}H^{(a)}$. The two image cosets of the subgroup $H^{(a)}$ are disjoint or equal according to whether $g^{2a}\notin H^{(a)}$ or $g^{2a}\in H^{(a)}$. In the disjoint case, every $v\in\mathcal D^{(a)}$ lies in exactly one of the two image cosets and has exactly $a$ preimages in $\mathcal D$; the twin-coset condition for $\mathcal D^{(a)}$, namely $(g^a)^2\notin H^{(a)}$, is the hypothesis itself. In the coincident case, every $v\in g^aH^{(a)}$ has exactly $a$ preimages in $gH$ and exactly $a$ in $g^{-1}H$, hence exactly $2a$ in $\mathcal D$. For the last claim: at $a=M$ we have $H^{(M)}=\{1\}$, so the scale-$M$ condition reads $g^{2M}\ne1$; conversely, if $g^{2a}=h^a$ for some $a\mid M$ and $h\in H$, then $g^{2M}=(g^{2a})^{M/a}=h^{M}=1$.
\end{proof}

\begin{lemma}[exact Chebyshev fibers on $x$-coordinate twin-coset domains]\label{lem:cheb-fibers}
Let $\mathcal D=gH\cup g^{-1}H$ be a twin coset of size $2M$, let $D:=\chi(\mathcal D)\subseteq\F_p$, $|D|=M$, and let $a\mid M$ satisfy $g^{2a}\notin H^{(a)}$. Then $D$ is $(T_a,a)$-smooth over $\F_p$, and
\[
T_a(D)\ =\ \chi\bigl(\mathcal D^{(a)}\bigr)
\]
is the $x$-coordinate domain of the twin coset $\mathcal D^{(a)}$, of size $M/a$. In particular, if $\operatorname{ord}(g)\nmid2M$, this holds at every scale $a\mid M$ simultaneously, and for $M=2^m$ the tower
\[
D\ \xrightarrow{\ T_2\ }\ T_2(D)\ \xrightarrow{\ T_2\ }\ T_4(D)\ \xrightarrow{\ T_2\ }\ \cdots
\]
consists of $x$-coordinate twin-coset domains of sizes $M,M/2,M/4,\dots$ --- the line-round tower of circle FRI \cite{HLP24}.
\end{lemma}

\begin{proof}
The semiconjugacy gives $T_a(D)=T_a(\chi(\mathcal D))=\chi(\mathcal D^{(a)})$; by \cref{lem:torus-fibers}(a), $\mathcal D^{(a)}$ is a twin coset, so $\chi$ is exactly two-to-one on it and $|T_a(D)|=|\mathcal D^{(a)}|/2=M/a$.

Fix $w\in T_a(D)$ and write $w=\chi(v)$ with $v\in\mathcal D^{(a)}=g^aH^{(a)}\sqcup g^{-a}H^{(a)}$. Since $\bigl(g^aH^{(a)}\bigr)^{-1}=g^{-a}H^{(a)}$, the set $\mathcal D^{(a)}$ is closed under inversion, and after replacing $v$ by $v^{-1}$ if necessary (which does not change $w$) we may assume $v\in g^aH^{(a)}$. Then $v\ne v^{-1}$: otherwise $v$ would lie in both of the disjoint cosets $g^aH^{(a)}$ and $g^{-a}H^{(a)}$. For $x=\chi(u)\in D$ with $u\in\mathcal D$,
\[
T_a(x)=w
\iff
\chi(u^a)=\chi(v)
\iff
u^a\in\{v,v^{-1}\}.
\]
By the fiber structure in the proof of \cref{lem:torus-fibers}, the equation $u^a=v$ has exactly $a$ solutions $u\in gH$ and none in $g^{-1}H$ (whose $a$-th powers fill the disjoint coset $g^{-a}H^{(a)}$), and the equation $u^a=v^{-1}$ has exactly $a$ solutions in $g^{-1}H$ --- the inverses of the former --- and none in $gH$. Hence
\[
E_w\ :=\ \{u\in\mathcal D:\ u^a\in\{v,v^{-1}\}\}
\]
has exactly $2a$ elements, is closed under inversion, and, being a subset of $\mathcal D$, contains no self-inverse element. Since $\chi(u)=\chi(u')$ forces $u'\in\{u,u^{-1}\}$, the image $\chi(E_w)$ --- which is precisely the fiber $T_a^{-1}(w)\cap D$ --- has exactly $|E_w|/2=a$ elements. As $w\in T_a(D)$ was arbitrary, $D$ is $(T_a,a)$-smooth with $|T_a(D)|=M/a=|D|/a$, as required. The tower statement follows from the all-scales equivalence in \cref{lem:torus-fibers} together with $T_{2^{j+1}}=T_2\circ T_{2^j}$.
\end{proof}

\begin{remark}[standard-position cosets are twin cosets with $\operatorname{ord}(g)\nmid2M$]\label{rem:standard-position}
A standard-position coset of circle FFT \cite{HLP24} is a coset $gH_0$, $H_0:=\langle g^2\rangle$, of size $2M$ with $\operatorname{ord}(g)=4M$ (so $|H_0|=2M$). Let $K\le H_0$ be the index-two subgroup, of order $M$. Then $g^{-1}K=g\,g^{-2}K$, and $g^{-2}$ generates $H_0$, hence $g^{-2}\notin K$, so
\[
gH_0\ =\ gK\ \sqcup\ g^{-1}K
\]
is exactly the twin coset attached to $(K,g)$: the condition $g^2\notin K$ holds because $g^2$ generates $H_0\supsetneq K$, and $\operatorname{ord}(g)=4M\nmid2M$. Thus every standard-position domain of \cite{HLP24}, at every folding level, satisfies the hypotheses of \cref{lem:torus-fibers,lem:cheb-fibers} at every dyadic scale simultaneously.
\end{remark}

\begin{lemma}[diagonal invariance of lists and correlated agreement]\label{lem:diag-invariance}
Let $C\subseteq\F^n$ be a linear code and $t\in(\F^{\times})^n$; write $t\circ v$ for the coordinatewise product and $t\circ C:=\{t\circ c:c\in C\}$. Then for every received word $U$, every radius $\delta$, and every pair of radii,
\[
\Lst(t\circ C,\delta,t\circ U)=t\circ\Lst(C,\delta,U),
\qquad
\eca(t\circ C,\cdot,\cdot)=\eca(C,\cdot,\cdot),
\qquad
\emca(t\circ C,\cdot)=\emca(C,\cdot).
\]
\end{lemma}

\begin{proof}
The map $v\mapsto t\circ v$ is an $\F$-linear bijection of $\F^n$ that preserves coordinatewise agreement: $(t\circ v)_j=(t\circ w)_j\iff v_j=w_j$. Hence it preserves $\dist$ and every restricted distance $\distS$, maps $C$ onto $t\circ C$, acts row-wise on pairs, and commutes with the line structure: $t\circ(f_1+\gamma f_2)=(t\circ f_1)+\gamma\,(t\circ f_2)$. It therefore carries every witness of the defining events of $\Lst$, $\eca$, and $\emca$ for $C$ bijectively onto a witness for $t\circ C$ with the same slope $\gamma$ and the same agreement supports, and conversely via $t^{-1}$.
\end{proof}

\begin{lemma}[torus uniformization of circle codes]\label{lem:circle-rs}
Let $p\equiv3\pmod4$, let $\F\supseteq\F_{p^2}$ be a finite field \textup(so $i\in\F$\textup), let $E\subseteq\mathcal C(\F_p)$ be a set of circle points, and let $E':=\iota(E)\subseteq\UU\subseteq\F_{p^2}^{\times}$. For $w\ge0$ define the $\F$-coefficient circle code
\[
\mathcal C_w(\F,E)\ :=\ \bigl\{(f(P))_{P\in E}\ :\ f\in\F[x,y]/(x^2+y^2-1),\ \deg f\le w\bigr\}.
\]
Then, under the coordinate bijection $\iota:E\to E'$ and with the diagonal vector $t:=(u^{-w})_{u\in E'}$,
\[
\mathcal C_w(\F,E)\ =\ t\circ\RS[\F,E',2w+1].
\]
Consequently, by \cref{lem:diag-invariance}, the circle code and the Reed--Solomon code $\RS[\F,E',2w+1]$ of odd dimension $2w+1$ on the torus domain $E'$ have identical list sizes at every radius and identical $\eca$ and $\emca$ at every (pair of) radii.
\end{lemma}

\begin{proof}
Since $i\in\F$ and $x+iy$ is invertible in the quotient ring (with inverse $x-iy$, as $(x+iy)(x-iy)=x^2+y^2=1$), the substitution $u:=x+iy$ defines a ring isomorphism
\[
\F[x,y]/(x^2+y^2-1)\ \xrightarrow{\ \sim\ }\ \F[u,u^{-1}],
\qquad
x\mapsto\frac{u+u^{-1}}2,\quad y\mapsto\frac{u-u^{-1}}{2i},
\]
with inverse $u\mapsto x+iy$; it matches evaluation at $P\in E$ on the left with evaluation at $\iota(P)\in E'$ on the right. Every class of degree $\le w$ has the canonical form $f_0(x)+y\,f_1(x)$ with $\deg f_0\le w$ and $\deg f_1\le w-1$ (reduce $y^2=1-x^2$), and $\{h_0(x)+y\,h_1(x)\}$ is a free-module canonical form of the quotient ring over $\F[x]$, so this space has dimension exactly $2w+1$. Under the isomorphism,
\[
x^j=2^{-j}u^{-j}(u^2+1)^j\in u^{-j}\,\F[u]_{\le2j},
\qquad
y\,x^j\in u^{-(j+1)}\,\F[u]_{\le2j+2},
\]
so the degree-$\le w$ subspace maps into $u^{-w}\F[u]_{\le2w}$; both spaces have dimension $2w+1$ and the map is injective, hence the image is exactly $u^{-w}\F[u]_{\le2w}$. Evaluating on $E'$ gives
\[
\mathcal C_w(\F,E)
=\bigl\{(u^{-w}h(u))_{u\in E'}\ :\ h\in\F[u],\ \deg h\le2w\bigr\}
=t\circ\RS[\F,E',2w+1].
\qedhere
\]
\end{proof}

Note the parity: the uniformized dimension $k_c=2w+1$ is always odd, while the useful folding scales $a$ on a dyadic torus twin coset are even, so $a\nmid k_c$ \emph{always}. The divisibility-free form of \cref{lem:phi-fiber} is therefore not a luxury here: the original \cref{lem:fiber} can never be applied to a uniformized circle code directly.

\subsection{Universal caps for circle codes and circle-FRI line rounds}
\label{sec:circle-caps}

\begin{corollary}[universal circle-row cap]\label{cor:circle-grand}
Let $p\equiv3\pmod4$ be prime, let $\mathcal D=gH\cup g^{-1}H\subseteq\UU$ be a twin coset of size $2M$ with $\operatorname{ord}(g)\nmid2M$ \textup(e.g.\ any standard-position coset of \cite{HLP24}, by \cref{rem:standard-position}\textup), and let $\F$ be any finite field with $q:=|\F|<2^{256}$.
\begin{enumerate}[label=\textup{(\alph*)}]
\item \textup{(Line rounds.)} Assume $\F\supseteq\F_p$. Let $D:=\chi(\mathcal D)$, $n:=M$, and $C:=\RS[\F,D,k]$ with $k\le2^{40}$ and $\rho:=k/n\in[\tfrac1{16},\tfrac12]$. If $\rho\ge\frac18$, assume $2^{10}\mid M$ and set $N:=2^{10}$; if $\rho<\frac18$, assume $2^{11}\mid M$ and set $N:=2^{11}$. Then
\[
\emca(C,\delta)\;>\;\frac1{2k}\Bigl(1-\frac nq\Bigr)\;\ge\;2^{-86}\;\gg\;2^{-128}
\qquad\text{for every }\ \delta\in\Bigl[1-\rho-\frac1N,\ 1-\rho\Bigr),
\]
and the same lower bound holds for $\eca(C,\delta)$ on the integer-admissible subinterval $[1-\frac{A_2}n,\,1-\rho-\frac1n]$, with $A_2$ as in \cref{lem:phi-fiber} for $a=n/N$. If moreover $\frac nN\mid k$ --- as in every $2$-power instantiation at the four challenge rates --- the MCA band extends to $[1-\rho-\frac2N,\,1-\rho)$. In particular
\[
\delta^*_C(2^{-128})\le1-\rho-2^{-10}\ \ (\rho\ge\tfrac18),
\qquad
\delta^*_C(2^{-128})\le1-\rho-2^{-11}\ \ (\rho<\tfrac18),
\]
improving to $1-\rho-2^{-9}$ \textup(resp.\ $1-\rho-2^{-10}$\textup) under $\frac nN\mid k$.
\item \textup{(Circle codes.)} Assume $\F\supseteq\F_{p^2}$. Let $C:=\mathcal C_w(\F,\iota^{-1}(\mathcal D))$ with $k_c:=2w+1\le2^{41}$, $n_c:=2M$, and $\rho_c:=k_c/n_c\in[\tfrac1{16},\tfrac12]$. If $\rho_c\ge\frac18$, assume $2^{9}\mid M$ and set $N:=2^{10}$; if $\rho_c<\frac18$, assume $2^{10}\mid M$ and set $N:=2^{11}$. Then
\[
\emca(C,\delta)\;>\;\frac1{2k_c}\Bigl(1-\frac{n_c}q\Bigr)\;\ge\;2^{-88}\;\gg\;2^{-128}
\qquad\text{for every }\ \delta\in\Bigl[1-\rho_c-\frac1N,\ 1-\rho_c\Bigr),
\]
with the corresponding $\eca$ bound on $[1-\frac{A_2}{n_c},\,1-\rho_c-\frac1{n_c}]$ \textup(here $a:=n_c/N\nmid k_c$ always, since $k_c$ is odd\textup). In particular $\delta^*_C(2^{-128})\le1-\rho_c-2^{-10}$ for $\rho_c\ge\frac18$, and $\delta^*_C(2^{-128})\le1-\rho_c-2^{-11}$ for $\rho_c<\frac18$.
\end{enumerate}
\end{corollary}

\begin{proof}
(a) Put $a:=n/N=M/N$, so $a\mid M$, and $g^{2a}\notin H^{(a)}$ by the all-scales equivalence of \cref{lem:torus-fibers}; \cref{lem:cheb-fibers} then makes $D$ a $(T_a,a)$-smooth domain over $\B:=\F_p\subseteq\F$. We check the hypotheses of \cref{thm:phi-cap}. First, $n<q$: since $M\mid p+1$ and $\operatorname{ord}(g)\mid p+1$, the case $M=p+1$ would give $\operatorname{ord}(g)\mid p+1\mid2M$, contradicting $\operatorname{ord}(g)\nmid2M$; hence $M\le\frac{p+1}2$ and $n\le\frac{p+1}2<p\le q$. Second, $\ell_2\le\rho N+2\le N/2+2\le N-1$. Third, for \eqref{eq:hyp-phi}: $\ell_2/N\in(\rho+\frac1N,\rho+\frac2N]$, which is contained in $[\frac18,\,0.502]$ when $\rho\ge\frac18$ and in $[\frac1{16},\,0.126]$ when $\rho<\frac18$; on each interval the concave $\Hb$ is minimized at an endpoint, and both right endpoints have larger $\Hb$-value than the left, so $\Hb(\ell_2/N)\ge\Hb(\frac18)\ge0.5435$, resp.\ $\Hb(\ell_2/N)\ge\Hb(\frac1{16})\ge0.3372$. Hence
\[
\log_2\binom N{\ell_2}\ \ge\ N\,\Hb(\ell_2/N)-\log_2(N+1)\ \ge\ 546,\ \ \text{resp.}\ \ 679,
\]
while $|\B|(q/k+1)\le q(q/k+1)\le q^2+q<2^{513}$. The error numerics are those of \cref{cor:grand}: $\rho\ge\frac1{16}$ gives $n\le16k\le2^{44}$, so $\frac1{2k}(1-\frac nq)\ge\frac1{2k(n+1)}>2^{-86}$. \Cref{thm:phi-cap} now yields the $\eca$ and $\emca$ lower bounds on $[1-\frac{A_2}n,1-\rho-\frac1n]$ and $[1-\frac{A_2}n,1-\rho)$ respectively; since $A_2\ge k+a+1$, the MCA interval contains $[1-\rho-\frac1N,1-\rho)$, and when $\frac nN\mid k$ one has $A_2=k+2a$, so it equals $[1-\rho-\frac2N,1-\rho)$.

(b) By \cref{lem:circle-rs}, all list sizes, $\eca$, and $\emca$ of $C$ coincide with those of $C':=\RS[\F,\mathcal D,k_c]$ on the torus twin coset $\mathcal D\subseteq\F_{p^2}$. Put $a:=n_c/N=2M/N$; then $a\mid M$ (indeed $M/a=N/2\in\mathbb Z$), and $g^{2a}\notin H^{(a)}$ as before, so $\mathcal D$ is $(X^a,a)$-smooth over $\B:=\F_{p^2}\subseteq\F$ by \cref{lem:torus-fibers}(a). Also $n_c=2M<p+1\le p^2\le q$: the case $2M=p+1$ would again give $\operatorname{ord}(g)\mid p+1=2M$. The verification of \eqref{eq:hyp-phi} is identical to (a), with the same two entropy cases and $|\B|=p^2\le q$; the error numerics use $n_c\le16k_c\le2^{45}$, so $\frac1{2k_c}(1-\frac{n_c}q)\ge\frac1{2k_c(n_c+1)}>2^{-88}$. \Cref{thm:phi-cap} applied to $C'$, transported back through \cref{lem:circle-rs,lem:diag-invariance}, gives the stated bands; $A_2\ge k_c+a+1$ again yields the containment of $[1-\rho_c-\frac1N,1-\rho_c)$.
\end{proof}

\begin{corollary}[deployed circle parameters: Mersenne-31 with QM31]\label{cor:circle-deployed}
Let $p=2^{31}-1$, let $\F=\F_{p^4}$ be the deployed QM31 extension \textup(so $q=p^4>2^{123.9}$ and $\F\supseteq\F_{p^2}$\textup), and let $\mathcal D$ be a standard-position twin coset with $M=2^{21}$ \textup(so $\operatorname{ord}(g)=2^{23}\nmid2^{22}=2M$\textup).
\begin{enumerate}[label=\textup{(\alph*)}]
\item \textup{(Line round, rate $\frac12$.)} With $D=\chi(\mathcal D)$, $|D|=2^{21}$, and $C=\RS[\F_{p^4},D,2^{20}]$,
\[
\emca(C,\delta)\;>\;\bigl(1-2^{-102}\bigr)\,2^{-21}\;>\;2^{-22}
\qquad\text{for every }\ \delta\in\Bigl[\frac12-2^{-7},\ \frac12\Bigr),
\]
and the same lower bound holds for $\eca(C,\delta)$ for every $\delta\in[\frac12-2^{-7},\,\frac12-2^{-21}]$.
\item \textup{(Circle code.)} With $C=\mathcal C_w(\F_{p^4},\iota^{-1}(\mathcal D))$, $w=2^{20}$, $k_c=2^{21}+1$, $n_c=2^{22}$, $\rho_c=\frac12+2^{-22}$,
\[
\emca(C,\delta)\;>\;2^{-23}
\qquad\text{for every }\ \delta\in\Bigl[\frac12-2^{-7},\ 1-\rho_c\Bigr),
\]
and the same lower bound holds for $\eca(C,\delta)$ for every $\delta\in[\frac12-2^{-7},\,\frac12-2^{-21}]$.
\end{enumerate}
\end{corollary}

\begin{proof}
Both parts use $N=256$ and, as in \cref{cor:deployed}, $\log_2\binom{256}{130}\ge247$.

(a) Here $a=n/N=2^{13}$ divides $k=2^{20}$, so $\ell_2=\rho N+2=130$, $A_2=k+2a=2^{20}+2^{14}$, and $1-A_2/n=\frac12-2^{-7}$. The standard-position hypothesis gives exact Chebyshev fibers at scale $a$ (\cref{rem:standard-position,lem:cheb-fibers}), with $\B=\F_p$. Hypothesis \eqref{eq:hyp-phi}: $|\B|(q/k+1)=p\,(p^4/2^{20}+1)<2^{31}\cdot2^{105}=2^{136}\le2^{247}\le\binom{256}{130}$. The error bound is $\frac1{2k}(1-n/q)$ with $n/q=2^{21}/p^4<2^{-102}$. \Cref{thm:phi-cap} gives the MCA band $[\frac12-2^{-7},\frac12)$ and the $\eca$ band $[\frac12-2^{-7},\frac12-2^{-21}]$, since $1-\rho-\frac1n=\frac12-2^{-21}$.

(b) Here $a=n_c/N=2^{14}\mid M=2^{21}$, $\ell_2=\lfloor(2^{21}+1)/2^{14}\rfloor+2=2^7+2=130$, $A_2=2^{14}\cdot130=2^{21}+2^{15}$, and $1-A_2/n_c=\frac12-2^{-7}$. \Cref{lem:circle-rs} transports everything to $C'=\RS[\F_{p^4},\mathcal D,2^{21}+1]$ on the torus twin coset, which is $(X^a,a)$-smooth over $\B=\F_{p^2}$ by \cref{lem:torus-fibers,rem:standard-position}. Hypothesis \eqref{eq:hyp-phi}: $|\B|(q/k_c+1)=p^2\,\bigl(p^4/(2^{21}+1)+1\bigr)<2^{62}\cdot2^{104}=2^{166}\le2^{247}$. The error bound: $n_c/q=2^{22}/p^4<2^{-101}$, so
\[
\frac1{2k_c}\Bigl(1-\frac{n_c}q\Bigr)\;>\;\frac{1-2^{-101}}{2^{22}+2}\;>\;2^{-22}\bigl(1-2^{-21}\bigr)\bigl(1-2^{-101}\bigr)\;>\;2^{-23}.
\]
\Cref{thm:phi-cap} gives the MCA band $[\frac12-2^{-7},\,1-\rho_c)$ and the $\eca$ band $[\frac12-2^{-7},\,1-\rho_c-\frac1{n_c}]=[\frac12-2^{-7},\,\frac12-2^{-21}]$.
\end{proof}

\begin{corollary}[circle certifying lines are genuinely extension-valued]\label{cor:circle-Fvalued}
The subfield-confinement lemma \cref{lem:confine} holds verbatim for an arbitrary evaluation set $D\subseteq\B$: its proof uses only that the evaluation points lie in $\B$, not that they are invertible. Consequently, at the parameters of \cref{cor:circle-deployed}:
\begin{enumerate}[label=\textup{(\alph*)}]
\item any pair of words on $D=\chi(\mathcal D)$ that is $\F_p$-valued up to a common scalar factor has MCA-bad-slope density at most $p/q=p^{-3}<2^{-92}$ at every radius; in particular, every pair attaining the bound $>2^{-22}$ of \cref{cor:circle-deployed}\textup{(a)} is genuinely $\F_{p^4}$-valued;
\item the diagonal vector $t=(u^{-w})_{u}$ of \cref{lem:circle-rs} is $\F_{p^2}$-valued, so any pair of circle words whose $\iota$-transports are $\F_{p^2}$-valued up to a common scalar factor has MCA-bad-slope density at most $p^2/q=p^{-2}<2^{-61}$ at every radius; in particular, every circle pair attaining the bound $>2^{-23}$ of \cref{cor:circle-deployed}\textup{(b)} contains a word whose $\iota$-transport is not $\F_{p^2}$-valued, even after removing a common scalar.
\end{enumerate}
As in \cref{rem:extension-pole-status}, the extension-pole floors \cref{cor:extension-pole-deep-list-floor,cor:extension-pole-quotient-remainder-floor} supply the matching constructive mechanism: the fiber received words of \cref{lem:phi-fiber} on circle domains are $\B$-valued with $\B$-valued lists, and any such list converts through a simple pole $\alpha\in\F\setminus\B$ into genuinely extension-valued bad lines, with denominator $q-b$ ($b=|\B|\in\{p,p^2\}$) in place of $q-n$.
\end{corollary}

\begin{proof}
Inspecting the proof of \cref{lem:confine}: it decomposes words and codewords over a $\B$-basis of $\F$ and evaluates polynomials at the points of $D\subseteq\B$; no step uses $D\subseteq\B^{\times}$, so the two-radius statement holds for arbitrary $D\subseteq\B$, in particular for the Chebyshev domain $\chi(\mathcal D)\subseteq\F_p$ and the torus domain $\mathcal D\subseteq\F_{p^2}$. (a) is then \cref{lem:confine} with $\B=\F_p$ inside $q=p^4$: the bad-slope density of a $\B$-confined pair is at most $|\B|/q=p^{-3}<2^{-92}<2^{-22}$. (b) A circle pair whose $\iota$-transports are $\F_{p^2}$-valued up to a common scalar corresponds, after multiplying by the $\F_{p^2}$-valued diagonal $t$, to an $\F_{p^2}$-confined pair for $\RS[\F_{p^4},\mathcal D,k_c]$; \cref{lem:diag-invariance} preserves the bad-slope set, and \cref{lem:confine} with $\B=\F_{p^2}$ bounds its density by $p^2/q=p^{-2}<2^{-61}<2^{-23}$. The final claim is \cref{cor:extension-pole-deep-list-floor} applied with $\B\in\{\F_p,\F_{p^2}\}$ to the $\B$-valued received words and lists of \cref{lem:phi-fiber}.
\end{proof}

\begin{remark}[interleaving, slacked fallback, and the two-tier picture for circle rows]\label{rem:circle-transfer}
The interleaving-transfer lemma \cref{lem:inter} is stated for arbitrary linear codes, so every bound of \cref{cor:circle-grand,cor:circle-deployed} transfers verbatim to $s$-interleaved circle rows, exactly as in \cref{cor:rows}. The slacked BCHKS/ABF fallback also carries over: \cref{lem:phi-fiber}(i) supplies a list of size $\ge\binom N{\ell_1}/|\B|$ in $\RS[\F,D,k]$ at radius $1-\frac{A_1}n$, and whenever $A_1\ge k+3$ \textup(automatic in the $2$-power line rounds with $a\ge4$, where $A_1=k+a$\textup) and $\binom N{\ell_1}\ge|\B|\,q$, the contrapositive of \cref{thm:B} gives $\eca\bigl(C,\,1-\frac{A_1}n+\frac2n,\,1-\rho-\frac1n\bigr)\ge\frac1{2n}$, as in \cref{prop:slacked}. Finally, the two-tier summary of \cref{rem:twotier} acquires a circle column: the deployed Mersenne-$31$ circle family carries the same kind of near-capacity negative band as the multiplicative rows --- gap $2^{-7}$ at the deployed sizes, $2^{-9}$--$2^{-11}$ universally --- over every challenge field $q<2^{256}$ containing the relevant base field.
\end{remark}

\subsection{The lacunarity obstruction for genus-one (ECFFT) domains}
\label{sec:ecfft}

\begin{remark}[what blocks the elliptic case]\label{rem:ecfft-obstruction}
The proof of \cref{lem:phi-fiber} works because $\varphi$ is a \emph{polynomial}: in the expansion $\prod_{b\in A}(\varphi-b)=\sum_j(-1)^je_j(A)\,\varphi^{\ell-j}$, consecutive terms drop by a full $a=\deg\varphi$ in degree, so the received word occupies the top two degree slots and everything else is a codeword sitting $2a$ degrees down --- gap $\approx2/N$. Consider instead a rational folding map $\psi=f/g$ of degree $a$, with $e:=\deg g<\deg f=a$ and $g$ nonvanishing on $D$, whose fibers inside $D$ are complete and of uniform size $a$. The natural locators are $\Lambda_A:=\prod_{b\in A}(f-b\,g)$, of degree $a\ell$, vanishing on the $\psi$-fibers over $A$; their expansion $\sum_j(-1)^je_j(A)\,f^{\ell-j}g^{\,j}$ has $j$-th term of degree $a\ell-j(a-e)$, so the slack-two window has width $2(a-e)$ in degree and the certified gap is $2(a-e)/n$. Polynomial maps have $e=0$ and recover gap $2/N$. The FFTree domains of ECFFT \cite{BCKL21} fold by $x$-coordinate maps of $2$-isogenies of elliptic curves, which are rational of degree type $(a,e)=(2,1)$; every composition of such maps is again the $x$-map of an isogeny, of type $(a,a-1)$. For these, $a-e=1$ at every scale: the window collapses to width $2$, and the construction certifies only the first staircase step --- a large list at agreement $k+2$, hence an $\eca$ floor at gap $2/n$ --- a nonvacuous but microscopic band, far from the $2^{-9}$--$2^{-11}$ bands above. The circle escapes this because its folding tower becomes polynomial after the $x$-projection: $T_2$ fixes $\infty$ and is totally ramified there. The elliptic $x$-maps also fix $\infty$ but are \emph{not} totally ramified there ($\psi^{-1}(\infty)=\{\infty,x(T)\}$ for a $2$-isogeny with kernel generator $T$), and the locator expansion sees exactly this defect as $e=a-1$. Obtaining a macroscopic universal cap for ECFFT rows therefore requires either a lacunary basis of isogeny locators, a different received-word family, or a matching safe-side theorem showing the collapse is essential. The second exit is realized by the rational graded floor of \cref{prop:graded-rational-floor} \textup(\cref{cor:ecfft-macroscopic}\textup); the residual band form is recorded as the genus-one/circle transport question \cref{prob:isogeny}.
\end{remark}


\section{Threshold formulation and certificate interface}\label{sec:prize-program}

The Proximity Prize asks for a threshold, not merely for a near-capacity
obstruction.  We therefore use the integer staircase below to state both sides
of the result: lower certificates mark agreements where MCA or list decoding is
unsafe, while upper certificates mark agreements where the corresponding error
is below the target.  The paper does not determine the staircase to one step in
the remaining band, but it reduces the unknown part to a single aperiodic
correlated-agreement problem and gives exact finite certificates for all
printed edges.

\subsection{Established inputs}

For the grand MCA challenge at target error $\varepsilon^*=2^{-128}$, the
unconditional output of \cref{cor:grand} is the following negative certificate:
for every challenge-envelope row satisfying the printed divisor hypothesis,
\[
        \delta\ge 1-\rho-2^{-9}
        \quad(\rho\in\{1/2,1/4,1/8\})
\]
and
\[
        \delta\ge 1-\rho-2^{-10}
        \quad(\rho=1/16)
\]
are impossible MCA radii.  Equivalently, any full determination of
$\delta_C^*(2^{-128})$ for these smooth multiplicative rows must search below
these gaps.  This conclusion uses only the self-contained deep-point conversion,
the quotient-fiber pigeonhole, and support-wise MCA monotonicity; it no longer
rests on importing the Crites--Stewart conversion.

The theorem also distinguishes two different grades of information.  First, the
\emph{threshold grade} is already enough for the prize inequality: the certified
error is $\gg2^{-128}$.  Second, the \emph{failure-profile grade} is still weak:
the present proof gives error of order $1/k$, not error $1-o(1)$.  Closing the
error-one profile at fixed gaps such as $1/64$ is a different, stronger problem.
It would improve our understanding of the failure landscape, but is not needed
to obtain the stated upper bound on $\delta_C^*(2^{-128})$.

The negative certificate also covers a wider row family.  The
same conclusion --- an unsafe near-capacity band at error $\gg2^{-128}$ ---
now holds for mixed-radix multiplicative rows without the integrality
condition $\rho N\in\mathbb Z$ (\cref{cor:mixed-radix}, gaps
$2^{-11}$/$2^{-12}$), for the $x$-coordinate line rounds of circle FRI over
any $p\equiv3\pmod4$, and for the bivariate circle codes over any challenge
field containing $i$ (\cref{cor:circle-grand}, gaps $2^{-9}$--$2^{-11}$),
including the deployed Mersenne-$31$/QM31 family at gap $2^{-7}$
(\cref{cor:circle-deployed}).  These additions use the same three
ingredients; the only new inputs are the exact Chebyshev-fiber geometry of
twin cosets and the torus uniformization, both elementary.

The treatment in \cref{sec:discharge} adds the safe side. \Cref{thm:deep-mca} certifies $\emca(C,\delta)\le(\lfloor\delta n\rfloor+1)/q$ whenever $3\lfloor\delta n\rfloor\le n-k$, for every row at once, and \cref{thm:johnson-list} certifies the interleaved-list challenge up to the Johnson radius uniformly in the arity; \cref{tab:status} records the per-obligation outcome, and \cref{thm:sandwich} assembles the resulting two-sided sandwich.

\subsection{Integer-staircase formulation}

Let
\[
        a(\delta):=\lceil(1-\delta)n\rceil,
        \qquad r(\delta):=n-a(\delta)
\]
be the agreement and integer-radius forms of the same radius.  For a fixed finite-slope line
field $F_{\rm line}$, write $q_{\rm line}=|F_{\rm line}|$ and let
\[
        B_{\rm mca}(a)
\]
be the exact maximum number of MCA-bad line parameters at agreement at least
$a$, for the line family used in the challenge.  Then
\[
        \emca(C,1-a/n)=B_{\rm mca}(a)/q_{\rm line}
\]
under the closed finite-slope convention.  For projective-slope or curve samplers, the same staircase formulation applies after replacing $q_{\rm line}$ and $B_{\rm mca}$ by the corresponding sampler denominator and numerator.  Thus the MCA prize problem is an
integer staircase problem: find the first agreement $a$ for which
\[
        B_{\rm mca}(a)\le 2^{-128}q_{\rm line}.
\]
The present paper supplies lower bounds on $B_{\rm mca}(a)$ near capacity by
routing quotient-fiber list mass through deep points.

\begin{proposition}[one-step threshold certificate]\label{prop:onestep}
Fix a row $C$ and a denominator $q_{\rm line}$.  Suppose that for some agreement
$a_0$ one has a lower certificate
\[
        B_{\rm mca}(a_0)>2^{-128}q_{\rm line}
\]
and an upper theorem
\[
        B_{\rm mca}(a_0+1)\le2^{-128}q_{\rm line}.
\]
Then, under the closed integer-radius convention, radius $1-a_0/n$ is unsafe and
radius $1-(a_0+1)/n$ is safe.  As a real-radius supremum, the threshold is
$1-a_0/n$ and is not attained if the closed ball at that endpoint is counted.
\end{proposition}

\begin{proof}
The function $B_{\rm mca}(a)$ is nonincreasing in $a$, since raising the
required agreement only removes witnesses.  At the closed radius $1-a/n$, the
MCA error is $B_{\rm mca}(a)/q_{\rm line}$.  The two displayed inequalities are
therefore exactly the unsafe/safe transition at consecutive integer agreement
levels.
\end{proof}

\subsection{Conditional MCA one-step criterion}

The following theorem is a deliberately explicit proof blueprint.  Each
hypothesis is a concrete mathematical task, not a black-box ``large field''
assumption.

\begin{theorem}[conditional MCA threshold theorem]\label{thm:conditional-mca}
Fix a smooth multiplicative row $C=\RS[F,D,k]$, a line field $F_{\rm line}$, and
a target $\varepsilon^*=2^{-128}$.  Suppose the following four inputs are proved
for this row, uniformly in the received line:
\begin{enumerate}[label=(\roman*),leftmargin=2.2em]
\item \textbf{Quotient-profile lower and upper ledger.}  For every divisor
$c\mid n$ and every agreement $a=mc+s$ with $0\le s<c$, the quotient-periodic
bad slopes, including remainder supports consisting of $m$ full fibers plus
$s$ residual points, are counted with overlaps across divisors accounted for.
Write the resulting divisor-lattice upper numerator as $B_Q^{\rm all}(a)$;
individual divisor contributions $B_Q(a;c)$ may also be printed, but the upper
ledger must count their union or a certified safe sum rather than only their
maximum.
\item \textbf{Aperiodic residue-line packing.}  After removing tangent and
quotient-periodic locators, the Hankel/residue-line incidence variety has at
most $B_{\rm ap}(a)$ bad line parameters.
\item \textbf{Extension-line accounting.}  If $F_{\rm line}$ is a proper
extension of the field generated by $D$, either all genuinely extension-valued
line witnesses are included in $B_{\rm ap}(a)$, or a separate extension term
$B_{\rm ext}(a)$ is proved.
\item \textbf{No missing line family.}  The line family in the theorem is the
same finite-slope or projective-slope family used by the challenge or protocol
under consideration.
\end{enumerate}
Assume moreover that the resulting upper numerator
\[
        B^+(a):=B_{\rm tan}(a)+B_Q^{\rm all}(a)
                +B_{\rm ap}(a)+B_{\rm ext}(a)
\]
and a corresponding lower numerator $B^-(a)$ satisfy
\[
        B^-(a_0)>\varepsilon^* q_{\rm line},
        \qquad
        B^+(a_0+1)\le\varepsilon^* q_{\rm line}
\]
for some $a_0$.  Then $\delta_C^*(\varepsilon^*)$ is determined at the
one-step accuracy of \cref{prop:onestep}.
\end{theorem}

\begin{proof}
The four hypotheses classify every possible MCA-bad line parameter into one of
four buckets: tangent, quotient-periodic, aperiodic, or genuinely extension-line.
The quotient bucket is counted by the divisor-lattice union ledger $B_Q^{\rm all}$,
so overlaps or simultaneous quotient descriptions cannot invalidate the upper
bound.  The displayed $B^+(a)$ is therefore an upper bound for $B_{\rm mca}(a)$, while
$B^-(a)$ is a certified lower bound.  The conclusion is exactly
\cref{prop:onestep}.
\end{proof}

The present paper supplies part of the lower numerator $B^-(a)$ in the
near-capacity band: quotient-fiber list mass forces MCA mass above the printed
universal caps.  It does not prove the aperiodic upper input (ii), and therefore
does not by itself determine the exact staircase.

\subsection{Conditional interleaved-list one-step criterion}

The grand list-decoding challenge asks for the largest $\delta$ such that an
interleaved list, divided by the field used for the relevant challenge or list
ledger, is below $2^{-128}$.  In the simplest same-field form this is
\[
        |\Lambda(C^{\equiv m},\delta)|\le 2^{-128}|F| .
\]
The quotient-fiber construction in this paper already gives lower floors for
ordinary lists, and ordinary lists embed into interleaved lists by placing the
same received word in one row and zero words in the others.  On the MCA side,
\cref{sec:aperiodic-hankel-certificates} supplies a certified eliminant route
for aperiodic Hankel packing, but it does not prove that the needed eliminants
always exist.  The missing positive list direction is an arbitrary-word
locator-fiber theorem.

\begin{theorem}[conditional interleaved-list threshold theorem]\label{thm:conditional-list}
Fix an interleaving arity $m$.  Suppose that for every received word $U$ and
agreement $a=k+\sigma$ one has a base-code list bound of the form
\[
        |\Lambda(C,1-a/n,U)|
        \le
        L_Q(a,U)+L_{\rm ap}(a),
\]
where $L_Q$ is an explicit quotient-profile list contribution and
$L_{\rm ap}$ is a uniform aperiodic locator-fiber bound.  Suppose also that the
same bound is available at the higher agreement levels that occur in the
codegree recursion, starting with $2a-k$.  Then the interleaved-list threshold is
obtained by applying the codegree recursion row by row and comparing the
resulting numerator to $2^{-128}|F|$.
\end{theorem}

\begin{proof}
For two rows, split the second-row list according to whether its agreement set
has size below or above $2a-k$.  Below $2a-k$, two first-row codewords would
agree on more than $k$ positions and hence are equal; above $2a-k$, use the
base-code list bound at agreement $2a-k$.  Iterating gives the stated $m$-row
recursion.  This is the same codegree mechanism used in the experimental ledger;
the only missing input is the arbitrary-word locator-fiber bound.
\end{proof}

\subsection{Residual ingredients for one-step threshold certificates}

The one-step criteria above consume the following residual ingredients, which we
record with their current status.  The lower-floor half of the periodic-support
ledger is supplied in scanner-ready form by
\cref{lem:heaviest-prefix-locator-floor,thm:quotient-remainder-deep-floor}: every
quotient-plus-remainder agreement $a=mc+s$ has a generated-field prefix-fiber
floor, an image-size fallback, an ambient-box fallback, and a converted CA/MCA
lower numerator.  The support-family upper half is exact at the divisor-lattice
level by \cref{thm:exact-divisor-block-support-ledger}: overlaps between quotient
descriptions at different divisor scales are counted by a common-block
coefficient extraction rather than by a raw safe sum.  The distinct-parameter
branch image is exact inside every declared quotient family by the lcm ledger of
\cref{thm:exact-quotient-image-lcm-ledger,cor:affine-line-lcm-ledger,cor:projective-line-lcm-ledger}; duplicate finite slopes, projective slopes, and
finite curve parameters are coalesced before the numerator is printed.  The
residual-band analysis is split into four certified pieces.  The regular overdetermined nonsingular branch
is canonical by \cref{thm:canonical-affine-rankdrop-ledger,cor:canonical-regular-closed-ball-hankel-packing,thm:canonical-curve-rankdrop-ledger}.  The underdetermined
and rank-drop singular branches are covered by the rank-pivot certificate theorem
\cref{thm:residual-rank-pivot-cover,thm:scanner-checkable-residual-aperiodic-ledger}.  The new compressed-residual layer adds the kernel-compressed eliminant construction of
\cref{thm:kernel-compressed-affine-eliminant,cor:kernel-compressed-projective-curve-eliminants} and the exact arbitrary-residual support-image lcm ledger of
\cref{thm:exact-arbitrary-residual-image-ledger,cor:residual-exact-image-after-paid-removal}.  Thus a residual chart must now be discharged by a small compressed
eliminant, by an exact residual image polynomial, by a Weil-restricted extension
projection bound, or by a named structural obstruction.  Extension-line
accounting is handled on the safe side by
\cref{thm:extension-line-dimension-degree-ledger}: after choosing a $B$-basis of
$F$, each extension chart prints a parameter-projection degree and dimension over
$B$, or an exact zero-dimensional point count.  The curve-sampler ledger is
upgraded for finite-parameter curves and projective lines in parallel: quotient branches have exact gcd/lcm
image ledgers, regular curve rank-drop branches have canonical maximal-minor gcd
ledgers, and residual curve branches have both coefficient-pivot rank charts and
kernel-compressed curve eliminants.  What remains to be settled is no longer a bookkeeping
ambiguity: prove that the compressed residual eliminants are small enough in the
near-capacity rows, or classify any large residual image as quotient,
tangent/common-line, base-field exceptional, genuinely extension-valued, or a new
obstruction.

\begin{enumerate}[label=\textbf{T\arabic*.},leftmargin=2.5em]
\item \textbf{Remainder quotient-profile theorem.}  Generalize the fiber lemma
from full quotient-fiber unions to all agreements $a=mc+s$.  The target is an
exact lower and upper quotient ledger over the divisor lattice of $D$.  The present quotient-ledger package supplies the scanner-ready lower half via heaviest prefix fibers and the
support-union upper half via the exact divisor-block coefficient formula; inside
a declared quotient family, the gcd/lcm image ledger also coalesces duplicate
finite, projective, and curve parameters.  What remains is the structural
exhaustion theorem saying that the declared quotient family captures all
quotient-periodic bad parameters in the intended regime, with the subfield
denominator $|B|$ printed separately from the line denominator $q$.

\item \textbf{Aperiodic Hankel-packing theorem.}  Use the syndrome/Hankel normal
form of \cite{Cho26b} to prove that, after quotient-periodic and tangent
locators are removed, the number of remaining bad line parameters is at most
$B_{\rm ap}(a)$.  The regular overdetermined nonsingular branch is now
canonical by \cref{thm:canonical-affine-rankdrop-ledger,cor:canonical-regular-closed-ball-hankel-packing}: a scanner takes the gcd of all nonzero maximal
Hankel minors and then an lcm over exact agreement levels.  The underdetermined
and rank-drop singular buckets now have the rank-pivot cover of
\cref{thm:residual-rank-pivot-cover}, projective-infinity and curve coefficient
pivots, and the kernel-compressed eliminant ideals of
\cref{thm:kernel-compressed-affine-eliminant,cor:kernel-compressed-projective-curve-eliminants}.  The remaining proof problem is constructive and structural:
produce small eliminants for these compressed residual charts, use the exact
residual image lcm of \cref{thm:exact-arbitrary-residual-image-ledger} when
finite enumeration is feasible, or classify every large residual component as
quotient, tangent/common-line, base-field exceptional, genuinely extension-valued,
or a new obstruction.

\item \textbf{Extension-line lift or counterexample.}  For $B\subset F$ and
$D\subset B$, either prove that genuinely $F$-valued lines are covered by the
aperiodic ledger with a stable numerator, or exhibit explicit extension-valued
counterexamples.  \Cref{cor:Fvalued} shows that deployed failures certified by
this paper must be genuinely extension-valued; \cref{cor:extension-pole-deep-list-floor,cor:extension-pole-quotient-remainder-floor}
provide an explicit lower-side mechanism from any large list by taking the pole
in $F\setminus B$.  On the safe side,
\cref{thm:extension-line-dimension-degree-ledger} gives the certificate format:
Weil-restrict each residual extension chart to $B$, print a parameter-projection
degree/dimension pair or an exact zero-dimensional point count, and divide by the
line field size.  The remaining structural problem is to prove that all genuinely
extension-valued residual charts have small projected dimension/degree or to
exhibit the chart that violates this expectation.

\item \textbf{Arbitrary-word locator-fiber theorem.}  Prove the base-code list
bound needed by \cref{thm:conditional-list}, first for monomial-prefix received
words and quotient-free dimensions, then for arbitrary received words with
quotient cores budgeted.

\item \textbf{Curve and projective-line ledgers.}  Extend the affine-line Hankel
classification to projective slopes and to finite-parameter power curves
$f_0+\gamma f_1+\cdots+\gamma^d f_d$.  The projective and curve ledger package gives exact gcd/lcm ledgers
for projective and finite-curve quotient branches, canonical rank-drop gcds for
regular overdetermined finite-curve buckets, projective-infinity residual charts,
coefficient-pivot curve rank charts, and kernel-compressed projective/curve
eliminant certificates.  The remaining work is to prove smallness or structural
classification for the residual projective and curve charts, including
exceptional vanishing-vector strata and genuinely extension-valued components.

\item \textbf{Certificate scanner.}  Implement a finite staircase scanner that,
for each challenge row, prints tangent, quotient, generated-field entropy,
aperiodic, extension, curve, and interleaved-list numerators at every agreement
$a$.  The desired output is an interval
\[
        a_{\rm unsafe}<a_C^*\le a_{\rm safe}
\]
that shrinks to one step when \cref{thm:conditional-mca} or
\cref{thm:conditional-list} applies.

\item \textbf{Formal finite core.}  Formalize the self-contained deep-point
conversion, the quotient-fiber pigeonhole, support-wise MCA monotonicity, and
the integer staircase lemma.  These are small enough for proof-assistant
checking and directly strengthen a formal certificate package.

\item \textbf{Circle-row completion and genus-one floors.}  Extend the
quotient-support and quotient-image upper ledgers (item~1) and the aperiodic
Hankel classification (item~2) from power locators on multiplicative divisor
lattices to Chebyshev locator towers on $x$-coordinate twin-coset domains,
where the divisor lattice of $D$ is replaced by the lattice of dyadic
Chebyshev scales; and settle the genus-one lower-floor question of
the genus-one/circle transport question \cref{prob:isogeny}, for which the lacunarity collapse of
\cref{rem:ecfft-obstruction} blocks the present route beyond the first
staircase step.
\end{enumerate}

\subsection{Scope of claims}

We state precisely what this paper establishes and what it does not. Established: a universal unsafe band at capacity for the grand MCA challenge, over smooth multiplicative rows in the stated field-size envelope (\cref{cor:grand,cor:mixed-radix}) and, by the later sections, over Chebyshev, circle, and genus-one rows as well (\cref{cor:circle-grand,cor:ecfft-macroscopic,cor:circle-anyfield}); an unconditional safe region below one third of the distance (\cref{thm:deep-mca}) with a list-side safe region up to the Johnson radius, uniform in the interleaving arity (\cref{thm:johnson-list}); a safe region up to one half of the distance modulo a single classical import (\cref{thm:mca-from-ca,cor:conditional-half}); and the resulting sandwich, which determines the grand threshold of every deployed row within a factor smaller than two modulo that import, and within a factor of three unconditionally (\cref{thm:sandwich,thm:sandwich2}, \cref{tab:scanner2}). Not established, and not claimed: the exact threshold, a positive theorem inside the remaining frontier band, protocol soundness, or a curve-MCA theorem. What separates the sandwich from a one-step determination is precisely the broad band formulation \cref{prob:band}.

\section{Structural ledgers, residual bands, and certificate machinery}\label{sec:discharge}

This section collects the technical infrastructure used by the two-sided threshold statements, organized by mathematical function.  Stabilizer and quotient-descent lemmas isolate the periodic supports and make their numerators field-size independent.  The deep-regime theorem controls MCA directly below one third of the distance and removes the singular Hankel branches there.  The restriction-of-scalars lift shows that extension-valued lines are not a separate class of witnesses.  The Johnson and recursion statements control the list side outside the band.  The curve theorem covers projective and finite-parameter samplers in the deep regime.  Finally, the scanner and finite-certificate grammar package all deployed verdicts as exact integer comparisons.  The only mathematical component not closed by these ledgers is the aperiodic band problem, the broad band formulation \cref{prob:band} below.

\subsection{Stabilizer exhaustion and exact quotient descent}\label{sec:discharge-T1}

Throughout this subsection, \emph{multiplicative case} means: $D=\alpha H$ is a multiplicative coset of a cyclic group $H$ of order $n$, acting freely on $D$ by multiplication; for $c\mid n$, $H_c\le H$ is the unique subgroup of order $c$, equal to the full group $\mu_c$ of $c$-th roots of unity in $\B$, and the fibers of $\pi_c:x\mapsto x^c$ on $D$ are exactly the $H_c$-orbits. \emph{Chebyshev case} means: $D=\chi(\mathcal D)$ for a twin coset $\mathcal D=gH\cup g^{-1}H$ of size $2M$ with $\operatorname{ord}(g)\nmid2M$ as in \cref{sec:circle-geometry}; for $c\mid M$, $K_c\le\UU$ is the order-$c$ kernel of the $c$-th power map, and for $S\subseteq D$ we write $\widetilde S:=\chi^{-1}(S)\cap\mathcal D$ for the (automatically inversion-closed) lift. By \cref{lem:cheb-fibers}, the $T_c$-fibers in $D$ are exactly the $\chi$-images of the $\langle K_c,\,u\mapsto u^{-1}\rangle$-orbits on $\mathcal D$. In both cases the \emph{scale lattice} is the divisor lattice of $n$ (resp.\ of $M$).

\begin{definition}[periodicity scale of a support]\label{def:periodicity-scale}
For $S\subseteq D$, the periodicity scale $c(S)$ is $|\{h\in H:\ hS=S\}|$ in the multiplicative case and $|\{\zeta\in K_M:\ \zeta\widetilde S=\widetilde S\}|$ in the Chebyshev case. A support is \emph{aperiodic} if $c(S)=1$.
\end{definition}

\begin{theorem}[stabilizer partition and structural exhaustion]\label{thm:stabilizer-partition}
In either case:
\begin{enumerate}[label=\textup{(\roman*)}]
\item The stabilizer in the definition is a subgroup of a cyclic group, so $c(S)$ divides $n$ \textup(resp.\ $M$\textup), and for every scale $c$,
\[
S\ \text{is a union of complete scale-$c$ fibers}
\quad\Longleftrightarrow\quad
c\mid c(S).
\]
In particular $c(S)$ is the unique maximal such scale, and the supports of each size are partitioned by the value of $c(S)$, with $c(S)=1$ the aperiodic class.
\item Fix a received line and an agreement threshold $a\ge k$, and call a bad parameter \emph{periodically witnessed} if at least one of its witness supports $S$ \textup(in the sense of \cref{def:ca,def:mca}\textup) has $c(S)\ge2$. Then every bad parameter is either periodically witnessed --- and hence counted, with multiplicity one per support, by the family of supports of size $\ge a$ that are unions of complete fibers at some scale $\ge2$ --- or all of its witness supports are aperiodic, in which case it belongs to the aperiodic bucket by definition. The two classes are exhaustive and disjoint.
\end{enumerate}
\end{theorem}

\begin{proof}
(i) The set $\{h\in H:hS=S\}$ is a subgroup of the cyclic group $H$ (resp.\ $\{\zeta\in K_M:\zeta\widetilde S=\widetilde S\}\le K_M$), hence equals $H_{c_0}$ (resp.\ $K_{c_0}$) for a unique divisor $c_0$, which is $c(S)$. In the multiplicative case, the $\pi_c$-fibers in $D$ are the $H_c$-orbits, and a set is a union of complete $H_c$-orbits if and only if it is $H_c$-invariant, i.e.\ $H_c\le H_{c_0}$, i.e.\ $c\mid c_0$. In the Chebyshev case, $S$ is a union of complete $T_c$-fibers if and only if $\widetilde S$ is a union of complete $\langle K_c,\mathrm{inv}\rangle$-orbits; since $\widetilde S$ is inversion-closed automatically, this holds if and only if $\widetilde S$ is $K_c$-invariant, i.e.\ $c\mid c_0$. The partition statement is immediate.

(ii) Exhaustiveness and disjointness are definitional given (i). The counting clause is \cref{lem:one-support-one-line} applied to the support family $\{S:\ |S|\ge a,\ c(S)\ge2\}$: each fixed support pays for at most one finite line parameter, so the number of periodically witnessed bad parameters is at most the size of that family.
\end{proof}

\begin{corollary}[exact size of the periodic support family; $q$-free numerator]\label{cor:periodic-support-count}
Let $P$ be the set of primes dividing $n$ \textup(resp.\ $M$\textup). The number of supports of exact size $A$ with $c(S)\ge2$ is
\[
\Pi(A)\ :=\ \sum_{\emptyset\ne T\subseteq P}(-1)^{|T|+1}\,
\mathbf 1\bigl[\operatorname{lcm}(T)\mid A\bigr]\,
\binom{\,n/\operatorname{lcm}(T)\,}{\,A/\operatorname{lcm}(T)\,},
\]
and consequently, for every received line and every agreement threshold $a\ge k$,
\[
\#\{\text{periodically witnessed bad finite parameters}\}\ \le\ \sum_{A=a}^{n}\Pi(A),
\]
with an extra factor $d$ for degree-$d$ power-curve samplers \textup(\cref{lem:one-support-d-curve}\textup). The numerator is independent of $q$ and of $|\B|$; the subfield order enters only the lower floors, as the periodic-support ledger demands.
\end{corollary}

\begin{proof}
$c(S)\ge2$ if and only if $H_pS=S$ for some $p\in P$, i.e.\ $S$ lies in the union over $p\in P$ of the families of complete-$p$-fiber unions. For $T\subseteq P$, invariance under every $H_p$, $p\in T$, is invariance under the subgroup they generate, which is $H_{\operatorname{lcm}(T)}$; the number of unions of complete $\operatorname{lcm}(T)$-fibers of total size $A$ is $\binom{n/\operatorname{lcm}(T)}{A/\operatorname{lcm}(T)}$ when $\operatorname{lcm}(T)\mid A$ and $0$ otherwise. Inclusion--exclusion gives $\Pi(A)$; the parameter bound is \cref{thm:stabilizer-partition}(ii) \textup(resp.\ \cref{lem:one-support-d-curve}\textup). The Chebyshev case is identical with $K_p\le K_M$.
\end{proof}

This complements \cref{thm:exact-divisor-block-support-ledger}: that formula counts a \emph{declared} remainder-profile family exactly; \cref{cor:periodic-support-count} counts the \emph{intrinsic} periodic family exactly, and \cref{thm:stabilizer-partition} proves that nothing quotient-periodic escapes it. What the periodic-support ledger still lacked was the passage from support-level to image-level counting. The next theorem supplies it exactly on the class of received data for which the quotient family was designed.

\begin{theorem}[exact quotient descent for fiber-constant data]\label{thm:fiber-descent}
Work in the multiplicative or Chebyshev case, write $\varphi_c$ for the scale-$c$ map \textup($X^c$, resp.\ $T_c$\textup) and $Q_c:=\varphi_c(D)$, $N_c:=n/c$. Let $c$ be a scale, let $p\in\F[X]_{<k}$, let $S\subseteq D$ be a union of complete $c$-fibers with $|S|\ge k$, and suppose $p$ agrees on $S$ with a fiber-constant word $W\circ\varphi_c$. Then there is a unique $g\in\F[Y]$ with
\[
p=g\circ\varphi_c,\qquad \deg g<\lceil k/c\rceil,
\]
and $g$ agrees with $W$ on $\varphi_c(S)$, a set of $|S|/c$ points of $Q_c$. Consequently, for a received line $f+zg_0$ with $f$ and $g_0$ both constant on $c$-fibers, and any agreement threshold $a\ge k$: a finite slope $z$ is MCA-bad with some witness support that is a union of complete $c$-fibers if and only if $z$ is MCA-bad for the quotient line $(f_c,g_{0,c})$ on $\RS[\F,Q_c,\lceil k/c\rceil]$ at the corresponding agreement $\ge\lceil a/c\rceil$, with witness $\varphi_c(S)$; and the same holds verbatim for lists: $p\mapsto g$ is a bijection between scale-$c$-witnessed list elements and quotient list elements.
\end{theorem}

\begin{proof}
\emph{Multiplicative case.} For $h\in H_c=\mu_c$ and $x\in S$ we have $hx\in S$ and $\varphi_c(hx)=\varphi_c(x)$, so $p(hx)=W(\varphi_c(hx))=W(\varphi_c(x))=p(x)$. Thus $p(hX)-p(X)$, of degree $<k$, has at least $|S|\ge k$ roots, hence vanishes identically: $p(hX)=p(X)$ for every $h\in\mu_c$. Comparing coefficients, $a_ih^i=a_i$ for all $h\in\mu_c$, so $a_i=0$ whenever $c\nmid i$ (take $h$ a generator of $\mu_c$), i.e.\ $p=g(X^c)$ with $c\deg g=\deg p<k$, so $\deg g<\lceil k/c\rceil$; uniqueness is clear. For $x\in S$, $g(\varphi_c(x))=p(x)=W(\varphi_c(x))$, and $\varphi_c(S)$ has $|S|/c$ elements since $S$ is a union of complete $c$-fibers.

\emph{Chebyshev case.} Work over the compositum $\F':=\F\cdot\F_{p^2}$ and put $P(u):=p(\chi(u))\in\F'[u,u^{-1}]$, a Laurent polynomial with exponents in $[-(k-1),k-1]$. For $\zeta\in K_c=\mu_c$ and $u\in\widetilde S$ we have $\zeta u\in\widetilde S$ and $T_c(\chi(\zeta u))=\chi(\zeta^cu^c)=T_c(\chi(u))$, so $P(\zeta u)=P(u)$ on $\widetilde S$; since $|\widetilde S|=2|S|\ge2k$ exceeds the degree $2(k-1)$ of $u^{k-1}\bigl(P(\zeta u)-P(u)\bigr)$, the identity $P(\zeta u)=P(u)$ holds in $\F'[u,u^{-1}]$. Hence only exponents divisible by $c$ occur, and $P$ is also inversion-symmetric because $\chi$ is, so $P$ lies in the symmetric Laurent ring $\F'[u^c+u^{-c}]$: $P(u)=G(u^c+u^{-c})$ for a unique $G$. Since $u^c+u^{-c}=2\,T_c(\chi(u))$, setting $g(Y):=G(2Y)$ gives $p(\chi(u))=g(T_c(\chi(u)))$ identically, and injectivity of the ring map $r(x)\mapsto r((u+u^{-1})/2)$ yields $p=g\circ T_c$ in $\F'[X]$; as $g$ is the unique solution of an $\F$-linear system with data in $\F$, in fact $g\in\F[Y]$, and $c\deg g=\deg p<k$ as before. Agreement on $\varphi_c(S)=T_c(S)$ follows as above, with $|T_c(S)|=|S|/c$ by \cref{lem:cheb-fibers}.

\emph{Lines and lists.} If $z$ has an MCA witness $S$ that is a union of complete $c$-fibers, the explaining codeword for the point $f+zg_0$ (a fiber-constant word) descends by the above; conversely a quotient explanation composes with $\varphi_c$. For noncontainment, if the pair $(f,g_0)$ were explained on $S$ by $(p_1,p_2)$, both $p_i$ descend (each agrees on $S$ with a fiber-constant word), so the quotient pair would be explained on $\varphi_c(S)$; and conversely quotient explanations compose. Hence badness with a scale-$c$ witness corresponds exactly. For lists, $p\mapsto g$ is well defined and injective (uniqueness of $g$), and $g\mapsto g\circ\varphi_c$ is its inverse from the quotient list, since the full agreement set of $g\circ\varphi_c$ with $W\circ\varphi_c$ is the full $\varphi_c$-preimage of the agreement set of $g$ with $W$.
\end{proof}

\begin{corollary}[structural exhaustion is exact]\label{cor:T1-status}
With the intrinsic \cref{def:periodicity-scale}, the structural exhaustion is a theorem: every bad parameter is quotient-periodic \textup(periodically witnessed, hence counted exactly at the support level by \cref{cor:periodic-support-count}, and exactly at the image level by \cref{thm:fiber-descent} whenever the received line is fiber-constant\textup) or belongs to the aperiodic bucket. The subfield order $|\B|$ appears only in lower floors; all upper numerators above are $q$-free. The entire residual content of the periodic-support ledger is thereby the aperiodic bucket, treated next.
\end{corollary}

\subsection{The deep regime, singular pencils, and the residual band}\label{sec:discharge-T2}

The following theorem is the safe-side counterpart of the paper's unsafe caps. It is self-contained, uses nothing beyond interpolation and double counting, and holds for arbitrary linear codes --- in particular for every multiplicative row, every Chebyshev line-round row, every bivariate circle code over any field \textup(including fields not containing $i$\textup), and every ECFFT row.

\begin{theorem}[deep-regime mutual correlated agreement, unconditional]\label{thm:deep-mca}
Let $C\subseteq\F^n$ be any $\F$-linear code of minimum Hamming weight $w$, let $q=|\F|$, let $\delta\in(0,1)$, and put $r:=\lfloor\delta n\rfloor$. Assume
\[
3r\ \le\ w-1 .
\]
Then for every pair $(f_1,f_2)\in\F^n\times\F^n$, at least one of the following alternatives holds:
\begin{enumerate}[label=\textup{(\alph*)}]
\item \textup{(Tangent pair.)} There are codewords $c_1,c_2\in C$ and a set $T'\subseteq\{1,\dots,n\}$ with $|T'|\le r$ such that $f_i=c_i$ off $T'$ for $i=1,2$; in this case $(f_1,f_2)$ has no CA-bad slope at radii $(\delta,\delta_{\mathrm{int}})$ for any $\delta_{\mathrm{int}}\ge\delta$, and at most $|T'|\le r$ MCA-bad slopes at radius $\delta$, namely a subset of $\{-e_{1,j}/e_{2,j}:\,j\in T',\ e_{2,j}\ne0\}$ where $e_i:=f_i-c_i$.
\item \textup{(Generic pair.)} $\#\{\gamma\in\F:\ \dist(f_1+\gamma f_2,\,C)\le\delta\}\ \le\ r+1$.
\end{enumerate}
Consequently
\[
\eca(C,\delta)\ \le\ \frac{r+1}q,
\qquad
\emca(C,\delta)\ \le\ \frac{r+1}q .
\]
For $C=\RS[\F,D,k]$ on any evaluation set $D$ of size $n$, the hypothesis reads $3\lfloor\delta n\rfloor\le n-k$.
\end{theorem}

\begin{proof}
Write $S_{\rm cl}:=\{\gamma:\dist(f_1+\gamma f_2,C)\le\delta\}$; note $\dist\le\delta$ means disagreement on at most $r$ coordinates, by integrality. If $|S_{\rm cl}|\le1$, alternative (b) holds ($r+1\ge1$). Otherwise fix $\gamma_1\ne\gamma_2$ in $S_{\rm cl}$ with codewords $p_1,p_2$ and error sets $E_1,E_2$ ($|E_i|\le r$), and set
\[
c_2:=\frac{p_1-p_2}{\gamma_1-\gamma_2},
\qquad
c_1:=p_1-\gamma_1c_2,
\qquad
e_i:=f_i-c_i .
\]
Off $T:=E_1\cup E_2$ ($|T|\le2r$) we can solve the two linear equations $f_1+\gamma_if_2=p_i$ coordinatewise, giving $f_2=c_2$ and $f_1=c_1$ there; hence $e_1,e_2$ are supported in $T$. Let $T':=\{j:(e_{1,j},e_{2,j})\ne(0,0)\}\subseteq T$.

\emph{Step 1: every close slope rides the recovered line.} Let $\gamma\in S_{\rm cl}$ with codeword $p$ and error set $E$, $|E|\le r$. Off $T\cup E$ we have $p=f_1+\gamma f_2=c_1+\gamma c_2$, so the codeword $p-(c_1+\gamma c_2)$ has weight at most $|T|+|E|\le3r\le w-1<w$ and is therefore zero. Hence
\[
f_1+\gamma f_2-(c_1+\gamma c_2)\ =\ e_1+\gamma e_2
\quad\text{has weight at most } r
\qquad\text{for every }\gamma\in S_{\rm cl}. \tag{$*$}
\]

\emph{Step 2: dichotomy.} Suppose $|T'|\ge r+1$. For $j\in T'$, the affine function $\gamma\mapsto e_{1,j}+\gamma e_{2,j}$ is not identically zero, so it vanishes for at most one $\gamma$. By $(*)$, each $\gamma\in S_{\rm cl}$ needs at least $|T'|-r\ge1$ vanishing coordinates inside $T'$. Double counting vanishings,
\[
|S_{\rm cl}|\,\bigl(|T'|-r\bigr)\ \le\ |T'|,
\qquad\text{hence}\qquad
|S_{\rm cl}|\ \le\ \frac{|T'|}{|T'|-r}\ \le\ r+1,
\]
which is alternative (b). Otherwise $|T'|\le r$ and $(f_1,f_2)$ is a tangent pair as in (a) with this $T'$.

\emph{Step 3: bad slopes of tangent pairs.} Let $(f_1,f_2)$ be tangent with data $(c_1,c_2,T')$, $|T'|\le r$. First, $\dist_2((f_1,f_2),C^{\equiv2})\le|T'|/n\le r/n\le\delta\le\delta_{\mathrm{int}}$, so no slope is CA-bad at radii $(\delta,\delta_{\mathrm{int}})$. For MCA, note $2r\le w-1$ (from $3r\le w-1$ and $r\ge0$), so two codewords agreeing on at least $n-2r$ coordinates are equal. Let $\gamma$ be MCA-bad with witness $S$, $|S|\ge n-r$. On $S\setminus T'$, which has at least $n-2r$ elements, $f_1+\gamma f_2=c_1+\gamma c_2$; hence the codeword explaining the point on $S$ equals $c_1+\gamma c_2$, and thus $e_1+\gamma e_2$ vanishes on $S\cap T'$. Moreover, any pair explanation $(p'_1,p'_2)$ on $S$ satisfies $p'_i=c_i$ on the $\ge n-2r$ points of $S\setminus T'$, forcing $p'_i=c_i$; so the pair fails to be explained on $S$ exactly when $S\cap T'\ne\emptyset$ (every $j\in T'$ has $(e_{1,j},e_{2,j})\ne0$). Therefore $\gamma$ is MCA-bad if and only if $e_{1,j}+\gamma e_{2,j}=0$ for some $j\in T'$: necessity was just shown; for sufficiency take $S:=(\{1,\dots,n\}\setminus T')\cup\{j\}$, of size $\ge n-r$, on which $c_1+\gamma c_2$ explains the point while the pair is not explained. Each $j\in T'$ contributes at most one slope, giving at most $|T'|\le r$ MCA-bad slopes.

Finally, generic pairs have at most $r+1$ close slopes, hence at most $r+1$ CA- or MCA-bad slopes (an MCA witness of size $\ge n-r$ implies $\dist\le r/n\le\delta$); tangent pairs have $0$ CA-bad and at most $r$ MCA-bad slopes. Taking the maximum over pairs gives the two displayed error bounds. For $\RS[\F,D,k]$, $w=n-k+1$, so $3r\le w-1$ is $3r\le n-k$.
\end{proof}

\begin{remark}[two radii; scope]
The CA statement holds in the two-radius form: for any $\delta_{\mathrm{int}}\ge\delta_{\mathrm{fld}}=\delta$ with $3\lfloor\delta n\rfloor\le w-1$, $\eca(C,\delta,\delta_{\mathrm{int}})\le(r+1)/q$, since the tangent alternative kills the far condition at every $\delta_{\mathrm{int}}\ge\delta$. The constant $\frac13$ can be improved toward $\frac12$ of the distance by the methods of \cite{BCHKS25}; we deliberately keep the three-line-per-step, formalization-ready proof --- it feeds the formal finite core --- and record the import separately. Note also that \cref{thm:deep-mca} applies to the bivariate circle code over $\F_p$ itself \textup(minimum weight $n_c-k_c+1$, by extending scalars to $\F_{p^2}$ and using \cref{lem:circle-rs}\textup), so the deep regime of the genus-one/circle transport question \cref{prob:isogeny}\textup{(b)} is closed as well.
\end{remark}

We next resolve a named obstruction of the residual-band analysis --- the singular Hankel buckets --- in the regime where \cref{thm:deep-mca} applies. First, the structural normal form of the identically singular branch.

\begin{lemma}[identically singular pencils are locator families]\label{lem:singular-pencil}
In the setting of \cref{sec:aperiodic-hankel-certificates}, fix an exact agreement $A$ with $t_A\ge j_A+1$ and set $M(Z):=H_{t_A,j_A}(u)+ZH_{t_A,j_A}(v)$. The following are equivalent:
\begin{enumerate}[label=\textup{(\roman*)}]
\item every $(j_A+1)\times(j_A+1)$ minor of $M(Z)$ vanishes identically in $Z$;
\item there is a nonzero polynomial vector $\widetilde v(Z)=\sum_{i=0}^{d}Z^iv_i\in\F[Z]^{j_A+1}$ with $M(Z)\widetilde v(Z)\equiv0$, equivalently a chain
\[
H_{t_A,j_A}(u)\,v_0=0,\qquad
H_{t_A,j_A}(u)\,v_{i}=-H_{t_A,j_A}(v)\,v_{i-1}\ (1\le i\le d),\qquad
H_{t_A,j_A}(v)\,v_d=0 .
\]
\end{enumerate}
In that case, writing $\lambda(Z)$ for the last coordinate of $\widetilde v(Z)$, every $Z=\gamma$ with $\lambda(\gamma)\ne0$ admits a monic degree-$j_A$ syndrome annihilator for $u+\gamma v$, i.e.\ every such slope passes the rank test at exact agreement $A$; the at most $\deg\lambda$ remaining slopes pass it after column truncation. The rank test is necessary for badness at agreement $A$ but not sufficient \textup(the annihilator must in addition be $D$-split\textup), so membership in the singular bucket never certifies badness by itself.
\end{lemma}

\begin{proof}
(i) says $\operatorname{rank}_{\F(Z)}M(Z)\le j_A$, which holds if and only if the kernel of $M(Z)$ over the rational function field is nonzero; clearing denominators and dividing by common factors produces a nonzero polynomial kernel vector, and expanding $M(Z)\widetilde v(Z)\equiv0$ in powers of $Z$ gives the chain. The converse direction is immediate. Specializing $Z=\gamma$ with $\lambda(\gamma)\ne0$ and normalizing the last coordinate to $1$ produces the vector $\widetilde\ell$ of the chart atlas with $A(\ell)+\gamma B(\ell)=0$, i.e.\ a monic annihilator; $\lambda\ne0$ has at most $\deg\lambda$ roots. The final sentence restates the convention of \cref{sec:aperiodic-hankel-certificates}.
\end{proof}

\begin{corollary}[the singular buckets are resolved in the deep regime]\label{cor:singular-deep}
Let $C=\RS[\F,D,k]$ and let $A$ be an exact agreement with $3(n-A)\le n-k$. Then for every received line --- in particular for every line in the underdetermined or identically singular buckets of \cref{sec:aperiodic-hankel-certificates} --- the number of finite slopes that are bad at agreement threshold $A$ is at most $n-A+1$ if the pair is not $\frac{n-A}n$-tangent, and at most $n-A$ MCA-bad \textup(zero CA-bad\textup) slopes if it is. In particular, in this regime the singular bucket contributes at most $n-A+1$ to $B_{\rm ap}(A)$ per line, and no line contributes $\Theta(q)$: the ``totally bad line'' scenario that the singular bucket was guarding against does not occur below one third of the distance.
\end{corollary}

\begin{proof}
Immediate from \cref{thm:deep-mca} with $r=n-A$ and $\delta=r/n$: badness at agreement threshold $A$ implies distance at most $r/n$.
\end{proof}

\begin{remark}[calibration: near-capacity unsafety is quotient-borne]\label{rem:quotient-borne}
The unsafe mass produced by this paper does not press on the aperiodic bucket. The fiber witnesses of \cref{lem:fiber,lem:phi-fiber} have agreement supports that are unions of complete fibers, hence periodicity scale $\ge a\ge2$; by \cref{thm:stabilizer-partition} they are charged to the quotient side of the ledger, and by \cref{thm:fiber-descent} even at the image level. So a polynomially bounded aperiodic numerator in the band remains consistent with every theorem in this paper. This motivates stating the entire residual band as a single conjecture.
\end{remark}

\begin{remark}[Original broad band formulation]\label{prob:band}
Fix a multiplicative or Chebyshev row $C=\RS[\F,D,k]$. Prove or refute: there is a polynomial $P$, independent of $q$, such that for every agreement $a$ in the open band
\[
k+\tfrac{2n}N\ <\ a\ <\ n-\Bigl\lfloor\frac{n-k}3\Bigr\rfloor ,
\]
the number of MCA-bad finite slopes of any received line, all of whose witness supports are aperiodic in the sense of \cref{def:periodicity-scale}, is at most $P(n)$. By \cref{cor:T1-status,thm:weil-lines} \textup(structural exhaustion and the extension-line lift\textup) and \cref{thm:deep-mca} \textup(deep regime discharged\textup), this conjecture is exactly the missing input \textup{(ii)} of \cref{thm:conditional-mca} throughout the band, for every challenge row.
\end{remark}

\subsection{The exact extension-line lift}\label{sec:discharge-T3}

\begin{theorem}[restriction of scalars for lines]\label{thm:weil-lines}
Let $\B\subseteq\F$ with $e:=[\F:\B]$ and $D\subseteq\B$, and fix a $\B$-basis $1=\beta_0,\beta_1,\dots,\beta_{e-1}$ of $\F$. Decompose words as $w=\sum_i\beta_iw_i$ with $w_i\in\B^n$, and let $\mu_{ij}(\gamma)\in\B$ be the multiplication tensor, $\gamma\beta_i=\sum_j\mu_{ij}(\gamma)\beta_j$, which is $\B$-linear and injective in $\gamma$. Then:
\begin{enumerate}[label=\textup{(\roman*)}]
\item A word $c$ lies in $\RS[\F,D,k]$ if and only if every component $c_j$ lies in $\RS[\B,D,k]$, and for all words $w$,
\[
\dist\bigl(w,\RS[\F,D,k]\bigr)\ =\ \dist_e\bigl((w_0,\dots,w_{e-1}),\,\RS[\B,D,k]^{\equiv e}\bigr).
\]
More generally, a set $S\subseteq D$ explains $w$ over $\F$ if and only if $S$ column-explains the component tuple over $\B$.
\item For a pair $(f_1,f_2)$ over $\F$ and $\gamma\in\F$, the components of $f_1+\gamma f_2$ are
\[
(f_1+\gamma f_2)_j\ =\ f_{1,j}+\sum_{i}\mu_{ij}(\gamma)\,f_{2,i},
\]
an $e$-dimensional $\B$-linear \textup(in $\gamma$\textup) family of $e$-interleaved words over $\B$.
\item For every radius \textup(and pair of radii\textup), $\gamma$ is CA-bad \textup(resp.\ MCA-bad\textup) for $(f_1,f_2)$ over $\F$ if and only if the parameter $\gamma$ of the sampler in \textup{(ii)} is bad in the corresponding interleaved sense over $\B$, with pair noncontainment read on the $2e$-tuple $(f_{1,0},\dots,f_{1,e-1},f_{2,0},\dots,f_{2,e-1})$ against $\RS[\B,D,k]^{\equiv2e}$, on the same supports. For $\gamma\in\B$ the sampler restricts to the diagonal $e$-interleaved line $(f_{1,j}+\gamma f_{2,j})_j$.
\end{enumerate}
Consequently the extension term $B_{\rm ext}(a)$ of \cref{thm:conditional-mca} is not an independent quantity: it equals, definitionally, a base-field interleaved-sampler numerator. In particular no ``extension-valued counterexample'' can exist outside the base ledgers: every genuinely $\F$-valued line is charged, exactly, to a base bucket.
\end{theorem}

\begin{proof}
(i) is the componentwise decomposition already used in the proof of \cref{lem:confine} \textup(valid for any $D\subseteq\B$, cf.\ \cref{cor:circle-Fvalued}\textup): the coefficients of a polynomial over $\F$ decompose in the basis, evaluation at $D\subseteq\B$ commutes with the decomposition, and pointwise agreement over $\F$ is simultaneous componentwise agreement, i.e.\ column agreement. (ii) is the definition of $\mu$. (iii) The point condition, the pair condition, and every support-restricted condition in \cref{def:ca,def:mca} are, by (i), equalities of components on sets of coordinates, and these translate verbatim; injectivity of $\gamma\mapsto\mu(\gamma)$ makes the parameterizations match one-to-one. The diagonal claim is $\mu_{ij}(\gamma)=\gamma\delta_{ij}$ for $\gamma\in\B$.
\end{proof}

\begin{corollary}[status of the extension-line lift]\label{cor:T3-status}
The extension-line question is discharged on the ``lift'' side of its dichotomy: extension lines are covered --- exactly, not merely up to a stable numerator --- by base-field ledgers, namely the interleaved-sampler ledgers of the list and curve sides via \cref{thm:weil-lines}. In the deep regime nothing at all remains: \cref{thm:deep-mca} applies over $\F$ directly, uniformly in the field of definition of the pair. On the lower side, \cref{cor:extension-pole-deep-list-floor,cor:extension-pole-quotient-remainder-floor} already produce genuinely extension-valued bad lines from any supplied list, and \cref{lem:confine} \textup(with \cref{cor:Fvalued,cor:circle-Fvalued}\textup) shows the certified deployed failures must be of this kind. What the extension-line lift leaves behind is exactly the band input of the list and curve ledgers, i.e.\ the broad band formulation \cref{prob:band}.
\end{corollary}
\subsection{Uniform interleaved-list bounds, exact recursion, and calibration}\label{sec:discharge-T4}

\begin{theorem}[elementary Johnson list bound, uniform in the interleaving arity]\label{thm:johnson-list}
Let $C=\RS[\F,D,k]$ on any evaluation set $D$ of size $n$, let $m\ge1$, and let $a$ be an agreement threshold with $a^2>(k-1)n$. Then for every received $m$-tuple $U$,
\[
\bigl|\Lst\bigl(C^{\equiv m},\,1-\tfrac an,\,U\bigr)\bigr|
\ \le\
\frac{n\,(a-k+1)}{a^2-(k-1)n}\,,
\]
where the list is taken in the column distance $\dist_m$. The same bound holds for any code family in which two distinct codewords column-agree on at most $k-1$ coordinates; in particular it holds for interleaved bivariate circle codes with $k$ replaced by $k_c$.
\end{theorem}

\begin{proof}
Let $\mathbf c^{(1)},\dots,\mathbf c^{(L)}$ be distinct listed codewords with column-agreement sets $S_1,\dots,S_L\subseteq D$, $|S_i|\ge a$. Two distinct interleaved codewords column-agree exactly where all their rows agree; some pair of rows consists of distinct polynomials of degree $<k$, so $|S_i\cap S_j|\le k-1$ for $i\ne j$. Let $\deg(x):=\#\{i:x\in S_i\}$, so $\sum_x\deg(x)\ge La$ and
\[
\sum_{x\in D}\binom{\deg(x)}2=\sum_{i<j}|S_i\cap S_j|\le\binom L2(k-1).
\]
By convexity of $t\mapsto t(t-1)$, $\sum_x\deg(x)(\deg(x)-1)\ge n\cdot\frac{La}n\bigl(\frac{La}n-1\bigr)$ \textup(the right side is nonpositive when $La\le n$, and the average-value bound applies otherwise\textup). Hence $La(La-n)\le L(L-1)(k-1)n$, i.e.\ $a(La-n)\le(L-1)(k-1)n$, i.e.
\[
L\bigl(a^2-(k-1)n\bigr)\ \le\ an-(k-1)n\ =\ n(a-k+1),
\]
which is the claim when $a^2>(k-1)n$.
\end{proof}

\begin{theorem}[exact divisor-lattice list recursion for fiber-constant words]\label{thm:list-recursion}
Multiplicative case; let $c_0\mid n$, let $U=W\circ\pi_{c_0}$ be constant on $c_0$-fibers, and let $a\ge k$. For $p\in\Lst(C,1-a/n,U)$ let $S_p:=\{x\in D:p(x)=U(x)\}$ be the full agreement set and put $c_p:=\gcd\bigl(c(S_p),c_0\bigr)$. For $c\mid c_0$ write $W_c(y):=W\bigl(y^{c_0/c}\bigr)$ on $Q_c$, so $U=W_c\circ\pi_c$. Then, for each $c\mid c_0$, the descent map of \cref{thm:fiber-descent} is a bijection from $\{p\in\Lst(C,1-a/n,U):\ c_p=c\}$ onto
\[
\Bigl\{\,g\in\Lst\bigl(\RS[\F,Q_c,\lceil k/c\rceil],\,1-\lceil a/c\rceil/N_c,\,W_c\bigr)\ :\ \gcd\bigl(c(\bar S_g),\,c_0/c\bigr)=1\,\Bigr\},
\]
where $\bar S_g$ is the full agreement set of $g$ with $W_c$ and its scale is computed in the quotient lattice. In particular
\[
\bigl|\Lst(C,1-\tfrac an,U)\bigr|=\sum_{c\mid c_0}\#\{\cdots\},
\]
and if moreover $a^2>kn$, then every bucket obeys the level-$c$ Johnson bound and
\[
\bigl|\Lst(C,1-\tfrac an,U)\bigr|\ \le\ \sum_{c\mid c_0}\frac{n\,(a-k+2c-1)}{a^2-kn}\ \le\ \tau(c_0)\,\frac{n\,(a-k+2c_0-1)}{a^2-kn}\,,
\]
$\tau$ the number-of-divisors function.
\end{theorem}

\begin{proof}
Fix $p$ with $c_p=c$. Then $c\mid c(S_p)$, so $S_p$ is a union of complete $c$-fibers by \cref{thm:stabilizer-partition}(i), and $U=W_c\circ\pi_c$ is constant on $c$-fibers; \cref{thm:fiber-descent} (with $S=S_p$, $|S_p|\ge a\ge k$) gives $p=g\circ\pi_c$ with $g$ in the displayed quotient list, and $S_p=\pi_c^{-1}(\bar S_g)$ because $p(x)=U(x)\iff g(x^c)=W_c(x^c)$. The power map is equivariant, $\pi_c(hx)=h^c\pi_c(x)$, and $Q_c$ is a coset of $H^{(c)}:=\{h^c:h\in H\}$ with kernel $H_c$; hence $hS_p=S_p\iff h^c\bar S_g=\bar S_g$ and $c(S_p)=c\cdot c(\bar S_g)$. Since $c\mid c_0$, $\gcd(c\cdot c(\bar S_g),c_0)=c\cdot\gcd(c(\bar S_g),c_0/c)$, so $c_p=c$ is exactly the displayed coprimality. Conversely, for $g$ in the displayed set, $p:=g\circ\pi_c$ has full agreement set $\pi_c^{-1}(\bar S_g)$ of size $c|\bar S_g|\ge c\lceil a/c\rceil\ge a$, and the same computation returns $c_p=c$; the two maps are mutually inverse. Summing the buckets gives the exact formula. For the Johnson clause: at level $c$ the code has dimension $\lceil k/c\rceil$ and length $N_c=n/c$, and
\[
\Bigl\lceil\frac ac\Bigr\rceil^2-\Bigl(\Bigl\lceil\frac kc\Bigr\rceil-1\Bigr)\frac nc\ \ge\ \frac{a^2}{c^2}-\frac{kn}{c^2}\ =\ \frac{a^2-kn}{c^2}\ >\ 0,
\]
using $\lceil k/c\rceil-1<k/c$, so \cref{thm:johnson-list} applies at every level; with $\lceil a/c\rceil-\lceil k/c\rceil+1\le(a-k+2c-1)/c$ the level-$c$ bound is at most $n(a-k+2c-1)/(a^2-kn)$.
\end{proof}

\begin{proposition}[calibration: the quotient list term is necessarily word-dependent and exponential]\label{prop:list-calibration}
Let the hypotheses of \cref{lem:phi-fiber}\textup{(ii)} hold and let $U=u_z$ be the pigeonholed received word at agreement $A_2$. Every listed codeword $c_A=-R_A(\varphi)$ lies in the scale-$a$ quotient image $\{g\circ\varphi:\deg g\le\ell_2-2\}$. Hence in any bound of the shape demanded by the interleaved-list ledger,
\[
|\Lst(C^+,1-\tfrac{A_2}n,U)|\ \le\ L_Q(A_2,U)+L_{\rm ap}(A_2),
\]
in which $L_Q$ upper-bounds the quotient-image sublist and $L_{\rm ap}$ the rest, one necessarily has
\[
L_Q(A_2,u_z)\ \ge\ \binom N{\ell_2}\Big/|\B| .
\]
No uniform subexponential $L_Q$ exists in the band; only the aperiodic term $L_{\rm ap}$ can be uniformly polynomial, which is consistent with the broad band formulation \cref{prob:band} and with \cref{rem:quotient-borne}.
\end{proposition}

\begin{proof}
Immediate from \cref{lem:phi-fiber}: the $\ge\binom N{\ell_2}/|\B|$ listed codewords are of the displayed composed form, hence are counted by $L_Q$.
\end{proof}

\begin{corollary}[the grand interleaved-list staircase outside the band]\label{cor:list-sandwich}
Fix any row satisfying the hypotheses of \cref{cor:grand}, \cref{cor:mixed-radix}, or \cref{cor:circle-grand}, and any interleaving arity $m\ge1$.
\begin{enumerate}[label=\textup{(\roman*)}]
\item \textup{(Unsafe near capacity.)} For every agreement $a\le k+n/N$ \textup(radius $\ge1-\rho-1/N$\textup),
\[
\bigl|\Lst\bigl(C^{\equiv m},\,1-\tfrac an\bigr)\bigr|\ \ge\ \binom N{\ell_1}\Big/|\B|\ >\ 2^{-128}\,q ,
\]
by \cref{lem:phi-fiber}\textup{(i)} and the entropy counts of the corresponding corollary, together with the embedding $c\mapsto(c,0,\dots,0)$ of base lists into interleaved lists.
\item \textup{(Safe in the Johnson regime.)} For every agreement $a$ with $a^2>(k-1)n$ and $n(a-k+1)/(a^2-(k-1)n)\le2^{-128}q$,
\[
\bigl|\Lst\bigl(C^{\equiv m},\,1-\tfrac an\bigr)\bigr|\ \le\ 2^{-128}\,q ,
\]
by \cref{thm:johnson-list}, uniformly in $m$.
\end{enumerate}
Thus the interleaved-list staircase is determined outside the same Johnson-to-capacity band as the MCA staircase; inside the band, \cref{prop:list-calibration} shows that any resolution of the interleaved-list ledger must carry an exponential, word-dependent quotient term, and \cref{thm:list-recursion} shows that for fiber-constant words this term is exactly a divisor-lattice recursion.
\end{corollary}

\begin{proof}
(i) At agreement $A_1\le k+n/N$ the base-field list floor is $\binom N{\ell_1}/|\B|$ by \cref{lem:phi-fiber}\textup{(i)}; list sizes are nondecreasing as the agreement threshold decreases, and the zero-padding embedding preserves column agreement, so the interleaved list is at least as large. The entropy counts of the cited corollaries give $\binom N{\ell_1}\ge2^{513}\ge|\B|\cdot2^{-128}q\cdot2^{0}$ throughout the challenge envelope, since $|\B|\,q\,2^{-128}\le q^2\,2^{-128}<2^{384}$. (ii) is \cref{thm:johnson-list}.
\end{proof}

\subsection{Curves and projective slopes in the deep regime}\label{sec:discharge-T5}

\begin{theorem}[deep-regime theorem for degree-$d$ power-curve samplers]\label{thm:deep-curve}
Let $C\subseteq\F^n$ be $\F$-linear with minimum weight $w$, let $d\ge1$, let $W_\gamma=f_0+\gamma f_1+\cdots+\gamma^df_d$ be a power-curve sampler, let $\delta\in(0,1)$, $r:=\lfloor\delta n\rfloor$, and assume
\[
(d+2)\,r\ \le\ w-1 .
\]
Then either there are codewords $c_0,\dots,c_d\in C$ and a set $T'$ with $|T'|\le r$ such that $f_i=c_i$ off $T'$ for every $i$ \textup(curve-tangent tuple; at most $d\,|T'|\le dr$ MCA-bad and no CA-bad parameters\textup), or
\[
\#\{\gamma\in\F:\ \dist(W_\gamma,C)\le\delta\}\ \le\ d\,(r+1).
\]
In particular the curve analogues of $\eca$ and $\emca$ \textup(with noncontainment read on the coefficient tuple against $C^{\equiv(d+1)}$, as in \cref{lem:one-support-d-curve}\textup) are at most $d(r+1)/q$ at radius $\delta$. Adjoining the projective parameter at infinity \textup(the leading word $f_d$\textup) adds at most one to every count. For $C=\RS[\F,D,k]$ the hypothesis reads $(d+2)\lfloor\delta n\rfloor\le n-k$.
\end{theorem}

\begin{proof}
Let $S_{\rm cl}$ be the displayed set. If $|S_{\rm cl}|\le d$ the second alternative holds. Otherwise pick distinct $\gamma_0,\dots,\gamma_d\in S_{\rm cl}$ with codewords $p_j$ and error sets $E_j$ ($|E_j|\le r$), and let $c_i:=\sum_j\lambda_{ij}p_j\in C$, where $(\lambda_{ij})$ is the inverse of the Vandermonde matrix $(\gamma_j^i)$; off $T:=\bigcup_jE_j$ ($|T|\le(d+1)r$) the same inversion applied to values gives $f_i=c_i$. Put $e_i:=f_i-c_i$, supported in $T$, and $T':=\{x:(e_{0,x},\dots,e_{d,x})\ne0\}$. For any $\gamma\in S_{\rm cl}$ with codeword $p$ and error set $E$: off $T\cup E$, $p=W_\gamma=\sum_i\gamma^ic_i$, so the codeword $p-\sum_i\gamma^ic_i$ has weight at most $(d+2)r<w$, hence is zero, and
\[
W_\gamma-\sum_i\gamma^ic_i=\sum_i\gamma^ie_i
\quad\text{has weight at most }r .
\]
For $x\in T'$ the nonzero polynomial $\gamma\mapsto\sum_ie_{i,x}\gamma^i$ has at most $d$ roots. If $|T'|\ge r+1$, each $\gamma\in S_{\rm cl}$ needs at least $|T'|-r$ vanishing coordinates in $T'$, so $|S_{\rm cl}|(|T'|-r)\le d|T'|$ and $|S_{\rm cl}|\le d|T'|/(|T'|-r)\le d(r+1)$. Otherwise $|T'|\le r$ and the tuple is curve-tangent. In the tangent case, $2r\le w-1$ makes two codewords agreeing on $\ge n-2r$ coordinates equal; arguing exactly as in Step~3 of \cref{thm:deep-mca} \textup(with $n-(d+2)r\ge n-w+1$ forcing both the point codeword and any tuple explanation to coincide with the $c_i$ off $T'$\textup), a parameter is MCA-bad if and only if $\sum_ie_{i,x}\gamma^i=0$ for some $x\in T'$, giving at most $d|T'|$ parameters, and the coefficient tuple is $\delta$-close to $C^{\equiv(d+1)}$ so no parameter is CA-bad. MCA-bad parameters are close, so the counts cover both notions; the parameter at infinity is a single additional value. For Reed--Solomon, $w=n-k+1$.
\end{proof}

Together with \cref{thm:weil-lines}, \cref{thm:deep-curve} also covers the multi-parameter samplers into which extension lines decompose, since those are handled over $\F$ directly by \cref{thm:deep-mca}; what the curve-sampler ledger still lacks is only the band-regime aperiodic pivot atlas, i.e.\ the curve instance of the broad band formulation \cref{prob:band}, together with the exceptional vanishing-vector strata already flagged in \cref{sec:aperiodic-hankel-certificates}.

\subsection{The certificate scanner, its soundness, and deployed output}\label{sec:discharge-T6}

\begin{proposition}[small-field degeneracy]\label{prop:small-field}
Let $C\subsetneq\F^n$ be any linear code and $q=|\F|$. For every $\delta\in(0,1)$, $\emca(C,\delta)\ge1/q$. Consequently $\delta^*_C(\varepsilon^*)=0$ whenever $\varepsilon^*<1/q$; in particular every row over a field of size $q<2^{128}$ --- including the deployed QM31 circle rows, $q=p^4<2^{124}$ --- has a degenerate grand-challenge threshold $\delta^*_C(2^{-128})=0$, and the meaningful staircase question there is posed at a larger target such as $\varepsilon^*=2^{-100}$. At equality $1/q=\varepsilon^*$ the construction below gives only a floor meeting the target exactly, not strict unsafety; all degeneracy statements therefore require $1/q>\varepsilon^*$.
\end{proposition}

\begin{proof}
Pick $f_2\notin C$, any $c_0\in C$ and $\gamma_0\in\F$, and set $f_1:=c_0-\gamma_0f_2$. With $S=D$: the point $f_1+\gamma_0f_2=c_0$ is explained on $S$, while the pair is explained on $S=D$ only if both words lie in $C$, which fails. Hence $\gamma_0$ is MCA-bad at every radius, and $\emca\ge1/q$.
\end{proof}

\begin{definition}[staircase scanner]\label{def:scanner}
The scanner takes a row header $(D\text{-type},n,k,|\B|,q,q_{\rm line}=q)$, a target $\varepsilon^*$, and the applicable divisor data, and for each agreement $a\in\{k+1,\dots,n\}$ evaluates the following verdicts, each tagged by its certifying statement:
\begin{enumerate}[label=\textup{(V\arabic*)},leftmargin=2.6em]
\item \textbf{Cap-unsafe.} If $a\le A_2$ for an instantiated cap \textup(\cref{thm:main,thm:phi-cap} via \cref{cor:grand,cor:mixed-radix,cor:circle-grand,cor:circle-deployed}\textup) and $\frac1{2k}(1-\frac nq)>\varepsilon^*$: mark $a$ MCA-unsafe.
\item \textbf{Floor-unsafe.} If a remainder floor applies at $a$ \textup(\cref{thm:quotient-remainder-deep-floor}\textup) with converted numerator exceeding $\varepsilon^*q$: mark $a$ MCA-unsafe; if a fiber list floor at $a$ exceeds $\varepsilon^*q$ \textup(\cref{lem:fiber,lem:phi-fiber} with the entropy count\textup): mark $a$ list-unsafe.
\item \textbf{Deep-safe.} If $3(n-a)\le n-k$ and $n-a+1\le\varepsilon^*q$: mark $a$ MCA-safe \textup(\cref{thm:deep-mca}\textup); with $(d+2)(n-a)\le n-k$ and $d(n-a+1)\le\varepsilon^*q$: curve-$d$-safe \textup(\cref{thm:deep-curve}\textup).
\item \textbf{Johnson-list-safe.} If $a^2>(k-1)n$ and $n(a-k+1)/(a^2-(k-1)n)\le\varepsilon^*q$: mark $a$ list-safe for every interleaving arity \textup(\cref{thm:johnson-list}\textup).
\end{enumerate}
The output per challenge type is the pair $a_{\rm unsafe}:=\max\{a\ \text{marked unsafe}\}$ and $a_{\rm safe}:=\min\{a\ \text{marked safe}\}$, together with the open interval between them. If $\varepsilon^*<1/q$ the scanner instead reports the degeneracy of \cref{prop:small-field}.
\end{definition}

\begin{theorem}[scanner soundness]\label{thm:scanner-sound}
Every verdict printed by the scanner of \cref{def:scanner} is correct: agreements marked unsafe satisfy $\emca(C,1-a/n)>\varepsilon^*$ \textup(resp.\ list size $>\varepsilon^*q$\textup), agreements marked safe satisfy $\emca(C,1-a/n)\le\varepsilon^*$ \textup(resp.\ list size $\le\varepsilon^*q$, for every $m$\textup), and consequently $a_{\rm unsafe}<a^*_C\le a_{\rm safe}$ in the staircase formulation of \cref{prop:onestep}.
\end{theorem}

\begin{proof}
Each verdict cites a theorem of this paper whose hypotheses are verified numerically by the scanner; monotonicity across agreements is \cref{lem:mca-monotone} for MCA and radius-monotonicity of lists. The bracketing of $a^*_C$ is \cref{prop:onestep}.
\end{proof}

\begin{table}[ht]
\centering
\footnotesize
\setlength{\tabcolsep}{4pt}
\begin{tabular}{@{}lllll@{}}
\toprule
Row / challenge & Radius range $\delta$ & Verdict & Certified numerator & Source \\
\midrule
\multicolumn{5}{@{}l}{\emph{KoalaBear sextic}: $n=2^{21}$, $k=2^{20}$, $q=p^6>2^{185.9}$, $\varepsilon^*=2^{-128}$, $\varepsilon^*q\approx2^{57.9}$}\\
\addlinespace[2pt]
MCA & $[\tfrac12-2^{-7},\,\tfrac12)$ & unsafe & $>2^{-22}\,q\approx2^{164}$ & \cref{cor:deployed} \\
MCA & $\delta\le349525/2^{21}\approx0.16666$ & safe & $\le349526<2^{-167}q$ & \cref{thm:deep-mca} \\
MCA & $(0.16667,\,0.49218)$ & open & --- & the broad band formulation \cref{prob:band} \\
List (any $m$) & $\delta\ge\tfrac12-2^{-8}$ & unsafe & $\ge\binom{256}{129}/p\ge2^{216}$ & \cref{lem:fiber}(i) \\
List (any $m$) & $\delta\le0.29289$ \ ($a\ge1482910$) & safe & $\le2^{20}$ & \cref{thm:johnson-list} \\
List (any $m$) & $(0.29290,\,0.49609)$ & open & --- & the broad band formulation \cref{prob:band} \\
\midrule
\multicolumn{5}{@{}l}{\emph{M31 circle line round}: $n=2^{21}$, $k=2^{20}$, $q=p^4<2^{124}$; at $2^{-128}$: $\delta^*=0$ (\cref{prop:small-field}); below: $\varepsilon^*=2^{-100}$}\\
\addlinespace[2pt]
MCA & $[\tfrac12-2^{-7},\,\tfrac12)$ & unsafe & $>2^{-22}\,q\approx2^{102}$ & \cref{cor:circle-deployed} \\
MCA & $\delta\le0.16666$ & safe & $\le349526<2^{-105}q$ & \cref{thm:deep-mca} \\
List (any $m$) & $\delta\ge\tfrac12-2^{-8}$ & unsafe & $\ge2^{216}$ & \cref{lem:phi-fiber}(i) \\
List (any $m$) & $\delta\le0.29289$ & safe & $\le2^{20}\le2^{-100}q$ & \cref{thm:johnson-list} \\
\bottomrule
\end{tabular}
\caption{Baseline scanner output (\cref{def:scanner}) for the two deployed rows, before the pincer and the optimized graded-prefix floors; it is retained to display the self-contained deep-safe and Johnson-list-safe certificates, and the final deployed threshold bands are the narrowed bands of \cref{tab:scanner2}. The open MCA band for the circle row at $\varepsilon^*=2^{-100}$ is the same interval $(0.16667,\,0.49218)$ as above.}
\label{tab:scanner}
\end{table}

The table realizes the interval output demanded by the staircase interface: for the KoalaBear MCA staircase, $a_{\rm unsafe}=1{,}064{,}960$ \textup($\delta=\frac12-2^{-7}$\textup) and $a_{\rm safe}=1{,}747{,}627$ \textup($\delta\approx0.16666$\textup); the interval shrinks to one step exactly when the broad band formulation \cref{prob:band} is settled, via \cref{thm:conditional-mca} whose other inputs are now discharged \textup(\cref{rem:conditional-status}\textup).

\subsection{The formal finite core}\label{sec:discharge-T7}

\begin{proposition}[finite certificates and decidability]\label{prop:finite-certificate}
Fix a row $(C,q,q_{\rm line})$ and an agreement $a$. The numerator $B_{\rm mca}(a)$ of \cref{sec:prize-program} is a well-defined integer, computable by finite enumeration over pairs, slopes, supports, and codewords; hence the one-step certificate of \cref{prop:onestep} is a finite object. A verifiable prize certificate consists of: the row header; the integer $a_0$; a \emph{lower witness} --- a pair $(f_1,f_2)$, a set of more than $\varepsilon^*q_{\rm line}$ slopes, and for each slope a codeword and a support, checkable by evaluation --- or a theorem tag from this paper with its numerically verified hypotheses; and an \emph{upper witness} --- a theorem tag with numerically verified hypotheses certifying $B_{\rm mca}(a_0+1)\le\varepsilon^*q_{\rm line}$. Verification of a certificate is polynomial in its length.
\end{proposition}

\begin{proof}
All quantifiers in \cref{def:ca,def:mca} range over the finite sets $\F^n$, $\F$, $2^D$, and $C$; the maximum defining $B_{\rm mca}(a)$ is therefore attained and computable. Checking a lower witness is evaluation and counting; checking a theorem tag is checking finitely many integer inequalities.
\end{proof}

\begin{table}[ht]
\centering
\small
\begin{tabular}{@{}ll@{}}
\toprule
Core statement & Proof ingredients \\
\midrule
\cref{thm:A} (deep-point conversion) & polynomial division, pigeonhole over $\F\setminus D$ \\
\cref{lem:fiber}, \cref{lem:phi-fiber} (fiber floors) & elementary symmetric functions, subfield pigeonhole \\
\cref{fact:chain}, \cref{lem:mca-monotone} & set manipulation \\
\cref{prop:onestep} (staircase) & monotone integer staircase \\
\cref{thm:deep-mca} (deep MCA) & linear interpolation, double counting of vanishings \\
\cref{thm:johnson-list} (Johnson lists) & convexity double counting of incidences \\
\cref{thm:stabilizer-partition}, \cref{thm:fiber-descent} & cyclic group actions, root counting \\
\cref{thm:weil-lines} (restriction of scalars) & basis decomposition of $\F/\B$ \\
\cref{def:scanner}, \cref{thm:scanner-sound} & finite case analysis over the statements above \\
\bottomrule
\end{tabular}
\caption{The formal finite core: each statement quantifies over explicitly bounded finite sets and uses only the listed elementary ingredients; the dependency order is top to bottom. Formalizing this table in a proof assistant is the outstanding transcription task.}
\label{tab:core}
\end{table}

The list deliberately includes the new safe-side theorems: a certificate package that contains both an unsafe certificate \textup(\cref{cor:grand} and its relatives\textup) and a machine-checked safe region \textup(\cref{thm:deep-mca,thm:johnson-list}\textup) documents the staircase from both sides in the sense of \cref{prop:onestep}.

\subsection{Tower transfer and the genus-one first step}\label{sec:discharge-T8}

The tower-transfer ledger half is already discharged: \cref{thm:stabilizer-partition,thm:fiber-descent,cor:periodic-support-count} were stated and proved for the Chebyshev case alongside the multiplicative case, with the divisor lattice of $D$ replaced by the dyadic scale lattice of the twin coset, exactly as the tower transfer requires; and the broad band formulation \cref{prob:band} is likewise stated for both cases. It remains to treat the genus-one floor. We first extend \cref{def:map-smooth} to rational maps.

\begin{definition}[rationally smooth domain]\label{def:rational-smooth}
Let $\B\subseteq\F$, let $\psi=f/g$ with $f,g\in\B[X]$ coprime, $a:=\deg f>e:=\deg g\ge0$, and let $D\subseteq\B$ with $g$ nonvanishing on $D$. The set $D$ is \emph{$(\psi,a)$-smooth} if every nonempty fiber $\psi^{-1}(w)\cap D$, $w\in\psi(D)$, has exactly $a$ elements; then $N:=|D|/a=|\psi(D)|$. The FFTree domains of ECFFT \cite{BCKL21} are $(\psi,2)$-smooth for the $x$-coordinate map $\psi$ of each $2$-isogeny in the chain, of degree type $(a,e)=(2,1)$.
\end{definition}

\begin{proposition}[rational-map fiber floor; unified $(a,e)$ form]\label{prop:rational-floor}
Let $D\subseteq\B$ be $(\psi,a)$-smooth of size $n$, with $\psi=f/g$ of degree type $(a,e)$, and let $2\le\ell\le N-1$. Put
\[
k'\ :=\ a\ell-2(a-e),
\qquad
u_z\ :=\ \bigl(f(x)^{\ell}+z\,f(x)^{\ell-1}g(x)\bigr)_{x\in D}\in\B^n .
\]
Then there exists $z\in\B$ with
\[
\bigl|\Lst\bigl(\RS[\F,D,k'+1],\,1-\tfrac{a\ell}n,\,u_z\bigr)\bigr|\ \ge\ \binom N{\ell}\Big/|\B| ,
\]
i.e.\ a list at agreement $a\ell=k'+2(a-e)$. For $e=0$ this is \cref{lem:phi-fiber}\textup{(ii)} \textup(window $2a$, gap $\approx2/N$\textup); for $e=a-1$ the window is $2$ and the certified agreement is $k'+2$ --- the first staircase step, in accordance with the collapse of \cref{rem:ecfft-obstruction}.
\end{proposition}

\begin{proof}
For $A\subseteq\psi(D)$ with $|A|=\ell$ put $\Lambda_A:=\prod_{b\in A}(f-bg)$, of degree $a\ell$ with leading coefficient $\mathrm{lc}(f)^\ell\ne0$, vanishing at every $x\in D$ with $\psi(x)\in A$ --- exactly $a\ell$ points by smoothness \textup(and $g(x)\ne0$\textup). Expanding,
\[
\Lambda_A=\sum_{j=0}^{\ell}(-1)^je_j(A)\,f^{\ell-j}g^{\,j},
\]
with $j$-th term of degree $a\ell-j(a-e)$. Set $z_A:=-\eone(A)\in\B$ and
\[
c_A:=\Bigl(-\sum_{j\ge2}(-1)^je_j(A)\,f^{\ell-j}(x)\,g^{\,j}(x)\Bigr)_{x\in D},
\]
a polynomial of degree at most $a\ell-2(a-e)=k'$, so $c_A\in\RS[\B,D,k'+1]$; on the $a\ell$ fiber points, $u_{z_A}=c_A$. For injectivity at fixed $z$: if $c_A=c_{A'}$ as words then as polynomials \textup($\deg\le k'\le a(N-1)<n$\textup), so $\sum_{j\ge2}(-1)^j\bigl(e_j(A)-e_j(A')\bigr)f^{\ell-j}g^{\,j}=0$; dividing by $g^{\ell}$ in $\F(X)$ gives a polynomial identity in $\psi$, and $\psi$ is nonconstant, hence transcendental over $\F$ inside $\F(X)$, so all coefficients vanish: $e_j(A)=e_j(A')$ for $j\ge2$, and $\eone$ agrees via $z$; then $A=A'$ from the root sets of $\prod_{b}(Y-b)$. Pigeonholing $z_A$ over $\B$ finishes as in \cref{lem:phi-fiber}.
\end{proof}

\begin{corollary}[every ECFFT row has an unsafe first staircase step]\label{cor:ecfft-onestep}
Let $D$ be an ECFFT domain that is $(\psi,2)$-smooth \textup(\cref{def:rational-smooth}\textup) over $\B$, let $C=\RS[\F,D,k]$ with $2\mid k$, $\rho=k/n\in[\tfrac1{16},\tfrac12]$, $n<q=|\F|<2^{256}$, and $n\ge2^{12}$. Then, with $\ell=(k+2)/2\le N-1$,
\[
\binom{n/2}{(k+2)/2}\ \ge\ |\B|\Bigl(\frac qk+1\Bigr),
\]
and consequently
\[
\eca\Bigl(C,\,1-\rho-\tfrac2n\Bigr)\ >\ \frac1{2k}\Bigl(1-\frac nq\Bigr),
\qquad
\emca(C,\delta)\ >\ \frac1{2k}\Bigl(1-\frac nq\Bigr)\ \ \text{for every }\delta\in\Bigl[1-\rho-\tfrac2n,\,1-\rho\Bigr),
\]
so $\delta^*_C(\varepsilon^*)\le1-\rho-2/n$ for every $\varepsilon^*\le\frac1{2k}(1-\frac nq)$; and, provided $k<2^{128}$, the list challenge is likewise unsafe at that step, since $\binom{n/2}{(k+2)/2}/|\B|\ge q/k>2^{-128}q$. For odd $k$ the same argument with $\ell=\lfloor k/2\rfloor+1$ gives the step at $1-\rho-1/n$.
\end{corollary}

\begin{proof}
Apply \cref{prop:rational-floor} with $(a,e)=(2,1)$: $k'=2\ell-2=k$, so the list lives in $\RS[\F,D,k+1]=C^+$ at radius $\delta_0:=1-(k+2)/n$, and $\lfloor\delta_0n\rfloor=n-k-2\le n-k-1$ is admissible for \cref{thm:A}. The entropy count: $\ell/(n/2)=(k+2)/n\in(\rho,\rho+2^{-11}]\subseteq[\tfrac1{16},0.502]$, so $\Hb(\ell/(n/2))\ge\Hb(\tfrac1{16})\ge0.3372$ by the endpoint argument of \cref{cor:mixed-radix}, and
\[
\log_2\binom{n/2}{\ell}\ \ge\ \frac n2\,\Hb\Bigl(\frac{\ell}{n/2}\Bigr)-\log_2\Bigl(\frac n2+1\Bigr)\ \ge\ 0.3372\cdot\frac n2-\log_2\Bigl(\frac n2+1\Bigr)\ \ge\ 690-12\ \ge\ 513,
\]
since the middle expression is increasing in $n$ and equals $690.6\ldots-11.0\ldots$ at $n=2^{12}$, while $|\B|(q/k+1)\le q^2+q<2^{513}$. The contradiction with \cref{thm:A} at $\eta=\frac12$, then \cref{fact:chain} and \cref{lem:mca-monotone}, run exactly as in \cref{thm:phi-cap}. The list-challenge clause follows because $\binom{n/2}{\ell}/|\B|\ge q/k+1>2^{-128}q$ whenever $k<2^{128}$, which covers every challenge-envelope row.
\end{proof}
\subsection{The two-sided picture and what remains}\label{sec:discharge-status}

\begin{theorem}[two-sided sandwich for the grand MCA challenge]\label{thm:sandwich}
Let $C=\RS[\F,D,k]$ be a row for which one of the capacity caps of this paper applies at target $\varepsilon^*=2^{-128}$ --- \cref{cor:grand}, \cref{cor:mixed-radix}, \cref{cor:circle-grand} or \cref{cor:circle-deployed}, or \cref{cor:ecfft-onestep} --- with printed gap $G$ \textup(respectively $2^{-9}$/$2^{-10}$, $2^{-11}$/$2^{-12}$, $2^{-9}$--$2^{-11}$ or $2^{-7}$, and $2/n$\textup). Put
\[
r^*\ :=\ \min\Bigl(\Bigl\lfloor\frac{n-k}3\Bigr\rfloor,\ \bigl\lfloor2^{-128}q\bigr\rfloor-1\Bigr).
\]
If $q>2^{128}$ then
\[
\frac{r^*}n\ \le\ \delta^*_C(2^{-128})\ \le\ 1-\rho-G ,
\]
so, whenever $\lfloor2^{-128}q\rfloor-1\ge\lfloor(n-k)/3\rfloor$ --- e.g.\ whenever $q\ge2^{128}n$, as at every deployed multiplicative row --- the grand MCA threshold of such a row is determined up to a multiplicative factor of three, up to the microscopic gap $G$; for $q$ nearer $2^{128}$ the lower edge degrades to $r^*/n$ and no factor-of-three claim is made. If $q<2^{128}$ then $\delta^*_C(2^{-128})=0$ \textup(\cref{prop:small-field}\textup); at the boundary $q=2^{128}$ the pair of \cref{prop:small-field} gives $\emca\ge2^{-128}$ at every radius, which meets the threshold condition with equality and does not force degeneracy.
\end{theorem}

\begin{proof}
Lower bound: $3r^*\le n-k$ and $r^*+1\le\lfloor2^{-128}q\rfloor\le2^{-128}q$, so \cref{thm:deep-mca} gives $\emca(C,r^*/n)\le(r^*+1)/q\le2^{-128}$, and $r^*/n<(n-k)/n=1-\rho$, so $\delta^*\ge r^*/n$ by \cref{def:dstar} (the case $r^*=0$ being trivial). Upper bound: the applicable cap certifies $\emca(C,\delta)>2^{-128}$ for every $\delta\in[1-\rho-G,\,1-\rho)$, so no such $\delta$ belongs to the supremum set, and by \cref{lem:mca-monotone} neither does any larger sub-capacity radius; hence $\delta^*\le1-\rho-G$. The degenerate clause is \cref{prop:small-field} with $\varepsilon^*=2^{-128}<1/q$; the boundary remark is the same construction, whose error floor $1/q$ equals the target exactly at $q=2^{128}$.
\end{proof}

For $\rho=\tfrac12$ and $q\ge2^{128}n$ --- e.g.\ the KoalaBear-sextic row, $q>2^{185}$ --- the sandwich reads $\delta^*_C(2^{-128})\in[0.1666\ldots,\,\tfrac12-G]$: the first two-sided determination of a grand-challenge MCA threshold, exact to within a factor of three.

\begin{table}[ht]
\centering
\small
\begin{tabular}{@{}p{2.7cm}p{7.9cm}l@{}}
\toprule
Component & Outcome & Where \\
\midrule
Periodic-support ledger & Stabilizer exhaustion is exact at the support level; for fiber-constant lines it descends to an exact quotient-image ledger; upper numerators are independent of $q$ and $|\B|$. & \cref{sec:discharge-T1} \\
Deep MCA core & MCA is safe below one third of the distance; singular and rank-drop Hankel branches are removed in that range; the remaining band is the broad band formulation \cref{prob:band}, with rank-pivot and compressed-eliminant certificates. & \cref{sec:discharge-T2,thm:deep-mca} \\
Extension-line lift & Extension-valued lines are handled exactly by restriction of scalars; there is no independent extension bucket. & \cref{sec:discharge-T3} \\
Interleaved-list side & Johnson safety is uniform in the interleaving arity; fiber-constant recursion is exact; the quotient term is necessarily exponential in the band. & \cref{sec:discharge-T4} \\
Curve/projective samplers & The deep regime is controlled for degree-$d$ curves and projective parameters; residual curve charts are part of the same aperiodic band problem. & \cref{sec:discharge-T5} \\
Scanner & The finite staircase scanner is specified, proved sound, and run on the deployed rows. & \cref{sec:discharge-T6} \\
Certificate core & The finite certificate format, decidability statement, and dependency chart are supplied; machine formalization is separate implementation work. & \cref{sec:discharge-T7} \\
Tower transfer & The ledger half is transferred to the Chebyshev tower; the genus-one first-step cap is proved, later strengthened to a macroscopic cap by the rational graded floors. & \cref{sec:discharge-T8,sec:answers-isogeny} \\
\bottomrule
\end{tabular}
\caption{Structural and certificate ingredients used in the two-sided threshold statements.}
\label{tab:status}
\end{table}

\begin{remark}[dependency status of the one-step criteria]\label{rem:conditional-status}
The one-step MCA criterion of \cref{thm:conditional-mca} has the following dependency status.  The periodic-support input is supplied by \cref{cor:periodic-support-count} at the support level and by \cref{thm:fiber-descent} at the image level for fiber-constant lines.  The deep-regime input is supplied by \cref{thm:deep-mca,cor:singular-deep}; outside that range it is exactly the broad band formulation \cref{prob:band}, represented by the residual rank-pivot and kernel-compressed certificate language of \cref{thm:residual-rank-pivot-cover,thm:kernel-compressed-affine-eliminant,thm:exact-arbitrary-residual-image-ledger}.  The extension-line input is supplied by \cref{thm:weil-lines}.  The sampler input is definitional once the protocol's line sampler is fixed.  For the one-step list criterion of \cref{thm:conditional-list}, the required base-code bound is unconditional throughout the Johnson regime and uniform in $m$ \textup(\cref{thm:johnson-list}, making the codegree recursion unnecessary there\textup), is exactly recursive for fiber-constant words \textup(\cref{thm:list-recursion}\textup), and is provably exponential in the band \textup(\cref{prop:list-calibration}\textup), so any resolution must take that shape.
\end{remark}

After these ingredients are assembled, the remaining mathematical gap has a compact form.  The central residue is the broad band formulation \cref{prob:band}, the aperiodic band conjecture, now represented by the finite rank-pivot, compressed-eliminant, and exact-image certificate language of \cref{thm:residual-rank-pivot-cover,thm:kernel-compressed-affine-eliminant,thm:exact-arbitrary-residual-image-ledger}; this is the analytic kernel separating the two-sided interval of \cref{thm:sandwich} from a one-step threshold determination.  The other residues are narrower: the failure-profile upgrade from error $1/k$ to error $1-o(1)$ \textup(the error-one question \cref{prob:errorone}\textup) and explicit deployed-size certifying pairs \textup(the explicit-witness question \cref{prob:explicit}\textup); the genus-one band, formerly a separate residue, is settled at macroscopic scale in \cref{sec:answers-isogeny}.  The next section attacks the main band directly, from both sides at once.

\section{Safe-side pincer and scale-optimized unsafe floors}\label{sec:pincer}

The residue identified at the end of \cref{sec:discharge} is a single open band per row. This section attacks the band from both sides and narrows it substantially, with three new results. On the safe side, we prove that \emph{mutual} correlated agreement coincides with plain correlated agreement, up to an explicit tangent term, all the way up to half the minimum distance (\cref{thm:mca-from-ca}); this eliminates the support-mismatch layer of the problem on that range and, combined with the imported unique-decoding proximity gap of Ben-Sasson--Carmon--Ishai--Kopparty--Saraf \cite{BCIKS20}, moves the certified safe frontier of every deployed row from one third to one half of the distance. We also prove a self-contained correlated-agreement bound at half the Johnson radius (\cref{thm:elementary-ca}), which improves the \emph{unconditional} safe frontier for every row of rate below $\tfrac14$. On the unsafe side, we extend the graded locator-prefix pigeonhole of \cref{lem:quotient-remainder-prefix} to every map-smooth scale --- in particular to the Chebyshev scales of the circle tower --- and optimize it over the scale (\cref{prop:graded-prefix-floor}); at deployed sizes this widens the certified unsafe band by a factor of four with a hand-checkable certificate, and by a factor of about sixteen with an exactly verified one, saturating an intrinsic entropy--subfield envelope that we identify in closed form (\cref{rem:entropy-frontier}). \Cref{thm:sandwich2} and \cref{tab:scanner2} assemble the narrowed sandwich.

\subsection{Mutual equals correlated agreement up to half the distance}\label{sec:pincer-half}

The deep-regime theorem (\cref{thm:deep-mca}) controls $\emca$ directly, but only below one third of the distance: its alternatives identify every close point-codeword with a value of a single interpolated line, and that identification costs a factor of three in the radius. The next theorem decouples the two layers of the problem. It shows that below \emph{half} the distance, the support-mismatch phenomenon that separates mutual from plain correlated agreement is carried entirely by the tangent slopes of \cref{thm:deep-mca}\textup{(a)}, of which there are at most $\lfloor\delta n\rfloor$; everything else is already visible to $\eca$.

\begin{theorem}[mutual from correlated agreement, up to half the distance]\label{thm:mca-from-ca}
Let $C\subseteq\F^n$ be any $\F$-linear code of minimum Hamming weight $w$, let $q=|\F|$, let $\delta\in(0,1)$, and put $r:=\lfloor\delta n\rfloor$. Assume
\[
2r\ \le\ w-1 .
\]
Then for every pair $(f_1,f_2)\in\F^n\times\F^n$:
\begin{enumerate}[label=\textup{(\alph*)}]
\item if $\dist_2\bigl((f_1,f_2),C^{\equiv2}\bigr)\le\delta$, the pair has \emph{no} CA-bad slope at radius $\delta$, and at most $r$ MCA-bad slopes at radius $\delta$: writing $(p_1,p_2)$ for a pair of codewords column-agreeing with $(f_1,f_2)$ outside a set $E$ with $|E|\le r$ and $e_i:=f_i-p_i$, the MCA-bad slopes form a subset of $\{-e_{1,j}/e_{2,j}:\ j\in E,\ e_{2,j}\ne0\}$;
\item if $\dist_2\bigl((f_1,f_2),C^{\equiv2}\bigr)>\delta$, every MCA-bad slope at radius $\delta$ is CA-bad at radii $(\delta,\delta)$.
\end{enumerate}
Consequently
\[
\emca(C,\delta)\ \le\ \max\Bigl(\,\eca(C,\delta),\ \frac rq\,\Bigr).
\]
For $C=\RS[\F,D,k]$ on any evaluation set $D$ of size $n$, the hypothesis reads $2\lfloor\delta n\rfloor\le n-k$: mutual and plain correlated agreement coincide, up to the explicit tangent term $r/q$, on the entire unique-decoding range of the interleaved code $C^{\equiv2}$.
\end{theorem}

\begin{proof}
Part \textup{(b)} is immediate from the definitions: an MCA-bad slope $\gamma$ has a witness set $S$ with $|S|\ge\lceil(1-\delta)n\rceil=n-r$ on which $f_1+\gamma f_2$ agrees with a codeword, so $\dist(f_1+\gamma f_2,C)\le r/n\le\delta$; if additionally the pair is $\delta$-far in column distance, $\gamma$ is CA-bad at radii $(\delta,\delta)$ by \cref{def:ca}.

For part \textup{(a)}, fix codewords $(p_1,p_2)$ column-agreeing with $(f_1,f_2)$ outside the column error set $E:=\{j:\ (e_{1,j},e_{2,j})\ne(0,0)\}$, $e_i:=f_i-p_i$, $|E|\le r$. The pair is $\delta$-close, so it has no CA-bad slope. Let $\gamma$ be MCA-bad with witness $(S,c)$: $|S|\ge n-r$, $c\in C$, $c=f_1+\gamma f_2$ on $S$, and no pair of codewords of $C$ explains $(f_1,f_2)$ on $S$.

First, $c=p_1+\gamma p_2$. Indeed, $z:=c-(p_1+\gamma p_2)\in C$ vanishes on $S\setminus E$: there $f_i=p_i$, so $f_1+\gamma f_2=p_1+\gamma p_2$, while $c=f_1+\gamma f_2$ on all of $S$. Since $|S\setminus E|\ge n-2r$, the codeword $z$ has weight at most $2r\le w-1$, hence $z=0$.

Next, $S\not\subseteq D\setminus E$: otherwise $(p_1,p_2)$ would explain the pair on $S$. So there is $j\in S\cap E$. On $S$ we have $0=(f_1+\gamma f_2)-c=(e_1+\gamma e_2)$, so at this $j$,
\[
e_{1,j}+\gamma\,e_{2,j}=0,\qquad (e_{1,j},e_{2,j})\ne(0,0).
\]
If $e_{2,j}=0$ the left side is the nonzero constant $e_{1,j}$, which cannot vanish; hence $e_{2,j}\ne0$ and $\gamma=-e_{1,j}/e_{2,j}$. The MCA-bad slopes of the pair therefore form a subset of the at most $|E|\le r$ listed values. Taking the maximum over pairs gives the displayed bound.
\end{proof}

\begin{remark}[scope and sharpness]\label{rem:half-scope}
Three comments. \textup{(i)} The hypothesis $2r\le w-1$ is exactly the unique-decoding radius of the interleaved code $C^{\equiv2}$ in column weight, and the proof uses it twice, both times through the statement ``a codeword of weight $\le 2r$ vanishes''; beyond half the distance a $\delta$-close pair can admit several column explanations with incomparable agreement sets, and witness codewords need not ride any explanation line, which is precisely where the band begins. \textup{(ii)} The theorem is an unconditional \emph{implication}, with no list bound and no field-size hypothesis; instantiated with the self-contained $\eca$ bound of \cref{thm:deep-mca} it returns the deep theorem (with the marginally better error $\max((r+1)/q,r/q)$ replaced by $(r+1)/q$), and its value is that it accepts \emph{any} stronger correlated-agreement input on the range $(w-1)/3<2r\le w-1$, where no self-contained input is available. \textup{(iii)} The same proof gives the two-radius statement: for $\delta_{\mathrm{int}}\ge\delta$, every MCA-bad slope of a $\delta_{\mathrm{int}}$-far pair is CA-bad at radii $(\delta,\delta_{\mathrm{int}})$, while a $\delta_{\mathrm{int}}$-close pair that is also $\delta$-close contributes at most $r$ slopes; only pairs in the annulus between the two radii require the no-loss form.
\end{remark}

\begin{corollary}[the band conjecture loses its support-mismatch layer below half the distance]\label{cor:band-reduction}
Let $C=\RS[\F,D,k]$ and let $\delta$ satisfy $2\lfloor\delta n\rfloor\le n-k$. Then any correlated-agreement bound $\eca(C,\delta)\le\varepsilon$ implies $\emca(C,\delta)\le\max(\varepsilon,\lfloor\delta n\rfloor/q)$. In particular, on the agreement subrange $a\ge\lceil(n+k)/2\rceil$ of the broad band formulation \cref{prob:band}, the aperiodic band conjecture reduces to its correlated-agreement form: column-close pairs contribute at most $\lfloor\delta n\rfloor$ explicit tangent slopes, and what must be bounded polynomially is only, for each column-far pair, the number of slopes $\gamma$ for which $f_1+\gamma f_2$ has an aperiodically witnessed agreement set of size at least $\lceil(1-\delta)n\rceil$.
\end{corollary}

\begin{proof}
Immediate from \cref{thm:mca-from-ca}: case \textup{(a)} pairs contribute at most $r=\lfloor\delta n\rfloor$ slopes, each explicitly a tangent ratio; case \textup{(b)} identifies every remaining MCA-bad slope, together with its witness support, as a CA-bad slope with the same witness.
\end{proof}

\subsection{A self-contained correlated-agreement bound at half the Johnson radius}\label{sec:pincer-ca}

\Cref{thm:mca-from-ca} converts the safe side of the band problem into a demand for correlated-agreement inputs. The strongest known inputs are the proximity-gap theorems of \cite{BCIKS20} (unique-decoding range with error $n/q$; Johnson range with error $\operatorname{poly}(n)/q$), whose proofs run through a Berlekamp--Welch decoder over the rational function field $\F(Z)$ and an induction over curves, and the recent strengthening \cite{BCHKS25}. Those we import. But the certificate interface of \cref{sec:prize-program} also prices \emph{self-contained} inputs, since the formal core (\cref{tab:core}) is built from elementary counting; the following theorem is the elementary member of the family. Its only ingredients are the pair-interleaved Johnson bound of \cref{thm:johnson-list} and a rider double count, and it reaches half the Johnson radius with an explicit list-size error.

\begin{theorem}[elementary correlated agreement at half the Johnson radius]\label{thm:elementary-ca}
Let $C=\RS[\F,D,k]$ on any evaluation set $D$ of size $n$, let $\delta\in(0,1)$, put $r:=\lfloor\delta n\rfloor$, and assume
\[
n-2r\ >\ 0
\qquad\text{and}\qquad
(n-2r)^2\ >\ (k-1)\,n .
\]
Then, with
\[
L_2\ :=\ \frac{n\,\bigl(n-2r-k+1\bigr)}{(n-2r)^2-(k-1)n}
\]
the pair-interleaved Johnson bound of \cref{thm:johnson-list} at agreement $n-2r$,
\[
\eca(C,\delta)\ \le\ \frac{1+(r+1)\,L_2}{q}\,,
\]
and the same bound holds for $\eca(C,\delta,\delta_{\mathrm{int}})$ for every $\delta_{\mathrm{int}}\ge\delta$.
\end{theorem}

\begin{proof}
Fix a pair $(f_1,f_2)$ with $\dist_2((f_1,f_2),C^{\equiv2})>\delta_{\mathrm{int}}\ge\delta$; only such pairs have CA-bad slopes. Let $B\subseteq\F$ be the set of CA-bad slopes; for each $\gamma\in B$ fix a codeword $c_\gamma\in C$ and an agreement set $S_\gamma$ with $|S_\gamma|\ge n-r$ and $c_\gamma=f_1+\gamma f_2$ on $S_\gamma$. If $|B|\le1$ we are done, so fix $\gamma_1\in B$ and consider any $\gamma\in B\setminus\{\gamma_1\}$. Define
\[
P_2:=\frac{c_{\gamma_1}-c_{\gamma}}{\gamma_1-\gamma}\in C,
\qquad
P_1:=c_{\gamma_1}-\gamma_1P_2\in C,
\]
so that, identically as codewords,
\begin{equation}\label{eq:rider-identity}
P_1+\gamma_1P_2=c_{\gamma_1},
\qquad
P_1+\gamma P_2=c_{\gamma}.
\end{equation}
On $S_{\gamma_1}\cap S_{\gamma}$, which has size at least $n-2r$, the two coordinatewise linear equations $f_1+\gamma_1f_2=c_{\gamma_1}$ and $f_1+\gamma f_2=c_{\gamma}$ have the unique solution $f_2=P_2$, $f_1=P_1$; hence the codeword pair $(P_1,P_2)$ column-agrees with $(f_1,f_2)$ on at least $n-2r$ coordinates, i.e.
\[
(P_1,P_2)\ \in\ \mathcal L\ :=\ \Lst\Bigl(C^{\equiv2},\,1-\tfrac{n-2r}{n},\,(f_1,f_2)\Bigr),
\qquad
|\mathcal L|\le L_2
\]
by \cref{thm:johnson-list} with $m=2$ and $a=n-2r$. Say that $\gamma$ \emph{rides} a member $(P_1,P_2)\in\mathcal L$ if $c_\gamma=P_1+\gamma P_2$ as codewords; by \eqref{eq:rider-identity}, every $\gamma\in B\setminus\{\gamma_1\}$ rides some member of $\mathcal L$.

Fix $(P_1,P_2)\in\mathcal L$ and let $E:=\{j:\ (f_{1,j},f_{2,j})\ne(P_{1,j},P_{2,j})\}$ be its column error set. Since the pair is $\delta$-far, $|E|\ge r+1$. If $\gamma$ rides $(P_1,P_2)$, then $(f_1-P_1)+\gamma(f_2-P_2)=(f_1+\gamma f_2)-c_\gamma$ vanishes on $S_\gamma$, hence on at least $|E|-r\ge1$ coordinates of $E$; and at each $j\in E$ the map $\gamma\mapsto(f_{1,j}-P_{1,j})+\gamma(f_{2,j}-P_{2,j})$ is a nonzero affine function, vanishing for at most one slope. Double counting vanishing incidences between riders of $(P_1,P_2)$ and coordinates of $E$,
\[
\#\{\text{riders of }(P_1,P_2)\}\cdot\bigl(|E|-r\bigr)\ \le\ |E|,
\qquad\text{so}\qquad
\#\{\text{riders}\}\ \le\ \frac{|E|}{|E|-r}\ \le\ r+1 .
\]
Summing over $\mathcal L$ and adding back $\gamma_1$: $|B|\le1+(r+1)L_2$. Taking the maximum over pairs completes the proof.
\end{proof}

\begin{remark}[provenance and comparison]\label{rem:ca-provenance}
The argument is elementary and in the folklore spirit ``a list bound implies a proximity gap with a list-size loss''; we do not claim novelty of the mechanism, only of the exact constants in this normalization and of its place in the formal core: its only inputs are \cref{thm:johnson-list} and counting, so it enters the dependency chart of \cref{tab:core} without enlarging the trusted base. The theorems of \cite{BCIKS20,BCHKS25} are strictly stronger where they apply --- they reach the full Johnson radius $1-\sqrt\rho$ with error $\operatorname{poly}(n)/q$ --- at the price of a substantially deeper proof. Note also the built-in factor of two: the recovered pair $(P_1,P_2)$ lives at column radius $2r$, so the list bound is invoked at agreement $n-2r$, which is why the theorem stops at half the Johnson radius rather than at the Johnson radius.
\end{remark}

\begin{corollary}[self-contained safe frontier at half the Johnson radius]\label{cor:self-contained-safe}
Let $C=\RS[\F,D,k]$ with $n\ge k+2$, and let $\delta$, $r$, $L_2$ be as in \cref{thm:elementary-ca}, so $n-2r>0$ and $(n-2r)^2>(k-1)n$. Then
\[
\emca(C,\delta)\ \le\ \frac{1+(r+1)L_2}{q}\,.
\]
For the challenge rows this improves the self-contained safe frontier exactly when $\rho<\tfrac14$: the frontier moves from $(1-\rho)/3$ to $\approx(1-\sqrt\rho\,)/2$, i.e.\ from $0.2917$ to $0.3232$ at $\rho=\tfrac18$ and from $0.3125$ to $0.3750$ at $\rho=\tfrac1{16}$, while at $\rho\in\{\tfrac12,\tfrac14\}$ the deep frontier $(1-\rho)/3$ remains the better self-contained certificate. At the deployed KoalaBear row the half-Johnson certificate reads: $r=307121$, $\delta=r/n\approx0.14645$, $L_2\le1001282$, error at most $(1+307122\cdot1001282)/q<2^{-147}$ --- dominated by \cref{thm:deep-mca} at this rate, but recorded as an independent certificate of the same verdict.
\end{corollary}

\begin{proof}
The hypothesis gives $n-2r>0$ and $(n-2r)^2>(k-1)n$, hence $n-2r>\sqrt{(k-1)n}$; since $n\ge k+2$ we also have $\sqrt{(k-1)n}\ge k$, so $n-2r>k$, i.e.\ $2r\le n-k-1$: the radius lies below half the distance. \Cref{thm:mca-from-ca} therefore applies, and with the input of \cref{thm:elementary-ca},
\[
\emca(C,\delta)\le\max\Bigl(\tfrac{1+(r+1)L_2}{q},\ \tfrac rq\Bigr)=\tfrac{1+(r+1)L_2}{q}.
\]
The frontier comparison is the elementary inequality $(1-\sqrt\rho)/2>(1-\rho)/3\iff\sqrt\rho<\tfrac12$; the KoalaBear numbers are $n=2^{21}$, $k=2^{20}$, $(n-2r)^2-(k-1)n=909700$ at $r=307121$, and $L_2=\lfloor n(n-2r-k+1)/909700\rfloor\le1001282$, with $(1+307122\cdot1001282)<2^{38.2}$ and $q>2^{185.9}$.
\end{proof}

\begin{corollary}[conditional safe frontier at half the distance]\label{cor:conditional-half}
Let $C=\RS[\F,D,k]$, $\rho=k/n$. Import the unique-decoding proximity gap of \cite[Thm.~4.1]{BCIKS20}: in the normalization of \cref{def:ca}, $\eca(C,\delta)\le n/q$ for every $\delta\le(1-\rho)/2$. (Only $\eca(C,\delta)\le\operatorname{poly}(n)/q$ is used below, so the conclusion is robust to the exact constant.) Then for every $\delta$ with $2\lfloor\delta n\rfloor\le n-k$,
\[
\emca(C,\delta)\ \le\ \max\Bigl(\frac nq,\ \frac{\lfloor\delta n\rfloor}q\Bigr)\ =\ \frac nq\,.
\]
Deployed instantiations: for the KoalaBear sextic row, $\emca(C,\delta)\le n/q<2^{-164}<2^{-128}$ for every $\delta\le\tfrac14$ (the endpoint $r=2^{19}$, $2r=n-k$, is included); for the circle line-round row over $\F_{p^4}$, $p=2^{31}-1$, $\emca\le n/q<2^{-102}<2^{-100}$ for every $\delta\le\tfrac14$. In both cases the certified safe frontier moves from one third to one half of the distance.
\end{corollary}

\begin{proof}
Combine the import with \cref{thm:mca-from-ca}; $\lfloor\delta n\rfloor\le n$ makes the maximum $n/q$. For the numbers: $n/q=2^{21}/p^6<2^{-164.9}$ with $p=2^{31}-2^{24}+1$, and $n/q=2^{21}/(2^{31}-1)^4<2^{-102.9}$. The circle case uses the diagonal equivalence of \cref{lem:circle-rs}, which preserves both correlated-agreement errors exactly, so the import applies verbatim to the equivalent Reed--Solomon code on the torus twin coset.
\end{proof}

\subsection{Widening the unsafe band: graded prefix floors on every scale}\label{sec:pincer-floors}

The deployed caps of \cref{cor:deployed,cor:circle-deployed} certify the unsafe band $[\tfrac12-2^{-7},\tfrac12)$ using only the top \emph{two} slots of the locator expansion at the single top scale ($N=256$). The quotient-remainder ledger already prices deeper prefixes at the multiplicative level (\cref{lem:quotient-remainder-prefix}); two observations remain to be made. First, the graded pigeonhole runs verbatim on \emph{every} map-smooth scale --- the locator expansion $\prod_{b\in M}(\varphi-b)=\sum_j(-1)^je_j(M)\varphi^{m-j}$ is graded by $\deg\varphi$ for any polynomial $\varphi$, in particular for the Chebyshev scales of the circle tower --- which extends the tower ledger from the first staircase step to every depth. Second, the floor can be \emph{optimized over the scale}: small scales trade prefix cost $|\B|$ per slot against fiber entropy $\binom{n/c}{m}$, and the optimum widens the certified unsafe band by an order of magnitude at deployed sizes. The route has an intrinsic ceiling, which we identify exactly in \cref{rem:entropy-frontier}.

\begin{proposition}[graded prefix floor for map-smooth scales]\label{prop:graded-prefix-floor}
Let $\B\subseteq\F$ be a subfield, let $D\subseteq\B$ with $|D|=n$, and let $\varphi\in\B[X]$ of degree $c\ge2$ whose fibers within $D$ are complete and of uniform size $c$ (the map-smooth hypothesis of \cref{lem:phi-fiber}); put $N:=n/c$ and $Q:=\varphi(D)\subseteq\B$, so $|Q|=N$. Let $1\le K\le n$ and let $m$ satisfy $\lceil K/c\rceil\le m\le N$; put
\[
w\ :=\ m-\lceil K/c\rceil\ \ge\ 0 .
\]
Then there is a $\B$-valued received word $U:D\to\B$ with
\[
\bigl|\Lst\bigl(\RS[\F,D,K],\,1-\tfrac{mc}{n},\,U\bigr)\bigr|
\ \ge\
\left\lceil\frac{\binom Nm}{|\B|^{\,w}}\right\rceil .
\]
For $\varphi=X^c$ on a multiplicative coset this is \cref{lem:quotient-remainder-prefix} with $s=0$ (note $w=\lfloor(mc-K)/c\rfloor=m-\lceil K/c\rceil$); the content here is that the same graded pigeonhole runs on every polynomial scale, in particular on the Chebyshev scales $T_c$ of the circle tower (\cref{lem:cheb-fibers}), whose locators are not lacunary in the monomial basis but are perfectly graded in powers of $T_c$.
\end{proposition}

\begin{proof}
For $M\subseteq Q$ with $|M|=m$, the locator $\Lambda_M:=\prod_{b\in M}(\varphi-b)\in\B[X]$ vanishes exactly on the union of the $\varphi$-fibers over $M$ inside $D$, a set of size $mc$, and expands as $\Lambda_M=\sum_{j=0}^m(-1)^je_j(M)\,\varphi^{m-j}$ with $e_j(M)\in\B$ the elementary symmetric functions of $M\subseteq\B$. For $z=(z_1,\dots,z_w)\in\B^w$ define the $\B$-valued word
\[
U_z\ :=\ \Bigl(\varphi^m+\textstyle\sum_{j=1}^{w}z_j\,\varphi^{m-j}\Bigr)\Big|_D .
\]
Assign to each $M$ the vector $z(M):=\bigl((-1)^je_j(M)\bigr)_{1\le j\le w}\in\B^w$; by pigeonhole some $z^*$ is assigned by at least $\binom Nm/|\B|^w$ of the $\binom Nm$ sets. For each such $M$ put
\[
c_M\ :=\ U_{z^*}-\Lambda_M\big|_D\ =\ -\sum_{j=w+1}^{m}(-1)^je_j(M)\,\varphi^{m-j}\Big|_D ,
\]
a polynomial of degree at most $c(m-w-1)=c\bigl(\lceil K/c\rceil-1\bigr)\le K-1$, hence a codeword of $\RS[\F,D,K]$; it agrees with $U_{z^*}$ on the $mc$ zeros of $\Lambda_M$ in $D$, so it lies in the displayed list. The map $M\mapsto c_M$ is injective: if $c_M=c_{M'}$ then $\sum_{j>w}(-1)^j\bigl(e_j(M)-e_j(M')\bigr)\varphi^{m-j}=0$ as a polynomial of degree $<n$ vanishing on $D$, hence identically; the powers $\varphi^{m-j}$ have pairwise distinct degrees $c(m-j)$ and are therefore linearly independent over $\F$, so $e_j(M)=e_j(M')$ for $j>w$, while for $j\le w$ equality holds by the choice of $z^*$; thus $\prod_{b\in M}(Y-b)=\prod_{b\in M'}(Y-b)$ and $M=M'$.
\end{proof}

\begin{corollary}[widened unsafe band at deployed sizes]\label{cor:widened-deployed}
\begin{enumerate}[label=\textup{(\roman*)}]
\item \textup{(KoalaBear sextic, MCA.)} With $\B$, $\F$, $D$, $n=2^{21}$, $k=2^{20}$ as in \cref{cor:deployed},
\[
\emca\bigl(\RS[\F_{p^6}\!,D,2^{20}],\,\delta\bigr)\ >\ 2^{-22}
\qquad\text{for every}\qquad
\delta\in\Bigl[\tfrac{15}{32},\,\tfrac12\Bigr),
\]
a certified gap of $2^{-5}$, four times the gap $2^{-7}$ of \cref{cor:deployed}, by a hand-checkable certificate at scale $c=16$.
\item \textup{(KoalaBear sextic, grand list challenge.)} For every $\delta\in[\tfrac{15}{32},1-\rho)$ there is a $\B$-valued word $U$ with $|\Lst(C,\delta,U)|>2^{-128}q$.
\item \textup{(Circle line-round row.)} With $p'=2^{31}-1$, $\F=\F_{p'^4}$, and the twin-coset $x$-domain of size $n=2^{21}$, $k=2^{20}$ as in \cref{cor:circle-deployed}, the same two statements hold verbatim on $[\tfrac{15}{32},\tfrac12)$ resp.\ $[\tfrac{15}{32},1-\rho)$, at the Chebyshev scale $\varphi=T_{16}$, against the target $2^{-100}$.
\end{enumerate}
\end{corollary}

\begin{proof}
All four claims instantiate \cref{prop:graded-prefix-floor} at scale $c=16$, $N=n/c=2^{17}$, $m=69632$, so $mc=k+2^{16}$ and $m/N=\tfrac{17}{32}$, and convert the floor as in \cref{cor:deployed}. The common entropy estimate is
\[
\log_2\binom{N}{m}\ \ge\ N\,\Hb\bigl(\tfrac{17}{32}\bigr)-\log_2(N+1)\ \ge\ 131072\cdot0.99718-17\ >\ 130{,}685 .
\]
\textup{(i)} Take $K=k+1$, so $w=m-\lceil(k+1)/16\rceil=69632-65537=4095$. The threshold for the deep-point conversion (\cref{thm:A} with $\eta=\tfrac12$, admissible since $n-mc\le n-k-1$) is a floor exceeding $q/k+1$:
\[
w\log_2p+\log_2(q/k+1)\ <\ 4095\cdot30.989+166\ <\ 127{,}116\ <\ 130{,}685 ,
\]
with over $3500$ bits of room. So $\eca(C,1-mc/n)>\tfrac1{2k}(1-n/q)>2^{-22}$, and \cref{fact:chain} with \cref{lem:mca-monotone} extend the same lower bound to $\emca(C,\delta)$ for every $\delta\in[1-mc/n,\tfrac12)=[\tfrac{15}{32},\tfrac12)$.
\textup{(ii)} Take $K=k$, so $w=4096$; the list threshold is $\binom Nm/|\B|^w>2^{-128}q$, i.e.\ $w\log_2p+\log_2q-128<4096\cdot30.989+58<127{,}000<130{,}685$; lists only grow with the radius, giving the full range.
\textup{(iii)} $T_{16}$ is a polynomial of degree $16$ over $\B=\F_{p'}$ with complete uniform fibers on the twin-coset $x$-domain at every dyadic scale (\cref{lem:cheb-fibers}), so \cref{prop:graded-prefix-floor} applies with the same $N$, $m$, $w$; the thresholds are weaker ($\log_2(q/k+1)<104$ and $\log_2q-100<24$), and the conversion runs through the diagonal equivalence of \cref{lem:circle-rs} exactly as in \cref{cor:circle-deployed}.
\end{proof}

\begin{remark}[exactly verified frontier]\label{rem:exact-frontier}
The hand-checkable rows above deliberately stop at the round gap $2^{-5}$. Optimizing $(c,m)$ under exact integer verification of $\binom{N}{m}>p^{\,w}(q/k+1)$ resp.\ $\binom Nm\cdot2^{128}>p^{\,w}q$ --- a finite certificate in the sense of \cref{def:scanner,prop:finite-certificate} --- widens the band further, to the following edges (all verified by exact arithmetic; $\Delta:=mc-k$):
\begin{center}\footnotesize
\begin{tabular}{@{}llllll@{}}
\toprule
row / challenge & scale $c$ & $m$ & $w$ & $\Delta$ & certified unsafe range \\
\midrule
KoalaBear, MCA & $16$ & $69748$ & $4211$ & $67392$ & $\delta\in[\,\tfrac{15331}{32768},\tfrac12)=[0.46786\ldots,\tfrac12)$ \\
KoalaBear, list & $32$ & $34874$ & $2106$ & $67392$ & $\delta\in[\,\tfrac{15331}{32768},1-\rho)$ \\
circle $p'{=}2^{31}{-}1$, MCA & $32$ & $34873$ & $2104$ & $67360$ & $\delta\in[\,\tfrac{30663}{65536},\tfrac12)=[0.46787\ldots,\tfrac12)$ \\
circle $p'{=}2^{31}{-}1$, list & $32$ & $34874$ & $2106$ & $67392$ & $\delta\in[\,\tfrac{15331}{32768},1-\rho)$ \\
\bottomrule
\end{tabular}
\end{center}
The certified gap is thus $\approx0.03214\approx2^{-4.96}$ throughout, versus $2^{-7}$ before.
\end{remark}

\begin{remark}[the entropy--subfield frontier]\label{rem:entropy-frontier}
The optimized route has a closed-form ceiling. At scale $c$ the floor beats the conversion threshold iff
\[
\log_2\binom{n/c}{m}\ \gtrsim\ w\log_2|\B|+\log_2(q/k),
\]
and writing $mc=:(\rho+g)n$ --- so $m/N=\rho+g$ and $w=gN+O(1)$ --- and dividing by $N$, this reads
\[
\Hb(\rho+g)\ \ge\ g\,\log_2|\B|+o(1)
\]
for challenge-envelope field sizes ($\log_2q\le256$, so the additive terms are $O(1)$ against $N\to\infty$). Define
\[
g^*(\rho,\beta)\ :=\ \sup\bigl\{\,g\in(0,1-\rho]:\ \Hb(\rho+g)\ge\beta g\,\bigr\},
\qquad
\beta:=\log_2|\B| .
\]
Every gap $g<g^*$ is certified at large enough deployed sizes, and \emph{no} gap beyond $g^*$ is certifiable by graded prefix floors at any scale: the frontier is intrinsic to the route, not to our optimization. For $\beta=31$ (both deployed base primes): $g^*(\tfrac12)\approx0.03216$, $g^*(\tfrac14)\approx0.02748$, $g^*(\tfrac18)\approx0.01920$, $g^*(\tfrac1{16})\approx0.01238$. The deployed KoalaBear row achieves $0.03214$ at $n=2^{21}$: the envelope is saturated to three digits already at deployed sizes. Two structural consequences: the left edge of the band in the broad band formulation \cref{prob:band} moves from $k+2n/N$ to $k+g^*n-o(n)$; and the frontier depends on the base field only through $\beta$, so larger base fields \emph{narrow} the certifiable unsafe band --- a quantitative form of the subfield-confinement phenomenon of \cref{sec:subfield}.
\end{remark}

\subsection{The narrowed sandwich}\label{sec:pincer-sandwich}

\begin{theorem}[narrowed two-sided sandwich]\label{thm:sandwich2}
Let $\delta^*_C$ be the challenge threshold of \cref{def:dstar}.
\begin{enumerate}[label=\textup{(\roman*)}]
\item \textup{(KoalaBear sextic MCA row, $\varepsilon^*=2^{-128}$.)} Self-contained:
\[
\delta^*_C(2^{-128})\ \in\ \Bigl[\,\tfrac{349525}{2^{21}},\ \tfrac{15331}{32768}\,\Bigr]\ \approx\ [\,0.16666,\ 0.46787\,] .
\]
With the imported unique-decoding proximity gap \textup{(\cref{cor:conditional-half})}:
\[
\delta^*_C(2^{-128})\ \in\ \Bigl[\,\tfrac14,\ \tfrac{15331}{32768}\,\Bigr],
\]
an frontier band of width $<0.218$, down from $0.326$ in \cref{thm:sandwich}.
\item \textup{(Circle line-round row, $\varepsilon^*=2^{-100}$.)} Self-contained: $\delta^*_C\in[\tfrac{349525}{2^{21}},\tfrac{30663}{65536}]$; with the import: $\delta^*_C\in[\tfrac14,\tfrac{30663}{65536}]$.
\item \textup{(General envelope rows.)} For any row with $q\ge2^{128}(n+1)$ and a map-smooth scale lattice: the safe edge is $\max\bigl(\lfloor(n-k)/3\rfloor,\,r_{\mathrm{hJ}}\bigr)/n$ self-contained, where $r_{\mathrm{hJ}}$ is the largest radius passing \cref{cor:self-contained-safe}, and $\lfloor(n-k)/2\rfloor/n$ with the import; the unsafe edge is $1-\rho-g$ for every exactly certified $g$, with envelope $g^*(\rho,\log_2|\B|)$ of \cref{rem:entropy-frontier}.
\end{enumerate}
\end{theorem}

\begin{proof}
Lower edges: \cref{thm:deep-mca} at $r=349525=\lfloor(n-k)/3\rfloor$ (errors $349526/q<2^{-167}$ resp.\ $<2^{-105}$, below the respective targets, as in \cref{tab:scanner}); \cref{cor:self-contained-safe} where it exceeds the deep edge; \cref{cor:conditional-half} for the imported edge $\tfrac14$ (errors $2^{-164}$ resp.\ $2^{-102}$, below the respective targets). Upper edges: \cref{cor:widened-deployed} and \cref{rem:exact-frontier} certify $\emca>2^{-22}\ge\varepsilon^*$ (resp.\ $>2^{-100}$ via the circle thresholds) on $[\,\text{edge},\tfrac12)$, so $\delta^*_C$ cannot exceed the printed edge. Part \textup{(iii)} collects the corresponding general statements.
\end{proof}

\begin{table}[t]
\centering\footnotesize
\setlength{\tabcolsep}{4pt}
\begin{tabular}{@{}llllll@{}}
\toprule
row & challenge & unsafe \textup{(widened)} & safe, self-contained & safe, with import & open band \\
\midrule
KoalaBear sextic & MCA, $2^{-128}$ & $[\,0.46787,\tfrac12)$ & $\delta\le0.16666$ & $\delta\le0.25$ & $(0.25,\ 0.46787)$ \\
KoalaBear sextic & list, $2^{-128}$ & $\delta\ge0.46787$ & $\delta\le0.29289$ & --- & $(0.29290,\ 0.46787)$ \\
circle, $p'{=}2^{31}{-}1$ & MCA, $2^{-100}$ & $[\,0.46788,\tfrac12)$ & $\delta\le0.16666$ & $\delta\le0.25$ & $(0.25,\ 0.46788)$ \\
circle, $p'{=}2^{31}{-}1$ & list, $2^{-100}$ & $\delta\ge0.46787$ & $\delta\le0.29289$ & --- & $(0.29290,\ 0.46787)$ \\
\bottomrule
\end{tabular}
\caption{Scanner output after the pincer, superseding \cref{tab:scanner}. ``Unsafe (widened)'' is certified by \cref{cor:widened-deployed} and \cref{rem:exact-frontier} (edges rounded outward to five digits: the exact edges are $\tfrac{15331}{32768}$ and $\tfrac{30663}{65536}$). ``Safe, self-contained'' is \cref{thm:deep-mca} resp.\ \cref{thm:johnson-list} as in \cref{tab:scanner}; at rates below $\tfrac14$, \cref{cor:self-contained-safe} would replace the MCA entries. ``Safe, with import'' is \cref{cor:conditional-half}; the list rows need no import, their safe edge being the self-contained Johnson bound. Open bands are stated for the stronger certificate in each column.}
\label{tab:scanner2}
\end{table}

The reduced kernel, restated once. After the pincer, the open MCA band of each deployed row is $(\tfrac14,\ \text{edge})$ with edge $\approx0.4679$, modulo one import whose own proof is independent of everything here; self-contained, it is $(0.1667,\ \text{edge})$, and at rates below $\tfrac14$ the left endpoint improves to half the Johnson radius. Below half the distance the mutual and plain correlated-agreement questions are now provably the same question up to an explicit tangent term (\cref{thm:mca-from-ca}), so the broad band formulation \cref{prob:band} survives only above half the distance, and there the certified unsafe mass presses down to the entropy--subfield envelope $g^*$ (\cref{rem:entropy-frontier}), which the deployed rows already saturate. Closing the band therefore requires one of exactly two kinds of new mathematics: a floor construction whose received words carry more than a graded locator prefix --- anything beating $g^*$ --- or a correlated-agreement theorem crossing half the distance toward capacity, to which \cref{thm:mca-from-ca} would immediately add the mutual layer. This is, for the first time in the program, a two-sided quantitative target with both envelopes made explicit.

\section{Further refinements and finite certificate closure}\label{sec:answers}

This section records three groups of consequences that are not part of the main unsafe/safe sandwich but complete the account.  First, graded rational floors extend the unsafe construction to genus-one scale maps, so the ECFFT obstruction of \cref{rem:ecfft-obstruction} is a one-slot obstruction rather than a macroscopic barrier (\cref{sec:answers-isogeny}), and a stereographic model treats circle codes over challenge fields that do not contain $i$, completing the circle-code transfer (\cref{sec:answers-stereo}).  Second, explicit head words, simple-pole certifying pairs, and the optimized failure profile clarify what is constructive in the failure mechanism and what remains beyond the present method (\cref{sec:answers-explicit,sec:answers-profile}).  Third, \cref{sec:answers-T7} closes the finite certificate grammar over every verdict used in the narrowed sandwich, prints the deployed certificates as integer tuples, bounds verification complexity, and isolates the single imported input to one optional clause.

\subsection{Graded floors on rational scales: the genus-one band is macroscopic}\label{sec:answers-isogeny}

The graded prefix floor of \cref{prop:graded-prefix-floor} pays one subfield factor per locator slot and buys $c$ degrees of agreement per slot on a polynomial scale of degree $c$. On a rational scale of degree type $(a,e)$ the slot spacing is $a-e$ rather than $a$; nothing else changes.

\begin{proposition}[graded prefix floor for rational scales]\label{prop:graded-rational-floor}
Let $\B\subseteq\F$ be a subfield and let $D\subseteq\B$ with $|D|=n$ be $(\psi,a)$-smooth in the sense of \cref{def:rational-smooth}, with $\psi=f/g$ of degree type $(a,e)$, $a>e\ge0$; put $N:=n/a$ and $Q:=\psi(D)\subseteq\B$. Let $1\le K\le n$ and let $m$ satisfy
\[
am\ \ge\ K,\qquad em\ \le\ K,\qquad m\ \le\ N,
\qquad\text{and put}\qquad
w\ :=\ \Bigl\lfloor\frac{am-K}{a-e}\Bigr\rfloor\ \le\ m .
\]
Then there is a $\B$-valued received word $U:D\to\B$ with
\[
\bigl|\Lst\bigl(\RS[\F,D,K],\,1-\tfrac{am}{n},\,U\bigr)\bigr|
\ \ge\
\left\lceil\frac{\binom Nm}{|\B|^{\,w}}\right\rceil .
\]
For $e=0$ this is \cref{prop:graded-prefix-floor} \textup(then $w=m-\lceil K/a\rceil$\textup); for $w\le1$ it recovers the slack-two window of \cref{prop:rational-floor}.
\end{proposition}

\begin{proof}
For $M\subseteq Q$ with $|M|=m$ put $\Lambda_M:=\prod_{b\in M}(f-bg)$. Each factor has degree exactly $a$ with leading coefficient $\mathrm{lc}(f)$, so $\deg\Lambda_M=am$ with leading coefficient $\mathrm{lc}(f)^m\ne0$; and $\Lambda_M\in\B[X]$ vanishes at $x\in D$ if and only if $\psi(x)\in M$, i.e.\ at exactly $am$ points of $D$, by smoothness and $g\ne0$ on $D$. Expanding,
\[
\Lambda_M\ =\ \sum_{j=0}^{m}(-1)^je_j(M)\,f^{\,m-j}g^{\,j},
\]
with $j$-th term of degree $am-j(a-e)$, strictly decreasing in $j$. For $z=(z_1,\dots,z_w)\in\B^w$ define the $\B$-valued word
\[
U_z\ :=\ \Bigl(f(x)^m+\sum_{j=1}^{w}z_j\,f(x)^{\,m-j}g(x)^{\,j}\Bigr)_{x\in D},
\]
which makes sense because $w\le m$: indeed $am-K\le m(a-e)$ is the hypothesis $em\le K$. By pigeonhole over the prefix map $M\mapsto z(M):=\bigl((-1)^je_j(M)\bigr)_{j\le w}\in\B^w$ there is $z^*$ attained by at least $\lceil\binom Nm/|\B|^w\rceil$ sets $M$. For each such $M$ set
\[
c_M\ :=\ U_{z^*}-\Lambda_M|_D\ =\ \Bigl(-\sum_{j>w}(-1)^je_j(M)\,f^{\,m-j}g^{\,j}\Bigr)_{x\in D},
\]
of degree at most $am-(w+1)(a-e)\le K-1$ by the choice of $w$, so $c_M\in\RS[\B,D,K]\subseteq\RS[\F,D,K]$; and $U_{z^*}$ agrees with $c_M$ on the $am$ fiber points of $M$, i.e.\ within relative radius $1-am/n$. Injectivity of $M\mapsto c_M$ on the pigeonhole class: equality of words of degree $\le K-1\le n-1$ forces equality of polynomials, hence $\sum_{j>w}(-1)^j\bigl(e_j(M)-e_j(M')\bigr)f^{\,m-j}g^{\,j}=0$; dividing by $g^m$ in $\F(X)$ gives a polynomial identity in $\psi$, and $\psi$ is nonconstant, hence transcendental over $\F$ inside $\F(X)$, so all coefficients vanish, exactly as in the proof of \cref{prop:rational-floor}. Thus $e_j(M)=e_j(M')$ for $j>w$; for $j\le w$ they agree through $z^*$; hence $\prod_{b\in M}(Y-b)=\prod_{b\in M'}(Y-b)$ and $M=M'$. The $c_M$ are therefore distinct listed codewords for the single word $U_{z^*}$.
\end{proof}

\begin{corollary}[macroscopic universal cap for genus-one rows]\label{cor:ecfft-macroscopic}
Let $D$ be $(\psi,2)$-smooth over $\B$ with $\psi$ of degree type $(2,1)$ \textup(\cref{def:rational-smooth}; in particular every FFTree domain of ECFFT \cite{BCKL21}\textup), let $C=\RS[\F,D,k]$ with $\rho:=k/n\in[\tfrac1{16},\tfrac12]$, $n\ge2^{14}$, $n<q=|\F|<2^{256}$, and $\beta:=\log_2|\B|\le64$. Put $m:=\lceil(k+1+2^{-9}n)/2\rceil$ and $N:=n/2$. Then
\[
\binom Nm\ >\ |\B|^{\,2m-k-1}\Bigl(\frac qk+1\Bigr)
\qquad\text{and}\qquad
\binom Nm\ >\ |\B|^{\,2m-k}\,2^{-128}q ,
\]
and consequently:
\begin{enumerate}[label=\textup{(\roman*)}]
\item \textup(MCA/CA cap\textup) $\eca\bigl(C,\delta\bigr)>\frac1{2k}\bigl(1-\frac nq\bigr)$ for every integer-admissible $\delta\in[1-\frac{2m}n,\,1-\rho-\frac1n]$, and $\emca(C,\delta)$ obeys the same lower bound for every $\delta\in[1-\rho-2^{-9},\,1-\rho)$;
\item \textup(list cap\textup) there is a $\B$-valued word $U$ with $|\Lst(C,1-\frac{2m}n,U)|>2^{-128}q$, so the grand list challenge is unsafe at target $2^{-128}$ on the same band.
\end{enumerate}
The gap $2^{-9}$ coincides with the universal multiplicative gap of \cref{cor:grand}: genus-one rows are no longer behind the polynomial-scale rows.
\end{corollary}

\begin{proof}
Write $w_1:=2m-k-1$ and $w_0:=2m-k$ for the two prefix depths ($K=k+1$ and $K=k$). From $k+1+2^{-9}n\le2m\le k+2^{-9}n+3$: first, $w_0\le2^{-9}n+3=N/256+3$, so $\beta w_0\le N/4+192$; second, $m/N\in[\rho+2^{-9},\,\rho+2^{-9}+2^{-12}]$ using $3/n\le3\cdot2^{-14}<2^{-12}$. On the interval $[\tfrac1{16}+2^{-9},\,\tfrac12+2^{-9}+2^{-12}]$ the concave $\Hb$ is minimized at an endpoint, and $\Hb(\tfrac1{16}+2^{-9})\ge0.344$, $\Hb(\tfrac12+2^{-9}+2^{-12})\ge0.999$, so $\Hb(m/N)\ge0.344$. Hence
\[
\log_2\binom Nm\ \ge\ N\,\Hb(m/N)-\log_2(N+1)\ \ge\ 0.344\,N-\log_2(N+1),
\]
while both right-hand sides are at most $2^{\beta w_0+256}\le2^{N/4+448}$ crudely \textup(using $q/k+1\le q<2^{256}$ and $2^{-128}q<2^{128}$\textup). Since $0.094\,N\ge0.094\cdot2^{13}=770>448+\log_2(N+1)$ for all $N\ge2^{13}$, both displayed inequalities hold. The hypotheses of \cref{prop:graded-rational-floor} hold for both $K\in\{k,k+1\}$: $2m\ge k+1\ge K$; $em=m\le k\le K$ because $2^{-9}n+3\le k$ for $\rho\ge\tfrac1{16}$, $n\ge2^{14}$; and $m\le N$ because $2m\le k+2^{-9}n+3\le n$. Part \textup{(ii)} is the floor at $K=k$. For \textup{(i)}, the floor at $K=k+1$ exceeds $q/k+1\ge\lceil2q\,\eca\rceil$ whenever $\eca(C,\delta)\le\frac12(\frac1k-\frac n{kq})$, contradicting \cref{thm:A} at every integer-admissible $\delta\ge\delta_0:=1-2m/n$ \textup(admissible since $n-2m\le n-k-1$\textup); the assembly through \cref{fact:chain} at $\delta_0$ and \cref{lem:mca-monotone} up to $1-\rho$ is verbatim the proof of \cref{thm:phi-cap}. The band containment is $1-2m/n\le1-\rho-2^{-9}$, i.e.\ $2m\ge k+2^{-9}n$.
\end{proof}

\begin{remark}[the trichotomy of \cref{rem:ecfft-obstruction}, resolved]\label{rem:trichotomy}
\Cref{rem:ecfft-obstruction} offered three exits from the lacunarity collapse: a lacunary basis of isogeny locators, a different received-word family, or a matching safe-side theorem showing the collapse essential. The second exit is the one that exists, and the collapse was not essential: the one-slot analysis priced the window width $2(a-e)$ as the \emph{reach} of the route, when it only sets the \emph{exchange rate} --- $a-e$ degrees of agreement per subfield factor, against an entropy budget of $\Hb$ bits per fiber. Optimizing the number of slots gives the closed-form frontier
\[
\Hb(\rho+g^*)\ =\ g^*\cdot\beta\cdot\frac a{a-e},
\]
so the degree type enters only through the ramification-defect ratio $a/(a-e)$: polynomial scales ($e=0$) have ratio $1$ and recover \cref{rem:entropy-frontier}; genus-one $x$-maps have $e=a-1$, ratio $a$, optimized at the first fold $a=2$, i.e.\ an effective subfield cost of $2\beta$. Numerically $g^*(\tfrac12,62)\approx0.0161\approx2^{-5.96}$ at $\beta=31$ and $g^*(\tfrac12,128)\approx0.0078\approx2^{-7}$ at $\beta=64$: macroscopic in every case, against the $2/n$ of the one-slot route. The safe side needs no new work: \cref{thm:deep-mca,thm:mca-from-ca,thm:elementary-ca,cor:conditional-half} are domain-agnostic, so every $(\psi,2)$-smooth row carries the full narrowed sandwich of \cref{thm:sandwich2}'s shape --- self-contained safe frontier at $(1-\rho)/3$, safe at $(1-\rho)/2$ modulo the single import, unsafe from gap $2^{-9}$ \textup(and, row by row, from the optimized frontier $g^*(\rho,2\beta)$\textup). This answers the genus-one/circle transport question \cref{prob:isogeny}\textup{(a)} as posed.
\end{remark}

\subsection{The stereographic model: circle codes over every challenge field}\label{sec:answers-stereo}

The torus uniformization of \cref{lem:circle-rs} needs $i\in\F$. The circle is a conic with the rational point $(1,0)$, so it also has a uniformization defined over $\F_p$ itself --- the stereographic projection --- and this one transports the code theory to \emph{every} challenge field. The price is that the resulting Reed--Solomon domain is no longer a multiplicative coset; but it is rationally smooth of the same $(2,1)$ type as the genus-one rows, so \cref{sec:answers-isogeny} covers it.

\begin{lemma}[stereographic uniformization of circle codes; no $i$ required]\label{lem:stereographic}
Let $p\equiv3\pmod4$, let $\mathcal D=gH\cup g^{-1}H\subseteq\UU$ be a twin coset of size $2M$ with $\operatorname{ord}(g)\nmid2M$ and $4\mid M$, let $E:=\iota^{-1}(\mathcal D)\subseteq\mathcal C(\F_p)$, and let $\F\supseteq\F_p$ be \emph{any} finite field. Define $s(x,y):=y/(1+x)$ and $T:=s(E)$.
\begin{enumerate}[label=\textup{(\roman*)}]
\item The four torsion points $(\pm1,0),(0,\pm1)$ lie in $\iota^{-1}(H)$ and hence not in $E$; $s$ is injective on $E$ with inverse $s\mapsto\bigl(\frac{1-s^2}{1+s^2},\frac{2s}{1+s^2}\bigr)$; and $T\subseteq\F_p\setminus\{0,\pm1\}$ with $|T|=2M$. \textup(Over $\F_p$ one has $1+s^2\ne0$ throughout, since $p\equiv3\pmod4$.\textup)
\item With the diagonal vector $d:=\bigl((1+s(P)^2)^{-w}\bigr)_{P\in E}$ and the coordinate bijection $s|_E:E\to T$,
\[
\mathcal C_w(\F,E)\ =\ d\circ\RS[\F,T,2w+1].
\]
Consequently, by \cref{lem:diag-invariance}, the bivariate circle code on $E$ and $\RS[\F,T,2w+1]$ have identical list sizes at every radius and identical $\eca$ and $\emca$ at every \textup(pair of\textup) radii --- with no hypothesis on $i\in\F$.
\item $T$ is $(\psi,2)$-smooth over $\B:=\F_p$ in the sense of \cref{def:rational-smooth} for the tangent-halving map $\psi(s):=\frac{s^2-1}{2s}$, of degree type $(2,1)$: the $\psi$-fibers within $T$ are the pairs $\{s(P),s(P\tau)\}=\{s,-1/s\}$, $\tau:=(-1,0)$, complete because $\tau\in\iota^{-1}(H)$ and of uniform size $2$ because $s^2\ne-1$ in $\F_p$; and $N=M$, $Q=\psi(T)\subseteq\F_p$.
\end{enumerate}
\end{lemma}

\begin{proof}
(i) $H$ is the unique subgroup of order $M$ of the cyclic group $\UU$ of order $p+1$; since $4\mid M$ it contains all elements of order dividing $4$, whose $\iota^{-1}$-images are exactly the four torsion points. From $\operatorname{ord}(g)\nmid2M$ we get $g,g^{-1}\notin H$, so $gH\cap H=g^{-1}H\cap H=\emptyset$ and $\mathcal D\cap H=\emptyset$; in particular $(-1,0)\notin E$, so $s$ is defined on $E$, and $(1,0),(0,\pm1)\notin E$, so $0,\pm1\notin T$. For the inverse formula: from $x^2+y^2=1$, $1+s^2=\frac{(1+x)^2+y^2}{(1+x)^2}=\frac2{1+x}$, whence $\frac{1-s^2}{1+s^2}=(1+x)-1=x$ and $\frac{2s}{1+s^2}=s(1+x)=y$; and $1+s^2\ne0$ for $s\in\F_p$ because $-1$ is a nonsquare when $p\equiv3\pmod4$.

(ii) Every class of degree $\le w$ has the canonical form $f_0(x)+y\,f_1(x)$ with $\deg f_0\le w$, $\deg f_1\le w-1$, a space of dimension $2w+1$ \textup(proof of \cref{lem:circle-rs}\textup). Substituting the inverse of (i) and clearing denominators,
\[
(1+s^2)^w\bigl(f_0(x)+y\,f_1(x)\bigr)\Big|_{x=\frac{1-s^2}{1+s^2},\,y=\frac{2s}{1+s^2}}
\]
is a polynomial in $s$ of degree at most $2w$: the monomial $x^j$ contributes $(1-s^2)^j(1+s^2)^{w-j}$ and $y\,x^j$ contributes $2s\,(1-s^2)^j(1+s^2)^{w-1-j}$. These $2w+1$ polynomials form a basis of $\F[s]_{\le2w}$: the first family is even, the second odd, so it suffices to check each family separately; and in the variable $v:=s^2$ the first family is $\{(1-v)^j(1+v)^{w-j}\}_{j\le w}$, the image of the monomial basis $\{Y^jZ^{w-j}\}$ under the invertible substitution $(Y,Z)=(1-v,1+v)$ \textup($p$ odd\textup), and similarly for the odd family. Hence $F\mapsto(1+s^2)^wF$ is a linear isomorphism from the circle space onto $\F[s]_{\le2w}$, and evaluating at $P\in E$ gives $F(P)=(1+s(P)^2)^{-w}\,f(s(P))$ with $f\in\F[s]_{\le2w}$ ranging over all of $\RS[\F,T,2w+1]$: the displayed code identity. The vector $d$ is $\F_p$-valued and everywhere nonzero, so \cref{lem:diag-invariance} applies.

(iii) The group law is $(x_1,y_1)\cdot(x_2,y_2)=(x_1x_2-y_1y_2,\,x_1y_2+x_2y_1)$, so $P\tau=(-x,-y)$ and
\[
s(P)\,s(P\tau)\ =\ \frac y{1+x}\cdot\frac{-y}{1-x}\ =\ \frac{-y^2}{1-x^2}\ =\ -1,
\]
using $x^2+y^2=1$ and $y\ne0$ on $E$. Since $\tau\in\iota^{-1}(H)$ \textup($2\mid M$\textup) and $\mathcal D H=\mathcal D$ coset-wise, $E\tau=E$, so the involution $s\mapsto-1/s$ preserves $T$. The fibers of $\psi$: $\psi(s)=b$ if and only if $s^2-2bs-1=0$, whose two roots multiply to $-1$; hence every fiber is $\{s,-1/s\}$, of size exactly $2$ since $s^2=-1$ is impossible in $\F_p$, and complete within $T$ by the closure just shown. Finally $f=X^2-1$ and $g=2X$ are coprime of degrees $2>1$, and $g$ is nonvanishing on $T$ because $0\notin T$.
\end{proof}

\begin{corollary}[circle codes over every challenge field; the genus-one/circle transport question \cref{prob:isogeny}\textup{(b)}]\label{cor:circle-anyfield}
Let $E$ be a standard-position circle domain as in \cref{lem:stereographic} with $M=2^{21}$ \textup(the deployed shape: $n_c=2M=2^{22}$, $k_c=2^{21}+1$\textup), and let $\F\supseteq\F_p$, $p=2^{31}-1$, be any challenge field, $q=|\F|$.
\begin{enumerate}[label=\textup{(\roman*)}]
\item If $q<2^{100}$ --- in particular $\F=\F_p$ and $\F=\F_{p^3}$, hence every field without $i$ among the extensions $\F_{p^r}$, $r\le4$, that deployed verifiers sample from, since $i\in\F_{p^r}$ if and only if $2\mid r$ --- the row is degenerately unsafe at every deployed target by \cref{prop:small-field}, and the genus-one/circle transport question \cref{prob:isogeny}\textup{(b)} is out of scope there as posed.
\item If $i\in\F$, \cref{cor:circle-grand,cor:circle-deployed} apply verbatim through \cref{lem:circle-rs}.
\item If $i\notin\F$ and $q\ge2^{100}$ \textup(the first such field is $\F_{p^5}$\textup), the stereographic model of \cref{lem:stereographic} transports the row to $\RS[\F,T,k_c]$ on the $(\psi,2)$-smooth tangent domain $T$, and \cref{prop:graded-rational-floor} with
\[
m=2^{20}+2^{15}+1,\qquad w=2m-(k_c+1)=2^{16},\qquad
\binom{2^{21}}{m}\ >\ p^{\,2^{16}}\cdot2^{256}
\]
\textup(verified exactly in integer arithmetic; the margin exceeds $2^{63000}$\textup) gives, for every $q\in(n_c,2^{256})$ at once, an MCA and list cap on
\[
\delta\ \in\ \bigl[\,1-\rho_c-2^{-6},\ 1-\rho_c\bigr),
\qquad
\text{error}\ >\ \tfrac1{2k_c}\bigl(1-\tfrac{n_c}q\bigr)\ >\ 2^{-23},
\]
by the assembly of \cref{cor:ecfft-macroscopic} \textup(here $\beta=\log_2p$, effective cost $2\beta$, asymptotic frontier $g^*\approx2^{-5.96}$\textup). The safe side is generic: \cref{thm:deep-mca} applies to the circle code directly \textup(as already noted after its proof\textup), and \cref{thm:mca-from-ca,thm:elementary-ca,cor:conditional-half} apply to the diagonal model. Every bivariate circle row over every challenge field therefore carries a two-sided sandwich, and the deep-point conversion \cref{thm:A} does extend to the bivariate setting --- through the stereographic model rather than a new proof.
\end{enumerate}
\end{corollary}

\begin{proof}
(i) is \cref{prop:small-field} with $q=p<2^{31}$ resp.\ $q=p^3<2^{93}$: $\emca\ge1/q>2^{-100}$ at every radius, so $\delta^*_C(\varepsilon^*)=0$ for every deployed target $\varepsilon^*\le2^{-100}$; the parity criterion is $\operatorname{ord}(i)=4\mid p^r-1$ if and only if $r$ is even, as $p\equiv3\pmod4$. (ii) is quoted. (iii): the transport is \cref{lem:stereographic}\textup{(ii)}, the smoothness is \textup{(iii)}, and the hypotheses of \cref{prop:graded-rational-floor} hold with $K\in\{k_c,k_c+1\}$: $2m=k_c+1+2^{16}\ge K$, $em=m\le k_c\le K$, $m\le N=2^{21}$. The single displayed integer inequality dominates both thresholds for every $q<2^{256}$, since $q/k_c+1\le q<2^{256}$ and $2^{-100}q<2^{156}$; the edge satisfies $1-2m/n_c=1-\rho_c-2^{-6}-2^{-21}+2^{-22}\le1-\rho_c-2^{-6}$, and $\lfloor\delta_0 n_c\rfloor=n_c-2m\le n_c-k_c-1$, so \cref{thm:A}, \cref{fact:chain}, and \cref{lem:mca-monotone} assemble exactly as in \cref{cor:ecfft-macroscopic}; the error numerics use $q\ge p^5$, so $\frac1{2k_c}\bigl(1-\frac{n_c}q\bigr)>2^{-22}(1-2^{-21})(1-2^{-133})>2^{-23}$.
\end{proof}

\subsection{Explicit witness words and explicit certifying pairs}\label{sec:answers-explicit}

Every floor so far selects its received word by pigeonhole over a prefix. At the first staircase step the pigeonhole can be removed entirely: if the support is assembled from \emph{antipodal pairs}, its first elementary symmetric function vanishes identically, and the received word collapses to a pure power.

\begin{theorem}[explicit head-and-pairs floor]\label{thm:explicit-head-floor}
Let $D$ be $(\varphi,c)$-smooth over $\B$ \textup(\cref{def:map-smooth}\textup) with $N:=n/c$ even, and suppose $Q:=\varphi(D)$ satisfies $-Q=Q$ and $0\notin Q$, so that $Q$ is partitioned into $N/2$ antipodal classes $\{y,-y\}$. Let $K\le n$ and let $m$ satisfy $2\le m\le N-2$ and $cm\le K-1+2c$.
\begin{enumerate}[label=\textup{(\roman*)}]
\item If $m$ is even, the \emph{pure power word} $u:=\varphi^m|_D$ satisfies
\[
\bigl|\Lst\bigl(\RS[\F,D,K],\,1-\tfrac{cm}n,\,u\bigr)\bigr|\ \ge\ \binom{N/2}{m/2},
\]
the listed codewords being $c_M:=u-\Lambda_M|_D$ for $M$ a union of $m/2$ antipodal classes.
\item If $m$ is odd, then for every $t\in Q$ the \emph{one-head word} $u:=\bigl(\varphi^m-t\,\varphi^{m-1}\bigr)|_D$ satisfies
\[
\bigl|\Lst\bigl(\RS[\F,D,K],\,1-\tfrac{cm}n,\,u\bigr)\bigr|\ \ge\ \binom{N/2-1}{(m-1)/2},
\]
with $M=\{t\}\cup(\text{$(m-1)/2$ antipodal classes distinct from $\{t,-t\}$})$.
\end{enumerate}
In both cases the word, the family, and every listed codeword are explicit closed forms, with no pigeonhole and no subfield division. Both deployed towers qualify: on a multiplicative coset, $Q$ is a coset of the even-order subgroup $\mu_N$, so $-Q=Q$ and $0\notin Q$; on a Chebyshev domain $D=\chi(\mathcal D)$ with $2c\mid M$, the antipodal classes of $Q=T_c(D)$ are exactly the $T_2$-fibers inside $Q$, complete of size $2$ by \cref{lem:cheb-fibers} at scales $c$ and $2c$ \textup(two elements with equal $T_2$-value satisfy $y'=-y$, and $0\notin Q$ since a $T_2$-fiber $\{0\}$ would have size $1$\textup).
\end{theorem}

\begin{proof}
In case \textup{(i)}, $M$ is a union of antipodal pairs, so $e_1(M)=0$ and $\Lambda_M=\varphi^m+\sum_{j\ge2}(-1)^je_j(M)\varphi^{m-j}$, whence $c_M:=u-\Lambda_M|_D=-\sum_{j\ge2}(-1)^je_j(M)\varphi^{m-j}|_D$ has degree at most $c(m-2)\le K-1$ and agrees with $u$ on the $cm$ fiber points of $M$. In case \textup{(ii)}, $e_1(M)=t+\sum(\text{pair sums})=t$, so the slot-$1$ coefficient of $\Lambda_M$ is $-t\varphi^{m-1}$ for every $M$ in the family, and $c_M:=u-\Lambda_M|_D$ again consists of the slots $j\ge2$ only, of degree at most $c(m-2)\le K-1$; note that $e_2(M)$ and beyond are permitted to vary, since they land in the codeword rather than in the received word. Injectivity: $c_M=c_{M'}$ as words forces $\Lambda_M=\Lambda_{M'}$ on all of $D$, and both are polynomials of degree $cm<cN=n$, hence equal as polynomials, with equal root multisets, hence $M=M'$ through the fibers. The family counts are the number of choices of antipodal classes.
\end{proof}

\begin{corollary}[explicit witnesses at the deployed rows]\label{cor:explicit-deployed}
At the KoalaBear row of \cref{cor:deployed} \textup($n=2^{21}$, $k=2^{20}$, $\B=\F_p$, $q=p^6$\textup) and at the Chebyshev line-round row of \cref{cor:circle-deployed}\textup{(a)} \textup($n=2^{21}$, $k=2^{20}$, $q=p'^4$, $p'=2^{31}-1$, target $2^{-100}$\textup):
\begin{enumerate}[label=\textup{(\roman*)}]
\item \textup(MCA, gap $2^{-8}$\textup) With $c=2^{12}$, $m=k/c+2=258$, $K=k+1$: the explicit word $u=X^{k+2^{13}}|_D$ \textup(resp.\ $u=T_c^{258}|_D$\textup) has list at agreement $k+2^{13}$ at least $\binom{256}{129}\ge2^{251}$, exceeding $q/k+1<2^{166}$ \textup(resp.\ $<2^{104}$\textup); by \cref{thm:A} with $\eta=\frac12$, exactly as in \cref{thm:main}, $\eca(C,\delta)>\frac1{2k}(1-\frac nq)$ on the integer-admissible part of $[\frac12-2^{-8},\frac12)$ and $\emca$ obeys the same bound throughout it.
\item \textup(lists, gap $2^{-7}$\textup) With $c=2^{14}$, $m=k/c+1=65$, $K=k$: for any $t\in Q$ the explicit word $u=\bigl(X^{k+2^{14}}-t\,X^{k}\bigr)|_D$ \textup(resp.\ its $T_c$ analogue\textup) has list at agreement $k+2^{14}$ at least $\binom{63}{32}=916312070471295267>2^{59}$, exceeding $2^{-128}q<2^{58}$ \textup(resp.\ $2^{-100}q<2^{24}$\textup): the first staircase step of the grand list challenge, at its original printed edge $\frac12-2^{-7}$, is certified by a fully explicit word.
\end{enumerate}
All four inequalities are verified exactly in integer arithmetic.
\end{corollary}

\begin{proof}
The hypotheses of \cref{thm:explicit-head-floor} hold with $N=512$ resp.\ $128$: parity as displayed, $cm\le K-1+2c$ \textup(with equality in \textup{(i)}\textup), and the Chebyshev scale conditions $2c\mid M=2^{21}$. The threshold comparisons are the displayed integer inequalities; the MCA conversion and monotone assembly are as quoted.
\end{proof}

The words above certify lists explicitly, and \cref{thm:A} converts large lists into CA failure --- but by contradiction, without exhibiting the failing pair. the explicit-witness question \cref{prob:explicit}, however, asks for the \emph{pair}. The proof of \cref{thm:A} is more generous than its statement: it converts through the completely explicit simple-pole lines $\bigl(U/(x-\alpha),\,-1/(x-\alpha)\bigr)$ --- precisely the residue-line normal form with extension denominator that \cref{cor:Fvalued} predicts --- and its collision count can be read as a lower bound holding for \emph{most} poles.

\begin{theorem}[explicit certifying pairs, up to the choice of a pole]\label{thm:explicit-pairs}
At the KoalaBear row of \cref{cor:deployed}, let $u:=X^{k+2^{13}}|_D$ be the pure power word of \cref{cor:explicit-deployed}\textup{(i)}, let $L_0:=\lceil(q-n)/k\rceil<2^{166}$, and fix any explicit subfamily of $L_0$ of its listed codewords $P_M\in\F[X]_{<k+1}$ \textup(e.g.\ the lexicographically least $L_0$ antipodal-class sets; $L_0\le\binom{256}{129}$\textup). For $\alpha\in\Omega:=\F\setminus D$ define the explicit pair
\[
f_\alpha(x):=\frac{u(x)}{x-\alpha},
\qquad
g_\alpha(x):=-\frac1{x-\alpha}
\qquad(x\in D).
\]
Then for at least half of all $\alpha\in\Omega$, the pair $(f_\alpha,g_\alpha)$ has at least
\[
\Bigl\lceil\frac{q-n}{3k}\Bigr\rceil
\]
distinct CA-bad slopes at radius $\delta=\frac12-2^{-8}$, namely among the explicit values $\{P_M(\alpha)\}$; its CA-bad-slope density exceeds $2^{-21.6}>2^{-22}$ there and at every larger integer-admissible radius, and its MCA-bad-slope density exceeds the same bound throughout $[\frac12-2^{-8},\frac12)$. All but at most $p$ of these poles have $\alpha\notin\B$, and for those the pair is genuinely $\F$-valued, in accordance with \cref{cor:Fvalued}. The same construction on the circle line-round row gives density $>2^{-21.6}>2^{-23}$ for at least half of all poles.
\end{theorem}

\begin{proof}
Each $P_M$ agrees with $u$ on at least $a=k+2^{13}$ points, so, as in the proof of \cref{thm:A}, for every $\alpha\in\Omega$ and every $M$ the slope $P_M(\alpha)$ carries the point $f_\alpha+P_M(\alpha)g_\alpha$ to within radius $1-a/n=\frac12-2^{-8}$ of $C$, while the pair is farther than every $\delta<1-k/n=\frac12$ from $C^{\equiv2}$ by the $(X-\alpha)G+1$ argument there; so every distinct value $P_M(\alpha)$ is a CA-bad slope at every integer-admissible $\delta\in[\frac12-2^{-8},\frac12)$, hence \textup(as failing CA is failing MCA, \cref{fact:chain}\textup) an MCA-bad slope throughout that interval. It remains to count distinct values for many $\alpha$ at once. Distinct sets $M$ give distinct polynomials $P_M$ \textup(\cref{thm:explicit-head-floor}\textup), and $P_M-P_{M'}$ is nonzero of degree at most $k$, so the total number of colliding pairs over all $\alpha\in\Omega$ is at most $k\binom{L_0}2$. Let $x^*:=\lceil2k\binom{L_0}2/(q-n)\rceil$; by Markov, fewer than $(q-n)/2$ poles have collision count at least $x^*$, so at least half have collision count at most $x^*-1$. For such $\alpha$, Cauchy--Schwarz as in \cref{thm:A} gives
\[
M(\alpha)\ \ge\ \frac{L_0^2}{L_0+2(x^*-1)}\ \ge\ \frac{L_0}{1+2k(L_0-1)/(q-n)}\ >\ \Bigl\lceil\frac{q-n}{3k}\Bigr\rceil-1,
\]
so, since $M(\alpha)$ is an integer,
$M(\alpha)\ge\lceil(q-n)/(3k)\rceil$. The last comparison and
$\lceil(q-n)/3k\rceil/q>2^{-21.6}$ are verified exactly in integer arithmetic
at both rows. The subfield count is $|\Omega\cap\B|\le|\B|=p$; for
$\alpha\notin\B$ the word $g_\alpha$ is not $\B$-valued \textup(its value at any
$x\in D\subseteq\B$ lies in $\B$ only if $x-\alpha\in\B$\textup), and no common
scalar repairs this while keeping $f_\alpha$ a $\B$-multiple, so the pair falls
outside the inert class of \cref{lem:confine}, as \cref{cor:Fvalued} requires
of any pair with this density.
\end{proof}

\begin{remark}[what remains of the explicit-witness question \cref{prob:explicit}]\label{rem:explicit-scope}
\Cref{thm:explicit-pairs} answers the explicit-witness question \cref{prob:explicit} at $\delta=\frac12-2^{-8}$, on the upper half of the stated band, in the predicted normal form, and for both densities the problem asks about; the non-constructiveness has collapsed from a received word --- an object of $n\log_2q$ bits with no known certificate --- to the choice of one field element in an explicit majority subset of $\F\setminus D$. Two residues remain. First, a certified \emph{individual} pair: derandomizing the pole. Second, the left half of the band: at gap exactly $2^{-7}$ the explicit family at scale $2^{13}$ has only $\binom{128}{65}<2^{125}$ members, short of the $\approx2^{166}$ the counting needs, so explicitness there --- and throughout the widened band of \cref{cor:widened-deployed}, where the floors are pigeonhole-selected deep prefixes --- would require explicit families with prescribed deep prefixes, i.e.\ explicit subsets of $\mu_N$ with many prescribed elementary symmetric functions. We record this as the surviving form of the explicit-witness question \cref{prob:explicit}.
\end{remark}

\subsection{The failure profile: the deep-point ceiling}\label{sec:answers-profile}

\Cref{thm:A} was used throughout with $\eta=\frac12$, giving error floors of $\frac1{2k}(1-\frac nq)$. The dichotomy improves when the list is far above threshold, and the improvement is exactly quantifiable.

\begin{proposition}[optimized failure profile]\label{prop:profile}
Let $C=\RS[\F,D,k]$, let $\delta$ be integer-admissible \textup($\lfloor\delta n\rfloor\le n-k-1$, $q>n$\textup), and put $T:=\frac1k\bigl(1-\frac nq\bigr)$. If some received word has $\bigl|\Lst(C^+,\delta,u)\bigr|\ge1+\kappa qT$ for a real $\kappa>0$, then
\[
\eca(C,\delta)\ \ge\ \frac{\kappa}{\kappa+1}\,T .
\]
\end{proposition}

\begin{proof}
Write $x:=\eca(C,\delta)$; if $x\ge T$ there is nothing to prove, so assume $x<T$ and set $\eta:=x/T\in[0,1)$. Then $x=\eta T$, so \cref{thm:A} applies and gives $L\le\lceil qx/(1-\eta)\rceil\le qx/(1-\eta)+1$, i.e.\ $\kappa qT\le L-1\le qx/(1-x/T)$. Rearranging, $\kappa(T-x)\le x$, i.e.\ $x\ge\kappa T/(\kappa+1)$.
\end{proof}

\begin{corollary}[deployed profile at the first staircase step]\label{cor:profile-deployed}
At the KoalaBear row, the fiber floor of \cref{lem:phi-fiber} at scale $2^{13}$ gives $L\ge\lceil\binom{256}{130}/p\rceil>2^{220}$ at agreement $k+2^{14}$, hence $\kappa\ge2^{54}$ and
\[
\eca\bigl(C,\delta\bigr)\ \ge\ \bigl(1-2^{-54}\bigr)\,\frac1k\Bigl(1-\frac nq\Bigr)
\]
for every integer-admissible $\delta\in[\frac12-2^{-7},\frac12-2^{-21}]$, with the same lower bound for $\emca(C,\delta)$ throughout $[\frac12-2^{-7},\frac12)$: twice the printed constant of \cref{cor:deployed}, saturating the dichotomy's ceiling to relative error $2^{-54}$. On the circle line-round row the same floor gives $\kappa\ge2^{116}$ and the factor $1-2^{-116}$.
\end{corollary}

\begin{proof}
$\kappa=\lfloor(L-1)k/(q-n)\rfloor\ge2^{54}$ \textup(resp.\ $2^{116}$\textup) is verified exactly; $\kappa/(\kappa+1)\ge1-2^{-54}$; the list persists at every larger radius, and \cref{fact:chain} with \cref{lem:mca-monotone} extend as usual.
\end{proof}

\begin{remark}[the ceiling, and what keeps the error-one question \cref{prob:errorone} open]\label{rem:profile-ceiling}
Two ceilings frame the error-one question \cref{prob:errorone}. First, the route ceiling: as $\kappa\to\infty$ the profile of \cref{prop:profile} tends to $T=\frac1k(1-\frac nq)$ and no list, however large, certifies more through \cref{thm:A} --- error beyond $\approx\frac1k$ requires strengthening the dichotomy itself, which is exactly the unnamed problem following the error-one question \cref{prob:errorone}. Second, at the specific radius $1-\rho-\frac1{64}$ of the error-one question \cref{prob:errorone} over \emph{prime} fields $2^{150}\le p\le2^{256}$, the present floors certify nothing at all: the only subfield is $\B=\F_p$, so every prefix frontier obeys $g^*\le1/\beta=1/\log_2p\le\frac1{150}<\frac1{64}$, and the radius lies strictly outside the entropy--subfield envelope. \textup(There is also no confinement barrier there: \cref{lem:confine} caps $\B$-valued pairs at $|\B|/q$, which is vacuous when $\B=\F$.\textup) So the error-one question \cref{prob:errorone} is genuinely beyond both mechanisms of this paper: it needs either non-prefix floors at macroscopic gaps over prime fields, or conversions past $\frac1k$. The certified state after this section: error at least $(1-2^{-54})\frac1k(1-\frac nq)$ at gap $2^{-7}$, error at least $\frac1{2k}(1-\frac nq)$ at the widened gaps, and error one still open everywhere.
\end{remark}

\subsection{Certificate closure of the two-sided sandwich}\label{sec:answers-T7}

\Cref{prop:finite-certificate} supplied the initial certificate format and \cref{tab:core} the formal core.  The pincer and the refinements above add new verdict types.  This subsection closes the certificate grammar over the statements used by the two-sided sandwich, prints the deployed certificates, and isolates the single imported proximity-gap input.

\begin{definition}[certificate grammar, second edition]\label{def:certificate-v2}
Extend \cref{def:scanner} by the following verdict types, each a tuple of integers together with a theorem tag:
\begin{enumerate}[label=\textup{(V\arabic*)},leftmargin=2.6em,start=5]
\item \textbf{Prefix-floor \textup(PF\textup) / rational-floor \textup(RF\textup).} A scale datum \textup(polynomial or Chebyshev $c$, or rational type $(a,e)$\textup), integers $(m,w,K)$, and the single inequality $\binom Nm>|\B|^w\cdot\Theta$, $\Theta\in\{q/k+1,\ \varepsilon^*q\}$: mark the band $[1-am/n,1-\rho)$ MCA-unsafe \textup(via \cref{prop:graded-prefix-floor,prop:graded-rational-floor}, \cref{thm:A}\textup) resp.\ $1-am/n$ list-unsafe.
\item \textbf{Explicit-head \textup(XH\textup).} Integers $(c,m,K)$ with $cm\le K-1+2c$ and the inequality $\binom{N/2-\epsilon}{\lfloor m/2\rfloor}>\Theta$, $\epsilon\in\{0,1\}$ by parity: as \textup{(V5)}, with the received word printed in closed form \textup(\cref{thm:explicit-head-floor}\textup).
\item \textbf{Mutual-from-correlated \textup(MC\textup).} An integer $r$ with $2r\le w_{\min}-1$ and a CA-verdict handle: transfer a CA-safe verdict at radius $r/n$ to MCA-safe with error $\max\{\varepsilon_{\rm CA},\,r/q\}$ \textup(\cref{thm:mca-from-ca}\textup). The CA handle is either \textbf{HJ} --- integers $(r,L_2)$ with $n-2r>0$, $(n-2r)^2>(k-1)n$, $(L_2+1)\bigl((n-2r)^2-(k-1)n\bigr)>n(n-2r-k+1)$, and $(1+(r+1)L_2)/q\le\varepsilon^*$, certifying the self-contained half-Johnson CA bound of \cref{thm:elementary-ca} --- or \textbf{CH} --- the inequality $2r\le n-k$ with error $n/q\le\varepsilon^*$, tagged \textup{\textsc{import}}:~\cite[Thm.~4.1]{BCIKS20} \textup(\cref{cor:conditional-half}\textup).
\item \textbf{Profile \textup(PR\textup).} A floor handle and an integer $\kappa$ with $L\ge1+\kappa q T$: certify error $\ge\frac{\kappa}{\kappa+1}T$ \textup(\cref{prop:profile}\textup).
\end{enumerate}
\end{definition}

\begin{theorem}[certificate closure of the two-sided sandwich]\label{thm:T7-closure}
With the grammar of \cref{prop:finite-certificate} extended by \cref{def:certificate-v2}:
\begin{enumerate}[label=\textup{(\roman*)}]
\item \textup(Soundness\textup) Every verdict of \cref{def:certificate-v2} is implied by its cited theorem once its integer inequalities hold; \cref{thm:scanner-sound} extends verbatim.
\item \textup(Completeness for the printed results\textup) Every clause of \cref{thm:sandwich2}, of \cref{tab:scanner2}, of \cref{cor:ecfft-macroscopic}, of \cref{cor:circle-anyfield}, and of \cref{cor:explicit-deployed,cor:profile-deployed} is certified by a certificate of \cref{prop:finite-certificate} or \cref{def:certificate-v2}; the deployed certificates are printed in \cref{tab:certs}.
\item \textup(Complexity\textup) A \textup{(V5)}/\textup{(V6)} certificate is checked with $O(N)$ multiplications of integers of $O(N+\log q)$ bits \textup(the binomial by its product formula, the power $|\B|^w$ by repeated squaring\textup); every other certificate with $O(1)$ arithmetic operations on integers of $O(\log q+\log n)$ bits. A full row is checked in $O(n)$ big-integer multiplications.
\item \textup(Import isolation\textup) Exactly one imported statement occurs anywhere in the grammar: the \textup{\textsc{import}} tag of the CH handle. Deleting it degrades only the safe edge, from $(1-\rho)/2$ to $\max\{(1-\rho)/3,\ \text{half-Johnson}\}$, and leaves every other verdict intact. The remaining statements --- \cref{tab:core} extended by the rows of \cref{tab:core2} --- quantify over explicitly bounded finite sets and use only the listed elementary ingredients; their dependency closure is import-free.
\end{enumerate}
\end{theorem}

\begin{proof}
(i) is quotation: each verdict's hypotheses are the hypotheses of its theorem, specialized to integers. (ii) is inspection of the quoted proofs against \cref{tab:certs}: each proof consists of its printed integer inequalities plus the cited assembly steps, all of which are verdicts or core statements. (iii): $\binom Nm=\prod_{j=1}^m\frac{N-m+j}j$ uses $2m\le2N$ multiplications and divisions of integers of at most $N+O(\log N)$ bits; thresholds and powers are $O(\log w)$ squarings of $O(w\log|\B|+\log q)$-bit integers, and $w\log|\B|=O(N)$ in every certificate above; the remaining verdicts compare $O(1)$ products. (iv): the only \textup{\textsc{import}} tag is CH; \cref{thm:sandwich2} lists its conditional clause separately, and the self-contained clauses cite only \cref{thm:deep-mca,thm:elementary-ca,thm:mca-from-ca} on the safe side. The new core rows are in \cref{tab:core2}; each proof above is elementary in the stated ingredients, as the reader may verify line by line.
\end{proof}

\begin{table}[ht]
\centering
\small
\begin{tabular}{@{}ll@{}}
\toprule
Added core statement & Proof ingredients \\
\midrule
\cref{thm:mca-from-ca} \textup(mutual from correlated\textup) & minimum weight, two-point interpolation \\
\cref{thm:elementary-ca} \textup(half-Johnson CA\textup) & pair-list double counting, convexity \\
\cref{prop:graded-prefix-floor}, \cref{prop:graded-rational-floor} & symmetric functions, graded pigeonhole \\
\cref{thm:explicit-head-floor} \textup(explicit heads\textup) & antipodal pairing, no pigeonhole \\
\cref{thm:explicit-pairs} \textup(explicit pairs\textup) & \cref{thm:A}'s collision count, Markov \\
\cref{prop:profile} \textup(optimized profile\textup) & one-line rearrangement of \cref{thm:A} \\
\cref{lem:stereographic} \textup(stereographic model\textup) & conic parametrization, basis change \\
\cref{def:certificate-v2}, \cref{thm:T7-closure} & finite case analysis over the above \\
\bottomrule
\end{tabular}
\caption{The formal finite core, continued: statements added by the pincer and the refinements, extending \cref{tab:core}. Each row quantifies over explicitly bounded finite sets; the dependency order is top to bottom, below the rows of \cref{tab:core}. The one import (\textsc{CH}) is not a core row.}
\label{tab:core2}
\end{table}

\begin{table}[ht]
\centering
\small
\begin{tabular}{@{}llll@{}}
\toprule
Row / verdict & Certificate & Verified inequality & \\
\midrule
KB MCA, PF & $(c,m,w)=(16,\,69748,\,4211)$ & $\binom{2^{17}}{69748}>p^{4211}(q/k+1)$ & exact \\
KB list, PF & $(32,\,34874,\,2106)$ & $\binom{2^{16}}{34874}>p^{2106}\,2^{-128}q$ & exact \\
Circle MCA, PF & $(32,\,34873,\,2104)$ & $\binom{2^{16}}{34873}>p'^{2104}(q/k+1)$ & exact \\
Circle list, PF & $(32,\,34874,\,2106)$ & $\binom{2^{16}}{34874}>p'^{2106}\,2^{-100}q$ & exact \\
KB MCA, XH & $(2^{12},\,258)$, $u=X^{k+2^{13}}|_D$ & $\binom{256}{129}>q/k+1$: $2^{251.6}$ vs $2^{165.94}$ & exact \\
KB list, XH & $(2^{14},\,65)$, one head & $\binom{63}{32}>2^{-128}q$: $2^{59.66}$ vs $2^{57.94}$ & exact \\
Circle MCA, XH & $(2^{12},\,258)$, $u=T_c^{258}|_D$ & $\binom{256}{129}>q/k+1$: $2^{251.6}$ vs $2^{104}$ & exact \\
Circle list, XH & $(2^{14},\,65)$, one head & $\binom{63}{32}>2^{-100}q$: $2^{59.66}$ vs $2^{24}$ & exact \\
Both rows, MC & $r=\lfloor(n-k)/2\rfloor=2^{19}$ & $2r\le n-k=2^{20}$ & --- \\
KB safe, HJ & $(r,L_2)=(307121,\,1001282)$ & $(1+(r+1)L_2)/q<2^{-147}$ & exact \\
KB safe, CH \textsc{import} & $r=2^{19}$ & $n/q<2^{-164}$ & --- \\
Circle safe, CH \textsc{import} & $r=2^{19}$ & $n/q<2^{-102}$ & --- \\
KB profile, PR & $\kappa=2^{54}$ & $\lceil\binom{256}{130}/p\rceil>1+2^{54}q T$ & exact \\
$i$-free circle, RF & $(m,w)=(2^{20}{+}2^{15}{+}1,\,2^{16})$ & $\binom{2^{21}}{m}>p'^{\,2^{16}}2^{256}$ & exact \\
\bottomrule
\end{tabular}
\caption{Printed certificates for the deployed rows (\cref{def:certificate-v2}; $T=\frac1k(1-\frac nq)$). Each line is one integer comparison; ``exact'' means it was verified in exact integer arithmetic. The two \textup{\textsc{import}} lines are optional imported half-distance clauses --- the only non-self-contained entries.}
\label{tab:certs}
\end{table}

\medskip
\noindent\textbf{Remaining residues.}
The refinements above reduce the remaining landscape to explicit residues.  The genus-one and circle-code questions are reduced to the same normalized safe-side package after the rational-floor and stereographic transports of \cref{cor:ecfft-macroscopic,cor:circle-anyfield}.  The explicit-witness question is resolved at the level of closed-form words and residue-line pairs in \cref{cor:explicit-deployed,thm:explicit-pairs}, with the remaining scope stated in \cref{rem:explicit-scope}.  The error-one profile is not settled here: \cref{cor:profile-deployed} raises the certified profile to $(1-2^{-54})\frac1k(1-\frac nq)$, while \cref{rem:profile-ceiling} explains why reaching $1-o(1)$ requires new input.  The frontier residue is the active package stated in \cref{subsec:capg-active-interface}: quotient/prefix flatness, the base-field-normalized split-pencil census, and primitive shift-pair control.  The certificate material above is finite and checkable; a proof-assistant transcription would be implementation work rather than a new mathematical ingredient.




\providecommand{\F}{\mathbb F}
\providecommand{\RS}{\mathrm{RS}}
\providecommand{\ImgFib}{\operatorname{ImgFib}}
\providecommand{\Paid}{\operatorname{Paid}}
\providecommand{\Dloc}{\mathcal D}
\providecommand{\PP}{\mathbb P}
\providecommand{\E}{\mathbb E}
\providecommand{\Var}{\operatorname{Var}}
\providecommand{\Span}{\operatorname{span}}
\providecommand{\UU}{\mathbb U}
\providecommand{\ceil}[1]{\left\lceil #1\right\rceil}
\providecommand{\floor}[1]{\left\lfloor #1\right\rfloor}

\section{Frontier package: statements, certificates, and remaining inputs}\label{sec:frontier-package}

This part of the paper contains the finite certificate calculus and the remaining safe-side interface.  The status is deliberately separated into proved certificates, imported generic safe anchors, and three final residual inputs.

\begin{center}\small
\begin{tabular}{@{}p{0.30\linewidth}p{0.26\linewidth}p{0.34\linewidth}@{}}
\toprule
component & status & role \\
\midrule
prefix floors, simple-pole conversion, and deployed integer certificates & proved, with exact integer arithmetic & supply the last unsafe agreements in the deployed rows \\

deep-regime MCA, tangent accounting, and MCA-from-CA pincer & proved, elementary & supply unconditional safe anchors and remove support-mismatch cells in the covered region \\

half-distance CA proximity gap & imported where explicitly marked & gives the generic half-distance safe edge used only in the broad sandwich \\

quotient/prefix flatness & residual input & controls worst prefix fibers and quotient recursions beyond the entropy--subfield envelope \\

base-field-normalized split-pencil census & residual input & controls the primitive interior aperiodic support census after slope elimination \\

primitive shift-pair control & residual input & controls the primitive part of the exact second-moment ledger \\
\bottomrule
\end{tabular}
\end{center}

The active finite deployed frontier is the following adjacent-pair prediction:
\[
\begin{array}{c|cc|c}
\text{row} & \text{proved unsafe agreement }a_0 &
\text{first safe }a_0+1 &
\text{margin (bits)}\\ \hline
\text{KoalaBear MCA} & 1116047 & 1116048 & 22.2\\
\text{KoalaBear list} & 1116046 & 1116047 & 22.0\\
\text{Mersenne-31 MCA} & 1116023 & 1116024 & 3.3\\
\text{Mersenne-31 list} & 1116022 & 1116023 & 3.1
\end{array}
\]
The unsafe entries are theorems, proved by the identity-prefix floor together with the flexible-budget conversion.  The safe entries are conjectural and are kept separate from the proved certificates.  Polynomial-loss estimates are sufficient only after an agreement reserve of order larger than \(\log n\); the finite adjacent pairs require explicit constants inside the displayed margins.

The final safe-side residual package consists of exactly three inputs:
\[
\boxed{\begin{gathered}
\text{quotient/prefix flatness},\\
\text{base-field-normalized split-pencil census},\\
\text{primitive shift-pair control}.
\end{gathered}}
\]
The quotient, rank-one, balanced-core, tangent, and prefix-collision sections below are reductions toward these inputs, not additional assumptions to be imposed in parallel.

\subsection{Active interface for the remaining inputs}
\label{subsec:capg-active-interface}

The proof below contains several reductions of the safe side.  They should be
read in the following order:
\[
\begin{gathered}
        \text{aperiodic/prefix split}
        \Longrightarrow
        \text{rank-one support census}
        \Longrightarrow
        \text{balanced-core and split-pencil census},\\
        \text{then the three active inputs in the boxed package above.}
\end{gathered}
\]
Only the last package is consumed by the final conditional closure.  The
intermediate problems are retained because their lemmas are part of the proof:
slope elimination explains how the aperiodic line problem becomes a support
census; the shifted-lattice normal form explains why the support census has a
split-pencil shape; the prefix-collision ledgers calibrate the quotient side.
They are not additional assumptions to be imposed in parallel with the final
package.

More precisely, whenever a closing theorem uses a safe-side input, the active
reading is:
\[
\boxed{
\begin{gathered}
\text{quotient/prefix flatness},\\
\text{base-field-normalized split-pencil census},\\
\text{primitive shift-pair control}.
\end{gathered}}
\]
The first input is formulated in \cref{prob:capg-active-Q}, the second in
\cref{prob:capg-active-BC}, and the third in
\cref{prob:capg-active-shiftpairs}.  In particular, the aperiodic band is read
with the normalized left edge of \cref{prob:capfr1-normalized-band}; below the
entropy--subfield envelope the identity-scale floor supplies genuine unsafe
mass.  Polynomial or \(e^{o(n)}\) losses are used only with an agreement
reserve.  At the printed adjacent pairs, every multiplier and additive error
must be included in the integer upper ledger and must fit inside the displayed
row margin.

\begin{proposition}[active closure interface]
\label{prop:capg-active-interface}
Read every safe-side closure statement below through
\[
        \cref{prob:capg-active-Q,prob:capg-active-BC,prob:capg-active-shiftpairs}.
\]
If these three inputs hold with polynomial or \(e^{o(n)}\) loss in the reserve
regime, then the asymptotic frontier follows with any reserve
\(R\gg\log n\).  If their finite forms hold with explicit constants below the
row margins of \cref{cor:capg-adjacent-pairs}, then the finite adjacent pairs
follow from the staircase compiler \cref{thm:capf-staircase} and the one-step
mechanism \cref{prop:onestep}.  No unspecified polynomial factor is promoted to
a finite adjacent-pair certificate.
\end{proposition}

\section{Threshold certificate compilers}\label{sec:capf-threshold}

\subsection{Integer staircases and corridors}

\begin{definition}[integer numerator staircase]\label{def:capf-staircase}
Fix an agreement interval
\[
        I=\{a_{\min},a_{\min}+1,\ldots,a_{\max}\}
\]
and a nonincreasing function
\[
        N:I\longrightarrow \mathbb Z_{\ge0}.
\]
For MCA or line-MCA rows, $N(a)$ is the number of bad sampled parameters at agreement at least $a$; the closed integer Hamming radius is $r=n-a$.  For list rows one may use the same agreement coordinate, so the list numerator is nonincreasing in $a$.  Given a denominator $Q$ and a target $\varepsilon^*$, put
\[
        B_*:=\floor{\varepsilon^* Q}.
\]
An agreement $a$ is \emph{safe} if $N(a)\le B_*$ and \emph{unsafe} if $N(a)>B_*$.
\end{definition}

\begin{theorem}[staircase localization]\label{thm:capf-staircase}
Let $N:I\to\mathbb Z_{\ge0}$ be nonincreasing.  Exactly one of the following holds:
\begin{enumerate}[label=\textup{(\roman*)}]
\item every $a\in I$ is safe;
\item every $a\in I$ is unsafe;
\item there is a unique first safe agreement $a_*\in I$ such that the safe set is $\{a\in I:a\ge a_*\}$ and the unsafe set is $\{a\in I:a<a_*\}$.
\end{enumerate}
In the interior case, when $a_*>a_{\min}$, one has
\[
        N(a_*-1)>B_*\ge N(a_*).
\]
\end{theorem}

\begin{proof}
The predicate $N(a)\le B_*$ is monotone nondecreasing in $a$: if it holds at $a$, then for $a'\ge a$ we have $N(a')\le N(a)\le B_*$.  A monotone Boolean function on a finite linearly ordered set is either identically false, identically true, or has a unique first true point.  The displayed inequalities are exactly the assertion that the point just before the first true point is false while the first true point is true.
\end{proof}

\begin{corollary}[closed-radius endpoint convention]\label{cor:capf-endpoint}
In the interior case of \cref{thm:capf-staircase}, the largest safe closed integer radius is
\[
        r_{\mathrm{safe}}=n-a_*,
\]
and the first unsafe closed integer radius above it is $r_{\mathrm{safe}}+1$.  As a real closed-ball threshold under the convention $a(\delta)=\ceil{(1-\delta)n}$ of \cref{sec:prize-program}, the supremum of safe radii is
\[
        \frac{n-a_*+1}{n},
\]
and this endpoint is not attained when $N(a_*-1)>B_*$.
\end{corollary}

\begin{proof}
The integer radius corresponding to agreement $a$ is $n-a$, so the first safe agreement $a_*$ gives the largest safe integer radius $n-a_*$.  A real radius $\delta$ is safe exactly when $\ceil{(1-\delta)n}\ge a_*$, which holds if and only if $(1-\delta)n>a_*-1$, i.e.\ for all $\delta<(n-a_*+1)/n$.  At $\delta=(n-a_*+1)/n$ the agreement requirement is $a_*-1$, which is unsafe in the interior case.
\end{proof}

\begin{theorem}[corridor from lower and upper certificates]\label{thm:capf-corridor}
Suppose that $L(a)\le N(a)\le U(a)$ are certified lower and upper bounds on the same staircase.  Define
\[
        A_{\mathrm{unsafe}}:=\{a\in I:L(a)>B_*\},
        \qquad
        A_{\mathrm{safe}}:=\{a\in I:U(a)\le B_*\}.
\]
If an interior first safe agreement $a_*$ exists, then
\[
        1+\max A_{\mathrm{unsafe}}\le a_*\le \min A_{\mathrm{safe}},
\]
with the conventions $1+\max\varnothing=a_{\min}$ and $\min\varnothing=a_{\max}$.
\end{theorem}

\begin{proof}
If $L(a)>B_*$, then $N(a)>B_*$, so $a$ is unsafe and the first safe point is strictly larger than $a$.  Taking the largest such certified unsafe point gives the lower bound.  If $U(a)\le B_*$, then $N(a)\le B_*$, so $a$ is safe and the first safe point is at most $a$.  Taking the smallest such certified safe point gives the upper bound.  The displayed conventions are exactly the empty-set cases.
\end{proof}

\begin{remark}[relation to the main interface]\label{rem:capf-staircase-main}
\Cref{prop:onestep} is the one-step case of \cref{thm:capf-staircase,thm:capf-corridor}: a matched pair of certificates $L(a_0)>B_*\ge U(a_0+1)$ pins $a_*=a_0+1$ exactly.  In the conditional theorems \cref{thm:conditional-mca,thm:conditional-list}, the hypothesized band inputs enter precisely as the upper staircase $U$, and the unsafe floors of \cref{sec:pincer} enter as the lower staircase $L$; \cref{thm:capf-corridor} is the exact bookkeeping of the sandwich between them at a fixed budget.
\end{remark}

\begin{theorem}[coarse steepness compiler]\label{thm:capf-steepness}
Let $J\subseteq I$ be an interval of consecutive agreements, let $M(a)>0$ be a positive model count, and suppose a row packet proves
\[
        \frac{M(a)}{E}\le N(a)\le E M(a)\qquad(E\ge1)
\]
for every $a\in J$.  For a single adjacent pair $a,a+1\in J$, if
\[
        \frac{M(a)}{M(a+1)}>E^2,
\]
then the two ambiguity intervals
\[
        \left[\frac{M(a)}E,EM(a)\right]
        \quad\text{and}\quad
        \left[\frac{M(a+1)}E,EM(a+1)\right]
\]
are disjoint, and a fixed budget $B_*$ can lie in at most one of them.

If the same separation holds for every adjacent pair in $J$, then the ambiguity intervals are pairwise disjoint in their natural order.  Consequently at most one agreement in $J$ can remain undecided by the coarse envelope for a fixed budget $B_*$; every other agreement has a certified safe or unsafe sign.
\end{theorem}

\begin{proof}
For the adjacent pair, the displayed inequality is equivalent to
\[
        \frac{M(a)}E>EM(a+1),
\]
which is precisely the assertion that the lower end of the $a$-interval is above the upper end of the $(a+1)$-interval.  If $B_*<M(b)/E$, then $N(b)\ge M(b)/E>B_*$ and $b$ is certified unsafe.  If $B_*\ge EM(b)$, then $N(b)\le EM(b)\le B_*$ and $b$ is certified safe.  Thus only budgets inside the ambiguity interval for $b$ can be undecided.

When the adjacent separation holds throughout $J$, induction gives
\[
        \frac{M(a)}E>EM(a+1)\ge \frac{M(a+1)}E>EM(a+2)\ge\cdots
\]
for consecutive points, using $E\ge1$, so all ambiguity intervals are pairwise disjoint.  A single integer budget can belong to at most one member of a pairwise disjoint collection of intervals, and the sign criterion of the earlier paragraph certifies all remaining agreements.
\end{proof}

\subsection{Paid ledger cells}

\begin{definition}[high-agreement tangent exact cell]\label{def:capf-tangent-cell}
For a Reed--Solomon row of length $n$ and dimension $k<n$, put $r(A)=n-A$ and
\[
        R_{\tan}:=\floor{\frac{n-k}{3}}.
\]
The high-agreement tangent paid cell is
\[
        \Paid_{\tan}^{\mathrm{hi}}(n,k,A):=r(A)+1
\]
when $0\le r(A)\le R_{\tan}$, and is unavailable outside this range unless a separate tangent/common-line theorem is supplied.
\end{definition}

\begin{proposition}[exact tangent cell]\label{prop:capf-tangent}
Let $C=\RS[F,D,k]$ with $|D|=n$ and $k<n$.  If $A=n-r$ and
\[
        3r\le n-k,
\]
then the maximum over received lines of the number of finite support-wise MCA-bad slopes at closed agreement threshold $A$ is exactly $r+1=n-A+1$, and the same maximum for no-loss CA-bad slopes is also exactly $r+1$.  Hence \cref{def:capf-tangent-cell} is an exact paid cell throughout the high-agreement range.
\end{proposition}

\begin{proof}
The upper bound is \cref{thm:deep-mca} applied at radius $r/n$: the Reed--Solomon minimum Hamming weight is $w=n-k+1$, so the hypothesis $3r\le w-1$ there is exactly $3r\le n-k$, and the theorem bounds the number of CA- or MCA-bad finite slopes of every received line by $r+1$.

For the matching lower witness, choose a set
\[
        T=\{t_0,\ldots,t_r\}\subseteq D
\]
of size $r+1$ \textup{(possible since $r+1\le n$)} and choose distinct field elements $\gamma_0,\ldots,\gamma_r\in F$ \textup{(possible since $|F|>n\ge r+1$ for $D\subseteq F^\times$, and $|F|\ge n\ge r+1$ in general)}.  Define two received words by
\[
        f_2(t_i)=1,
        \qquad
        f_1(t_i)=-\gamma_i,
        \qquad
        f_1(x)=f_2(x)=0\quad (x\in D\setminus T).
\]
For the slope $\gamma_i$, the word $f_1+\gamma_i f_2$ vanishes at $t_i$ and off $T$, hence agrees with the zero codeword on
\[
        S_i:=D\setminus (T\setminus\{t_i\}),
        \qquad |S_i|=n-r=A .
\]

We first record the two degree inequalities used below.  If $r=0$, then $n-r-1=n-1\ge k$ and $n-2r-1=n-1\ge k$.  If $r\ge1$, then $r+1\le 3r\le n-k$ gives $n-r-1\ge k$, and $2r\le 3r-r\le n-k-r\le n-k-1$ gives $n-2r-1\ge k$.

\emph{Support-wise badness.}  Suppose $P_1,P_2\in F[X]_{<k}$ satisfy $P_1=f_1$ and $P_2=f_2$ on $S_i$.  On $D\setminus T\subseteq S_i$ both received words vanish, and $|D\setminus T|=n-r-1\ge k$, so each $P_j$ has at least $k$ roots and is the zero polynomial.  But $t_i\in S_i$ and $f_2(t_i)=1\ne0$, a contradiction.  Hence the pair $(f_1,f_2)$ is not simultaneously explained on $S_i$, while the point $f_1+\gamma_i f_2$ is explained there by the zero codeword: each $\gamma_i$ is support-wise MCA-bad at agreement threshold $A$, and the $r+1$ slopes $\gamma_i$ are distinct.

\emph{No-loss badness.}  Let $S\subseteq D$ be any support with $|S|\ge A=n-r$ and suppose $P_1,P_2\in F[X]_{<k}$ explain $(f_1,f_2)$ on $S$.  Since $|T|=r+1$, inclusion--exclusion gives $|S\cap T|\ge|S|+|T|-n\ge1$, and
\[
        |S\cap(D\setminus T)|\ge |S|-|T|\ge n-2r-1\ge k .
\]
Hence both $P_j$ vanish at $\ge k$ points and are identically zero, contradicting $f_2(t)=1$ at any point $t\in S\cap T$.  So no support of size $\ge A$ simultaneously explains the pair, i.e.\ the pair is no-loss far at radius $r/n$; since each $f_1+\gamma_if_2$ is explained on $S_i$, the same $r+1$ slopes are no-loss CA-bad.  Combining the witness with the upper bound of \cref{thm:deep-mca} proves exactness.
\end{proof}

\begin{remark}[tangent-cell calibration]\label{rem:capf-tangent-calibration}
The exactness makes the high-agreement staircase pointwise computable: in the tangent range the exact numerator is $N(A)=n-A+1$, which is strictly decreasing, so \cref{thm:capf-staircase} localizes with a one-step transition.  As a worked integer example, take the calibration row $n=512$, $k=256$ with sampling denominator $Q=17^{32}$ and target $2^{-128}$.  Exact arithmetic gives $B_*=\floor{17^{32}/2^{128}}=6$ \textup{($17^{32}$ has $131$ bits)}, and the tangent range is $r\le R_{\tan}=85$, i.e.\ $A\ge427$.  Inside this range, $N(A)=n-A+1\le6$ first holds at $a_*=507$: agreement $506$ is unsafe with $N=7>6$ and agreement $507$ is safe with $N=6\le6$.  The largest safe closed integer radius for the tangent-range staircase is $r_{\mathrm{safe}}=5$, and the real closed-ball supremum of \cref{cor:capf-endpoint} is $6/512$, not attained.  This is a calibration example for the compiler, not a deployed row.
\end{remark}

\begin{definition}[quotient paid status cell]\label{def:capf-quotient-status}
Let $\mathcal C$ be a finite set of declared quotient fiber sizes $c\mid n$ and let $A$ be an agreement threshold.  The conservative safe-sum support count is
\[
U_{\mathrm{sum}}(n,A,\mathcal C):=
\sum_{c\in\mathcal C}\sum_{B=A}^{n}
\binom{n/c}{\floor{B/c}}
\binom{n-c\floor{B/c}}{B-c\floor{B/c}},
\]
the quotient-remainder count of \cref{sec:discharge}: for each $c$ and each size $B$, choose $\floor{B/c}$ complete $c$-fibers among the $n/c$ fibers and then a residual set of size $B-c\floor{B/c}$ among the remaining $n-c\floor{B/c}$ points.  The quotient paid cell at $(n,A,\mathcal C)$ is status-valued:
\[
\Paid_{\mathrm{quot}}=
\begin{cases}
\mathrm{EXACT\_IMAGE}, & \text{an image-lcm certificate as in \cref{thm:exact-quotient-image-lcm-ledger}},\\
\mathrm{EXACT\_SUPPORT}, & \text{an exact block-profile union as in \cref{thm:exact-divisor-block-support-ledger}},\\
\mathrm{SAFE\_SUM}(U_{\mathrm{sum}}), & \text{otherwise}.
\end{cases}
\]
Unsafe quotient lower cells should carry zone tags, for example norm-exact, collision-not settled here, or cap-floor.  A collision-open zone is interval-valued and must not be silently replaced by a point estimate.
\end{definition}

\begin{proposition}[safe-sum validity]\label{prop:capf-quotient-safe-sum}
Let $C=\RS[F,D,k]$ and $A\ge k$.  If a quotient packet declares that every finite bad slope in the cell is witnessed on at least one support counted by the quotient-remainder model above, then $U_{\mathrm{sum}}(n,A,\mathcal C)$ is a valid conservative upper numerator for the union of the declared quotient finite-slope cells at agreement at least $A$, for support-wise MCA-bad and for no-loss CA-bad slopes alike.
\end{proposition}

\begin{proof}
The displayed binomial product counts, for each declared $c$ and size $B$, the supports consisting of $\floor{B/c}$ complete $c$-fibers plus a residual set outside those fibers; summing over $B\ge A$ and $c\in\mathcal C$ is a union bound over the declared support family, which may double-count supports with several quotient descriptions but cannot undercount them.  By \cref{lem:one-support-one-line}, each fixed support pays at most one finite bad slope, in either badness sense.  The declared coverage hypothesis then converts the support union bound into a finite-slope numerator.  This is the safe-sum instantiation of \cref{prop:divisor-union-support-ledger}; when the exact union or the image-lcm coalescing is available, the sharper grades of \cref{def:capf-quotient-status} replace it.
\end{proof}

\begin{proposition}[extension-pole conversion cell]\label{prop:capf-extension}
Let $B\subseteq F$ be the generated field inside the line field, with $q_{\mathrm{gen}}=|B|$ and $q_{\mathrm{line}}=|F|$.  If $q_{\mathrm{gen}}=q_{\mathrm{line}}$, then there is no proper extension-only line-parameter cell:
\[
        \Paid_{\mathrm{ext}}^{\mathrm{only}}=0.
\]
Assume $q_{\mathrm{gen}}<q_{\mathrm{line}}$, put $m=q_{\mathrm{line}}-q_{\mathrm{gen}}$, and let $\mathcal P$ be a base list of size $L\ge1$.  Suppose a pole-conversion certificate assigns to every $P\in\mathcal P$ and every pole $\alpha\in F\setminus B$ a finite extension parameter $\lambda_\alpha(P)$ such that:
\begin{enumerate}[label=\textup{(\roman*)}]
\item for each fixed pole $\alpha$, distinct values among $\{\lambda_\alpha(P):P\in\mathcal P\}$ give distinct extension-only bad finite line parameters for the lifted row;
\item for every $P\ne P'$ in $\mathcal P$,
\[
        \#\{\alpha\in F\setminus B:\lambda_\alpha(P)=\lambda_\alpha(P')\}\le \kappa .
\]
\end{enumerate}
Then some pole produces at least
\[
        \mathrm{ExtPole}(q_{\mathrm{line}},q_{\mathrm{gen}},\kappa,L)
        :=
        \left\lceil
        \frac{L(q_{\mathrm{line}}-q_{\mathrm{gen}})}
        {q_{\mathrm{line}}-q_{\mathrm{gen}}+\kappa(L-1)}
        \right\rceil
\]
distinct extension-only bad finite line parameters.  When the conversion is evaluation at the pole, one may take $\kappa=k$ for a degree-$\le k$ auxiliary list and $\kappa=k-1$ for an ordinary degree-$<k$ Reed--Solomon list.  For safe-side extension charts over $B$, a closed parameter set of degree $\Delta$ and dimension $e$ contributes at most $\Delta q_{\mathrm{gen}}^{\,e}$ affine $B$-parameters, by \cref{thm:extension-line-dimension-degree-ledger}.
\end{proposition}

\begin{proof}
If $B=F$, then $F\setminus B=\varnothing$, so no line parameter is extension-only.  Now assume $m>0$.  For a pole $\alpha$, let $M_\alpha$ be the number of distinct values among $\{\lambda_\alpha(P):P\in\mathcal P\}$, let $n_{\alpha,\lambda}$ be the multiplicity of the value $\lambda$, and let $C_\alpha:=\sum_\lambda\binom{n_{\alpha,\lambda}}2$ be the number of colliding unordered pairs at $\alpha$.  Summing the pairwise certificate \textup{(ii)} over unordered pairs of list elements gives
\[
        \sum_{\alpha\in F\setminus B}C_\alpha
        \le \kappa\binom L2 ,
\]
so some pole $\alpha$ has $C_\alpha\le \kappa\binom L2/m$.  At this pole, Cauchy--Schwarz gives
\[
        L^2
        =\Bigl(\sum_\lambda n_{\alpha,\lambda}\Bigr)^2
        \le
        M_\alpha\sum_\lambda n_{\alpha,\lambda}^2
        =M_\alpha(L+2C_\alpha)
        \le
        M_\alpha\Bigl(L+\frac{\kappa L(L-1)}m\Bigr),
\]
hence
\[
        M_\alpha\ge \frac{Lm}{m+\kappa(L-1)} .
\]
Since $M_\alpha$ is an integer, the displayed ceiling is certified, and condition \textup{(i)} converts the distinct pole values into distinct extension-only bad line parameters.  For the evaluation conversion, two distinct polynomials of degree $\le k$ \textup{(resp.\ $<k$)} agree at most $k$ \textup{(resp.\ $k-1$)} times on $F\setminus B$, giving the stated $\kappa$.  The safe-side clause is \cref{thm:extension-line-dimension-degree-ledger} verbatim.
\end{proof}

\begin{remark}[relation to the main extension-pole floors]\label{rem:capf-extension-main}
\Cref{prop:punctured-simple-pole-floor} and \cref{cor:extension-pole-deep-list-floor} realize the conversion certificate concretely: the conversion is evaluation of the base list at a simple pole $\alpha\in F\setminus B$, condition \textup{(i)} holds with genuinely extension-valued witness lines, and condition \textup{(ii)} holds with $\kappa=k$ there.  The proof of \cref{cor:extension-pole-deep-list-floor} internally derives the denominator $m+k(L-1)$ before weakening it to $m+kL$ for display; \cref{prop:capf-extension} promotes the sharper denominator to the official compiler value.  The coarser denominator $q_{\mathrm{line}}-q_{\mathrm{gen}}+\kappa L$ remains a valid conservative fallback and may be printed when only the weaker collision certificate is available.
\end{remark}

\subsection{Budget windows and bounded quotient census}

\begin{theorem}[budget windows and dodge selection]\label{thm:capf-windows}
Let $N$ be the actual integer numerator for a fixed object, and suppose a certificate gives $L\le N\le K$.  Put $B(q)=\floor{q/2^{128}}$.  Then
\[
        B(q)<L \Rightarrow \text{certified unsafe},
        \qquad
        B(q)\ge K \Rightarrow \text{certified safe for this object},
\]
and the only unresolved budgets are
\[
        L\le B(q)<K.
\]
Equivalently, the unresolved field-size interval is
\[
        L2^{128}\le q\le K2^{128}-1.
\]
If the lower proof is $L=K-g$, then the largest tolerable missing mass while still proving unsafety is
\[
        g_{\max}=K-B(q)-1.
\]
\end{theorem}

\begin{proof}
The safety condition for the fixed object is $N\le B(q)$.  If $B(q)<L\le N$, then $N>B(q)$ and the object certifies unsafety.  If $N\le K\le B(q)$, then the object is safe under this certificate.  The complement is precisely $L\le B(q)<K$.  The floor definition turns this double inequality into the displayed field-size interval.  Finally, if $L=K-g$, unsafety requires $K-g>B(q)$, equivalently $g<K-B(q)$, and the largest such integer is $K-B(q)-1$.
\end{proof}

For a $2$-power quotient order $N'$, let $n_1=N'/2$.  The characteristic-zero antipodal quotient count used by the quotient census of \cite{Cho26b} is
\[
        A_2(N',\ell')
        :=
        \sum_{\substack{u\ge0,\ t=\ell'-2u\ge0\\ u\le n_1-t}}
        \binom{n_1}{t}2^t .
\]

\begin{lemma}[rate-$1/2$ census identity]\label{lem:capf-census-identity}
For every even $N'\ge2$, with $n_1=N'/2$,
\[
        A_2(N',n_1+1)=\frac{3^{n_1}-1}{2}.
\]
\end{lemma}

\begin{proof}
At $\ell'=n_1+1$ the summation constraints reduce to a single parity class.  Indeed $t=n_1+1-2u$, so $t\equiv n_1+1\pmod2$; the constraint $u\le n_1-t$ becomes $u\le n_1-(n_1+1-2u)=2u-1$, i.e.\ $u\ge1$, which is equivalent to $t\le n_1-1$; and $t\ge0$ closes the range.  Hence
\[
        A_2(N',n_1+1)
        =\sum_{\substack{0\le t\le n_1-1\\ t\equiv n_1+1\ (\mathrm{mod}\ 2)}}\binom{n_1}t2^t
        =\sum_{\substack{0\le t\le n_1\\ t\not\equiv n_1\ (\mathrm{mod}\ 2)}}\binom{n_1}t2^t,
\]
where the last step notes that $t=n_1$ has the excluded parity, and $t=n_1+1$ is out of range.  From $\sum_t\binom{n_1}t2^t=3^{n_1}$ and $\sum_t\binom{n_1}t(-2)^t=(-1)^{n_1}$, the odd-$t$ class sums to $(3^{n_1}-(-1)^{n_1})/2$ and the even-$t$ class to $(3^{n_1}+(-1)^{n_1})/2$.  If $n_1$ is even the class $t\not\equiv n_1$ is the odd class and $(-1)^{n_1}=1$; if $n_1$ is odd it is the even class and $(-1)^{n_1}=-1$.  In both cases the sum is $(3^{n_1}-1)/2$.
\end{proof}

\begin{theorem}[bounded quotient census, exact arithmetic certificate]\label{thm:capf-census}
Let $B_{\max}=2^{128}-1$.  In the relaxed even-scale model, where $N'$ ranges over the even integers with $\rho N'\in\mathbb Z$, the first admissible scales at which $A_2(N',\rho N'+1)>B_{\max}$ are
\[
\begin{array}{c|cccc}
\rho & 1/2 & 1/4 & 1/8 & 1/16\\ \hline
N'_{\mathrm{relaxed}} & 164 & 176 & 248 & 384 .
\end{array}
\]
For dyadic quotient orders $N'=2^a$, the corresponding first scales are
\[
\begin{array}{c|cccc}
\rho & 1/2 & 1/4 & 1/8 & 1/16\\ \hline
N'_{\mathrm{dyadic}} & 256 & 256 & 256 & 512 .
\end{array}
\]
Thus the exact integer quotient endgame at the $2^{-128}$ budget sees only bounded quotient scales, independent of the deployed blocklength, once the row has been reduced to this canonical quotient-count model.
\end{theorem}

\begin{proof}
The count $A_2(N',\ell')$ is an explicit finite integer sum, so the statement is an exact arithmetic certificate.  For each fixed $\rho$, the certificate evaluates $A_2(N',\rho N'+1)$ by exact integer arithmetic at \emph{every} admissible scale up to and including the printed first crossing, confirms $A_2\le B_{\max}$ at every admissible scale strictly below the crossing, and confirms $A_2>B_{\max}$ at the crossing; no floating-point estimate and no unproved monotonicity is used.  Writing $\operatorname{bits}(X)=\floor{\log_2X}+1$, so that $\operatorname{bits}(X)\le128\iff X\le B_{\max}$, the immediate predecessor and crossing values are
\[
\begin{array}{c|cc|cc}
\rho & N'_{\mathrm{prev}} & \operatorname{bits} A_2 & N'_{\mathrm{first}} & \operatorname{bits} A_2\\ \hline
1/2  & 162 & 128 & 164 & 129\\
1/4  & 172 & 127 & 176 & 130\\
1/8  & 240 & 127 & 248 & 131\\
1/16 & 368 & 124 & 384 & 129
\end{array}
\qquad
\begin{array}{c|cc|cc}
\rho & N'_{\mathrm{prev}} & \operatorname{bits} A_2 & N'_{\mathrm{first}} & \operatorname{bits} A_2\\ \hline
1/2  & 128 & 101 & 256 & 202\\
1/4  & 128 &  95 & 256 & 190\\
1/8  & 128 &  68 & 256 & 135\\
1/16 & 256 &  87 & 512 & 172
\end{array}
\]
for the relaxed even model \textup{(left)} and the dyadic model \textup{(right)}.  At rate $1/2$ the scanned values agree with the closed form of \cref{lem:capf-census-identity}, which is strictly increasing in $N'$, so the rate-$1/2$ rows also follow from the identity alone.
\end{proof}

\begin{definition}[dyadic planted quotient profile]\label{def:capf-qprofile}
For $n=2^m$, $k=\rho n$, and integer slack $\sigma=\ceil{\eta n}$, define
\[
K_Q(n,k,\sigma):=
\max_{\substack{M=2^e\mid\gcd(n,k)\\ \sigma<M\\ k/M\le n/M-1}}
\binom{n/M-1}{k/M},
\]
with value $0$ if the active set is empty.
\end{definition}

\begin{proposition}[bounded active quotient order]\label{prop:capf-qprofile}
If $\sigma\ge n/256$, then every active scale in $K_Q$ has quotient order $N=n/M<256$; because $N$ is dyadic in \cref{def:capf-qprofile}, this means $N\le128$.  If $\sigma\ge n/128$, then every active scale has $N<128$, hence dyadically $N\le64$.
\end{proposition}

\begin{proof}
Activity requires $\sigma<M=n/N$, hence $N<n/\sigma$.  The two displayed strict bounds follow from the two lower bounds on $\sigma$, and the dyadic constraint improves the strict inequalities to $N\le128$ and $N\le64$ respectively.
\end{proof}

\begin{proposition}[polynomial residual absorption]\label{prop:capf-integrality}
Fix a candidate scale and suppose a future theorem bounds an unpaid residual contribution by
\[
        n^C\Phi(n,q)
\]
for some explicit proxy $\Phi$.  If the verified entropy calculation gives
\[
        -\log_2\bigl(n^C\Phi(n,q)\bigr)>0,
\]
then the residual integer contribution is zero.
\end{proposition}

\begin{proof}
The residual contribution is a nonnegative integer bounded above by a real number strictly smaller than $1$.  The only such integer is $0$.
\end{proof}

\section{L1 list-side compilers and sunflower residuals}\label{sec:capf-l1}

Let $D\subset F$ be an evaluation domain and let $C=\RS[F,D,k]$ be a Reed--Solomon code of length $n=|D|$.  For a received word $U:D\to F$ and an agreement threshold $s$, write
\[
        \ImgFib_U^{<k}(s):=
        \{P\in F[X]_{<k}: \#\{x\in D:U(x)=P(x)\}\ge s\},
\]
and, for an auxiliary degree bound $d$ and a subdomain $T\subseteq D$,
\[
        \ImgFib_W^{\le d}(s;T):=
        \{G\in F[X]_{\le d}: \#\{x\in T:W(x)=G(x)\}\ge s\}.
\]
These are the actual list objects used below; they count codewords rather than raw supports.

\subsection{Planted quotient-core lower counts}

\begin{theorem}[planted quotient-core lower count]\label{thm:capf-planted}
Let $D=\alpha H\subset F^\times$ be a multiplicative coset of a cyclic subgroup $H$ of order $n$, let $C=\RS[F,D,k]$, and suppose $M\mid\gcd(n,k)$ satisfies
\[
        1\le \sigma<M,
        \qquad
        k/M\le n/M-1.
\]
Then there is a received word $U:D\to F$ with at least
\[
        P_M(n,k):=\binom{n/M-1}{k/M}
\]
distinct codewords agreeing with $U$ on at least $k+\sigma$ positions.  Consequently
\[
\max_U \bigl|\ImgFib_U^{<k}(k+\sigma)\bigr|
\ge
\max_{\substack{M\mid\gcd(n,k)\\ \sigma<M\\ k/M\le n/M-1}}
\binom{n/M-1}{k/M}.
\]
\end{theorem}

\begin{proof}
Let $K\le H$ be the unique subgroup of order $M$ and let $Q:=\{x^M:x\in D\}$ be the quotient image, so that the map $\pi:D\to Q$, $\pi(x)=x^M$, has fibers equal to the $K$-cosets inside $D$, each of size exactly $M$, and $|Q|=N:=n/M$.  Choose one quotient value $b_0\in Q$ and a set $T\subseteq\pi^{-1}(b_0)$ of size $\sigma$; this is possible since $\sigma<M$.  Put
\[
        E_T(X):=\prod_{t\in T}(X-t).
\]
For every subset $A\subseteq Q\setminus\{b_0\}$ of size $k/M$, define
\[
        L_A(X):=E_T(X)\prod_{b\in A}(X^M-b).
\]
Every point of the fiber $\pi^{-1}(b)$ with $b\in A$ is a root of $X^M-b$, and every point of $T$ is a root of $E_T$; these point sets are disjoint since $b_0\notin A$.  Hence $L_A$ vanishes at at least $M\cdot(k/M)+\sigma=k+\sigma$ points of $D$.

Since
\[
        \prod_{b\in A}(X^M-b)=X^k+\text{terms of degree at most }k-M,
\]
and $\deg E_T=\sigma<M$, the part of $L_A$ of degree at least $k$ is independent of $A$, namely
\[
        U_0(X):=X^kE_T(X).
\]
Thus $L_A=U_0+R_A$ with $\deg R_A<k$.  Let $U$ be the received word $x\mapsto U_0(x)$ on $D$ and let $P_A:=-R_A\in F[X]_{<k}$.  For every root $x\in D$ of $L_A$,
\[
        0=L_A(x)=U(x)-P_A(x),
\]
so $P_A$ agrees with $U$ on at least $k+\sigma$ positions.

It remains to check distinctness.  If $P_A$ and $P_{A'}$ give the same codeword on $D$, then $R_A-R_{A'}$ has degree $<k<n$ and vanishes on all $n$ points of $D$, so $R_A=R_{A'}$ as polynomials, hence $L_A=L_{A'}$, and dividing by the fixed nonzero polynomial $E_T$ gives
\[
        \prod_{b\in A}(X^M-b)=\prod_{b\in A'}(X^M-b).
\]
The substitution homomorphism $F[Y]\to F[X]$, $G(Y)\mapsto G(X^M)$, is injective \textup{(the monomials $X^{Mi}$ are linearly independent)}, so $\prod_{b\in A}(Y-b)=\prod_{b\in A'}(Y-b)$ in $F[Y]$, and unique factorization into the distinct monic linear factors gives $A=A'$.  \textup{(This argument needs no separability of $X^M-b$ and is valid in every characteristic.)}  There are $\binom{N-1}{k/M}$ choices of $A$, which proves the claim.
\end{proof}

\begin{corollary}[complete polynomial-folding planted core]\label{cor:capf-polyfold-planted}
Let $H\subset F$ be finite, let $\phi\in F[X]$ have degree $M>0$, and assume the map $\phi:H\to Q$ is onto with every fiber $\phi^{-1}(b)\cap H$ of size exactly $M$.  Suppose
\[
        M\mid k,
        \qquad
        1\le\sigma<M,
        \qquad
        k/M\le |Q|-1 .
\]
Then there is a received word $U:H\to F$ with at least
\[
        \binom{|Q|-1}{k/M}
\]
distinct degree-$<k$ Reed--Solomon codewords agreeing with $U$ on at least $k+\sigma$ positions.
\end{corollary}

\begin{proof}
Choose $b_0\in Q$ and a subset $T\subseteq \phi^{-1}(b_0)\cap H$ of size $\sigma$.  Put
\[
        E_T(X):=\prod_{t\in T}(X-t),
        \qquad
        \ell:=k/M.
\]
For every $A\subseteq Q\setminus\{b_0\}$ of size $\ell$, define
\[
        L_A(X):=E_T(X)\prod_{b\in A}(\phi(X)-b).
\]
Its roots in $H$ include the $\sigma$ planted points of $T$ and the $M\ell=k$ points of the complete fibers over $A$, disjointly, hence at least $k+\sigma$ points.  Since
\[
        \prod_{b\in A}(Y-b)=Y^\ell+\text{terms of degree at most }\ell-1,
\]
we have
\[
        L_A(X)=E_T(X)\phi(X)^\ell+R_A(X),
        \qquad
        \deg R_A\le \sigma+M(\ell-1)<k.
\]
Let $U$ be the word $x\mapsto E_T(x)\phi(x)^\ell$ on $H$ and put $P_A:=-R_A\in F[X]_{<k}$; then $U=P_A$ at every root of $L_A$ in $H$.

For distinctness, note $|H|=M|Q|$ by the complete-fiber hypothesis, and $k=M\ell\le M(|Q|-1)<|H|$.  If $P_A$ and $P_{A'}$ define the same codeword on $H$, then $R_A-R_{A'}$ has degree $<k<|H|$ and vanishes on $H$, hence $R_A=R_{A'}$ and
\[
        \prod_{b\in A}(\phi(X)-b)=\prod_{b\in A'}(\phi(X)-b).
\]
The substitution map $F[Y]\to F[X]$, $G(Y)\mapsto G(\phi(X))$, is injective because $\deg\phi>0$, so $\prod_{b\in A}(Y-b)=\prod_{b\in A'}(Y-b)$ and $A=A'$.  The number of choices is $\binom{|Q|-1}{\ell}$.
\end{proof}

\begin{remark}[relation to the main fiber lemmas and extension persistence]\label{rem:capf-planted-main}
\Cref{thm:capf-planted} is the exact-count planted form of the quotient-core construction of \cite{Cho26b}: a single received word carries the whole binomial sublist, with no pigeonhole loss and no subfield denominator.  By contrast, \cref{lem:fiber} pigeonholes locator subsets over their slope values to extract one heavy fiber at slack two; the planted form is the right tool on the list side, where no slope average is available.  \Cref{cor:capf-polyfold-planted} is the map-smooth analogue: its complete-uniform-fiber hypothesis is exactly the condition of \cref{def:map-smooth} used by \cref{lem:phi-fiber}, so the planted core applies verbatim to Chebyshev and circle-derived foldings.  Finally, all these list floors persist over every field extension: $\RS[B,D,k]\subseteq\RS[F,D,k]$ for $D\subseteq B\subseteq F$, so the planted sublist survives unchanged when the row is lifted, exactly as in the list-side persistence discussion of \cref{sec:subfield}.
\end{remark}

\begin{corollary}[list unsafe test]\label{cor:capf-list-unsafe}
For a list denominator $q_{\mathrm{list}}$ and target $2^{-128}$, the planted quotient-core packet proves a list row unsafe whenever
\[
\max_M\binom{n/M-1}{k/M}>
\floor{\frac{q_{\mathrm{list}}}{2^{128}}},
\]
where the maximum ranges over the active divisors of \cref{thm:capf-planted}.
\end{corollary}

\begin{proof}
By \cref{thm:capf-planted}, the maximum list size at agreement $k+\sigma$ is at least the displayed planted count.  If that count exceeds the allowed numerator $\floor{q_{\mathrm{list}}/2^{128}}$, the normalized list quantity exceeds $2^{-128}$.
\end{proof}

\begin{proposition}[dyadic planted crossings]\label{prop:capf-dyadic-planted}
For $n=2^m$, $k=\rho n$, and
\[
        \rho\in\{1/2,\,1/4,\,1/8,\,1/16\},
\]
the first dyadic quotient orders $N=n/M$ at which the planted count exceeds $2^{128}-1$ are
\[
\begin{array}{c|cccc}
\rho & 1/2 & 1/4 & 1/8 & 1/16\\ \hline
N & 256 & 256 & 256 & 512 ,
\end{array}
\]
and the exact bit lengths of $\binom{N-1}{\rho N}$ at the predecessor and first crossing are
\[
\begin{array}{c|cc|cc}
\rho & N_{\mathrm{prev}} & \operatorname{bits} & N_{\mathrm{first}} & \operatorname{bits}\\ \hline
1/2  & 128 & 124 & 256 & 251\\
1/4  & 128 & 100 & 256 & 204\\
1/8  & 128 &  67 & 256 & 136\\
1/16 & 256 &  83 & 512 & 169.
\end{array}
\]
\end{proposition}

\begin{proof}
The planted count at quotient order $N$ is $\binom{N-1}{\rho N}$, and $\operatorname{bits}(X)\le128\iff X\le2^{128}-1$; the displayed bit lengths were checked by exact integer arithmetic on binomial coefficients, not floating-point logarithms.  It remains to justify that the predecessor check identifies the first dyadic crossing.  Along dyadic orders the counts are monotone: from a subset $S\subseteq\{1,\ldots,N-1\}$ of size $\rho N$, form a subset of $\{1,\ldots,2N-1\}$ of size $2\rho N$ by adjoining a fixed set of $\rho N$ elements of $\{N,\ldots,2N-1\}$.  This map is injective, so
\[
        \binom{N-1}{\rho N}\le \binom{2N-1}{2\rho N}.
\]
Hence once the predecessor is below the budget, every smaller dyadic order is below it as well.
\end{proof}

\begin{corollary}[planted list budget windows]\label{cor:capf-list-windows}
If a planted exact count is $P$ and the certified lower count is $L\le P$, then the only unresolved list budgets are
\[
        L\le \floor{q_{\mathrm{list}}/2^{128}}<P,
\]
equivalently $L2^{128}\le q_{\mathrm{list}}\le P2^{128}-1$.
\end{corollary}

\begin{proof}
This is \cref{thm:capf-windows} with $q=q_{\mathrm{list}}$, $K=P$.
\end{proof}

\subsection{Agreement-threshold Johnson bounds}

\begin{theorem}[$\kappa$-intersecting agreement bound]\label{thm:capf-johnson-list}
Let $\Omega$ be a finite set of size $n_0$, let $W:\Omega\to F$ be a received word, and let $\mathcal F$ be a family of words $\Omega\to F$ such that any two distinct members of $\mathcal F$ agree on at most $\kappa$ points of $\Omega$.  For an agreement threshold $s$ with $0<s\le n_0$, let $L$ be the number of members of $\mathcal F$ agreeing with $W$ on at least $s$ points.
\begin{enumerate}[label=\textup{(\roman*)}]
\item If $2s>n_0+\kappa$, then $L\le1$.
\item If $s^2>\kappa n_0$, then
\[
        L
        \le
        \frac{n_0(s-\kappa)}{s^2-\kappa n_0}.
\]
\end{enumerate}
In particular, for $C=\RS[F,\Omega,k]$ and $\kappa=k-1$,
\[
        \bigl|\ImgFib_W^{<k}(s)\bigr|\le\frac{n_0(s-k+1)}{s^2-(k-1)n_0}
        \qquad\text{when }s^2>(k-1)n_0,
\]
and for the auxiliary degree-$\le d$ code on a petal domain, $\kappa=d$.
\end{theorem}

\begin{proof}
Let $P_1,\ldots,P_L$ be the listed members and $A_i:=\{x\in\Omega:P_i(x)=W(x)\}$, so $|A_i|\ge s$ and, for $i\ne j$, $|A_i\cap A_j|\le\kappa$ by the intersecting hypothesis \textup{(two words agreeing with $W$ at a point agree with each other there)}.

\textup{(i)}  For $L\ge2$, inclusion--exclusion gives $|A_1\cap A_2|\ge2s-n_0>\kappa$, a contradiction.

\textup{(ii)}  Assume $L\ge1$ and put $m_x:=\#\{i:x\in A_i\}$ and $I:=\sum_{x\in\Omega}m_x\ge Ls$.  Counting incident pairs,
\[
        \sum_{x\in\Omega}\binom{m_x}2=\sum_{i<j}|A_i\cap A_j|\le\binom L2\kappa .
\]
First suppose $Ls\le n_0$.  Then $L\le n_0/s$, and since $s\le n_0$,
\[
        \frac{n_0}s\le\frac{n_0(s-\kappa)}{s^2-\kappa n_0}
        \iff
        s^2-\kappa n_0\le s^2-\kappa s
        \iff
        \kappa s\le\kappa n_0,
\]
which holds; note $s-\kappa>0$ because $s^2>\kappa n_0\ge\kappa s$.  Now suppose $Ls>n_0$.  The function $\varphi(y)=y(y-1)$ is convex and nondecreasing on $[1,\infty)$, so Jensen and $I/n_0\ge Ls/n_0>1$ give
\[
        n_0\cdot\frac{Ls}{n_0}\Bigl(\frac{Ls}{n_0}-1\Bigr)
        \le
        n_0\,\varphi\Bigl(\frac I{n_0}\Bigr)
        \le
        \sum_x m_x(m_x-1)
        \le
        L(L-1)\kappa .
\]
Multiplying out, $Ls(Ls-n_0)\le n_0L(L-1)\kappa$, hence $s(Ls-n_0)\le n_0(L-1)\kappa$, i.e.
\[
        L(s^2-\kappa n_0)\le n_0(s-\kappa),
\]
which is the claim.
\end{proof}

\begin{remark}[relation to the main Johnson theorem]\label{rem:capf-johnson-main}
For $\kappa=k-1$ this is exactly the $m=1$ case of \cref{thm:johnson-list}, restated in agreement-threshold list coordinates; the numerator $n_0(s-k+1)$ is the sharp one and strictly improves the numerator $n_0(n_0-k+1)$ that a plain Cauchy--Schwarz argument yields.  The abstract $\kappa$-intersecting form is recorded because the sunflower reduction below applies it with $\kappa=d$ to an auxiliary code that is not a row of the challenge.
\end{remark}

\subsection{Sunflower and petal interfaces}

The statements in this subsection are compiler reductions: they become theorems about an actual sunflower decomposition only when the decomposition supplies the stated layer maps.  This is the safe capf form: the paid pieces are proved, while the mixed-petal and growing-defect classification remains a named residual branch.  The first proposition records one concrete situation in which the required layer injection is automatic.

\begin{definition}[concrete sunflower core layer]\label{def:capf-concrete-sunflower}
Let $\Omega\subset F$ be a finite evaluation set for $\RS[F,\Omega,k]$.  Fix pairwise disjoint subsets
\[
        Y,\quad R_0,\quad T_1,\ldots,T_M\subseteq\Omega,
        \qquad |Y|=k-1,
        \qquad |T_i|=\ell=\sigma+1,
\]
a reference polynomial $P_\star\in F[X]_{<k}$, and labels $c_1,\ldots,c_M\in F$.  Put
\[
        L_Y(X)=\prod_{y\in Y}(X-y),
\]
and assume the received word $U:\Omega\to F$ satisfies
\[
        U=P_\star\quad\text{on }Y\cup R_0,
        \qquad
        U=P_\star+c_iL_Y\quad\text{on }T_i .
\]
For $D_0\subseteq Y$ write $d=|D_0|$, $L_{D_0}(X)=\prod_{y\in D_0}(X-y)$, and $T:=\bigcup_iT_i$.  The concrete $(D_0,R_0)$ layer consists of those $P\in F[X]_{<k}$ that agree with $U$ on at least $k+\sigma$ points of $\Omega$, agree with $U$ on all of $(Y\setminus D_0)\cup R_0$, and have no agreements with $U$ outside $(Y\setminus D_0)\cup R_0\cup T$.
\end{definition}

\begin{proposition}[concrete sunflower layer injection]\label{prop:capf-concrete-sunflower}
In the setting of \cref{def:capf-concrete-sunflower}, the map
\[
        P\longmapsto
        G_P:=\frac{P-P_\star}{L_{Y\setminus D_0}}
\]
is well-defined and injects the concrete $(D_0,R_0)$ layer into $F[X]_{\le d}$.  Moreover $G_P$ agrees on at least
\[
        a=\sigma+d+1-|R_0|
\]
points of the petal domain $T$ with the auxiliary word
\[
        U_{D_0}(x):=c_iL_{D_0}(x),\qquad x\in T_i,
\]
and if the layer retains the residual set $R_0$, then in addition $G_P(x)=0$ for every $x\in R_0$.
\end{proposition}

\begin{proof}
Since $P$ and $P_\star$ agree on $Y\setminus D_0$ \textup{(both equal $U=P_\star$ there)}, the difference $P-P_\star$ is divisible by $L_{Y\setminus D_0}$, which has degree $k-1-d$; the quotient $G_P$ therefore has degree at most $(k-1)-(k-1-d)=d$.  The quotient determines $P=P_\star+G_PL_{Y\setminus D_0}$, so the map is injective.

The layer has at least $k+\sigma$ agreements in total, of which at most $|Y\setminus D_0|+|R_0|=(k-1-d)+|R_0|$ lie off the petal domain, by the containment condition of the definition.  Hence at least
\[
        (k+\sigma)-(k-1-d)-|R_0|=\sigma+d+1-|R_0|=a
\]
agreement points lie in $T$.  If $x\in T_i$ is such a point, then $x\notin Y$, so $L_{Y\setminus D_0}(x)\ne0$, and
\[
        G_P(x)L_{Y\setminus D_0}(x)
        =P(x)-P_\star(x)
        =U(x)-P_\star(x)
        =c_iL_Y(x)
        =c_iL_{D_0}(x)L_{Y\setminus D_0}(x);
\]
cancelling the nonzero factor gives $G_P(x)=c_iL_{D_0}(x)=U_{D_0}(x)$.  Finally, if $x\in R_0$, then $P(x)=U(x)=P_\star(x)$ and $x\notin Y$, so $G_P(x)=0$.
\end{proof}

\begin{definition}[fixed sunflower auxiliary layer]\label{def:capf-sunflower-layer}
Fix a missed-core set $D_0$ of size $d$ with monic locator $L_{D_0}(X)=\prod_{y\in D_0}(X-y)$, a retained residual subset $R_0$ of size $r$, pairwise disjoint petals $T_1,\ldots,T_M$ of common size $\sigma+1$ disjoint from $D_0\cup R_0$, and labels $c_i\in F$.  Let $T:=\bigcup_iT_i$ and define
\[
        U_{D_0}(x):=c_iL_{D_0}(x),\qquad x\in T_i.
\]
A fixed $(D_0,R_0)$ sunflower layer is a finite set $\mathcal L(D_0,R_0)$ together with an injective map
\[
        \Psi:\mathcal L(D_0,R_0)\hookrightarrow F[X]_{\le d}
\]
such that every $G=\Psi(P)$ agrees with $U_{D_0}$ on at least
\[
        a:=\sigma+d+1-r
\]
points of $T$.  Extra conditions forcing zeros on $R_0$ are called retained zero constraints.
\end{definition}

\begin{proposition}[sunflower core-defect reduction]\label{prop:capf-pma}
For every fixed sunflower auxiliary layer in the sense of \cref{def:capf-sunflower-layer},
\[
        |\mathcal L(D_0,R_0)|
        \le
        \bigl|\ImgFib_{U_{D_0}}^{\le d}(a;T)\bigr|,
\]
the auxiliary degree-$\le d$ list on the petal domain $T$ at agreement $a=\sigma+d+1-r$.  Dropping the retained zero constraints on $R_0$ only enlarges the right-hand side, so it is safe for upper bounds.
\end{proposition}

\begin{proof}
By definition, $\Psi$ injects the fixed layer into the set of degree-$\le d$ polynomials agreeing with $U_{D_0}$ on at least $a$ points of $T$, which is the displayed auxiliary list.  Adding zero constraints cuts out a subset of it; removing them can only increase the counted set.
\end{proof}

\begin{corollary}[few-petal Johnson boundary]\label{cor:capf-pma-johnson}
In the notation of \cref{prop:capf-pma}, the auxiliary petal domain has size $|T|=M(\sigma+1)$, and two distinct degree-$\le d$ polynomials agree on at most $d$ points, so \cref{thm:capf-johnson-list} applies with $\kappa=d$.  If $a^2>d|T|$, the fixed $(D_0,R_0)$ auxiliary list has size at most
\[
        \frac{|T|\,(a-d)}{a^2-d\,|T|}.
\]
Equivalently, for $d>0$, the few-petal Johnson-covered range is exactly
\[
        M\le
        \floor{\frac{a^2-1}{d(\sigma+1)}}.
\]
\end{corollary}

\begin{proof}
The bound is \cref{thm:capf-johnson-list}\textup{(ii)} with $n_0=|T|$, $s=a$, $\kappa=d$ \textup{(if $a>|T|$ the list is empty and the bound is vacuous)}.  Since all quantities are integers, the strict inequality $a^2>M(\sigma+1)d$ is equivalent to $M(\sigma+1)d\le a^2-1$, which is the stated floor.
\end{proof}

\begin{definition}[full-petal fixed-excess counting chart]\label{def:capf-full-petal-chart}
A full-petal counting chart with cofactor excess $e=d-\ell$ is a parameterization of a full-petal layer with the following certificate; the certificate is for the whole layer being charged, and if several independent charts are used the bounds below must be summed over the chart index.
\begin{enumerate}[label=\textup{(\roman*)}]
\item the zero-excess layer $e=0$ is determined by an unordered pair of petals and one scalar in $F$;
\item for every $e\ge1$, the layer is determined by a subset of the $M$ petals and at most $e+1$ scalars in $F$.
\end{enumerate}
\end{definition}

\begin{proposition}[distinct-label zero-excess full-petal chart]\label{prop:capf-zero-excess-chart}
Work in a concrete sunflower layer \textup{(\cref{def:capf-concrete-sunflower})} with $R_0=\varnothing$, petal size $\ell=\sigma+1$, and $d=\ell$, and assume the labels $c_1,\ldots,c_M$ are pairwise distinct.  Consider a residual-free full-petal sublayer: every listed codeword $P$ agrees with $U$ on the full union of a set $I(P)$ of complete petals, and the listedness inequality forces $|I(P)|\ge2$.  Then the map
\[
        P\longmapsto(\{i,j\},\alpha_i),
\]
where $i<j$ are the two smallest indices in $I(P)$ and $\alpha_i$ is the scalar constructed in the proof, is injective, so the sublayer has size at most
\[
        \binom M2\,|F| .
\]
Moreover the injection recovers the missed-core locator $L_{D_0}$ itself, hence $D_0$; consequently the same bound $\binom M2|F|$ holds for the \emph{union} of these sublayers over all missed-core sets $D_0\subseteq Y$ of size $\ell$.  In particular the chart supplies the zero-excess certificate of \cref{def:capf-full-petal-chart}.
\end{proposition}

\begin{proof}
Write $L:=L_{D_0}$ and let $L_i$ be the monic locator of the petal $T_i$; both are monic of degree $\ell$.  For a codeword $P$ in the sublayer write, as in \cref{prop:capf-concrete-sunflower},
\[
        P-P_\star=G\,L_{Y\setminus D_0},
        \qquad \deg G\le d=\ell .
\]
If $i\in I(P)$, then $P$ agrees with $U=P_\star+c_iL_Y$ at every point of $T_i$, so by \cref{prop:capf-concrete-sunflower} the polynomial $G-c_iL$, of degree at most $\ell$, vanishes at all $\ell$ points of $T_i$; hence
\[
        G-c_iL=\alpha_iL_i
        \qquad\text{for some }\alpha_i\in F .
\]
Subtracting the identities for the two smallest touched petals $i<j$,
\[
        (c_j-c_i)L=\alpha_iL_i-\alpha_jL_j .
\]
Comparing leading coefficients of the three monic degree-$\ell$ locators gives
\[
        c_j-c_i=\alpha_i-\alpha_j,
        \qquad\text{so}\qquad
        \alpha_j=\alpha_i-(c_j-c_i).
\]
Given the pair $\{i,j\}$ and the scalar $\alpha_i$, the label difference $c_j-c_i\ne0$ determines $\alpha_j$, then
\[
        L=\frac{\alpha_iL_i-\alpha_jL_j}{c_j-c_i}
\]
determines $L$, then $G=c_iL+\alpha_iL_i$, and finally $P=P_\star+G\,L_{Y\setminus D_0}$, where $D_0$ is recovered as the root set of the monic squarefree polynomial $L$ inside $Y$ and $L_{Y\setminus D_0}=L_Y/L$.  Thus the map is injective, even jointly over the choice of $D_0$, and there are at most $\binom M2|F|$ values.
\end{proof}

\begin{proposition}[distinct-label fixed-excess full-petal chart]\label{prop:capf-distinct-excess-chart}
Work in a concrete sunflower layer \textup{(\cref{def:capf-concrete-sunflower})} with $R_0=\varnothing$, petal size $\ell=\sigma+1$, and pairwise distinct labels $c_1,\ldots,c_M$.  Fix $e\ge1$ and put $d=\ell+e$.  Consider the charged pairs $(P,D_0)$, where $D_0\subseteq Y$ has size $d$ and $P$ lies in a residual-free all-full sublayer with missed core $D_0$: there is a touched set $I(P,D_0)\subseteq[M]$ such that $P$ agrees with $U$ on every point of $T_i$ for $i\in I(P,D_0)$, has no petal agreements for $i\notin I(P,D_0)$, and has no agreements outside $(Y\setminus D_0)\cup\bigcup_{i\in I(P,D_0)}T_i$.  Then, for this fixed $e$, the number of charged pairs is at most
\[
        2^M |F|^{e+1}.
\]
Consequently the union of these residual-free all-full distinct-label sublayers over all missed cores $D_0$ also has size at most $2^M|F|^{e+1}$.  Equivalently, in this distinct-label residual-free all-full regime the higher-excess certificate of \cref{def:capf-full-petal-chart} is automatic, not an extra hypothesis.
\end{proposition}

\begin{proof}
Let $(P,D_0)$ be such a charged pair, put $L:=L_{D_0}$, and let $L_i$ be the monic locator of the petal $T_i$.  As in \cref{prop:capf-concrete-sunflower}, write
\[
        P-P_\star=G\,L_{Y\setminus D_0},
        \qquad \deg G\le d=\ell+e .
\]
If $i\in I:=I(P,D_0)$, full-petal agreement on $T_i$ gives
\[
        G-c_iL=L_iB_i
\]
for a uniquely determined polynomial $B_i$ of degree at most $e$, because the left side has degree at most $\ell+e$ and vanishes on the $\ell$ points of $T_i$.

The listedness threshold is $k+\sigma=k+\ell-1$.  Since the sublayer is residual-free and all-full, all agreements are contained in the $(k-1-d)$ retained core points and the $|I|\ell$ full petal points.  Hence
\[
        (k-1-d)+|I|\ell\ge k+\ell-1,
        \qquad\text{so}\qquad
        (|I|-1)\ell\ge d .
\]
In particular $I$ has at least two elements.  Let $i_0$ be the least element of $I$.  For every $j\in I\setminus\{i_0\}$, subtracting the full-petal identities gives
\[
        (c_{i_0}-c_j)L+L_{i_0}B_{i_0}=L_jB_j,
\]
so, since $c_{i_0}\ne c_j$,
\[
        L\equiv -(c_{i_0}-c_j)^{-1}L_{i_0}B_{i_0}\pmod {L_j}.
\]
The petal locators are pairwise coprime, so the touched set $I$ and the single polynomial $B_{i_0}$ determine, by the Chinese remainder theorem, the residue class of $L$ modulo
\[
        Q_I:=\prod_{j\in I\setminus\{i_0\}}L_j,
        \qquad \deg Q_I=(|I|-1)\ell\ge d .
\]
There is at most one monic polynomial of degree $d$ in this residue class: if $\deg Q_I>d$ this is immediate because the residue representative has degree $<\deg Q_I$, and if $\deg Q_I=d$ the difference of two monic degree-$d$ representatives has degree $<d$ while being divisible by $Q_I$.  Thus $L=L_{D_0}$ is recovered from $(I,B_{i_0})$.

Once $L$ is known, the actual charged pair recovers $D_0$ as the root set of the monic squarefree locator $L$ inside $Y$, then recovers $L_{Y\setminus D_0}=L_Y/L$, then recovers
\[
        G=c_{i_0}L+L_{i_0}B_{i_0},
        \qquad
        P=P_\star+G L_{Y\setminus D_0}.
\]
Hence the map $(P,D_0)\mapsto (I,B_{i_0})$ is injective.  There are at most $2^M$ possible touched sets and at most $|F|^{e+1}$ possible polynomials $B_{i_0}$ of degree at most $e$, proving the bound.
\end{proof}

\begin{remark}[higher-excess full petals after the distinct-label chart]\label{rem:capf-higher-excess-cert}
\Cref{prop:capf-distinct-excess-chart} pays residual-free all-full higher-excess petals when the petal labels are pairwise distinct.  For repeated labels, retained residual sets, partial or mixed petals, or charts that do not satisfy the residual-free all-full containment hypothesis, Chinese-remainder reconstructions may still certify the $2^Mq^{e+1}$ chart of \cref{def:capf-full-petal-chart}, but the reconstruction data must be printed as a certificate.  The part does not infer those remaining higher-excess charts from listedness alone.
\end{remark}

\begin{proposition}[fixed-excess full-petal compiler]\label{prop:capf-fixed-excess}
For every fixed $E\ge0$, any single certified full-petal atlas satisfying \cref{def:capf-full-petal-chart} has total size at most
\[
        \binom M2 q + 2^M\sum_{e=1}^{E}q^{e+1},
        \qquad q=|F| ,
\]
the certified layer bounds being $\binom M2q$ at $e=0$ and $2^Mq^{e+1}$ at each $e\ge1$.
\end{proposition}

\begin{proof}
For $e=0$ the certificate chooses an unordered pair of petals and one scalar: at most $\binom M2q$ possibilities.  For fixed $e\ge1$ it chooses a subset of the $M$ petals in at most $2^M$ ways and $e+1$ scalars in at most $q^{e+1}$ ways.  In the residual-free pairwise-distinct-label all-full case, \cref{prop:capf-distinct-excess-chart} proves exactly this certificate; in other cases it remains a printed atlas certificate.  Summing over $1\le e\le E$ proves the claim.
\end{proof}

\begin{corollary}[fixed-excess polynomiality]\label{cor:capf-fixed-excess-poly}
At the intended L1 lower cutoff, suppose $M=O(\log n)$, $q\le n^Q$ for a fixed $Q$, and $E$ is fixed.  Then the fixed-excess full-petal contribution certified by \cref{prop:capf-fixed-excess} is polynomial in $n$.
\end{corollary}

\begin{proof}
The term $\binom M2q$ is polynomial because $M=O(\log n)$ and $q\le n^Q$.  For fixed $e\le E$, the term $2^Mq^{e+1}$ is $n^{O(1)}\cdot n^{Q(e+1)}$, hence polynomial, and a fixed finite sum of polynomially bounded terms is polynomial.
\end{proof}

\begin{remark}[L1 evidence is not a theorem]\label{rem:capf-l1-evidence}
The small-field auxiliary scans and worst-word sweeps that motivated the sunflower split, together with the explicit growing-defect full-petal examples, should be retained only as evidence and route cuts.  They support the residual decomposition but do not prove polynomial growth for the mixed-petal or growing-defect branches.
\end{remark}

\begin{remark}[L1 primitive image-fiber reduction]\label{prob:capf-primitive-image-fiber}
Prove a polynomial bound for the stabilizer-primitive part of $\ImgFib_U^{<k}(k+\sigma)$ once quotient-periodic codewords are charged separately and the generated-field entropy reserve and lower cutoff clear.
\end{remark}

\begin{remark}[mixed-petal and growing-defect residuals]\label{prob:capf-l1-residuals}
Classify or bound mixed-petal sunflower extras in the sub-Johnson region after the few-petal Johnson layer \textup{(\cref{cor:capf-pma-johnson})}, the distinct-label all-full fixed-excess layer \textup{(\cref{prop:capf-distinct-excess-chart})}, and any certified fixed-excess full-petal layers \textup{(\cref{prop:capf-fixed-excess})} are paid.  Determine whether the full-petal branch with $d-\ell\to\infty$ remains polynomial or produces a new obstruction ledger.
\end{remark}

\section{M1 residue-line tools and Conjecture-F reductions}\label{sec:capf-m1}

The safe-side program asks for upper bounds on bad line parameters after tangent, quotient, and extension cells have been charged.  In exact-agreement notation,
\[
        j=n-A,
        \qquad
        t=A-k,
        \qquad
        j+t=n-k,
\]
and the M1 problem is to control aperiodic residue-line packings in this $(j,t)$ window --- the residual kernel of the broad band formulation \cref{prob:band}.  The results below are reusable local algebra and combinatorics; they do not prove the full worst-case M1 local limit.

\subsection{Split-locator moment calculus}

Let $F=\F_q$, let $D\subset F^\times$ have size $n$, and let $k<A\le n$.  Put $t=A-k\ge1$ and $j=n-A$.  For $R\subseteq D$ with $|R|=j$, define
\[
        \ell_R(X)=\prod_{r\in R}(X-r)
\]
and the locator-syndrome vector of a word $w:D\to F$,
\[
        S_R(w)=
        \Bigl(\sum_{x\in D}w(x)\,\ell_R(x)\,x^m\Bigr)_{m=0}^{t-1}
        \in F^t .
\]
The shifted convention $m=1,\ldots,t$ is equivalent on $D\subset F^\times$: it is the composite of $S_R$ with the linear automorphism $w(x)\mapsto xw(x)$ of $F^D$, which preserves the uniform measure.  For independent uniform words $u,v:D\to F$, call $R$ \emph{aligned} if $S_R(v)\ne0$ and $S_R(u)\in F\,S_R(v)$, and let $N_A(u,v)$ be the number of aligned locators.

\begin{theorem}[exact first moment]\label{thm:capf-first-moment}
\[
        \E N_A
        =
        \binom{n}{j}\,(1-q^{-t})\,q^{1-t}.
\]
\end{theorem}

\begin{proof}
Fix $R$ and put $E:=D\setminus R$, $|E|=n-j=A$.  The linear map $w\mapsto S_R(w)$ has matrix entries $\ell_R(x)x^m$; the columns indexed by $x\in R$ vanish, and on $E$ we have $\ell_R(x)\ne0$, so scaling the column of $x$ by $\ell_R(x)^{-1}$ leaves the $t\times A$ Vandermonde block $(x^m)_{0\le m<t,\ x\in E}$ on distinct points.  Since $A=k+t\ge t$, any $t$ of its columns form an invertible Vandermonde matrix, so the map is surjective onto $F^t$.  A surjective linear image of a uniform vector is uniform, and $u,v$ are independent, so $(S_R(u),S_R(v))$ is uniform on $F^t\times F^t$.  The number of pairs $(a,b)$ with $b\ne0$ and $a\in Fb$ is $(q^t-1)q$, giving alignment probability $(q^t-1)q/q^{2t}=(1-q^{-t})q^{1-t}$; summing over the $\binom nj$ locators proves the formula.
\end{proof}

\begin{theorem}[two-locator joint rank]\label{thm:capf-joint-rank}
For $R,T\subseteq D$ of size $j$, put $c=|R\cap T|$.  The linear map
\[
        w\longmapsto (S_R(w),S_T(w))\in F^t\times F^t
\]
has rank
\[
        2t-\max(0,\,t-j+c).
\]
\end{theorem}

\begin{proof}
A linear relation among the $2t$ coordinate functionals is a pair of polynomials $A_0,B_0\in F[X]_{<t}$ such that $\sum_x w(x)\bigl(\ell_R(x)A_0(x)+\ell_T(x)B_0(x)\bigr)=0$ for all $w$, i.e.\ such that
\[
        \ell_R(X)A_0(X)+\ell_T(X)B_0(X)
\]
vanishes on all of $D$.  Its degree is at most $j+t-1\le n-1<n$, so it is the zero polynomial.  Let $G=\gcd(\ell_R,\ell_T)$, the locator of $R\cap T$, and write $\ell_R=Gr$, $\ell_T=Gs$ with $\deg G=c$, $\deg r=\deg s=j-c$, and $\gcd(r,s)=1$.  The identity becomes $rA_0+sB_0=0$, so $s\mid A_0$ and $r\mid B_0$: $A_0=sH$, $B_0=-rH$ for a polynomial $H$ with
\[
        \deg H\le t-1-(j-c).
\]
Hence the relation space has dimension $\max(0,\,t-j+c)$, and the rank is $2t$ minus this nullity.
\end{proof}

\begin{corollary}[covariance support]\label{cor:capf-covariance}
The alignment indicators for $R$ and $T$ are independent whenever
\[
        d(R,T):=j-|R\cap T|\ge t .
\]
All covariance is supported in the radius-$(t-1)$ Johnson neighborhood.
\end{corollary}

\begin{proof}
The condition $d(R,T)\ge t$ is $t-j+|R\cap T|\le0$, so by \cref{thm:capf-joint-rank} the map $w\mapsto(S_R(w),S_T(w))$ is surjective onto $F^{2t}$; hence $S_R(w)$ and $S_T(w)$ are independent uniform vectors for each of the independent words $u,v$, and the four syndrome vectors are mutually independent.  The alignment event for $R$ depends only on $(S_R(u),S_R(v))$ and that for $T$ only on $(S_T(u),S_T(v))$, so the two events are independent.
\end{proof}

\begin{theorem}[exact second moment]\label{thm:capf-second-moment}
Put $h(c):=\max(0,\,t-j+c)$ and, for $0\le h\le t$,
\[
P_h:=
\frac{q\,(q^h-1)\,q^{2(t-h)}+q^2\,(q^{t-h}-1)^2}{q^{4t-2h}} .
\]
Then
\[
\E N_A^2
=
\sum_{c=\max(0,2j-n)}^{j}
\binom{n}{j}\binom{j}{c}\binom{n-j}{j-c}\,
P_{h(c)} .
\]
\end{theorem}

\begin{proof}
Fix an ordered pair of locators $(R,T)$ with $|R\cap T|=c$ and put $h=h(c)$; the number of such ordered pairs is $\binom nj\binom jc\binom{n-j}{j-c}$, and $c$ ranges over $\max(0,2j-n)\le c\le j$.  It suffices to show that the probability that both $R$ and $T$ are aligned equals $P_h$.

\emph{Normal form of the joint image.}  Let $V\subseteq F^t\oplus F^t$ be the image of $w\mapsto(S_R(w),S_T(w))$; by \cref{thm:capf-joint-rank}, $\dim V=2t-h$, and both coordinate projections $\pi_1,\pi_2:V\to F^t$ are surjective \textup{(each single-locator map is surjective, as in the first-moment proof)}.  Put
\[
        W_1:=\{a\in F^t:(a,0)\in V\},
        \qquad
        W_2:=\{b\in F^t:(0,b)\in V\}.
\]
Since $\pi_1$ is surjective with kernel $\{0\}\oplus W_2$, we get $\dim W_2=(2t-h)-t=t-h$, and symmetrically $\dim W_1=t-h$.  The correspondence $a+W_1\mapsto b+W_2$ for $(a,b)\in V$ is a well-defined linear isomorphism $F^t/W_1\to F^t/W_2$ between $h$-dimensional spaces: well-defined because $(a,b),(a,b')\in V$ force $b-b'\in W_2$, injective and surjective by the symmetric statement and the dimension count.  Choose splittings $F^t=Z_i\oplus W_i$ with $\dim Z_i=h$, identify $Z_1\cong Z_2\cong F^h$ through the isomorphism, and correct the $W$-components by a linear section of $V$ over $Z_1$.  In the resulting coordinates,
\[
        V=U_h:=\{(z,p;\,z,q):z\in F^h,\ p,q\in F^{t-h}\},
\]
with the first block representing $S_R$ and the second $S_T$; alignment is invariant under these blockwise invertible changes of coordinates because each acts by an invertible linear map on the corresponding syndrome vector, preserving the relation $S(u)\in F\,S(v)$ and the condition $S(v)\ne0$.

\emph{Counting.}  The pair of words $(u,v)$ maps to a uniform element of $U_h\times U_h$, of size $q^{2(2t-h)}=q^{4t-2h}$.  Write the syndrome of $v$ as $(z,p;z,p')$.  If $z\ne0$ \textup{($q^h-1$ choices, with $p,p'$ free in $q^{2(t-h)}$ ways)}, then both $S_R(v)$ and $S_T(v)$ are nonzero, and simultaneous alignment of $u$ forces a common scalar: $S_R(u)=\lambda_R S_R(v)$ and $S_T(u)=\lambda_T S_T(v)$ share the first block $z_u=\lambda_Rz=\lambda_Tz$, so $\lambda_R=\lambda_T=:\lambda$, and each $\lambda\in F$ gives exactly one admissible syndrome of $u$, which lies in $U_h$; this contributes $q(q^h-1)q^{2(t-h)}$ aligned pairs.  If $z=0$, alignment requires $p\ne0$ and $p'\ne0$ \textup{($(q^{t-h}-1)^2$ choices)}, the two scalars are unconstrained and each choice $(\lambda_R,\lambda_T)\in F^2$ gives a distinct admissible syndrome of $u$ inside $U_h$; this contributes $q^2(q^{t-h}-1)^2$.  Dividing the total by $q^{4t-2h}$ gives $P_h$.
\end{proof}

\begin{corollary}[exact variance ledger]\label{cor:capf-exact-variance}
Let $p:=(1-q^{-t})q^{1-t}$.  Then
\[
\Var(N_A)=
\sum_{c=\max(0,2j-n)}^{j}
\binom{n}{j}\binom{j}{c}\binom{n-j}{j-c}
\bigl(P_{h(c)}-p^2\bigr),
\]
and the terms with $j-c\ge t$ vanish.
\end{corollary}

\begin{proof}
By \cref{thm:capf-first-moment} each locator is aligned with probability $p$, and the Vandermonde identity $\sum_c\binom jc\binom{n-j}{j-c}=\binom nj$ decomposes $(\E N_A)^2=\binom nj^2p^2$ over the same ordered overlap classes as \cref{thm:capf-second-moment}; subtracting termwise gives the display.  If $j-c\ge t$, then $h(c)=0$ and $P_0=q^2(q^t-1)^2/q^{4t}=p^2$, so those covariance terms are zero, in accordance with \cref{cor:capf-covariance}.
\end{proof}

\begin{corollary}[dependency-degree concentration]\label{cor:capf-dependency}
Let
\[
        D_t(n,j):=
        \sum_{d=0}^{\min(t-1,\,j,\,n-j)}
        \binom jd\binom{n-j}{d}
\]
be the size of the closed radius-$(t-1)$ Johnson ball.  Then
\[
        \Var(N_A)\le \E N_A\cdot D_t(n,j),
\]
so $\E N_A\ge D_t(n,j)\,n^s$ implies $\Var(N_A)/(\E N_A)^2\le n^{-s}$.  The coarser sufficient condition $\E N_A\ge t\,n^{2(t-1)+s}$ implies the same conclusion.
\end{corollary}

\begin{proof}
Write $N_A=\sum_RI_R$ with alignment indicators $I_R$.  By \cref{cor:capf-covariance}, $\operatorname{Cov}(I_R,I_T)=0$ unless $d(R,T)\le t-1$; the number of such $T$ for fixed $R$ is $\sum_{d\le t-1}\binom jd\binom{n-j}d=D_t(n,j)$ \textup{(choose the $d$ points of $R$ to remove and the $d$ points outside to add)}.  For every pair, $\operatorname{Cov}(I_R,I_T)\le\E(I_RI_T)\le\E I_R$.  Hence
\[
        \Var(N_A)=\sum_{R,T}\operatorname{Cov}(I_R,I_T)
        \le \sum_R D_t(n,j)\,\E I_R
        =\E N_A\cdot D_t(n,j),
\]
and dividing by $(\E N_A)^2$ gives the exact condition.  The coarse condition follows from $D_t(n,j)\le\sum_{d=0}^{t-1}n^{2d}\le t\,n^{2(t-1)}$.
\end{proof}

\begin{remark}[random-word scope]\label{rem:capf-moment-scope}
The moment calculus is an average over random word pairs.  It does not bound worst-case pairs; those require the aperiodic M1 local limit of the broad band formulation \cref{prob:band}, a Conjecture-F incidence input, or exchange-rigidity input.
\end{remark}

\subsection{Conjecture-F reductions}

Let $K$ be a field, let $H\subset K$ be finite, and put $P_H(X)=\prod_{h\in H}(X-h)$.  Let $\Dloc_j(H)$ be the set of monic squarefree degree-$j$ divisors of $P_H$, i.e.\ of locators $L_S(X)=\prod_{h\in S}(X-h)$ with $S\subseteq H$, $|S|=j$.  Conjecture-F asks for upper bounds on $|\PP(W)\cap\Dloc_j(H)|$ for linear subspaces $W\le K[X]_{\le j}$; the lemmas below pay its reducible strata.

\begin{lemma}[common-GCD reduction]\label{lem:capf-gcd}
Let $W\le K[X]_{\le j}$ be a nonzero vector space, let $P=\PP(W)$, and let $E=P\cap\Dloc_j(H)$.  If every locator in $E$ is divisible by a fixed monic divisor $G$ of $P_H$ with $\deg G=w$, then division by $G$ injects $E$ into
\[
        \Dloc_{j-w}\bigl(H\setminus Z(G)\bigr)
\]
inside a projective linear space of dimension at most $\dim P$.
\end{lemma}

\begin{proof}
Division by the fixed polynomial $G$ is a linear injection from the subspace $W\cap GK[X]$ into $K[X]_{\le j-w}$; write $W_G$ for its image, so $\dim\PP(W_G)\le\dim P$.  If $[L]\in E$ with $L$ its monic representative, then $L/G$ is monic of degree $j-w$ and, since $L$ is squarefree with roots in $H$ and contains all roots of $G$, the quotient $L/G$ is squarefree with roots exactly $Z(L)\setminus Z(G)\subseteq H\setminus Z(G)$.  Thus $[L]\mapsto[L/G]$ is the claimed injection into $\PP(W_G)\cap\Dloc_{j-w}(H\setminus Z(G))$.
\end{proof}

\begin{lemma}[quotient-pullback recursion]\label{lem:capf-quot-pullback}
Assume $H=\mu_n$ and $\operatorname{char}K\nmid n$.  If $M\mid\gcd(n,j)$, then
\[
        g(Y)\longmapsto g(X^M)
\]
maps $\Dloc_{j/M}(\mu_{n/M})$ bijectively onto the $M$-periodic stratum of $\Dloc_j(\mu_n)$, i.e.\ onto the locators whose root sets are unions of complete $\mu_M$-cosets.  Intersecting with a projective linear space descends to an intersection with a projective linear space of no larger dimension in the quotient locator space.
\end{lemma}

\begin{proof}
Because $\operatorname{char}K\nmid n$, the polynomial $X^n-1$ is squarefree, and the $M$-power map $\mu_n\to\mu_{n/M}$ is surjective with all fibers of size $M$.  A monic squarefree divisor $g$ of $Y^{n/M}-1$ of degree $j/M$ pulls back to $g(X^M)$, a monic divisor of $X^n-1$ of degree $j$ whose root set is the union of the fibers over $Z(g)$, hence $M$-periodic and squarefree.  Conversely, an $M$-periodic squarefree divisor $L$ of $X^n-1$ has root set a union of fibers; its image $S\subseteq\mu_{n/M}$ has size $j/M$, and $g:=\prod_{b\in S}(Y-b)$ satisfies $g(X^M)=L$ \textup{(both are monic with the same roots inside the squarefree $X^n-1$)}.  This proves the bijection.  For the descent, the space $\{g:g(X^M)\in W\}$ is linearly isomorphic to $W\cap K[X^M]$ by injectivity of the substitution, hence has dimension at most $\dim W$.
\end{proof}

\begin{lemma}[dimension-one voting]\label{lem:capf-dim1}
Assume $j\ge1$.  Let $P=\PP(\Span(L_0,L_1))$ be a projective line in $K[X]_{\le j}$, and assume no $h\in H$ is a common zero of $L_0$ and $L_1$.  Then
\[
        |P\cap\Dloc_j(H)|\le\floor{\frac{|H|}{j}}.
\]
\end{lemma}

\begin{proof}
For each $h\in H$, the functional $[a:b]\mapsto aL_0(h)+bL_1(h)$ is nonzero by the common-zero hypothesis, so it has a unique kernel point on the projective line; say that $h$ votes for this point.  A locator point $[aL_0+bL_1]\in P\cap\Dloc_j(H)$ has exactly $j$ roots in $H$ and receives a vote precisely from each of them.  There are $|H|$ votes in total, so the number of locator points is at most $|H|/j$, hence at most the floor.
\end{proof}

\begin{lemma}[hyperplane concurrency]\label{lem:capf-concurrency}
Let $W\le K[X]_{\le j}$ be nonzero and assume no $h\in H$ is a common root of all elements of $W$.  For each $h$, let
\[
        E_h:=\{[L]\in\PP(W):L(h)=0\},
\]
a projective hyperplane of $\PP(W)$.  Then $\PP(W)\cap\Dloc_j(H)$ is exactly the set of projective points lying on at least $j$ of the hyperplanes $E_h$.
\end{lemma}

\begin{proof}
If $[L]\in\PP(W)\cap\Dloc_j(H)$, then $L$ has exactly $j$ distinct roots in $H$, so $[L]$ lies on the corresponding $j$ hyperplanes.  Conversely, if a nonzero $L\in W$ vanishes at $j$ or more distinct points of $H$, then $\deg L\le j$ forces $L$ to have exactly $j$ roots, all simple and all among those points; after monic normalization, $L\in\Dloc_j(H)$.
\end{proof}

\begin{theorem}[projective-plane pair bound]\label{thm:capf-dim2}
Under the hypotheses of \cref{lem:capf-concurrency}, assume $j\ge2$ and $\dim\PP(W)=2$.  Then
\[
        |\PP(W)\cap\Dloc_j(H)|
        \le
        \frac{\binom{|H|}{2}}{j-1},
\]
and if the evaluation lines $E_h$ are pairwise distinct, the sharper bound with denominator $\binom j2$ holds.
\end{theorem}

\begin{proof}
Here the hyperplanes $E_h$ are lines in the projective plane $\PP(W)$; distinct lines meet in exactly one point.  Group the $E_h$ into coincidence classes.  No class has multiplicity $\ge j$: if $E_{h_1}=\cdots=E_{h_j}$ for distinct $h_1,\ldots,h_j$, this common line is a two-dimensional projective space of polynomials of degree $\le j$ vanishing at the $j$ fixed points $h_1,\ldots,h_j$; but that vanishing space is spanned by the single locator $\prod_i(X-h_i)$, hence at most zero-dimensional projectively --- a contradiction.

Now let $[L]$ be a locator point and let $m_1,\ldots,m_s$ be the multiplicities of the line classes through $[L]$ among the $j$ lines $E_h$, $h\in Z(L)$, so $\sum_i m_i=j$ and each $m_i\le j-1$; in particular $s\ge2$.  Call a pair $\{h,h'\}\subseteq Z(L)$ a \emph{cross pair} if $E_h\ne E_{h'}$.  Fixing any one class, of multiplicity $m$, the number of cross pairs at $[L]$ is at least $m(j-m)$, which for integers $1\le m\le j-1$ is minimized at the endpoints with value $j-1$.  Each cross pair $\{h,h'\}$ determines the unique intersection point $E_h\cap E_{h'}$, so it can be charged to at most one locator point.  Since there are at most $\binom{|H|}2$ pairs of points of $H$ in total, the number of locator points is at most $\binom{|H|}2/(j-1)$.  When all $E_h$ are pairwise distinct, every locator point is incident to at least $j$ distinct lines, hence carries at least $\binom j2$ cross pairs, giving the sharper denominator.
\end{proof}

\begin{theorem}[fixed-dimensional Conjecture-F bound]\label{thm:capf-fixeddim}
Let $W\le K[X]_{\le j}$ have vector-space dimension $d+1$ with $0\le d\le j$, and assume $W$ is gcd-trivial on $H$, i.e.\ no point of $H$ is a common root of all elements of $W$.  Then
\[
        |\PP(W)\cap\Dloc_j(H)|\le \binom{|H|}{d}.
\]
More generally, if $C_0\subseteq H$ is the common root set of $W$ on $H$ and $d\le j-|C_0|$, then
\[
        |\PP(W)\cap\Dloc_j(H)|\le \binom{|H|-|C_0|}{d}.
\]
\end{theorem}

\begin{proof}
First assume gcd-triviality and fix a total order on $H$.  For a locator point $[L]$ with root set $R\subseteq H$, $|R|=j$, the vanishing subspace
\[
        W(-R):=\{G\in W:G|_R=0\}
\]
equals $K\cdot L$: it contains $L$, and any $G\in W\subseteq K[X]_{\le j}$ vanishing at the $j$ points of $R$ is a scalar multiple of the degree-$j$ locator of $R$.

Process the roots of $R$ in the fixed order, imposing one vanishing condition at a time, and record the roots at which the dimension strictly drops; call this set $S_0\subseteq R$.  At every stage the current subspace equals $W(-R')$ for the set $R'$ of processed roots, and a skipped root leaves the subspace unchanged, so by induction the current subspace equals $W(-S')$ for the set $S'$ of \emph{kept} roots so far; in particular $W(-S_0)=W(-R)=K\cdot L$ after all of $R$ is processed.  Each vanishing condition drops the dimension by at most one, and the total drop is $(d+1)-1=d$, so $|S_0|=d$ exactly.  The assignment $[L]\mapsto S_0$ is therefore a map into the $d$-subsets of $H$, and it is injective because $S_0$ determines $W(-S_0)=K\cdot L$, i.e.\ the point $[L]$.  This gives the bound $\binom{|H|}d$.

For the general version, every element of $W$ vanishes on $C_0$ and has degree $\le j$, so $W\subseteq G_{C_0}K[X]$ for $G_{C_0}:=\prod_{h\in C_0}(X-h)$; moreover every locator in the intersection vanishes on $C_0$, hence is divisible by $G_{C_0}$.  By \cref{lem:capf-gcd}, division by $G_{C_0}$ injects the intersection into a gcd-trivial instance on $H\setminus C_0$ with locator degree $j-|C_0|$ and the same vector-space dimension $d+1$ \textup{(gcd-trivial because a new common root in $H\setminus C_0$ would have been a common root of $W$)}.  The first part applies whenever $d\le j-|C_0|$ and gives $\binom{|H|-|C_0|}d$.
\end{proof}

\begin{remark}[what remains to be settled]\label{rem:capf-conjf-open}
\Cref{thm:capf-fixeddim} shows that the primitive Conjecture-F core is a dimension-growing problem: common-root, quotient-periodic, dimension-one, and projective-plane strata are paid by \cref{lem:capf-gcd,lem:capf-quot-pullback,lem:capf-dim1,thm:capf-dim2}, and every fixed-dimensional stratum is paid by \cref{thm:capf-fixeddim}, but the growing-dimensional incidence theorem needed for the aperiodic band is not settled here and is part of the broad band formulation \cref{prob:band}.
\end{remark}

\subsection{Quotient seam arithmetic and Hankel tools}

\begin{lemma}[GAP-2 quotient seam]\label{lem:capf-gap2}
Let $n,k,A$ be fixed and put $j=n-A$, $t=A-k$, and $r=n-k$, so $j+t=r$.  If $M\mid\gcd(n,k)$, then $M\mid r$, and hence
\[
        M\mid j \quad\Longleftrightarrow\quad M\mid t.
\]
On such buckets all quotient parameters $n/M,k/M,A/M,j/M,t/M$ are integral.  If $M\mid n$ and $M\mid j$ but $M\nmid k$, then the support descends but the dimension and window do not: a non-rate-preserving seam.
\end{lemma}

\begin{proof}
Since $M\mid n$ and $M\mid k$, we have $M\mid n-k=r=j+t$, so $j\equiv-t\pmod M$, proving the equivalence, and then $A=n-j$ and $t$ descend whenever $j$ does.  If $M\mid n$ and $M\mid j$ but $M\nmid k$, then $A=n-j$ is divisible by $M$ while $k/M$ and $t/M=(A-k)/M$ are not integral.
\end{proof}

For a finite ordered set $X$ of nonzero field elements and $m\ge1$, write $V_m(X)$ for the matrix with rows indexed by $x\in X$, columns $0\le r<m$, and entries $x^r$.

\begin{proposition}[Hankel factorization]\label{prop:capf-hankel}
Let $H\subseteq F^\times$ be finite.  For a word $w:H\to F$, define the rectangular syndrome Hankel block
\[
        H_{t,s}(w)_{a,b}:=\sum_{x\in H}w(x)x^{a+b}
        \qquad(0\le a<t,\ 0\le b<s).
\]
Then
\[
        H_{t,s}(w)=V_t(H)^{\mathsf T}\operatorname{diag}\bigl(w(x)\bigr)V_s(H).
\]
If $X=\{x_0,\ldots,x_{m-1}\}$ is a set of distinct nonzero points and
\[
        A_h(X):=V_m(X)^{\mathsf T}\operatorname{diag}(x^h:x\in X)\,V_m(X),
\]
then
\[
        \det A_h(X)=\det\bigl(V_m(X)\bigr)^2\prod_{x\in X}x^h\ \ne0 .
\]
\end{proposition}

\begin{proof}
The $(a,b)$ entry of the triple product is $\sum_{x}x^aw(x)x^b$, the defining Hankel entry.  The determinant identity is multiplicativity of the determinant together with $\det(V_m(X)^{\mathsf T})=\det V_m(X)$, and the Vandermonde determinant on distinct points is nonzero.
\end{proof}

\begin{proposition}[Johnson exchange mixing]\label{prop:capf-johnson-exchange}
Let $1\le j\le n-1$, and let $J(n,j)$ be the Johnson graph on $j$-subsets of an $n$-set, with degree $d=j(n-j)$.  If $\mathcal A$ is a family of $j$-subsets of density $\delta$ and $T_1$ is obtained from a uniform $T_0$ by one uniform exchange step, then
\[
        \Pr[T_0\in\mathcal A,\ T_1\in\mathcal A]
        \le
        \delta^2+\Bigl(1-\frac{n}{d}\Bigr)\delta(1-\delta),
\]
and point dictators $\mathcal A=\{T:x_0\in T\}$ attain equality.
\end{proposition}

\begin{proof}
Let $P$ be the normalized adjacency operator of $J(n,j)$ and write $1_{\mathcal A}=\delta\mathbf 1+f$ with $f\perp\mathbf1$ and $\|f\|_2^2=\delta(1-\delta)$ under the uniform measure.  The eigenvalues of the Johnson graph are classically
\[
        \theta_i=(j-i)(n-j-i)-i,
        \qquad 0\le i\le \min(j,n-j),
\]
the association-scheme spectrum of $J(n,j)$.  For $i<\min(j,n-j)$,
\[
        \theta_i-\theta_{i+1}
        =(j-i)(n-j-i)-(j-i-1)(n-j-i-1)+1
        =n-2i>0,
\]
so the sequence is strictly decreasing and the largest nontrivial normalized eigenvalue is $\lambda_1=\theta_1/d=1-n/d$.  Hence
\[
        \Pr[T_0,T_1\in\mathcal A]
        =\langle1_{\mathcal A},P1_{\mathcal A}\rangle
        =\delta^2+\langle f,Pf\rangle
        \le\delta^2+\lambda_1\|f\|_2^2 ,
\]
which is the display; only the largest nontrivial eigenvalue is needed for this upper bound.  The centered indicator of a point dictator lies in the $i=1$ eigenspace, so equality is attained there.
\end{proof}

\begin{proposition}[anticode packing]\label{prop:capf-anticode}
Let $0\le s\le j\le n$.  If a family $\mathcal F$ of $j$-subsets of an $n$-set satisfies $|A\setminus B|>s$ for all distinct $A,B\in\mathcal F$, then
\[
        |\mathcal F|\binom js\le \binom n{j-s}.
\]
\end{proposition}

\begin{proof}
For each $A\in\mathcal F$, consider its $(j-s)$-subsets, $\binom j{j-s}=\binom js$ of them.  These collections are disjoint across $\mathcal F$: a common $(j-s)$-subset of $A\ne B$ would give $|A\cap B|\ge j-s$, hence $|A\setminus B|\le s$, contrary to hypothesis.  There are at most $\binom n{j-s}$ subsets of size $j-s$ in total.
\end{proof}

\subsection{SPI deficiency-one chart}

Let
\[
        M(Z)=M_0+ZM_1
\]
be a $t\times(j+1)$ matrix pencil with affine-linear entries over $F$, and assume $t=j$.  For $0\le i\le j$, let $M_i(Z)$ be the maximal minor obtained by deleting column $i$, put
\[
        c_i(Z)=(-1)^iM_i(Z),
        \qquad
        L_Z(X)=\sum_{i=0}^{j}c_i(Z)X^i,
\]
and recall that whenever $\operatorname{rank}M(z)=j$, the cofactor vector $(c_0(z),\ldots,c_j(z))$ is nonzero and spans $\ker M(z)$.  In the syndrome application, $M(Z)$ is the exact-agreement syndrome pencil of \cref{sec:aperiodic-hankel-certificates} in the deficiency-one shape $t=j$, the kernel vector is the coefficient vector of a candidate locator of degree $\le j$, and a slope $z$ \emph{passes the split test} if the normalized kernel locator divides $X^n-1$; as in \cref{lem:singular-pencil}, passing the split test is necessary for badness at the exact agreement but not sufficient.

\begin{theorem}[deficiency-one eliminant or residual]\label{thm:capf-spi}
Assume $H$ is the full root set of $X^n-1$ and $\operatorname{char}F\nmid n$.
\begin{enumerate}[label=\textup{(\arabic*)},start=0]
\item If every maximal minor of $M(Z)$ vanishes identically, the pencil is identically rank-deficient; this chart is a named residual branch, to be classified by the polynomial-kernel chain of \cref{lem:singular-pencil}, and the statements below do not apply to it.
\end{enumerate}
Otherwise fix a maximal minor that is not identically zero.  Every finite slope $z$ then lies in exactly one of the classes:
\begin{enumerate}[label=\textup{(\roman*)}]
\item rank drop: all maximal minors vanish at $z$; this set is contained in the root set of the chosen minor, a nonzero polynomial of degree at most $t$;
\item low-degree full-rank chart: $\operatorname{rank}M(z)=j$ and $c_j(z)=0$; the Cramer locator has degree $<j$ and is charged to a contained higher-agreement branch;
\item top chart: $c_j(z)\ne0$; the Cramer locator is proportional to a monic degree-$j$ locator, and $z$ passes the split test if and only if $L_z(X)\mid X^n-1$.
\end{enumerate}
If $c_j\equiv0$, the top chart is empty and every full-rank slope lies in class \textup{(ii)}.  Assume $c_j\not\equiv0$ and perform pseudo-division over $F[Z]$:
\[
        c_j(Z)^\Delta\,(X^n-1)=Q(Z,X)\,L_Z(X)+R(Z,X),
        \qquad
        \Delta=n-j+1,
\]
with $\deg_XR<j$ and $\deg_ZR_m\le(n-j+1)t$ for every $X$-coefficient $R_m(Z)$ of $R$.  Then exactly one of the following holds:
\begin{enumerate}[label=\textup{(\alph*)}]
\item some $R_m$ is not identically zero; then the polynomial $c_j(Z)R_m(Z)$, of degree at most
\[
        t+(n-j+1)t,
\]
is a finite eliminant whose root set contains every top-chart slope passing the split test together with every slope of classes \textup{(i)} and \textup{(ii)};
\item all $R_m$ vanish identically; then every top-chart slope passes the split test, and the chart must be retained as a named residual branch rather than discarded.
\end{enumerate}
\end{theorem}

\begin{proof}
Class \textup{(0)} and the trichotomy are immediate from the definitions: each maximal minor is the determinant of a $j\times j$ matrix with affine-linear entries in $Z$, hence a polynomial of degree at most $j=t$; if all vanish at $z$ then $\operatorname{rank}M(z)<j$, and otherwise the kernel is one-dimensional and spanned by the nonzero cofactor vector, whose top coordinate $c_j(z)$ decides between \textup{(ii)} and \textup{(iii)}.  In class \textup{(iii)}, $L_z$ has degree exactly $j$; since $X^n-1$ is squarefree with root set $H$, its monic normalization is $H$-split if and only if it divides $X^n-1$, which is the split test.

For the pseudo-division, dividing the degree-$n$ polynomial $X^n-1$ by the degree-$j$ polynomial $L_Z$ with leading coefficient $c_j(Z)$ performs $\Delta=n-j+1$ elimination steps, each multiplying the running remainder by $c_j(Z)$ once and subtracting a $c_i(Z)$-multiple of a shift of $L_Z$; since every $c_i$ has $Z$-degree at most $t$, each step raises the $Z$-degree of the remainder coefficients by at most $t$, giving $\deg_ZR_m\le\Delta t=(n-j+1)t$ from the degree-$0$ start.

In case \textup{(a)}, let $z$ be a top-chart slope passing the split test.  Substituting $Z=z$, the left-hand side $c_j(z)^\Delta(X^n-1)$ is divisible by $L_z$, and so is $Q(z,X)L_z(X)$; hence $L_z(X)\mid R(z,X)$, and $\deg_XR(z,X)<j=\deg L_z$ forces $R(z,X)=0$, so $R_m(z)=0$ for every $m$, in particular for the chosen nonzero one.  Slopes of classes \textup{(i)} and \textup{(ii)} satisfy $c_j(z)=0$ \textup{(at rank drop, all minors vanish, in particular $M_j$)}.  Hence all these slopes are roots of $c_jR_m$, whose degree is at most $t+(n-j+1)t$.  In case \textup{(b)}, substituting any top-chart slope gives $c_j(z)^\Delta(X^n-1)=Q(z,X)L_z(X)$ with $c_j(z)\ne0$, so $L_z$ divides $X^n-1$: the split test is identically passed on the top chart, which is then a genuine residual family.
\end{proof}

\begin{remark}[relation to the main eliminant machinery, and calibration]\label{rem:capf-spi-calibration}
\Cref{thm:capf-spi} complements the certificate framework of \cref{sec:aperiodic-hankel-certificates}: \cref{thm:aperiodic-affine-eliminant} and \cref{thm:kernel-compressed-affine-eliminant} bound chart contributions \emph{given} a nonzero eliminant, while \cref{thm:capf-spi} \emph{produces} one for every deficiency-one chart with explicit degree, or else names the residual branch \textup{(identically rank-deficient pencils, handled by \cref{lem:singular-pencil}, and identically split top charts)}.  As calibration: for $n=512$, $k=256$, $A=384$, one has $t=j=128$, and the global one-parameter cap of case \textup{(a)} is
\[
        t+(n-j+1)t=128+385\cdot128=49408,
\]
a bounded eliminant problem.  The theorem does not classify higher-deficiency SPI charts; those remain inside the broad band formulation \cref{prob:band}.
\end{remark}

\section{Identity-scale frontier and quotient base cases}\label{sec:capff1-frontier}

The quotient floors of the main text work at every divisor scale.  At the identity scale $c=1$ no quotient remains, and the pigeonhole is simply over the leading coefficients of ordinary split locators.  This terminal scale gives the base unsafe certificates; the MCA rows are then strengthened further by the flexible-budget conversion in \cref{prop:capg-moved-frontier}.

\subsection{The identity prefix floor}

\begin{lemma}[identity-scale prefix floor]\label{lem:capff1-identity-prefix-floor}
Let $\B\subseteq\F$, let $D\subseteq\B$ have $|D|=n$, and let $1\le K\le n$.  For every integer $m$ with $K\le m\le n$, put $w=m-K$.  Then there is a $\B$-valued received word $U:D\to\B$ such that
\[
        \bigl|\Lst(\RS[\F,D,K],1-m/n,U)\bigr|
        \ge
        \left\lceil {\binom n m\over |\B|^{\,w}}\right\rceil .
\]
Consequently the graded-prefix route of \cref{prop:graded-prefix-floor} extends to the identity scale $c=1$.
\end{lemma}

\begin{proof}
For each $m$-subset $M\subseteq D$, write
\[
        \Lambda_M(X)=\prod_{b\in M}(X-b)
        =X^m+\sum_{j=1}^{m}(-1)^j e_j(M)X^{m-j}.
\]
For $z=(z_1,\ldots,z_w)\in\B^w$, define
\[
        U_z:=\Bigl(X^m+\sum_{j=1}^{w}z_jX^{m-j}\Bigr)\Big|_D.
\]
The map $M\mapsto((-1)^je_j(M))_{1\le j\le w}$ has codomain $\B^w$, so one fiber contains at least $\binom nm/|\B|^w$ subsets.  Fix such a prefix value $z^*$, and for each $M$ in the fiber put
\[
        c_M:=U_{z^*}-\Lambda_M|_D
        =-\sum_{j=w+1}^{m}(-1)^je_j(M)X^{m-j}\Big|_D.
\]
Since $m-w-1=K-1$, the word $c_M$ is a codeword of $\RS[\F,D,K]$.  On the $m$ points of $M$, the locator vanishes, so $c_M$ agrees with $U_{z^*}$ on at least $m$ positions.

It remains only to check injectivity.  If $c_M=c_{M'}$, then the difference of the two displayed polynomials has degree at most $K-1<n$ and vanishes on all points of $D$, hence is the zero polynomial.  Therefore the coefficients $e_j(M)$ and $e_j(M')$ agree for all $j>w$; the prefix-fiber condition gives agreement for $j\le w$.  The two monic locator polynomials are equal, and hence $M=M'$.
\end{proof}

\subsection{Identity-prefix certificates}

Throughout this subsection $n=2^{21}$, $k=2^{20}$.  For the KoalaBear row, $p=2^{31}-2^{24}+1$ and $q=p^6$; for the Mersenne-31 line-round circle row, $p'=2^{31}-1$ and $q=(p')^4$.  The MCA rows use the floor at $K=k+1$ and then the deep-point conversion to the dimension-$k$ code, exactly as in \cref{cor:widened-deployed,rem:exact-frontier}; the list rows use $K=k$.


\begin{proposition}[identity-scale list frontier, exact]\label{prop:capff1-identity-frontier}
The exact identity-scale list certificates are:
\begin{center}\footnotesize
\begin{tabular}{@{}llllll@{}}
\toprule
row / challenge & $c$ & $m$ & $w$ & certified unsafe range & margins (bits)\\
\midrule
KoalaBear, list & $1$ & $1116046$ & $67470$ &
$\delta\in[\,\tfrac{490553}{1048576},1-\rho)=[0.4678278\ldots,1-\rho)$ & $+9.2\ /\ -22.0$\\
Mersenne-31 line round, list & $1$ & $1116022$ & $67446$ &
$\delta\in[\,\tfrac{490565}{1048576},1-\rho)=[0.4678392\ldots,1-\rho)$ & $+28.1\ /\ -3.1$\\
\bottomrule
\end{tabular}
\end{center}
Each printed list value of $m$ passes by exact integer comparison, and $m+1$
fails by exact integer comparison.  The MCA rows use the same identity-prefix
floor, but their final finite frontier is the stronger flexible-budget
certificate of \cref{prop:capg-moved-frontier}.
\end{proposition}

\begin{proof}
The construction is \cref{lem:capff1-identity-prefix-floor} with $K=k$.  The
displayed list floors exceed the corresponding list budgets by the exact
integer inequalities
\[
\binom n{1116046}\,2^{128}>p^{67470}p^6,
\qquad
\binom n{1116022}\,2^{100}>(p')^{67446}(p')^4.
\]
In both rows the inequality with $m$ replaced by $m+1$ and $w$ by $w+1$ is
false.  The edge fractions are $1-m/n$.
\end{proof}


\begin{corollary}[deployed sandwich edges]\label{cor:capff1-sandwich-edges}
The deployed unsafe edges supplied by the identity-prefix and flexible-budget certificates give \[ \delta^*_C(2^{-128})\le {981105\over2097152} \quad\text{for KoalaBear MCA}, \] while the list-row edges are the list entries of \cref{prop:capff1-identity-frontier} and the Mersenne-31 MCA edge is the MCA entry of \cref{prop:capg-moved-frontier}.  These are the upper endpoints used in the frontier sandwich.
\end{corollary}

\begin{proof}
Substitute \cref{prop:capff1-identity-frontier} into the unsafe side of \cref{thm:sandwich2}; the safe-side inputs are unchanged.
\end{proof}

\subsection{The frontier conjecture}

\begin{definition}[frontier function, unchanged]\label{def:capff1-gstar}
For $\rho\in(0,1)$ and $\beta>0$, put
\[
        g^*(\rho,\beta)=\sup\{g\in(0,1-\rho]:\Hb(\rho+g)\ge \beta g\}.
\]
For KoalaBear, $\beta=\log_2(2^{31}-2^{24}+1)=30.988685\ldots$ and
\[
1-\rho-g^*(\rho,\beta)=0.4678266,\ 0.7225056,\ 0.8557947,\ 0.9251114
\]
for $\rho=1/2,1/4,1/8,1/16$.  For Mersenne-31, the corresponding values are
\[
0.4678383,\ 0.7225161,\ 0.8558023,\ 0.9251165.
\]
At rate $1/2$, the identity-scale deployed gaps are within $2.7\cdot10^{-6}$ of the KoalaBear MCA ceiling, $1.2\cdot10^{-6}$ of the KoalaBear list ceiling, $1.5\cdot10^{-6}$ of the Mersenne-31 MCA ceiling, and $1.0\cdot10^{-6}$ of the Mersenne-31 list ceiling.  Larger-block extrapolations from floating-point scans should be recorded as numerical evidence unless separately certified by exact integer comparisons.
\end{definition}

\begin{conjecture}[frontier conjecture, asymptotic and finite forms]\label{prob:capff1-frontier}
Prove or refute the asymptotic safe-side match
\[
        \delta^*_C(\varepsilon^*)=1-\rho-g^*(\rho,\log_2|\B|)+o(1)
\]
for smooth multiplicative and circle line-round rows at the deployed challenge
targets.  The finite deployed form predicts that the first safe agreement is
one more than the best exact unsafe agreement supplied by
\cref{prop:capg-moved-frontier} and the list entries of
\cref{prop:capff1-identity-frontier}:
\[
\begin{array}{c|cc|c}
\text{row} & \text{proved unsafe agreement }a_0 &
\text{first safe }a_0+1 &
\text{margin (bits)}\\ \hline
\text{KoalaBear MCA} & 1116047 & 1116048 & 22.2\\
\text{KoalaBear list} & 1116046 & 1116047 & 22.0\\
\text{Mersenne-31 MCA} & 1116023 & 1116024 & 3.3\\
\text{Mersenne-31 list} & 1116022 & 1116023 & 3.1
\end{array}
\]
with the real closed-ball endpoint convention of \cref{cor:capf-endpoint}.
\end{conjecture}

\begin{remark}[closure mechanisms after $c=1$]\label{rem:capff1-closure}
The two ``free'' mechanisms in the terminal-scale scan should be kept, but with the endpoint changed.  Subgroup-invariant subset selection at a nontrivial scale is absorbed into the composite scale exactly as before; iterating this can only move toward the identity prefix floor, where no nontrivial quotient remains.  The $c=2$, $\sigma=1$ planted hybrid is also dominated by the identity row at the deployed parameters: its best net agreement is at most the $c=2$ agreement minus one, whereas \cref{prop:capff1-identity-frontier} improves the $c=2$ agreement by $5,2,3,2$ steps in the four rows.  Thus planting remains useful as a structural lower construction, but it is not the deployed frontier certificate once identity prefixes are allowed.
\end{remark}

\subsection{Proved base cases of the identity-scale collision problem}

At the identity scale the census behind \cref{lem:capff1-identity-prefix-floor} runs over the deployed domain itself, and for the KoalaBear rows that domain is a coset of the order-$2^{21}$ $2$-adic subgroup of $\F_p^\times$ with $2^{21}>\ceil{\sqrt p}=46160$.  The prefix-fiber counts are therefore in Weil range, and their base cases are provable unconditionally.  Throughout, $e_p(x):=\exp(2\pi ix/p)$ and $B:=\ceil{\sqrt p}$.

\begin{theorem}[bounded-prefix equidistribution over large cosets, with planted cores]\label{thm:capff1-collision-base}
Let $p$ be prime, $H\le\F_p^\times$ of order $N$, $d=(p-1)/N$, and $S=\alpha H$ a coset.  Let $P\subseteq S$ be a fixed set of size $\ell\ge0$, let $1\le w<N-\ell$ satisfy $wd<p$, and let $\ell\le m\le N$, $m':=m-\ell$.  For $\vec z\in\F_p^w$ put
\[
        \mathcal N_w^P(m,\vec z)\ :=\ \#\bigl\{A:\ P\subseteq A\subseteq S,\ |A|=m,\ \bigl(e_1(A),\ldots,e_w(A)\bigr)=\vec z\,\bigr\}.
\]
Then for every $\vec z$,
\[
        \Bigl|\,\mathcal N_w^P(m,\vec z)-\frac1{p^w}\binom{N-\ell}{m'}\,\Bigr|
        \ \le\ p^{w/2}\binom{wB+\ell+m'-1}{m'},
\]
and for $w=1$ the factor $p^{1/2}$ may be dropped.
\end{theorem}

\begin{proof}
\emph{Step 1 (removing the core).}  Write $A=P\cup A'$ with $A'\subseteq S':=S\setminus P$, $|A'|=m'$.  From $\Lambda_A=\Lambda_P\Lambda_{A'}$,
\[
        e_j(A)=e_j(A')+\sum_{i=1}^{\min(j,\ell)}e_i(P)\,e_{j-i}(A')\qquad(1\le j\le w),
\]
so the map $(e_1,\ldots,e_w)(A')\mapsto(e_1,\ldots,e_w)(A)$ is affine and unipotent triangular over $\F_p^w$, hence a bijection.  Each fiber of the planted census is therefore a fiber of the plain census on $S'$ at a transformed value, and it suffices to prove the displayed bound for $\mathcal N_w(m',\cdot)$ on $S'$.

\emph{Step 2 (Fourier reduction).}  By inversion on $\F_p^w$, $\mathcal N_w(m',\vec z)=p^{-w}\sum_{\vec t}e_p(-\vec t\cdot\vec z)E(\vec t)$ with $E(\vec t)=\sum_{|A'|=m'}e_p(\sum_jt_je_j(A'))$; the term $\vec t=0$ gives the main term, so it suffices to bound $|E(\vec t)|$ for $\vec t\ne0$.

\emph{Step 3 (Newton conversion).}  Fix $\vec t\ne0$ and $j^*=\max\{j:t_j\ne0\}$.  Since $j\le w<p$, the Newton identities express each $e_j$ as a polynomial in the power sums with $e_j=(-1)^{j-1}p_j/j+R_j(p_1,\ldots,p_{j-1})$, so $\sum_jt_je_j(A')=Q(P(A'))$ for $P(A')=(p_1,\ldots,p_{j^*})(A')$ and a polynomial $Q$ that is linear in its last variable with constant coefficient $(-1)^{j^*-1}t_{j^*}/j^*\ne0$.  Fourier inversion on $\F_p^{j^*}$ gives
\[
        E(\vec t)=p^{-j^*}\sum_{\vec s}\widehat Q(\vec s)\,G(\vec s),\qquad
        \widehat Q(\vec s)=\sum_{\vec v}e_p(Q(\vec v)-\vec s\cdot\vec v),\quad
        G(\vec s)=\sum_{|A'|=m'}e_p(\vec s\cdot P(A')),
\]
and $\widehat Q(0)=0$ because the inner sum over the last coordinate of $\vec v$ is a complete nontrivial character sum.

\emph{Step 4 (power sums on the punctured coset).}  For $\vec s\ne0$, $G(\vec s)=e_{m'}\bigl((e_p(g_s(a)))_{a\in S'}\bigr)$ with $g_s(x)=\sum_{i\le j^*}s_ix^i$ of degree $i^*\le w$; its power sums are $\pi_j=\sum_{a\in S'}e_p(jg_s(a))$, $1\le j\le m'$.  Since $\deg g_s\le w<N-\ell=|S'|$ and $g_s$ is nonconstant, and $j\not\equiv0\pmod p$, the coset sum obeys Weil's bound through the power parametrization $S=\{\alpha y^d:y\in\F_p^\times\}$:
\[
        \Bigl|\sum_{a\in S}e_p(jg_s(a))\Bigr|
        =\frac1d\Bigl|\sum_{y\in\F_p}e_p\bigl(jg_s(\alpha y^d)\bigr)-e_p(jg_s(0^{\ast}))\Bigr|
        \le\frac1d\bigl[(i^*d-1)\sqrt p+1\bigr]\le w\sqrt p\le wB,
\]
where $0^\ast$ abbreviates the $y=0$ evaluation point and $\deg_y=i^*d\le wd<p$ ensures applicability.  Removing the $\ell$ points of $P$ costs at most $\ell$, so $|\pi_j|\le wB+\ell$ for all $j$.

\emph{Step 5 (Newton domination).}  The recursion $m'e_{m'}=\sum_{i=1}^{m'}(-1)^{i-1}e_{m'-i}\pi_i$ is majorized coefficientwise by the recursion of $[u^r](1-u)^{-(wB+\ell)}$, whence $|G(\vec s)|\le\binom{wB+\ell+m'-1}{m'}$ for every $\vec s\ne0$.

\emph{Step 6 (Parseval).}  $\sum_{\vec s}|\widehat Q(\vec s)|\le p^{j^*/2}\bigl(\sum_{\vec s}|\widehat Q(\vec s)|^2\bigr)^{1/2}=p^{3j^*/2}$, so $|E(\vec t)|\le p^{j^*/2}\binom{wB+\ell+m'-1}{m'}$, and summing the $p^w-1$ nonzero characters against $p^{-w}$ gives the theorem.  For $w=1$ no linearization is needed: $E(t)=e_{m'}\bigl((e_p(ta))_{a\in S'}\bigr)$ directly, with $|\pi_j|\le\sqrt p+\ell$.
\end{proof}

The margins below are certified by the standard entropy sandwich, whose two halves have one-line proofs: for $0<m<M$, taking $\lambda=m/M$ in $1=(\lambda+(1-\lambda))^M\ge\binom Mm\lambda^m(1-\lambda)^{M-m}$ gives $\binom Mm\le2^{M\Hb(m/M)}$, and since $\binom Mm\lambda^m(1-\lambda)^{M-m}$ is the largest of the $M+1$ binomial probabilities, $\binom Mm\ge2^{M\Hb(m/M)}/(M+1)$.  Every margin printed in \cref{cor:capff1-collision-head} is an inequality between two explicit numbers of the form $M\Hb(m/M)\pm\log_2(M+1)\pm\tfrac{3w}2\log_2p$, verifiable by interval arithmetic at $10^{-6}$ precision; the accompanying script is a cross-check, not the proof.

\begin{corollary}[the head of the identity staircase is exact]\label{cor:capff1-collision-head}
For the KoalaBear rows of \cref{prop:capff1-identity-frontier} \textup{($N=2^{21}$, $d=1016$, $B=46160$, $\ell=0$)}:
\begin{enumerate}[label=\textup{(\alph*)}]
\item at the deployed subset sizes $m=1116047$ \textup{(MCA)} and $m=1116046$ \textup{(list)}, \cref{thm:capff1-collision-base} is nonvacuous for exactly $w\le21$, with margins
\[
+1810770,\ +1620296,\ +1462138,\ \ldots,\ +63888,\ +11956
\]
bits at $w=1,2,3,\ldots,20,21$ for the MCA row \textup{(the list row agrees within $4$ bits)}, and failure at $w=22$ by $38279$ bits;
\item at the head rungs $m=K+w$ it is nonvacuous for exactly $w\le22$ \textup{(margin $+32394$ bits at $w=22$, failure by $14138$ at $w=23$)};
\item consequently every prefix fiber in these ranges equals $\binom Nm p^{-w}(1+\varepsilon)$ with $|\varepsilon|\le2^{-\mathrm{margin}}$: the pigeonhole step of \cref{lem:capff1-identity-prefix-floor} is \emph{exact} at head depths, two-sidedly and with no polynomial factor, giving the first exactly priced strata of the identity staircase;
\item the planted census is priced identically with $wB$ replaced by $wB+\ell$, which costs about $\log_2(1+m/B)\approx4.65$ bits of margin per planted point: cores of size $\ell\le2^{11}$ preserve every rung $w\le21$, and the $w=1$ rung tolerates $\ell$ up to $\approx2^{18}$, so the planted lower constructions of \cref{rem:capff1-closure} are two-sided at head depths.
\end{enumerate}
Not claimed: any depth beyond the printed $w_0$ at these subset sizes; and the $U$-twisted per-word census at general received words \textup{(the fibers priced here are the ones entering the floor and the witness words it constructs; \cref{prop:capff1-witness-exact} makes that pricing exact)}.  The circle rows are treated in \cref{thm:capff1-collision-circle} below.
\end{corollary}

\begin{proposition}[composite directions collapse to coarser scales]\label{prop:capff1-composite}
Let $e\mid N$, $e>1$, and let $\vec s\ne0$ be supported on indices divisible by $e$, so that $g_s(x)=h(x^e)$ with $\deg h\le w/e$.  Then, with $S_e:=\{a^e:a\in S\}$ a coset of the order-$N/e$ subgroup,
\[
        \sum_{\substack{A\subseteq S\\|A|=m}}e_p\Bigl(\sum_{a\in A}g_s(a)\Bigr)
        \ =\ [x^m]\prod_{b\in S_e}\bigl(1+x\,e_p(h(b))\bigr)^{e}.
\]
\end{proposition}

\begin{proof}
The generating function $\sum_m x^m\sum_{|A|=m}e_p(\sum_{a\in A}g_s(a))$ equals $\prod_{a\in S}(1+xe_p(g_s(a)))$.  The value $g_s(a)=h(a^e)$ is constant on each fiber of $a\mapsto a^e$, and each $b\in S_e$ has exactly $e$ preimages in $S$, so the product collapses as displayed.
\end{proof}

Two consequences.  First, \cref{thm:capff1-collision-base} never assumed primitivity: for composite $\vec s$ the reduced-scale power sums still obey $e\cdot|\pi^{\mathrm{red}}_j|\le e\cdot(w/e)\sqrt p=w\sqrt p$ by the same Weil argument at the coarser coset, so the head-depth pricing above covers all directions.  Second, beyond the head, composite directions genuinely carry the character mass of their coarser scale \textup{(quantified in \cref{rem:capff1-calibration})}, so any strengthening uniform in the direction must be graded along the divisor lattice --- the same absorption structure as \cref{rem:capff1-closure}.

\begin{remark}[the corrected collision target]\label{rem:capff1-collision-gap}
The proof of \cref{thm:capff1-collision-base} pays one Weil cost $\sqrt p$ per power sum, and this is exactly what stops it at $w_0=21$--$22$ while the frontier rows of \cref{prop:capff1-identity-frontier} sit at depth $w\approx67466$.  Computation rules out the tempting strengthening: a uniform character bound $|E(\vec t)|\le\binom Nm p^{-\kappa w}$ over all $\vec t\ne0$ is false, in two distinct ways.  Composite directions carry coarse-scale mass by \cref{prop:capff1-composite}, so no single-scale exponent can hold uniformly.  And even over primitive directions, maxima are pinned near the Parseval level: $\sum_{\vec t}|E(\vec t)|^2$ equals $p^w$ times the collision count of the prefix map \emph{exactly}, so square-root cancellation per slot is the typical size, while the passage from character maxima to fiber bounds sums $p^w$ signed terms and would need $\kappa\ge1$.  Character maxima therefore face a square-root barrier at depth, and the correct remaining target is the max-fiber inequality itself, at every primitive scale in the dense regime $\binom Nmp^{-w}\ge1$:
\[
        \max_{\vec z}\ \mathcal N_w(m,\vec z)\ \le\ \operatorname{poly}(n)\cdot\binom Nm p^{-w},
\]
with composite scales priced through their reductions.  Its consequences must be read against \cref{prop:capff1-closing}: a polynomial-loss max-fiber input, with the normalized aperiodic-band input \cref{prob:capfr1-normalized-band}, closes the asymptotic frontier away from an $O(\log n)$ agreement reserve, and does \emph{not} decide the printed adjacent pairs of \cref{prob:capff1-frontier}, whose finite margins are $22.2$, $22.0$, $3.3$, $3.1$ bits.  Deciding a printed pair requires the constant-factor form --- max within $(1+o(1))$ of the mean, with all constants displayed --- which is precisely the regime the exact censuses of \cref{rem:capff1-calibration} probe, and which \cref{cor:capff1-collision-head} proves outright at head depths.
\end{remark}

\begin{remark}[computational calibration; experimental]\label{rem:capff1-calibration}
The measurements behind \cref{rem:capff1-collision-gap} use two exact instruments.  \textup{(i)}  Because the Newton map $(e_1,\ldots,e_w)\leftrightarrow(p_1,\ldots,p_w)$ is a bijection on value vectors for $w<p$, the power-sum prefix map has identical fibers and elementwise-additive phases, so ratios $E(\vec s)/\binom Nm$ are computable at every depth by the normalized recursion $q_j^{(i)}=(1-\tfrac ji)q_j^{(i-1)}+\tfrac ji\varepsilon_iq_{j-1}^{(i-1)}$: a convex combination of unimodularly weighted values, hence $|q|\le1$ throughout and the update is non-expansive, so floating error cannot amplify; spot values were reproduced to all printed digits by exact arithmetic in $\mathbf Z[x]/(x^p-1)$.  \textup{(ii)}  Group-algebra dynamic programming gives exact joint fiber censuses at small scales.  Findings, on subgroups with $N/\sqrt p\approx3$--$4$ \textup{(the deployed identity scale has $N/\sqrt p\approx45$, the favourable direction)}: primitive character maxima decay toward the Parseval level \textup{(measured exponents per slot $2.26,1.13,0.80,0.67,0.53,0.46,0.43,0.39,0.31$ at depths $2$--$10$ for $p=257$, $N=64$, $m=34$)}, while directions forced onto index multiples of $e=2,4,8$ jump to $4.2\times10^{-6}$, $1.7\times10^{-4}$, $1.6\times10^{-2}$ at depth $8$ --- the divisor hierarchy of \cref{prop:capff1-composite} in plain sight.  The exact censuses show the max-fiber target holding with room where the ledger operates: max-to-mean ratios $1.0012,\ 1.21,\ 2.67$ at depths $1$--$3$ for $(p,N,m)=(17,16,8)$; $1.0022,\ 1.21,\ 4.10$ for $(41,20,10)$; $1+3\times10^{-12}$, $1+9\times10^{-9}$ at depths $1$--$2$ for $(257,64,34)$, minima equally pinned.  Ratios grow only at the Poisson boundary, where the mean fiber falls to $O(1)$ and the maximum of $p^w$ near-independent cells is logarithmic \textup{(maximum $11$ against mean $2.7$ over $41^3$ cells, matching the Poisson tail)}; the deployed crossings sit in the ultra-dense bulk at mean $\approx q/k$.  These tables are calibration evidence only and enter no proof; the falsifiable prediction is that any violation of the max-fiber target must place a super-polynomial fiber inside the dense bulk of a primitive scale.
\end{remark}

\subsection{Circle rows, exact witness lists, and quotient scales}

Three additions complete the base-case picture: the identity witness lists are \emph{exactly} the prefix fibers, so the census statements above are list-size statements about the code; the circle rows acquire the same proved base cases through a Weil bound on the conic, with doubled constants; and the quotient census cells inherit a per-scale table.

\begin{proposition}[witness lists are exactly the prefix fibers]\label{prop:capff1-witness-exact}
In the setting of \cref{lem:capff1-identity-prefix-floor}, for \emph{every} $z\in\B^w$,
\[
        \Lst\bigl(\RS[\F,D,K],\,1-m/n,\,U_z\bigr)\ =\ \{\,c_M:\ M\ \text{in the prefix fiber of }z\,\},
\]
so the list size equals the fiber size exactly; moreover no codeword agrees with $U_z$ on more than $m$ points, so $U_z$ is list-empty at every strictly smaller radius.
\end{proposition}

\begin{proof}
The inclusion $\supseteq$ and the distinctness of the $c_M$ are \cref{lem:capff1-identity-prefix-floor}.  Conversely, let $c\in\RS[\F,D,K]$ agree with $U_z$ on at least $m$ points, and let $\hat c\in\F[X]$, $\deg\hat c<K$, be its polynomial.  Then
\[
        P(X):=X^m+\sum_{j=1}^{w}z_jX^{m-j}-\hat c(X)
\]
is monic of degree exactly $m$ \textup{(as $\deg\hat c\le K-1=m-w-1<m-w$)} and vanishes at every agreement point, so it has at least $m$ roots in $D$; having degree $m$, it equals $\Lambda_{M'}$ for the set $M'\subseteq D$ of its $m$ distinct roots.  Comparing coefficients in degrees $m-1,\ldots,m-w$ places $M'$ in the prefix fiber of $z$, and comparing the remaining coefficients gives $\hat c=-\sum_{j>w}(-1)^je_j(M')X^{m-j}$, i.e.\ $c=c_{M'}$.  Since $P\ne0$ has at most $m$ roots, agreement above $m$ is impossible.
\end{proof}

Combined with \cref{thm:capff1-collision-base,cor:capff1-collision-head} \textup{(the sign twist $z_j\mapsto(-1)^jz_j$ is a coordinatewise bijection of the value space, so fiber counts transfer verbatim)}, every identity witness word at KoalaBear head depths $w\le w_0$ has list size $\binom nm p^{-w}(1+\varepsilon)$ --- all $p^w$ of them, not only the pigeonhole maximizer.

\begin{lemma}[power-sum bound on twin-coset $x$-domains]\label{lem:capff1-torus-weil}
Let $p\equiv3\pmod4$, let $\UU$ be the circle group of order $p+1$, let $\mathcal D=gH\cup g^{-1}H$ be a twin coset \textup{($H\le\UU$ of order $M$, $g^2\notin H$)}, and let $D=\chi(\mathcal D)\subseteq\F_p$, $|D|=M$, $d'=(p+1)/M$.  Then for every nonconstant $f\in\F_p[X]$ with $\deg f\cdot d'<p$,
\[
        \Bigl|\sum_{x\in D}e_p(f(x))\Bigr|\ \le\ 2\deg f\cdot\sqrt p.
\]
\end{lemma}

\begin{proof}
Since $\mathcal D$ contains no self-inverse element, $\chi$ is exactly two-to-one on it, so $\sum_{x\in D}e_p(f(x))=\tfrac12\sum_{u\in\mathcal D}e_p(f(\chi(u)))$, and it suffices to bound each coset sum by $2\deg f\cdot\sqrt p$.  Fix the coset $gH$; as $\UU$ is cyclic and $d'\mid p+1$, the map $v\mapsto gv^{d'}$ carries $\UU$ onto $gH$ with every fiber of size $d'$, so
\[
        \sum_{u\in gH}e_p(f(\chi(u)))=\frac1{d'}\sum_{v\in\UU}e_p\bigl(R(v)\bigr),\qquad R(v):=f\bigl(\chi(gv^{d'})\bigr).
\]
Writing $x_g:=\chi(g)$, $y_g:=(g-g^{-1})/(2i)\in\F_p$ and using $u^{a}\pm u^{-a}$ Chebyshev identities, the angle-addition formula gives
\[
        \chi(gv^{d'})=x_g\,T_{d'}(x_v)-y_g\,y_v\,W_{d'}(x_v),
\]
where $T_{d'},W_{d'}\in\F_p[X]$ have degrees $d'$ and $d'-1$ \textup{(with $v^{a}-v^{-a}=(v-v^{-1})W_a(x_v)$)}.  Hence $R$ is an $\F_p$-polynomial function on the affine conic $x^2+y^2=1$, whose $\F_p$-points are exactly the $p+1$ elements of $\UU$: the conic has no rational points at infinity, since $x^2+y^2=0$ has no nonzero solution over $\F_p$.  Over $\overline{\F_p}$ the conic is a projective line with coordinate $u$, and near its two \textup{(conjugate, non-rational)} points at infinity $u\in\{0,\infty\}$ the expansion $\chi(gv^{d'})=\tfrac g2v^{d'}(1+O(v^{-2d'}))$ shows $R$ has exactly two poles, each of order $\deg f\cdot d'$ with nonvanishing leading coefficient.  $R$ is nonconstant: the $d'$-power map is onto $\UU$, and $\chi(g\,\UU)=\chi(\UU)$ has $(p+1)/2+1>\deg f$ elements, so a constant $R$ would force $f$ constant.  Its pole orders lie in $[1,p)$, so $R$ is not of the form $h^p-h+c$, whose pole orders are multiples of $p$.  Weil's theorem for additive character sums on complete nonsingular curves, in Bombieri's form, now gives
\[
        \Bigl|\sum_{v\in\UU}e_p(R(v))\Bigr|\le\bigl(2\cdot0-2+2+2\deg f\cdot d'\bigr)\sqrt p=2\deg f\cdot d'\sqrt p,
\]
and dividing by $d'$ bounds the coset sum.  The coset $g^{-1}H$ is identical.
\end{proof}

\begin{theorem}[circle version, with planted cores]\label{thm:capff1-collision-circle}
Let $D=\chi(\mathcal D)$ be a twin-coset $x$-domain as in \cref{lem:capff1-torus-weil}, $P\subseteq D$ with $|P|=\ell\ge0$, $1\le w<M-\ell$ with $wd'<p$, $\ell\le m\le M$, $m'=m-\ell<p$, and $B=\ceil{\sqrt p}$.  Then for every $\vec z\in\F_p^w$,
\[
        \Bigl|\,\mathcal N_w^P(m,\vec z)-\frac1{p^w}\binom{M-\ell}{m'}\,\Bigr|
        \ \le\ p^{w/2}\binom{2wB+\ell+m'-1}{m'},
\]
and for $w=1$ the factor $p^{1/2}$ may be dropped.
\end{theorem}

\begin{proof}
Steps 1--3, 5, and 6 of the proof of \cref{thm:capff1-collision-base} are domain-agnostic and apply verbatim.  In Step 4, the power sums are $\pi_j=\sum_{x\in D\setminus P}e_p(jg_s(x))$ with $jg_s$ nonconstant of degree $i^*\le w$ and $j\not\equiv0\pmod p$; \cref{lem:capff1-torus-weil} bounds the sum over $D$ by $2i^*\sqrt p\le2wB$, and removing the $\ell$ points of $P$ costs at most $\ell$.  The domination step then runs with $2wB+\ell$ in place of $wB+\ell$.
\end{proof}

\begin{corollary}[circle head rungs; all four rows priced]\label{cor:capff1-circle-head}
For the Mersenne-31 line-round rows of \cref{prop:capff1-identity-frontier} \textup{($M=2^{21}$, $d'=2^{10}$, $B=46341$, $\ell=0$)}:
\begin{enumerate}[label=\textup{(\alph*)}]
\item at the deployed sizes $m=1116023$ \textup{(MCA)} and $m=1116022$ \textup{(list)}, \cref{thm:capff1-collision-circle} is nonvacuous for exactly $w\le10$, with margins $+1619022$, $+1322222$, $\ldots$, $+169807$, $+60243$ bits at $w=1,\ldots,10$ and failure at $w=11$ by $41996$ bits;
\item at the head rungs $m=K+w$ it is nonvacuous for exactly $w\le11$ \textup{(margin $+28847$ at $w=11$)};
\item by \cref{prop:capff1-witness-exact}, on all four deployed rows the identity witness lists at head depths are exactly $\binom Mm p^{-w}(1+\varepsilon)$, $|\varepsilon|\le2^{-\mathrm{margin}}$; the doubled Weil constant halves the proved depth relative to the KoalaBear rows but leaves the $w=1$ margin above $1.6\times10^6$ bits;
\item planted cores cost about $\log_2(1+m/(2B))\approx3.7$ bits of margin per point on these rows.
\end{enumerate}
All margins are certified by the entropy sandwich exactly as before.
\end{corollary}

\begin{corollary}[per-scale table for the bounded quotient census]\label{cor:capff1-quotient-scales}
At the KoalaBear quotient scales, with $N_c=2^{21}/c$ and the per-scale certificate sizes $m_c$, \cref{thm:capff1-collision-base} prices the scale-$c$ prefix census with proved depths
\[
\begin{array}{c|ccccc}
c & 2 & 4 & 8 & 16 & 32\\ \hline
m_c & 558019 & 279007 & 139501 & 69748 & 34871\\
w_0(c) & 10 & 5 & 2 & 1 & \text{---}\\
w=1\text{ margin (bits)} & +810165 & +331058 & +111115 & +18230 & -14582
\end{array}
\]
so the census cells of the budget windows at scales $c\le16$ carry exactly priced heads, while the theorem is vacuous from $c=32$ onward \textup{($N_c/\sqrt p\approx1.42$: inside Weil range but outside the entropy margin)}.  For the circle rows, dyadic quotients of twin cosets are again twin cosets by \cref{lem:torus-fibers}, so \cref{thm:capff1-collision-circle} applies at every quotient scale with $(M,d')$ replaced by $(M/c,\,2^{10}c)$ and correspondingly halved reach.
\end{corollary}

\subsection{A repaired conditional closing statement}

\begin{proposition}[what the two named inputs actually close]\label{prop:capff1-closing}
Let $a_0$ be one of the four identity-scale unsafe agreements in \cref{prob:capff1-frontier}.  An exact finite one-step theorem at the deployed size follows from the exact upper certificate
\[
        U(a_0+1)\le B_*<L(a_0),
\]
where $L$ is the lower staircase supplied by \cref{prop:capff1-identity-frontier}, $U$ is the complete upper ledger for the same row, and $B_*$ is the challenge budget.  Then \cref{prop:onestep}, equivalently \cref{thm:capf-staircase,thm:capf-corridor}, gives $a_*=a_0+1$.

By contrast, quotient-fiber equidistribution only up to a factor $\operatorname{poly}(n)$, together with the normalized aperiodic-band input \cref{prob:capfr1-normalized-band}, is enough for the asymptotic frontier away from an $O(\log n)$ agreement reserve, but it does not by itself prove the finite adjacent pair.  At the deployed size the adjacent margins are only $22.2$, $22.0$, $3.3$, and $3.1$ bits in the four deployed rows; an unspecified polynomial factor can exceed those margins.
\end{proposition}

\begin{proof}
The first paragraph is exactly the one-step staircase compiler.  For the second paragraph, a factor $n^C$ costs $C\log_2 n=21C$ bits at $n=2^{21}$, while the printed adjacent margins are constants.  Therefore a polynomial-loss quotient input cannot be read as a finite adjacent upper certificate unless its degree and all constants are included and fit inside the printed margin.  Asymptotically, however, if one asks for agreement at least $a_0+R$ with $R\gg\log n$, the entropy--subfield inequality is separated from the frontier by $\Omega(R)$ bits along the same prefix ledger, and a fixed polynomial factor is absorbed.  This proves the stated distinction between finite one-step closing and asymptotic closing.
\end{proof}


\providecommand{\Avg}{\operatorname{Avg}}
\providecommand{\Supp}{\operatorname{Supp}}
\providecommand{\poly}{\operatorname{poly}}
\providecommand{\Gap}{\operatorname{Gap}}

\section{Frontier prize and quotient flatness}\label{sec:capfpr-prize}

The identity-scale floor of \cref{lem:capff1-identity-prefix-floor,prop:capff1-identity-frontier} changes the shape of the remaining problem.  The unsafe side is no longer the main source of leverage: the best printed lower certificate is already the entropy prefix floor, and the gap between the deployed rows and the asymptotic ceiling of \cref{def:capff1-gstar} is at the scale of a few parts in $10^{-6}$.  The prize-winning statement is therefore a safe-side theorem.  It must prove that no upper mechanism beats the identity-prefix average except through structures already paid by the capf ledger.

\begin{remark}[First-form frontier formulation]\label{prob:capfpr-frontier-prize}
Prove the asymptotic identity-prefix frontier
\[
        \delta_C^*(\varepsilon^*)
        =1-\rho-g^*(\rho,\log_2|\B|)+o(1)
\]
for the smooth multiplicative and circle line-round rows at the deployed challenge budgets.  In finite deployed form, prove the one-step safe side at the four adjacent candidates of \cref{prob:capff1-frontier}: namely, prove safety at agreements
\[
1116048,\quad 1116047,\quad 1116024,\quad 1116023
\]
for the KoalaBear MCA, KoalaBear list, Mersenne-31 MCA, and Mersenne-31 list rows respectively, while the preceding agreements remain unsafe by \cref{prop:capff1-identity-frontier}.
\end{remark}

\begin{remark}[separation between asymptotic and finite prizes]\label{rem:capfpr-finite-vs-asymp}
The asymptotic prize and the finite adjacent prize should be kept separate.  A factor $n^{C}$ in either missing upper input costs $C\log_2 n$ bits; at $n=2^{21}$ this is $21C$ bits, larger than several adjacent fail margins in \cref{prop:capff1-identity-frontier,prop:capff1-closing}.  Thus polynomial-loss inputs can close the frontier away from an $O(\log n)$ reserve, but a printed adjacent pair needs exact constants and integer rounding.
\end{remark}

\subsection{First split of the safe side}\label{subsec:capfpr-AQ}

The first useful decomposition separates prefix/quotient fibers from the aperiodic residue.  This split is falsifiable and is used by the subsequent reductions, but it is not the final active interface; after slope elimination and the split-pencil normalization it is read through \cref{subsec:capg-active-interface}.

\begin{remark}[First-form quotient/prefix formulation]\label{prob:capfpr-Q}
Let $D\subseteq\B\subseteq\F$, $|D|=n$, $K=\rho n$, $m=K+w$, and define the identity prefix map
\[
        \Phi_{m,w}:\binom Dm\longrightarrow\B^w,
        \qquad
        M\longmapsto\bigl((-1)^h e_h(M)\bigr)_{1\le h\le w}.
\]
After quotient-periodic components and declared paid scales have been separated, prove the primitive max-fiber estimate
\[
        \max_{z\in\B^w}|\Phi_{m,w}^{-1}(z)|
        \le \exp(o(n))\,{\binom n m\over |\B|^w}
\]
in the dense frontier window.  The finite deployed version must replace $\exp(o(n))$ by an explicit multiplier and additive error whose exact bit cost fits into the one-step upper ledger at $a_0+1$.
\end{remark}

\begin{remark}[First-form aperiodic formulation]\label{prob:capfpr-A}
After removing the tangent cell, quotient cells, extension-pole cells, common-GCD and quotient-pullback strata, planted/sunflower paid layers, fixed-dimensional Conjecture-F strata, projective-plane strata, and bounded SPI charts, every remaining primitive aperiodic M1/L1 configuration in the band of the broad band formulation \cref{prob:band} has at most model size up to a subexponential factor:
\[
        N_{\rm ap}(A)\le \exp(o(n))\,N_{\rm model}(A).
\]
For an exact deployed adjacent theorem, replace the right-hand side by an explicit integer $U_A(a_0+1)$ and part it into the staircase inequality of \cref{prop:capff1-closing}.
\end{remark}

\begin{proposition}[first-form closing template]\label{prop:capfpr-AQ-closing}
The asymptotic forms of \textup{(Q)} and \textup{(A)}, together with the paid cells already proved in \cref{sec:capf-threshold,sec:capf-l1,sec:capf-m1}, imply the asymptotic statement of the first-form frontier statement \cref{prob:capfpr-frontier-prize} away from an $O(\log n)$ agreement reserve.  The finite deployed adjacent theorem follows only from finite forms \textup{(Q)}$_{\rm fin}$ and \textup{(A)}$_{\rm fin}$ satisfying
\[
        U_{\rm paid}(a_0+1)+U_Q(a_0+1)+U_A(a_0+1)\le B_*<L(a_0),
\]
where $L(a_0)$ is the identity-prefix lower numerator at the preceding unsafe agreement.
\end{proposition}

\begin{proof}
The finite implication is exactly \cref{prop:capff1-closing}.  For the asymptotic implication, \textup{(Q)} bounds the heaviest primitive prefix or quotient fiber by the entropy average up to $\exp(o(n))$, while \textup{(A)} bounds all non-prefix aperiodic residuals by their model entropy up to $\exp(o(n))$.  The paid cells have already been separated by the capf compilers.  If the agreement is at least $O(\log n)$ beyond the identity-frontier crossing, the entropy deficit against the budget is $\Omega(\log n)$, so every fixed polynomial loss is absorbed.  This gives the stated asymptotic reserve, but not an exact finite adjacent pair.
\end{proof}

\subsection{Collision-gap formulation of \textup{(Q)}}\label{subsec:capfpr-collision-gap}

Let $\mu_{m,w}$ be the probability distribution of $\Phi_{m,w}(M)$ for a uniformly random $m$-subset $M\subseteq D$.  For $r\ge2$ define the normalized collision gaps
\[
        \Gamma_r(m,w)
        := |\B|^{w(r-1)}\sum_{z\in\B^w}\mu_{m,w}(z)^r .
\]
The uniform distribution has $\Gamma_r=1$.  The max-fiber assertion in \textup{(Q)} is an $L^\infty$ statement, while the numbers $\Gamma_r$ are high-moment probes of the same obstruction.

\begin{remark}[Fixed-moment collision hierarchy]\label{prob:capfpr-collision-gap}
In the primitive dense frontier window, prove a quantitative hierarchy strong enough to imply the first-form quotient input \cref{prob:capfpr-Q}.  A useful asymptotic target is either
\[
        \Gamma_r(m,w)\le n^{C r}
        \quad\text{for all }2\le r\le c w,
\]
with constants uniform in the row, or an equivalent tail estimate implying
\[
        \max_z \mu_{m,w}(z)
        \le n^{C}\,|\B|^{-w}.
\]
For exact deployment, print the finite $r$-moment or direct-tail constants and show that the induced max-fiber multiplier fits under the adjacent fail margin.
\end{remark}

\begin{remark}[why fixed low moments are not enough]\label{rem:capfpr-low-moments}
A bound on $\Gamma_r$ for fixed $r$ proves useful bulk equidistribution but does not by itself imply the max-fiber estimate.  Indeed
\[
        \max_z\mu_{m,w}(z)
        \le \Gamma_r(m,w)^{1/r}|\B|^{-w(1-1/r)},
\]
so recovering the average scale $|\B|^{-w}$ up to a polynomial factor from moments alone requires $r$ comparable to $w$, or else a separate tail or inverse theorem.  This is why the collision-gap programme should measure the whole joint tail, not only the second moment.
\end{remark}

\subsection{Stable Fourier measurement protocol}\label{subsec:capfpr-normalized-recursion}

The normalized elementary-symmetric recursion gives a numerically stable way to measure Fourier coefficients of the prefix distribution without ever forming the enormous coefficients $e_m$.

\begin{proposition}[non-expansive normalized recursion]\label{prop:capfpr-normalized-recursion}
Let $\varepsilon_1,\ldots,\varepsilon_N\in\mathbb C$ satisfy $|\varepsilon_i|\le1$.  Define
\[
        q_j^{(i)}
        := {1\over \binom i j}
           \sum_{1\le r_1<\cdots<r_j\le i}
           \varepsilon_{r_1}\cdots\varepsilon_{r_j},
        \qquad 0\le j\le i,
\]
with $q_0^{(i)}=1$.  Then, for $1\le j\le i$,
\[
        q_j^{(i)}=
        \left(1-{j\over i}\right)q_j^{(i-1)}
        +{j\over i}\,\varepsilon_i q_{j-1}^{(i-1)}.
\]
In particular $|q_j^{(i)}|\le1$ for all $i,j$.
\end{proposition}

\begin{proof}
The elementary symmetric recursion is
\[
        e_j^{(i)}=e_j^{(i-1)}+\varepsilon_i e_{j-1}^{(i-1)}.
\]
Dividing by $\binom i j$ and using
\[
        {\binom {i-1}j\over\binom ij}=1-{j\over i},
        \qquad
        {\binom {i-1}{j-1}\over\binom ij}={j\over i}
\]
gives the formula.  If $|q_j^{(i-1)}|,|q_{j-1}^{(i-1)}|\le1$, the displayed formula is a convex combination of two complex numbers of modulus at most one.  The claim follows by induction.
\end{proof}

For a frequency vector $\tau=(\tau_1,\ldots,\tau_w)$ over a prime field and an additive character $\psi$, put
\[
        \varepsilon_x=\psi\!\left(\sum_{r=1}^{w}\tau_r x^r\right),
        \qquad x\in D.
\]
Then
\[
        q_m^{(n)}(\varepsilon_x:x\in D)
        ={1\over\binom n m}\sum_{|M|=m}\psi\!\left(\sum_{r=1}^{w}\tau_r\sum_{x\in M}x^r\right),
\]
which is the normalized Fourier coefficient of the power-sum prefix distribution.  When $w<\operatorname{char}\F$, Newton identities give a triangular bijection between the first $w$ power sums and the first $w$ elementary symmetric prefixes, so the same measurement probes the first-form quotient input \cref{prob:capfpr-Q}.

\begin{remark}[Weil barrier and the reason to measure joint cancellation]\label{rem:capfpr-weil-barrier}
The per-power-sum Weil bound is the wrong endpoint for \textup{(Q)}.  It pays a visible $\sqrt p$ cost per active power-sum direction, and the fable base-case theorem \cref{thm:capff1-collision-base} shows exactly how far that method reaches at head depths.  The computations described in \cref{rem:capff1-calibration} indicate that the second-moment refinement over a full period returns the same square-root scale: individual character maxima are controlled by the wrong statistic.  The normalized recursion should therefore be used to measure joint tails and collision gaps, not to claim a uniform character bound strong enough to imply \textup{(Q)} by itself.
\end{remark}

\subsection{Heuristic picture and falsification tests}\label{subsec:capfpr-heuristics}

The current heuristic is as follows.  In the dense primitive bulk, the prefix map behaves like an almost-random constant-weight syndrome map.  Quotient-periodic directions account for the visible heavy Fourier modes, and those are precisely the directions absorbed into coarser scales by \cref{prop:capff1-composite}.  After those scales are removed, the remaining primitive fibers should have only polynomial max-to-mean inflation.  On the M1 side, the split-locator ledger of \cref{thm:capf-first-moment,thm:capf-second-moment,cor:capf-exact-variance} says that random residual correlations are supported only in a bounded Johnson neighborhood; the missing safe-side theorem is an inverse statement that every worst-case excess correlation has a paid structural cause.

A practical falsification programme should therefore print, for each tested scale:
\begin{enumerate}[label=\textup{(\roman*)},leftmargin=2.4em]
\item max-to-mean prefix fiber ratios in the dense bulk, separated by primitive and composite frequency support;
\item $\Gamma_r$ values for increasing $r$, with enough growth in $r$ to test the max-fiber prediction rather than only the second moment;
\item normalized-recursion Fourier tail plots or tables, with divisor-lattice buckets for composite directions;
\item the exact paid scale to which every visible structured spike is absorbed.
\end{enumerate}
A counterexample to the frontier should then be classifiable: either it produces a super-polynomial primitive prefix fiber, refuting \textup{(Q)}, or it produces a super-polynomial primitive aperiodic residual family after the paid cells, refuting \textup{(A)}.

\providecommand{\B}{\mathbb B}
\providecommand{\F}{\mathbb F}
\providecommand{\RS}{\mathrm{RS}}
\providecommand{\Dloc}{\mathcal D}
\providecommand{\PP}{\mathbb P}
\providecommand{\Bad}{\operatorname{Bad}}
\providecommand{\Syn}{\operatorname{Syn}}

\section{Safe-side proof routes}\label{sec:capfaq-strategies}

The frontier package leaves explicit upper inputs, formalized in \cref{sec:frontier-package,sec:frontier-next-steps}.  This section records proof routes and intermediate reductions for those inputs.  The symbols \textup{(A)} and \textup{(B/Q)} are used locally for the first aperiodic and quotient split; the final interface is the three-part package stated in \cref{subsec:capg-active-interface}.

\subsection{First-form contracts}\label{subsec:capfaq-contracts}

\paragraph{First-form aperiodic contract.}
Fix a multiplicative or circle-line row $C=\RS[\F,D,k]$, an agreement value $A$, and write
\[
        j=n-A,
        \qquad t=A-k,
        \qquad j+t=n-k .
\]
For a received affine line $f+zg$, let
\[
        \Bad_{\rm ap}^{\rm M1}(f,g;A)
\]
be the set of finite slopes $z$ admitting a support-wise noncontained explanation at agreement at least $A$ after the tangent cell, declared quotient cells, extension-pole cells, common-GCD strata, quotient-pullback strata, and printed fixed-dimensional Conjecture-F strata have been paid.  Every remaining witness support is required to be aperiodic in the sense of \cref{def:periodicity-scale}.

The asymptotic \textup{(A)} theorem should prove a row-uniform polynomial bound
\[
        \sup_{f,g}\bigl|\Bad_{\rm ap}^{\rm M1}(f,g;A)\bigr|
        \le n^{C_A}
\]
throughout the frontier band of the broad band formulation \cref{prob:band}, with $C_A$ independent of $q$.  The deployed finite theorem must be sharper: for each printed adjacent candidate $a_0+1$, it must supply an explicit integer
\[
        U_A(a_0+1)
        \ge
        \sup_{f,g}\bigl|\Bad_{\rm ap}^{\rm M1}(f,g;a_0+1)\bigr|
\]
whose bit cost is included in the exact upper staircase inequality
\[
        U_{\rm paid}(a_0+1)+U_Q(a_0+1)+U_A(a_0+1)\le B_* .
\]
Thus a polynomial aperiodic theorem is an asymptotic closure input, while a finite adjacent closure input is an exact-numerator certificate.

\paragraph{First-form quotient contract.}
For the quotient-remainder prefix map of \cref{def:prefix-fiber-certificate}, write
\[
        \Phi_{c,m,s}^{K}:\mathcal X_{c,m,s}\longrightarrow \B^{w_c(s,A_0-K)},
        \qquad
        \mu(\xi)=\bigl|\Phi_{c,m,s}^{K}{}^{-1}(\xi)\bigr|,
\]
where $\mathcal X_{c,m,s}$ is the finite set of quotient-remainder supports counted by $M_{c,m,s}$.  In the identity-scale case of \cref{lem:capff1-identity-prefix-floor}, this is simply
\[
        \Phi_{1,m,0}^{K}(M)=\bigl((-1)^he_h(M)\bigr)_{1\le h\le m-K},
        \qquad |M|=m .
\]
The quotient-fiber upper theorem should bound the heaviest prefix fiber
\[
        H_{c,m,s}^{K}:=\max_{\xi}\mu(\xi)
\]
from above, not only from below by pigeonhole.  The asymptotic grade may allow
\[
        H_{c,m,s}^{K}
        \le n^{C_Q}\,{M_{c,m,s}\over |\operatorname{im}\Phi_{c,m,s}^{K}|}
\]
or the corresponding ambient-box denominator when an image-size theorem is also supplied.  This loses only $C_Q\log_2 n$ bits and is therefore usable away from an $O(\log n)$ reserve.

The deployed finite grade must instead print constants.  For each adjacent candidate it should give a certified inequality of the form
\[
        H_{c,m,s}^{K}
        \le
        \Lambda_Q(c,m,s,K)\,{M_{c,m,s}\over |\operatorname{im}\Phi_{c,m,s}^{K}|}+E_Q(c,m,s,K)
\]
with exact integer rounding and with
\[
        \log_2\Lambda_Q+
        \log_2\left(1+{E_Q\,|\operatorname{im}\Phi_{c,m,s}^{K}|\over M_{c,m,s}}\right)
\]
small enough to fit inside the printed adjacent fail margin after the paid and aperiodic terms have been added.  Since the smallest deployed adjacent margin in \cref{prop:capff1-identity-frontier,prop:capff1-closing} is only a few bits, an unspecified $\operatorname{poly}(n)$ multiplier is not a finite certificate.

\subsection{The normalized aperiodic input}\label{subsec:capfaq-A-strategies}

\paragraph{A1. Deterministic local-limit upgrade from the split-locator moment ledger.}
The moment ledger of \cref{thm:capf-first-moment,thm:capf-second-moment,cor:capf-exact-variance,cor:capf-dependency} gives the correct random-word model but not a worst-case theorem.  The proposed upgrade is an inverse theorem.

\begin{enumerate}[label=\textup{A1.\arabic*},leftmargin=3.0em]
\item Convert every bad slope to at least one co-support locator $L_T$ with $|T|=j$ satisfying the fixed-locator Hankel equation
\[
        H_{t,j}(\Syn(f+zg))\boldsymbol\ell_T=0 .
\]
Use a deterministic tie-breaking rule so that each slope pays one canonical primitive locator after all paid branches are removed.
\item Prove a density dichotomy for the chosen locator set inside the Johnson graph on $j$-subsets: either it has random-scale local intersections, in which case the dependency-degree calculation forces a polynomial number of distinct slopes, or it has too many short Johnson exchanges.
\item In the second case, apply the exchange identities of \cref{prop:capf-johnson-exchange} and the anticode obstruction of \cref{prop:capf-anticode}.  Long exchange chains should force one of the already paid structures: common GCD, quotient-pullback periodicity, tangent containment, extension-pole containment, or an identically deficient SPI chart.
\item The deliverable is a deterministic replacement for the random-pair variance sentence: after paid structures are removed, every worst-case pair behaves locally like the random split-locator model up to $n^{O(1)}$ exceptional choices.
\end{enumerate}

The advantage of this route is that it reuses the strongest capf calculations already printed.  The risk is the inverse step: one must prove that excess local correlation has a structural cause, not merely that it is rare on average.

\paragraph{A2. Growing-dimensional primitive Conjecture-F incidence.}
For a fixed slope $z$, the locator equation cuts a projective linear space
\[
        \PP(W_z)\subseteq \PP(K[X]_{\le j})
\]
and the explanations at that slope are precisely the squarefree locator points in
\[
        \PP(W_z)\cap \Dloc_j(D)
\]
after the usual locator normalization.  The fixed-dimensional and projective-plane cases are already paid in \cref{thm:capf-fixeddim,thm:capf-dim2,lem:capf-dim1}; the missing part is the primitive growing-dimensional case.

A proof can be organized as follows.

\begin{enumerate}[label=\textup{A2.\arabic*},leftmargin=3.0em]
\item Use \cref{lem:capf-gcd,lem:capf-quot-pullback,lem:capf-gap2} to remove common-root and quotient-periodic components before the incidence theorem is invoked.  This ensures that the incidence problem is genuinely primitive.
\item Prove a Hilbert-function or polynomial-method bound for
\[
        |\PP(W)\cap\Dloc_j(D)|
\]
when $W$ has growing dimension but contains no paid subspace forced by quotient pullback, large common divisor, or low-dimensional projection.
\item Apply the bound uniformly to the pencil $W_z$ as $z$ varies.  If too many $z$ give nonempty primitive intersections, eliminate $z$ from the minors defining $W_z$; the resulting identically singular or identically split charts are precisely the named SPI residuals of \cref{thm:capf-spi}, and the non-identical charts contribute by their printed degree caps.
\item The target output is either
\[
        \#\{z:\PP(W_z)\cap\Dloc_j(D)_{\rm prim}\ne\varnothing\}\le n^{C_A}
\]
for asymptotic closure, or an exact degree-sum table for the deployed adjacent rows.
\end{enumerate}

This route is the cleanest conceptual proof of \textup{(A)} because it attacks exactly the frontier sentence in the proof guide: the growing-dimensional primitive incidence problem inside the broad band formulation \cref{prob:band}.

\paragraph{A3. Exchange-rigidity and Johnson-spectral route.}
Instead of bounding each projective incidence independently, one can exploit the fact that many bad slopes produce many locators with constrained exchanges.  The proof target is a rigidity theorem of the following form: a large primitive aperiodic locator family in the M1 band either expands in the Johnson graph like a random family, or its adjacent exchanges preserve enough locator-syndrome data to force quotient periodicity.

The concrete steps are:
\begin{enumerate}[label=\textup{A3.\arabic*},leftmargin=3.0em]
\item Define the exchange graph on canonical bad co-supports, connecting $T$ and $T'$ when $|T\triangle T'|=2$.
\item Use \cref{prop:capf-hankel,prop:capf-johnson-exchange} to express equality or near-equality of contributed slopes across an edge as an explicit low-rank condition on the two exchanged points.
\item Use a Johnson-graph spectral or container estimate to show that a super-polynomial primitive family contains many exchange paths unless it is contained in a small anticode.
\item Use \cref{prop:capf-anticode} to pay the anticode case; use the exchange identities along paths to force a periodic stabilizer, a common divisor, or an SPI residual in the remaining case.
\end{enumerate}

This route is attractive for finite certificates because the spectral/container constants can be stated as explicit integer inequalities on the deployed Johnson graph.

\paragraph{A4. Finite SPI atlas route.}
For the four adjacent deployed rows, a full asymptotic primitive incidence theorem may be more than is needed.  A finite route is to promote \cref{thm:capf-spi} into a chart atlas with a degree ledger:
\[
        \Bad_{\rm ap}^{\rm M1}(f,g;a_0+1)
        \subseteq
        \bigcup_{\mathfrak c\in\mathfrak C_{\rm ap}} Z(E_{\mathfrak c})
        \cup \mathfrak R_{\rm named} .
\]
Here $E_{\mathfrak c}$ is an explicit eliminant in the slope parameter and every named residual branch is either paid by a prior theorem or carried separately in $U_A$.  The certificate is then
\[
        U_A(a_0+1)=\sum_{\mathfrak c}\deg E_{\mathfrak c}+|\mathfrak R_{\rm named}|,
\]
with no asymptotic notation.  This is the right form for an exact one-step theorem, because it can be parted directly into the staircase inequality.

\subsection{Quotient-fiber equidistribution}\label{subsec:capfaq-Q-strategies}

\paragraph{Q0. Collision-gap measurement before proof.}
Before committing to any proof architecture, measure the primitive prefix distribution using the non-expansive recursion of \cref{prop:capfpr-normalized-recursion}.  The required outputs are not individual character maxima alone, but the joint quantities of the collision-gap programme \cref{prob:capfpr-collision-gap}: high moments $\Gamma_r$, max-to-mean ratios, and divisor-lattice decompositions of the Fourier tail.  The experiment is useful only if every table records whether the frequency is primitive or collapses to a coarser quotient scale in the sense of \cref{prop:capff1-composite}.  A table of single-character Weil ratios should be treated as calibration, not as evidence for a finite adjacent certificate.

The theoretical target suggested by the measurements is an inverse theorem: if the primitive dense-bulk max fiber is larger than $n^C\binom nm|\B|^{-w}$, then the responsible support family has a quotient stabilizer, a common block structure, or a low-dimensional exceptional incidence already paid by the capf ledger.  Without such an inverse theorem, the normalized recursion remains an audit instrument rather than a proof of \textup{(Q)}.

\paragraph{Q1. Constant-weight syndrome local central limit theorem.}
The identity prefix map is a constant-weight syndrome map: an $m$-subset of $D$ is sent to the first $w=m-K$ elementary-symmetric coefficients of its locator.  More generally, \cref{def:prefix-fiber-certificate} gives the quotient-remainder version.  The desired theorem is a local central limit theorem for these fibers.

The proof route is:
\begin{enumerate}[label=\textup{Q1.\arabic*},leftmargin=3.0em]
\item Use additive-character inversion over $\B^w$:
\[
        \mu(\xi)
        = |\B|^{-w}\sum_{\alpha\in\B^w}\psi(-\alpha\cdot\xi)
          \sum_{X\in\mathcal X_{c,m,s}}
          \psi\bigl(\alpha\cdot\Phi_{c,m,s}^{K}(X)\bigr).
\]
The trivial character gives the average fiber size.
\item Bound every nontrivial symmetric exponential sum in the last display.  The needed estimate is not merely cancellation on average over $\xi$; it must be uniform in the character $\alpha$ and in the deployed window of $m,w$.
\item Convert the character-sum bound into an explicit multiplier $\Lambda_Q$ and additive error $E_Q$.  For asymptotic closure one may allow $\Lambda_Q=n^{O(1)}$; for deployed adjacent closure the multiplier and additive term must be evaluated in bits and compared against the printed fail margins.
\end{enumerate}

This route has the clearest finite audit trail: every loss is visible as a character-sum constant.

\paragraph{Q2. Algebraic-monodromy route for elementary-symmetric fibers.}
The map from squarefree locators to their high-degree prefix coefficients is a finite morphism after the degree and support model are fixed.  One can try to prove that, after quotient decompositions are removed, this morphism has full primitive monodromy and hence nearly uniform fibers.

The proof package would contain:
\begin{enumerate}[label=\textup{Q2.\arabic*},leftmargin=3.0em]
\item a geometric irreducibility theorem for the primitive fiber product of two prefix fibers;
\item a branch-locus bound for exceptional prefix values;
\item a finite-field point-count theorem with constants, specialized to the multiplicative and circle-line domains;
\item an explicit treatment of quotient-periodic components, so that the theorem does not re-count cells already handled by \cref{thm:exact-quotient-image-lcm-ledger}.
\end{enumerate}

This is conceptually strong, but the finite deployed use requires effective constants.  A qualitative Lang--Weil or monodromy theorem would only give an asymptotic \textup{(B/Q)} input unless all constants are printed.

\paragraph{Q3. Switch-chain stability and inverse equidistribution.}
There is also a purely combinatorial route.  Run the usual switch chain on quotient-remainder supports: replace one selected point or fiber by one unselected point or fiber while keeping the agreement size fixed.  A very heavy prefix fiber is stable under unexpectedly many switches.  The inverse statement should say that such stability forces one of the paid quotient structures.

\begin{enumerate}[label=\textup{Q3.\arabic*},leftmargin=3.0em]
\item Prove an edge-isoperimetric inequality for the support model $\mathcal X_{c,m,s}$.
\item Show that a prefix fiber whose size exceeds the average by more than the permitted multiplier must have many internal switches.
\item Translate internal switches into polynomial identities among elementary symmetric coefficients.  Long chains of such switches force either a quotient stabilizer, a common block structure, or a low-dimensional exceptional family.
\item Pay those forced structures using the quotient image ledger, leaving only a uniformly distributed primitive remainder.
\end{enumerate}

This route aligns well with the exchange-rigidity strategy for \textup{(A)} and may permit one combined inverse theorem: excess mass in either the M1 local-limit problem or the quotient prefix fibers must be quotient-periodic.

\paragraph{Q4. Exact finite quotient-ledger certificate.}
For the finite adjacent theorem, the useful output is a certificate rather than a broad asymptotic theorem.  For each row and each candidate agreement $a_0+1$, print a quotient-fiber ledger containing:
\begin{enumerate}[label=\textup{Q4.\arabic*},leftmargin=3.0em]
\item the exact map $\Phi_{c,m,s}^{K}$ or its identity-scale specialization;
\item the denominator used, either $|\operatorname{im}\Phi|$ or the ambient box $|\B|^w$;
\item a proved upper bound for $H_{c,m,s}^{K}$, with all constants and rounding;
\item the resulting quotient numerator $U_Q(a_0+1)$ after parameter-image coalescing;
\item a bit-margin line showing
\[
        U_{\rm paid}(a_0+1)+U_Q(a_0+1)+U_A(a_0+1)\le B_* .
\]
\end{enumerate}
If the line contains a hidden $n^{C}$ or an unspecified absolute constant, it should be downgraded to an asymptotic certificate and not used for the adjacent deployed pair.

\subsection{Proof order}\label{subsec:capfaq-order}

The lowest-risk development order is:
\begin{enumerate}[label=\textup{(\arabic*)},leftmargin=2.5em]
\item Prove the asymptotic \textup{(B/Q)} theorem first, even with an $n^{C_Q}$ multiplier.  This validates the repaired frontier statement away from an $O(\log n)$ reserve and tests the quotient-fiber formalism on the identity-scale floor.
\item Prove an asymptotic \textup{(A)} theorem by the A2 incidence route or the A3 exchange-rigidity route.  This isolates the primitive aperiodic residual after the quotient theorem has been made precise.
\item Return to finite deployment only after the asymptotic structure is stable.  Replace $n^{O(1)}$ in \textup{(B/Q)} by a printed $\Lambda_Q,E_Q$ pair, and replace $n^{O(1)}$ in \textup{(A)} by an exact SPI or exchange-degree ledger.
\item Run the one-step compiler of \cref{thm:capf-staircase,thm:capf-corridor}: finite closure is claimed only after the exact sum of paid, quotient, and aperiodic numerators is below $B_*$ at $a_0+1$.
\end{enumerate}

\begin{remark}[what not to promote]\label{rem:capfaq-not-to-promote}
The following statements should remain visibly weaker than the finite adjacent theorem: random-word M1 concentration without a worst-case inverse theorem; support-family quotient counts without parameter-image coalescing; quotient-fiber equidistribution with an unspecified polynomial multiplier; and SPI charts whose identically deficient branches are not paid or named.  Each of these is useful evidence or an asymptotic component, but none is by itself an exact deployed upper certificate.
\end{remark}

\providecommand{\Rone}{\mathrm{R1}}
\providecommand{\Ker}{\operatorname{Ker}}
\providecommand{\rank}{\operatorname{rank}}
\providecommand{\Mode}{\operatorname{Mode}}
\providecommand{\Fib}{\operatorname{Fib}}
\providecommand{\Aff}{\operatorname{Aff}}
\providecommand{\wt}{\operatorname{wt}}
\providecommand{\Cen}{\operatorname{cen}}
\providecommand{\wdeg}{\operatorname{wdeg}}

\section{Rank-one and split-pencil reduction of the safe side}\label{sec:capfr1-programme}

This section refines the safe-side package.  The first reduction has the form
\[
        \boxed{\textup{(Q)}+\Rone}
\]
where \textup{(Q)} is the prefix-fiber census and \(\Rone\) is a rank-one support census obtained by eliminating the slope parameter.  Thus \(\Rone\) is not a proof of \textup{(A)} by itself; rather, after the paid tangent, quotient, extension-pole, common-GCD, quotient-pullback, fixed-dimensional, sunflower, and SPI cells are removed, \(\Rone\) is the concrete support-counting statement that implies the remaining aperiodic-band input.  The additional lattice reduction below sharpens this one step further: \(\Rone\) is reduced to a balanced two-generator lattice census, and then to a split-pencil problem for determinantal representations of the domain polynomial.  Thus the intermediate aperiodic target becomes a split-pencil census; the final closing statement later adds the base-field normalization and the primitive shift-pair ledger.

\subsection{Band normalization and the master flatness target}\label{subsec:capfr1-band-flatness}

The identity-prefix floor changes how the broad band formulation \cref{prob:band} should be read.  Below the entropy--subfield envelope, the identity-scale construction produces exponentially many generic support witnesses, and generic supports are aperiodic in the sense of \cref{def:periodicity-scale} except for a negligible quotient-periodic subfamily.  Therefore the aperiodic band input should be normalized with its left edge at the entropy frontier plus a reserve.

\begin{remark}[Normalized aperiodic band formulation]\label{prob:capfr1-normalized-band}
Let \(g^*(\rho,\beta)\) be as in \cref{def:capff1-gstar}.  Prove the aperiodic upper input only in the range
\[
        A\ge k+g^*(\rho,\beta)n+R(n),
\]
where the first asymptotic target is \(R(n)\gg\log n\), and the finite deployed target is the printed one-step value \(a_0+1\) with a fully explicit integer numerator.  The sub-band \(A\lesssim k+g^*n\) is unsafe by the identity-prefix mechanism and should not be assigned to \textup{(A)}.
\end{remark}

The common geometric object behind \textup{(Q)} and \(\Rone\) is split-locator flatness.  Let \(\Dloc_j(D)\) denote the monic squarefree degree-\(j\) locators whose roots lie in \(D\).

\begin{remark}[master split-locator flatness]\label{prob:capfr1-master-flatness}
Find an explicit polynomial \(P\) such that, for every affine subspace \(\mathcal A\) of the monic degree-\(j\) polynomial space over \(\B\) of codimension \(s\), after quotient and common-divisor components have been removed,
\[
        |\mathcal A\cap\Dloc_j(D)|
        \le
        P(n)\left(\binom n j |\B|^{-s}+1\right).
\]
The prefix-fiber input \textup{(Q)} is the case where \(\mathcal A\) fixes the top prefix of a locator.  The aperiodic/rank-one input is the case where \(\mathcal A\) is cut out by the support interpolation equations below.
\end{remark}

The first term in \cref{prob:capfr1-master-flatness} crosses one at exactly the entropy--subfield envelope.  This explains why polynomial-loss versions can close the asymptotic frontier with an \(O(\log n)\) reserve but cannot certify the one-step deployed pair without exact constants.

\subsection{The Q-side census: moment kernels and packing}\label{subsec:capfr1-Q-kernel}

Assume in this subsection that \(w<\operatorname{char}\F\), so Newton identities identify the first \(w\) elementary-symmetric prefixes with the first \(w\) power sums.  For \(m\)-subsets define
\[
        \Psi_{m,w}(M)=\left(\sum_{x\in M}x^h\right)_{1\le h\le w}.
\]
Define the moment kernel
\[
        K_w(D)=\left\{f:D\to\B:\ \sum_{x\in D} f(x)x^h=0\text{ for }1\le h\le w\right\}.
\]

\begin{proposition}[moment-kernel dictionary]\label{prop:capfr1-moment-kernel}
Let \(M,M'\subseteq D\) be \(m\)-subsets.  If \(\Psi_{m,w}(M)=\Psi_{m,w}(M')\), then
\[
        f=1_M-1_{M'}
\]
belongs to \(K_w(D)\), takes values in \(\{-1,0,1\}\), and has equally many \(+1\) and \(-1\) coordinates.  Conversely, if \(f\in K_w(D)\) is \(\{-1,0,1\}\)-valued with \(|f^{-1}(1)|=|f^{-1}(-1)|=e\), then for every common core \(C\subseteq D\setminus\operatorname{supp}(f)\) of size \(m-e\), the two sets
\[
        M=C\cup f^{-1}(1),
        \qquad
        M'=C\cup f^{-1}(-1)
\]
have the same first \(w\) power sums.
\end{proposition}

\begin{proof}
The equality \(\Psi_{m,w}(M)=\Psi_{m,w}(M')\) says precisely that
\[
        \sum_{x\in D}(1_M(x)-1_{M'}(x))x^h=0
        \qquad(1\le h\le w),
\]
which is the definition of membership in \(K_w(D)\).  The values and balance are immediate from the fact that \(M\) and \(M'\) have the same size.  The converse is the same identity read backwards after adding an arbitrary common core on which the two indicators agree.

\end{proof}

\begin{proposition}[balanced-ternary collision ledger]\label{prop:capfr1-collision-ledger}
Assume that the moment map defining \(K_w(D)\) has rank \(w\) over the base field.  Then \(K_w(D)\) has dimension \(n-w\), and the second moment of the power-sum prefix fibers is the weighted balanced-ternary enumerator
\[
        \sum_z |\Psi_{m,w}^{-1}(z)|^2
        =
        \binom n m+
        \sum_{e\ge \delta_w}\binom{n-2e}{m-e}\,T_w(e),
\]
where \(T_w(e)\) is the number of ordered pairs \((A,A')\) of disjoint \(e\)-subsets of \(D\) satisfying
\[
        \sum_{x\in A}x^h=\sum_{x\in A'}x^h
        \qquad(1\le h\le w),
\]
equivalently the number of \(\{0,\pm1\}\)-valued words in \(K_w(D)\) with exactly \(e\) positive and \(e\) negative coordinates.  Consequently \(\Gamma_2\) is exactly a balanced-ternary weight enumerator of the moment kernel.
\end{proposition}

\begin{proof}
The rank claim follows because the Vandermonde minor on any \(w\) distinct nonzero points of \(D\) is nonsingular.  For the moment identity, count ordered pairs \((M,M')\) in one prefix fiber.  If \(M=M'\), there are \(\binom n m\) pairs.  Otherwise write
\[
        A=M\setminus M',\qquad A'=M'\setminus M,
\]
so \(|A|=|A'|=e\), the two sets are disjoint, and the common core \(C=M\cap M'\) has size \(m-e\).  Equality of power-sum prefixes is exactly the displayed condition on \((A,A')\), because the common core cancels.  Conversely, any such \((A,A')\) and any common core \(C\subseteq D\setminus(A\cup A')\) of size \(m-e\) give a colliding ordered pair \((C\cup A,C\cup A')\).  By \cref{lem:capfr1-kernel-distance}, nontrivial collisions have \(e\ge\delta_w\).  The number of possible cores is \(\binom{n-2e}{m-e}\), giving the formula.
\end{proof}

\begin{remark}[certified size of the packing improvement]\label{rem:capfr1-packing-bits}
At the four deployed frontier rows, the all-depth packing bound of \cref{prop:capfr1-prefix-packing} improves the trivial bound \(\binom n m\) by about \(2.13\cdot10^5\) bits, while still remaining about \(1.88\cdot10^6\) bits above the conjectured average fiber.  Thus it is a genuine unconditional \textup{(Q)} theorem, but the finite adjacent problem still requires the high-moment or extremality input of \cref{prob:capfr1-mode-null}.
\end{remark}

\begin{lemma}[BCH-style distance of the moment kernel]\label{lem:capfr1-kernel-distance}
If \(0\ne f\in K_w(D)\) is supported on fewer than \(w+1\) points, then this is impossible.  Consequently two distinct sets in one \(\Psi_{m,w}\)-fiber differ on at least
\[
        \delta_w:=\left\lceil {w+1\over2}\right\rceil
\]
points on each side.
\end{lemma}

\begin{proof}
Let \(S=\operatorname{supp}(f)\), \(|S|=s\le w\).  The equations \(\sum_{x\in S}f(x)x^h=0\) for \(1\le h\le s\) form a Vandermonde system after multiplying the rows by the nonzero diagonal \(x\mapsto x\) and using that the elements of \(D\) are distinct and nonzero in the smooth rows.  The only solution is \(f|_S=0\), contradiction.  If two fiber sets differ by \(e\) points on each side, the difference vector has support \(2e\).  The first assertion gives \(2e\ge w+1\), hence \(e\ge\delta_w\).
\end{proof}

\begin{proposition}[all-depth packing bound for prefix fibers]\label{prop:capfr1-prefix-packing}
For every value \(z\) of the prefix map and every \(m,w\),
\[
        |\Psi_{m,w}^{-1}(z)|
        \le
        {\binom{n}{m-\delta_w+1}\over \binom{m}{\delta_w-1}},
        \qquad
        \delta_w=\left\lceil {w+1\over2}\right\rceil .
\]
The same bound holds for the elementary-prefix map \(\Phi_{m,w}\) when \(w<\operatorname{char}\F\).
\end{proposition}

\begin{proof}
Let \(\mathcal F\) be one fiber.  By \cref{lem:capfr1-kernel-distance}, distinct \(M,M'\in\mathcal F\) satisfy \(|M\cap M'|\le m-\delta_w\).  Therefore no \((m-\delta_w+1)\)-subset of \(D\) can be contained in two different members of \(\mathcal F\).  Counting pairs \((R,M)\) with \(M\in\mathcal F\), \(R\subseteq M\), and \(|R|=m-\delta_w+1\), gives
\[
        |\mathcal F|\binom{m}{\delta_w-1}
        \le
        \binom n {m-\delta_w+1}.
\]
This is the displayed inequality.  The elementary-prefix statement follows from the triangular Newton equivalence for \(w<\operatorname{char}\F\).
\end{proof}

\begin{remark}[what the packing bound does and does not do]\label{rem:capfr1-packing-gap}
The packing bound is a genuine all-depth \textup{(Q)} theorem and should be kept.  It is still far above the conjectured average fiber in the deployed frontier window, so it is not a finite adjacent certificate.  The stronger locator-subtraction rigidity for equal elementary prefixes often gives the sharper side-distance \(e\ge w+1\); turning that sharpening into an exact second-moment or high-moment ledger is one of the low-risk Q1 deliverables.
\end{remark}

\subsection{Slope elimination and the rank-one support census}\label{subsec:capfr1-slope-elimination}

Let \(T\subseteq D\) have size \(m\), and set \(w=m-K\).  Write
\[
        \Lambda_T(X)=\prod_{x\in T}(X-x).
\]
For a word \(h:T\to\F\), define the interpolation functionals
\[
        S_r(h,T)=\sum_{x\in T}{h(x)x^r\over \Lambda_T'(x)},
        \qquad 0\le r\le w-1.
\]

\begin{lemma}[support interpolation test]\label{lem:capfr1-interpolation-test}
A word \(h:T\to\F\) is the restriction to \(T\) of a polynomial of degree \(<K\) if and only if
\[
        S_r(h,T)=0\qquad(0\le r\le w-1).
\]
\end{lemma}

\begin{proof}
Let \(P_T\) be the unique interpolating polynomial of degree \(<m\) with \(P_T(x)=h(x)\) on \(T\).  For every integer \(d\), Lagrange interpolation applied to \(X^d\) shows
\[
        \sum_{x\in T}{x^d\over\Lambda_T'(x)}=0\quad(0\le d\le m-2),
        \qquad
        \sum_{x\in T}{x^{m-1}\over\Lambda_T'(x)}=1.
\]
If \(\deg P_T<K\), then for \(0\le r\le w-1=m-K-1\), every monomial of \(X^rP_T(X)\) has degree at most \(m-2\), so \(S_r(h,T)=0\).

Conversely, write \(P_T=\sum_{i=0}^{m-1}c_iX^i\).  The equation \(S_0=0\) has leading term \(c_{m-1}\) and lower-degree terms vanish by the displayed identity, so \(c_{m-1}=0\).  Having shown \(c_{m-1}=\cdots=c_{m-r}=0\), the equation \(S_r=0\) has leading remaining term \(c_{m-1-r}\), and all lower terms vanish.  Induction for \(r=0,\ldots,w-1\) gives \(c_i=0\) for every \(i\ge K\).  Thus \(\deg P_T<K\).
\end{proof}

For a received affine line \((u,v)\), put
\[
        \alpha(T)=\bigl(S_r(u,T)\bigr)_{0\le r<w},
        \qquad
        \beta(T)=\bigl(S_r(v,T)\bigr)_{0\le r<w}.
\]
Call \(T\) a \emph{rank-one support} for \((u,v)\) if the two vectors \(\alpha(T),\beta(T)\in\F^w\) span a subspace of dimension at most one.  It is a \emph{common support} if \(\alpha(T)=\beta(T)=0\).

\begin{proposition}[slope elimination]\label{prop:capfr1-slope-elimination}
Fix \((u,v)\) and a support \(T\) of size \(m\).  A finite slope \(z\) has \((u+zv)|_T\) extending to a degree-\(<K\) polynomial if and only if
\[
        \alpha(T)+z\beta(T)=0.
\]
If \(T\) is not common, then it determines at most one such slope.  Consequently, after common-support/tangent cells are paid, the number of bad finite slopes is bounded by the number of non-common rank-one supports.
\end{proposition}

\begin{proof}
By linearity of the functionals and \cref{lem:capfr1-interpolation-test}, the restriction \((u+zv)|_T\) extends to degree \(<K\) exactly when
\[
        S_r(u,T)+zS_r(v,T)=0
        \qquad(0\le r<w),
\]
which is the vector equation.  If \(\beta(T)=0\), then either \(\alpha(T)=0\), the common case, or no slope works.  If \(\beta(T)\ne0\), then a solution, if it exists, is forced by any nonzero coordinate of \(\beta(T)\), and the vector equation is possible exactly when \(\alpha(T)
\) and \(\beta(T)\) are proportional.  Thus a non-common support determines at most one slope.  Choosing one canonical non-common support for each unpaid bad slope gives an injection from slopes to non-common rank-one supports.
\end{proof}

\begin{remark}[Rank-one support census after slope elimination]\label{prob:capfr1-rank-one-census}
After all paid cells and common supports have been removed, prove the primitive census
\[
        \#\{T\subseteq D:\ |T|=m,\ T\text{ non-common rank-one for }(u,v)\}
        \le
        \exp(o(n))\max\left(1,\binom n m q^{-(w-1)}\right),
\]
uniformly over primitive affine lines \((u,v)\).  The finite deployed form must replace \(\exp(o(n))\) by exact constants and an integer numerator usable at \(a_0+1\).
\end{remark}

\begin{corollary}[intermediate rank-one closing reduction]\label{cor:capfr1-Q-R1-closing}
The asymptotic frontier of \cref{prob:capfpr-frontier-prize} follows from the asymptotic form of \textup{(Q)} and the asymptotic form of \(\Rone\), together with the paid capf compilers.  The finite adjacent deployed theorem follows from finite \textup{(Q)} and finite \(\Rone\) numerators satisfying the one-step inequality
\[
        U_{\rm paid}(a_0+1)+U_Q(a_0+1)+U_{\Rone}(a_0+1)
        \le B_*<L(a_0).
\]
\end{corollary}

\begin{proof}
By \cref{prop:capfr1-slope-elimination}, every unpaid primitive bad finite slope contributes a distinct non-common rank-one support.  Thus \(\Rone\) supplies the aperiodic numerator \(U_A\) required by \cref{prop:capfpr-AQ-closing}.  The rest is exactly the closing implication already proved there and in \cref{prop:capff1-closing}.
\end{proof}

\begin{lemma}[pair descent for close slopes]\label{lem:capfr1-pair-descent}
Fix an affine line \((u,v)\) with \(v\notin C\), and fix an agreement \(A>n/2\).  Call a finite slope \(z\) close if \(u+zv\) agrees with some codeword \(c_z\) on at least \(A\) points.  If \(z_1\ne z_2\) are close, then
\[
        v_{12}:={c_{z_1}-c_{z_2}\over z_1-z_2}\in C
\]
agrees with \(v\) on at least \(2A-n\) points.  Coloring each close pair \(\{z_i,z_j\}\) by \(v_{ij}\), every color class is a clique.  Inside a color class with color \(v'\), the close slopes are exactly close slopes of the derived line
\[
        \bigl(u-c_{z_0}+z_0v',\ v-v'\bigr)
\]
at the same agreement, for any fixed member \(z_0\) of the clique.
\end{lemma}

\begin{proof}
Let \(T_i\) be an agreement set of size at least \(A\) for \(u+z_iv\) and \(c_{z_i}\).  Since \(A>n/2\), the intersection \(T_1\cap T_2\) has size at least \(2A-n\).  On this intersection,
\[
        (z_1-z_2)v=c_{z_1}-c_{z_2},
\]
which proves the first assertion.  If \(v_{12}=v_{23}=v'\), then
\[
        c_{z_1}-c_{z_3}=(z_1-z_2)v'+(z_2-z_3)v'=(z_1-z_3)v',
\]
so \(v_{1,3}=v'\), and color classes are cliques.  For a fixed clique member \(z_0\), the relation \(c_{z_i}=c_{z_0}+(z_i-z_0)v'\) gives
\[
        u+z_iv-c_{z_i}
        =
        (u-c_{z_0}+z_0v')+z_i(v-v'),
\]
so the agreement sets are unchanged after passing to the derived line.
\end{proof}

\begin{remark}[why pair descent alone does not close the band]\label{rem:capfr1-pair-descent-gap}
If two distinct slopes \(z_1,z_2\) are close on agreement supports of size at least \(A\), then subtracting the two explanations shows that the direction word \(v\) agrees with a codeword on at least \(2A-n\) positions.  This is useful only when that descended agreement is itself in a controlled list-decoding range.  At the KoalaBear adjacent MCA value \(A=1116048\), one has
\[
        2A-n=134944,
        \qquad
        3A-2n=-846160,
\]
so the first-order descent is far too weak to eliminate the underdetermined band.  This is precisely why the rank-one support census is the right replacement for a naive descent argument.
\end{remark}

\subsection{Lattice resolution of \(\Rone\): near-rational slopes, balanced core, and split pencils}\label{subsec:capfr1-lattice-split}

The rank-one support census admits a second reduction.  For one received word \(U:D\to\F\), support counting can be expressed as the split locus of a rank-two \(\F[X]\)-lattice.  Its near-rational stratum resolves exactly; outside that stratum every surviving support belongs to an explicit balanced two-generator family.  Finally, the divisibility condition in that family is automatic, so the terminal aperiodic input is a split-pencil problem for determinantal representations of \(\Lambda_D\).  Throughout this subsection \(K\) denotes the code dimension, \(m\) is the agreement support size, \(w=m-K\), and \(\omega=n-m\).  Let \(\widehat U\) be the degree-\(\le n-1\) interpolant of \(U\) on \(D\), and put \(\Lambda_D(X)=\prod_{x\in D}(X-x)\).

\begin{proposition}[the support census as a lattice locus]\label{prop:capfr1-lattice-census}
Define
\[
        M_U:=\{(W,N)\in\F[X]^2:\ W(x)U(x)=N(x)\text{ for all }x\in D\},
\]
and give \(\F[X]^2\) the shifted degree
\[
        \wdeg(W,N):=\max\bigl(\deg W,\deg N-(K-1)\bigr).
\]
Then:
\begin{enumerate}[label=\textup{(\alph*)}]
\item \(M_U\) is a free rank-two \(\F[X]\)-module with basis \((1,\widehat U),(0,\Lambda_D)\).  A shifted weak-Popov basis \(g_1,g_2\) has predictable degrees, and if its row degrees are \(d_1\le d_2\), then \(d_1+d_2=n-K+1\).
\item The number \(\Cen(U;m)\) of valid agreement supports of size \(m\) is exactly
\[
        \#\{(W,N)\in M_U:\ W\text{ monic},\ \deg W=\omega,\ \wdeg(W,N)\le\omega,\ W\mid\Lambda_D,\ W\mid N\}.
\]
Each element gives the codeword \(c=N/W\) and the support \(T=D\setminus Z(W)\).
\item Equivalently,
\[
        \Cen(U;m)=\sum_{c\in C}\binom{\operatorname{agr}(U,c)}{m},
\]
the \(m\)-th binomial moment of the agreement profile of \(U\).
\end{enumerate}
\end{proposition}

\begin{proof}
For \((W,N)\in M_U\), the difference \(N-W\widehat U\) vanishes on every point of \(D\), hence is a multiple of \(\Lambda_D\).  Thus \((W,N)=W(1,\widehat U)-C(0,\Lambda_D)\) for some \(C\in\F[X]\), and the two displayed generators form a basis.  Since \(\F[X]\) is Euclidean, shifted row reduction gives a weak-Popov basis; the predictable-degree property is the standard leading-position argument for weak-Popov rows, and the sum of the shifted row degrees equals the shifted degree of the determinant \(\Lambda_D\), namely \(n-K+1\).

If \(T\) is a valid support with codeword \(c\), take \(W=\Lambda_{D\setminus T}\) and \(N=Wc\).  Then \((W,N)\in M_U\), \(W\) is monic of degree \(\omega\), satisfies \(\wdeg(W,N)\le\omega\), and \(W\mid\Lambda_D,N\).  Conversely, such an element gives a polynomial \(c=N/W\) of degree \(<K\), because \(\wdeg(W,N)\le\omega\) in the census range, and the identity \(W(U-c)=0\) on \(D\) shows that \(U=c\) on \(D\setminus Z(W)\), a support of size \(m\).  Since \(m\ge K\), the codeword attached to a valid support is unique.  Counting by supports or by codewords gives the binomial-moment formula.
\end{proof}

\begin{theorem}[near-rational dichotomy]\label{thm:capfr1-near-rational-dichotomy}
Assume \(2w\le n-K\), and let \(d_1=d_1(U)\) be the minimal shifted degree of a nonzero element of \(M_U\).
\begin{enumerate}[label=\textup{(\roman*)}]
\item If \(d_1\le w\), then either \(U\) is within Hamming distance exactly \(d_1\) of a unique codeword \(c\), the minimal vector is the error-locator pair \((\Lambda_E,\Lambda_Ec)\) with \(E=\operatorname{Disagr}(U,c)\), and
\[
        \Cen(U;m)=\binom{n-d_1}{m},
\]
or \(\Cen(U;m)=0\).
\item If \(d_1\ge w+1\), every census element lies in the two-generator family
\[
        Ag_1+Bg_2,
        \qquad
        \deg A\le \omega-d_1,\quad \deg B\le \omega-d_2,
\]
whose parameter space has \(\F\)-dimension \(\omega-w+1\).
\end{enumerate}
\end{theorem}

\begin{proof}
Let \(g_1=(W_1,N_1)\), \(g_2=(W_2,N_2)\) be a shifted weak-Popov basis with row degrees \(d_1\le d_2\).  If \(d_1\le w\), then \(d_2=n-K+1-d_1\ge n-K+1-w=\omega+1\).  Predictable degrees imply that any census element of shifted degree \(\omega\) must be a multiple \(Ag_1\).  Hence the census is empty unless \(W_1\) is, up to scalar, a monic squarefree divisor of \(\Lambda_D\), and unless \(N_1/W_1\) is a polynomial \(c\) of degree \(<K\).  In the nonempty case \(c\) agrees with \(U\) off the roots of \(W_1\), and minimality forces the disagreement set to be exactly those roots; otherwise the true error-locator pair would have smaller shifted degree.  The codeword is unique because two codewords both within distance \(w\) of \(U\) would differ in at most \(2w\le n-K\) positions, below the Reed--Solomon distance \(n-K+1\).  The valid supports are exactly the \(m\)-subsets of the \(n-d_1\) agreement points, yielding the displayed binomial.  If any of the required split/divisibility/codeword conditions fails, no census element exists.

If \(d_1\ge w+1\), both rows can appear.  Predictable degrees bound the multipliers by \(\deg A\le\omega-d_1\) and \(\deg B\le\omega-d_2\).  The number of coefficients in \(A,B\) is \((\omega-d_1+1)+(\omega-d_2+1)=2\omega+2-(n-K+1)=\omega-w+1\).
\end{proof}

\begin{corollary}[per-line reduction to the balanced core]\label{cor:capfr1-balanced-line}
For an affine line \((u,v)\) at band agreement \(m\), if two finite slopes \(z_1,z_2\) have \(d_1(u+z_iv)\le w\) and nonzero census, then \((u,v)\) is jointly close to a codeword pair on at least \(n-3w\) points and hence is not MCA-bad at agreement \(m\) in the normalized band.  Consequently, outside at most one near-rational slope,
\[
        N_{\mathrm{MCA\text{-}bad}}(u,v;m)
        \le
        1+
        \#\{z:\ d_1(u+zv)\ge w+1,\ \Cen(u+zv;m)>0\}.
\]
Thus the aperiodic census reduces to the balanced core along pencils.
\end{corollary}

\begin{proof}
By \cref{thm:capfr1-near-rational-dichotomy}, a slope with \(d_1\le w\) and nonzero census has \(u+z_iv=c_i+\eta_i\) with \(\wt(\eta_i)\le w\).  For two such slopes,
\[
        v={c_1-c_2\over z_1-z_2}+\eta_v,
        \qquad \wt(\eta_v)\le2w,
\]
and then \(u\) is within distance \(3w\) of a codeword.  Hence the pair has a common codeword-pair explanation on at least \(n-3w\) positions.  In the deployed normalized band this density is far above the target agreement density \(m/n\), so no finite slope is MCA-bad once this common explanation is paid.  Therefore at most one near-rational slope can survive, and the rest must have \(d_1\ge w+1\).
\end{proof}

\begin{remark}[Balanced-core reduction]\label{prob:capfr1-balanced-core}
For every word \(U\) with \(d_1(U)\ge w+1\), prove
\[
        \#\{(A,B):\ \deg A\le\omega-d_1,\ \deg B\le\omega-d_2,
        \ W_{A,B}\mid\Lambda_D,
        \ W_{A,B}\mid N_{A,B}\}
        \le
        e^{o(n)}\max\left(1,\binom n m q^{-w}\right),
\]
where \((W_{A,B},N_{A,B})=Ag_1+Bg_2\).  This statement implies the twisted support census for every word, hence \cref{prob:capfr1-rank-one-census}.  The family has dimension \(\omega-w+1\), and its split condition is the same divisor variety that appears in \cref{thm:capf-spi}, now with a global two-generator origin.
\end{remark}

\begin{lemma}[automatic divisibility on split supports]\label{lem:capfr1-autodiv}
If \((W,N)\in M_U\) and \(W\) is squarefree with all roots in \(D\), then \(W\mid N\).
\end{lemma}

\begin{proof}
For every root \(x\) of \(W\), the defining relation gives \(N(x)=W(x)U(x)=0\).  Since \(W\) is squarefree, vanishing at all roots of \(W\) is equivalent to divisibility by \(W\).
\end{proof}

\begin{lemma}[unimodularity of lattice bases]\label{lem:capfr1-unimodular}
For any \(\F[X]\)-basis \((W_1,N_1),(W_2,N_2)\) of \(M_U\), one has \(\gcd(W_1,W_2)=1\) and
\[
        W_1N_2-W_2N_1=\gamma\Lambda_D
\]
for some \(\gamma\in\F^\times\).
\end{lemma}

\begin{proof}
Because \((1,\widehat U),(0,\Lambda_D)\) is also a basis, the change-of-basis matrix has determinant in \(\F^\times\).  Hence the determinant of \((W_1,N_1),(W_2,N_2)\) is a nonzero scalar multiple of \(\Lambda_D\).  Also \((1,\widehat U)=A_0(W_1,N_1)+B_0(W_2,N_2)\) for some \(A_0,B_0\in\F[X]\); taking first coordinates gives \(1=A_0W_1+B_0W_2\), so \(W_1,W_2\) are coprime.
\end{proof}

\begin{proposition}[determinantal representation of the adversary]\label{prop:capfr1-detrep}
Every \(\F[X]\)-basis \((W_1,N_1),(W_2,N_2)\) of an interpolation lattice \(M_U\) gives a determinantal representation
\[
        W_1N_2-W_2N_1=\gamma\Lambda_D,
        \qquad
        \gamma\in\F^\times,
        \qquad
        \gcd(W_1,W_2)=1.
\]
Conversely, any quadruple satisfying the displayed identity and coprimality defines a word \(U\) on \(D\) by
\[
        U(x)=N_i(x)/W_i(x)
\]
using any index \(i\) for which \(W_i(x)\ne0\), and the two pairs form a basis of the resulting lattice \(M_U\).  Therefore the balanced aperiodic adversary may be studied as a degree-capped family of determinantal representations of \(\Lambda_D\); the shifted weak-Popov profile \((d_1,d_2)\), \(d_1+d_2=n-K+1\), is a normalization imposed after this representation is chosen.
\end{proposition}

\begin{proof}
The forward direction is \cref{lem:capfr1-unimodular}.  For the converse, coprimality ensures that at least one of \(W_1(x),W_2(x)\) is nonzero at each \(x\in D\).  If both are nonzero, the determinant identity, evaluated at \(x\), gives \(W_1(x)N_2(x)=W_2(x)N_1(x)\), because \(\Lambda_D(x)=0\); hence the two ratios agree and \(U\) is well defined.  If \(W_1(x)=0\), then \(W_2(x)\ne0\), and the same identity gives \(N_1(x)=0\), so \((W_1,N_1)\in M_U\); the argument for the second row is identical.  The two rows generate a full-rank submodule of \(M_U\) with determinant \(\gamma\Lambda_D\), the same determinant as \(M_U\) up to a unit, hence the generated submodule equals \(M_U\).  Shifted row reduction can then be applied to choose the weak-Popov basis and its degree profile without changing the underlying determinantal parent.
\end{proof}

\begin{remark}[Split-pencil reduction]\label{prob:capfr1-split-pencil}
Let \((W_1,N_1,W_2,N_2)\) be a determinantal representation of \(\Lambda_D\) as in \cref{prop:capfr1-detrep}, with balanced profile \(d_1\ge w+1\).  Prove
\[
        \#\left\{(A,B):\begin{array}{l}
        \deg A\le\omega-d_1,
        \quad \deg B\le\omega-d_2,\\
        AW_1+BW_2\mid\Lambda_D,
        \quad \deg(AW_1+BW_2)=\omega
        \end{array}\right\}
        \le
        e^{o(n)}\max\left(1,\binom n\omega q^{1-w}\right),
\]
counting monic normalizations.  By \cref{lem:capfr1-autodiv}, this is the full balanced-core census: the condition \(W\mid N\) has been eliminated, and the remaining split condition is the single congruence \(\Lambda_D\equiv0\pmod{AW_1+BW_2}\).  For a smooth multiplicative coset this is \(X^n-\beta\equiv0\pmod{AW_1+BW_2}\); the circle rows replace \(X^n-\beta\) by the corresponding twin-coset domain polynomial.  Thus the terminal aperiodic input is an explicit split-pencil problem.
\end{remark}


\begin{remark}[active reading of the split-pencil reduction]
\label{rem:capfr1-split-pencil-active-reading}
The split-pencil reduction above is a proof reduction, not the final census
scale.  The active split-pencil input is the base-field-normalized interior
census of \cref{prob:capg-active-BC}; the boundary profile is delegated to
quotient/prefix flatness, and the primitive shift-pair ledger is kept as a
separate input.
\end{remark}

\begin{corollary}[intermediate split-pencil package]\label{cor:capfr1-terminal-package}
This intermediate package shows that, after the first-form quotient input is separated, the aperiodic numerator is supplied by a split-pencil census.  The final all-strata closure uses the base-field-normalized version in \cref{prob:capg-active-BC} and the primitive shift-pair input in \cref{prob:capg-active-shiftpairs}; polynomial-loss versions still do not decide the printed adjacent pairs.
\end{corollary}

\begin{proof}
By \cref{lem:capfr1-autodiv,prop:capfr1-detrep}, the balanced core of the balanced-core reduction \cref{prob:capfr1-balanced-core} is bounded by the split-pencil census.  By \cref{cor:capfr1-balanced-line}, the balanced core bounds the remaining line-by-line aperiodic supports after near-rational slopes have been paid.  By \cref{prop:capfr1-slope-elimination}, the rank-one support count bounds unpaid bad finite slopes.  Thus the split-pencil input supplies the aperiodic numerator required in \cref{cor:capfr1-Q-R1-closing}.  The separation between asymptotic reserve and finite adjacent constants is the same as in \cref{prop:capff1-closing,prop:capfpr-AQ-closing}.
\end{proof}

\subsection{Finite extremality and audits}\label{subsec:capfr1-finite-audits}

The finite adjacent pair is not a soft local-limit problem.  The printed fail margins are constant-sized; an unspecified factor \(n^C\) costs \(21C\) bits at \(n=2^{21}\) and therefore exceeds three of the four deployed adjacent margins.  The finite version should be reduced to two exact extremality statements.

\begin{remark}[mode-at-null / exchange-compression for \textup{(Q)}]\label{prob:capfr1-mode-null}
Let \(N_w(z)=|\Phi_{m,w}^{-1}(z)|\) at the identity scale, with quotient rungs separated.  Prove either
\[
        N_w(z)
        \le
        N_w(0)
        \qquad\text{for all }z,
\]
or an exchange-compression theorem implying the same deployed max-fiber bound after quotient-pulled-back fibers are charged to their rungs.  Then prove the exact finite inequality
\[
        N_w(0)
        \le
        \kappa\binom n m |\B|^{-w}
\]
with \(\kappa\) below the printed margin for the row.
\end{remark}

\begin{remark}[Finite-rung audit]\label{prob:capfr1-rung-audit}
For each of the four deployed adjacent values \(a_0+1\), run the exact integer scanner over every divisor/rung and every slack profile used by the quotient ledger.  A periodic support count is far above budget at support level, so the quotient bucket must be paid at image level by descent and lcm coalescing.  The audit must print per-rung bit losses; if any rung is tight or inverted, the one-step finite conjecture is threatened from the periodic side before \textup{(Q)} or \(\Rone\) is even invoked.
\end{remark}

\subsection{Updated proof ladder and risk register}\label{subsec:capfr1-risk}

The work queue is now ordered by leverage.
\begin{enumerate}[label=\textup{(\arabic*)},leftmargin=2.4em]
\item \textbf{Immediate.} Use the normalized left edge; run the rung-margin audit; prove and use the slope-elimination reduction; keep the fable2 anchors \cref{cor:capff1-collision-head,cor:capff1-circle-head,prop:capff1-witness-exact,prop:capff1-composite,cor:capff1-quotient-scales} as proved base cases for \textup{(Q)}.
\item \textbf{Low-to-medium risk.} Build the exact prefix-collision ledger using \cref{prop:capfr1-moment-kernel,prop:capfr1-prefix-packing}; stratify collisions by common core, difference size, and quotient scale.  This gives a typical-fiber theorem and calibrates the high-moment hierarchy, but not max-fiber control by itself.
\item \textbf{Main Q effort.} Prove heavy-fiber symmetry descent: few heavy fibers plus the twist action \(M\mapsto\zeta M\) force quotient-pulled-back prefix values, reducing worst-case flatness to a null-fiber or quotient-rung census.
\item \textbf{Main aperiodic effort.} Prove the split-pencil census of the split-pencil reduction \cref{prob:capfr1-split-pencil}.  This would imply the balanced-core census of the balanced-core reduction \cref{prob:capfr1-balanced-core}, then \(\Rone\), then the normalized aperiodic input.  Johnson-exchange rigidity, bounded-deficiency SPI induction, and determinantal monodromy now all target this single two-generator pencil.
\item \textbf{Finite prize.} Replace all polynomial losses by exact constants.  The Mersenne-31 list row, with only a few bits of room, should be treated as requiring exact extremality rather than a soft analytic estimate.
\end{enumerate}

\begin{center}
\small
\begin{tabular}{|p{0.28\linewidth}|p{0.13\linewidth}|p{0.47\linewidth}|}
\hline
item & risk & failure meaning\\
\hline
band normalization, slope elimination, rung audit & low & editorial or arithmetic correction\\
\hline
prefix collision ledger and packing & low--medium & technical refinement, not frontier failure\\
\hline
heavy-fiber symmetry descent for \textup{(Q)} & medium & defect recursion may lose constants; fatal only for finite one-step\\
\hline
split-pencil / balanced-core census & high & a true primitive aperiodic split-locator mechanism may exist\\
\hline
mode-at-null / exchange-compression & open & finite adjacent theorem remains conditional\\
\hline
\end{tabular}
\end{center}

The strategic conclusion is that the Proximity Prize has a clean form.  Asymptotically, prove split-locator flatness at the entropy--subfield boundary.  Finitely, prove exact prefix-fiber extremality, base-field-normalized split-pencil bounds, and primitive shift-pair control at the four adjacent rows.  A counterexample cannot fail silently: it must either produce a super-polynomial primitive prefix fiber, refuting \textup{(Q)}, or a super-polynomial primitive split-pencil family, refuting \cref{prob:capfr1-split-pencil} and hence \(\Rone\).

\providecommand{\F}{\mathbb F}
\providecommand{\RS}{\mathrm{RS}}
\providecommand{\ceil}[1]{\left\lceil #1\right\rceil}
\providecommand{\floor}[1]{\left\lfloor #1\right\rfloor}

\section{Unified safe-side frontier package}\label{sec:capfp-prize}

After \cref{prop:capff1-identity-frontier}, the unsafe side is essentially finished: the identity-prefix floor reaches the entropy transition $\Hb(\rho+g)=\beta g$ directly, quotient scales and planting are subordinate mechanisms, and the deployed rows sit within $3\times10^{-6}$ of the ceiling $g^*(\rho,\beta)$.  The prize is now a safe-side theorem: prove that nothing beats the identity floor.  This section states that goal as a package of named conjectures, registers which of its base cases are already theorems, proves two new unconditional statements --- a rigidity/dictionary result for the prefix map, with the first nontrivial all-depth fiber bound, and the pair-descent interface for the aperiodic band --- and records why the band resists every first-order mechanism.

\subsection{The conjecture package}

\begin{remark}[Frontier prize formulation]\label{prob:capfp-F}
With $\beta=\log_2|\B|$, prove
\[
        \delta^*_C(\varepsilon^*)\ =\ 1-\rho-g^*(\rho,\beta)+o(1),
        \qquad
        g^*(\rho,\beta)=\sup\{g:\Hb(\rho+g)\ge\beta g\},
\]
for the deployed smooth multiplicative and circle line-round rows.  The finite deployed form is the prediction of the four adjacent pairs of \cref{prob:capff1-frontier}: the first safe agreement is $a_0+1$ for the unsafe agreements $a_0=1116047$, $1116046$, $1116023$, $1116022$.
\end{remark}

\begin{remark}[Boundary quotient/prefix formulation]\label{prob:capfp-Q}
For $m=(\rho+g)n$, $K=\rho n$, $w=m-K$, let
\[
        \Phi_{m,w}:\binom Dm\to\B^w,\qquad M\mapsto\bigl(e_1(M),\ldots,e_w(M)\bigr)
\]
be the prefix map \textup(equivalently the power-sum prefix, by Newton\textup).  The identity floor is the average-fiber lower bound $\max_z|\Phi_{m,w}^{-1}(z)|\ge\binom nm|\B|^{-w}$.  Prove the matching upper bound
\[
        \max_z\bigl|\Phi_{m,w}^{-1}(z)\bigr|\ \le\ e^{o(n)}\,\binom nm\,|\B|^{-w}
\]
throughout the band, with quotient-periodic structure absorbed at its own scale as in \cref{prop:capff1-composite}.  The finite form replaces $e^{o(n)}$ by an explicit constant ledger fitting under the deployed budget at $a_0+1$.
\end{remark}

\begin{remark}[Fixed-moment collision hierarchy]\label{prob:capfp-gamma}
Let $\mu_{m,w}$ be the distribution of $\Phi_{m,w}$ on a uniform $m$-subset, and for $r\ge2$ define
\[
        \Gamma_r(m,w)\ :=\ |\B|^{w(r-1)}\sum_{z}\mu_{m,w}(z)^r,
\]
so that $\Gamma_r=1$ exactly at uniformity, $\Gamma_r\ge1$ always, $\Gamma_2=\sum_\tau|\widehat\mu_{m,w}(\tau)|^2$ by Parseval, and $\Gamma_r^{1/r}$ increases in $r$ to the max-to-mean fiber ratio.  Prove $\Gamma_r(m,w)\le n^{C_r}$ for each fixed $r$ in the band, and quantify the growth of $C_r$.  The precise conversion to \cref{prob:capfp-Q} is
\[
        \frac{\max_z|\Phi^{-1}(z)|}{\text{average}}\ \le\ \inf_r\ \Gamma_r^{1/r}\,|\B|^{w/r},
\]
so polynomial max-fiber control requires the hierarchy up to $r^*\approx w\beta/\log_2n$ --- about $10^5$ at the deployed frontier depth.  A second-moment bound alone cannot close the finite problem; the hierarchy is where joint cancellation beyond the sharp per-frequency Weil cost must appear.
\end{remark}

\begin{remark}[Aperiodic formulation before slope elimination]\label{prob:capfp-A}
After removing the paid strata --- tangent, quotient, extension-pole, planted/sunflower, fixed-dimensional Conjecture-F, projective-plane, common-GCD, quotient-pullback, and bounded SPI charts --- prove that every remaining aperiodic M1/L1 configuration at band agreement $a$ contributes at most $e^{o(n)}$ times its model size:
\[
        N_{\mathrm{aper}}(a)\ \le\ e^{o(n)}\,N_{\mathrm{model}}(a),
\]
with $N_{\mathrm{model}}(a)$ below budget beyond the frontier; this is the band form of \cref{prob:band}.
\end{remark}

\subsection{Proved anchors of the package}

\begin{proposition}[base cases are theorems]\label{prop:capfp-anchors}
At the deployed rows, for every prefix depth $w\le w_0$ of \cref{cor:capff1-collision-head,cor:capff1-circle-head} \textup($w_0=21$--$22$ KoalaBear, $10$--$11$ Mersenne-31\textup) and every $2\le r\le\infty$,
\[
        1\ \le\ \Gamma_r(m,w)\ \le\ \bigl(1+2^{-\mathrm{margin}}\bigr)^{r-1},
\]
with the printed margins; the identity witness lists at these depths are exactly the fibers \textup(\cref{prop:capff1-witness-exact}\textup); composite directions are graded along the divisor lattice \textup(\cref{prop:capff1-composite}\textup); and the quotient census cells at scales $c\le16$ carry the table of \cref{cor:capff1-quotient-scales}.
\end{proposition}

\begin{proof}
Lower bound: power means.  Upper: with $R:=|\B|^w\max_z\mu(z)\le1+2^{-\mathrm{margin}}$ by the two-sided pricing, $\Gamma_r=|\B|^{w(r-1)}\sum\mu^r\le(|\B|^w\max\mu)^{r-1}\sum\mu=R^{r-1}$; the case $r=\infty$ is $R$ itself.  The remaining claims are the quoted results.
\end{proof}

\subsection{The moment kernel: rigidity, dictionary, and an all-depth bound}

The following unconditional statements convert the $L^2$ layer of \cref{prob:capfp-Q} into a coding-theoretic object and make the inverse strategy precise.  Assume $0\notin D$ \textup(true for the multiplicative rows; for circle $x$-domains, $0\in D$ only if $\pm i$ lies in the twin coset, in which case the statements hold verbatim on $D\setminus\{0\}$\textup).

\begin{proposition}[moment-kernel dictionary]\label{prop:capfp-kernel}
Let $K_w:=\{f:D\to\B:\ \sum_{x\in D}f(x)x^j=0,\ 1\le j\le w\}$.  Then:
\begin{enumerate}[label=\textup{(\alph*)}]
\item $K_w$ is a linear code of dimension $n-w$ and minimum Hamming weight at least $w+1$;
\item \textup(fiber rigidity\textup)\ distinct members $M\ne M'$ of one $\Phi_{m,w}$-fiber satisfy $|M\triangle M'|\ge w+1$; equivalently $|M\cap M'|\le m-\delta$ with $\delta:=\ceil{(w+1)/2}$;
\item \textup(collision stratification\textup)\ with $T_w(\ell)$ the number of ordered pairs of disjoint $\ell$-subsets $A,A'$ with $p_j(A)=p_j(A')$ for $j\le w$ --- equivalently the number of $\{0,\pm1\}$-valued words of $K_w$ with $\ell$ entries of each sign ---
\[
        \sum_z\bigl|\Phi_{m,w}^{-1}(z)\bigr|^2
        \ =\ \binom nm+\sum_{\ell\ge\delta}\binom{n-2\ell}{m-\ell}\,T_w(\ell);
\]
\item consequently $\Gamma_2$ is a weighted balanced-ternary weight enumerator of $K_w$, and an exponentially heavy fiber forces exponentially many balanced ternary codewords of $K_w$ concentrated on the fiber's members.
\end{enumerate}
\end{proposition}

\begin{proof}
\textup{(a)}  The defining map $V:f\mapsto(\sum f(x)x^j)_{j\le w}$ has, on any $s\le w$ coordinates $x_1,\ldots,x_s\in D$, the top minor $\det(x_i^j)_{1\le i,j\le s}=(\prod_ix_i)\cdot\prod_{i<i'}(x_{i'}-x_i)\ne0$, since the $x_i$ are distinct and nonzero; so $V$ is injective on words of support size $\le w$ and surjective, giving both claims.
\textup{(b)}  $f=\mathbf1_{M}-\mathbf1_{M'}\in K_w$ is nonzero of weight $|M\triangle M'|$; apply \textup{(a)}; the weight is even and $\ge w+1$, hence $\ge2\delta$.
\textup{(c)}  Stratify pairs by $A=M\setminus M'$, $A'=M'\setminus M$ \textup(sizes equal, say $\ell$, with $\ell\ge\delta$ or $\ell=0$ by \textup{(b)}\textup); prefix matching is $p_j(A)=p_j(A')$ for $j\le w$ since the common part cancels, i.e.\ $\mathbf1_A-\mathbf1_{A'}\in K_w$; the common part is an arbitrary $(m-\ell)$-subset of the remaining $n-2\ell$ points.
\textup{(d)}  Restate \textup{(c)}; for the last claim, a fiber of size $S$ contributes $S(S-1)$ ordered collision pairs, each producing a distinct balanced ternary codeword together with its core.
\end{proof}

\begin{corollary}[unconditional all-depth fiber bound]\label{cor:capfp-packing}
For every $m,w$ and every $z$,
\[
        \bigl|\Phi_{m,w}^{-1}(z)\bigr|\ \le\ \binom{n}{m-\delta+1}\Big/\binom{m}{\delta-1},
        \qquad \delta=\ceil{(w+1)/2}.
\]
At the four deployed frontier rows \textup($w=67466$, $67470$, $67444$, $67446$\textup) this improves the trivial bound $\binom nm$ by $213501$, $213510$, $213448$, $213453$ bits respectively --- the first nontrivial fiber bound at full frontier depth --- while remaining about $1.877\times10^6$ bits above the conjectured average: an exact measure of the distance from rigidity alone to \cref{prob:capfp-Q}.
\end{corollary}

\begin{proof}
By rigidity, two distinct fiber members share at most $m-\delta$ points, so no $(m-\delta+1)$-subset lies in two members; each member contains $\binom m{m-\delta+1}=\binom m{\delta-1}$ of the $\binom n{m-\delta+1}$ available.  The bit counts are certified by the entropy sandwich as before.
\end{proof}

\subsection{The \textup{(A)} interface: pair descent and band obstructions}

\begin{lemma}[pair descent]\label{lem:capfp-pair-descent}
Fix a line $(u,v)$ with $v\notin C$, an agreement $a>n/2$, and call a finite slope $z$ \emph{close} if $u+zv$ agrees with some codeword $c_z$ on at least $a$ points.  If $z_1\ne z_2$ are close, then $v_{12}:=(c_{z_1}-c_{z_2})/(z_1-z_2)$ is a codeword \textup(possibly zero\textup) agreeing with $v$ on at least $2a-n$ points.  Color each pair of close slopes by its codeword $v_{ij}$; then each color class is a clique, and the close slopes of a clique in color $v'$ are, slope for slope, close slopes of the derived line $(u-c_{z_0}+z_0v',\,v-v')$ at the same agreement, for any fixed clique member $z_0$.  Consequently, with $\lambda:=\#\{v'\in C:\ v'\text{ agrees with }v\text{ on }\ge2a-n\text{ points}\}$,
\[
        L(u,v;a)\ \le\ 1+\lambda\cdot\max_{v'}\,L\bigl(u-c_{z_0}+z_0v',\,v-v';a\bigr).
\]
\end{lemma}

\begin{proof}
Subtracting the two agreement relations on $T_1\cap T_2$ \textup($|T_1\cap T_2|\ge2a-n$\textup) gives $(z_1-z_2)v=c_{z_1}-c_{z_2}$ there, whence the first claim.  For transitivity, if $v_{12}=v_{23}=v'$ then $c_{z_1}-c_{z_3}=(z_1-z_2)v'+(z_2-z_3)v'=(z_1-z_3)v'$, so $v_{1,3}=v'$: color classes are cliques.  Within the clique of $v'$, $c_{z_i}=c_{z_0}+(z_i-z_0)v'$, so $u+z_iv-c_{z_i}=(u-c_{z_0}+z_0v')+z_i(v-v')$, and the agreement set is unchanged.  Every close slope other than a fixed one shares an edge, hence a color, with it; counting cliques by color gives the displayed inequality.
\end{proof}

\begin{remark}[why the band is the band]\label{rem:capfp-band-obstruction}
The inventory of first-order mechanisms against \cref{prob:capfp-A}, at rate $\rho=\tfrac12$ and the deployed sizes:
\textup{(i)}\ $r$-wise differences that eliminate both $u$ and $v$ need $r\ge3$ and bite only when $3(n-a)<n-k$, i.e.\ $a>\tfrac56n$ --- the deep regime already covered.
\textup{(ii)}\ \Cref{lem:capfp-pair-descent} is unconditional, but its hypothesis-free use needs the direction list to be empty or small at agreement $2a-n$, and $2a-n$ exceeds the direction word's own entropy frontier only when $a\ge n(1+\rho+g^*)/2\approx0.76609\,n$; below that, generic directions carry astronomically many codewords at agreement $2a-n$, and adversarial directions sit near the code outright.
\textup{(iii)}\ Subtracting the near codeword descends to a direction of weight at most $2(n-a)$, after which every agreement account passes through $|T_i|-2(n-a)\ge3a-2n$, which is \emph{negative} throughout the band --- at the KoalaBear adjacent pair, $3a-2n=-846160$ while $2a-n=134944$.
Thus every pencil-type first-order mechanism degenerates identically at the quantity $3a-2n$, and the band $a/n\in(\rho+g^*,\tfrac56)$ admits no unconditional exclusion by descent alone.  This is the precise sense in which \textup{(A)} is a named input rather than a missing lemma: the viable routes are stability classification of near-extremal locator families, growing-dimensional incidence bounds for the residual Conjecture-F strata, and bounded-deficiency SPI induction, each of which must produce structure that the paid compilers then price.
\end{remark}

\begin{remark}[measurement protocol and measured hierarchy values; experimental]\label{rem:capfp-gamma-measured}
The Fourier coefficients $\widehat\mu_{m,w}(\tau)$ of the power-sum prefix distribution are the normalized elementary symmetric coefficients of the unimodular family $\varepsilon_x=e_p(\sum_r\tau_rx^r)$, computable at any scale by the non-expansive convex recursion of \cref{rem:capff1-calibration}; the collision gaps then follow from censuses or from Parseval.  Exact census values \textup(group-algebra dynamic programming; scripts \texttt{gamma\_census.py}, \texttt{prize\_packing.py}\textup):
\[
\begin{array}{l|ccc|ccc|cc}
 & \multicolumn{3}{c|}{(p,N,m)=(17,16,8)} & \multicolumn{3}{c|}{(41,20,10)} & \multicolumn{2}{c}{(257,64,34)}\\
w & 1 & 2 & 3 & 1 & 2 & 3 & 1 & 2\\ \hline
\Gamma_2 & 1.000000 & 1.001539 & 1.137687 & 1.000001 & 1.004136 & 1.502534 & 1.000000 & 1.000000\\
\Gamma_3 & 1.000000 & 1.004679 & 1.419892 & 1.000002 & 1.012429 & 2.755825 & 1.000000 & 1.000000\\
\Gamma_4 & 1.000001 & 1.009499 & 1.914863 & 1.000003 & 1.024965 & 5.865781 & 1.000000 & 1.000000\\
\end{array}
\]
The hierarchy sits at $1+o(1)$ throughout the dense bulk and grows only at the Poisson boundary \textup(mean fiber $O(1)$\textup), matching \cref{rem:capff1-calibration}.  The falsifiable law is $\log_2\Gamma_r(m,w)=O_r(\log n)$ across the band, against the alternative $\Theta(n)$; by \cref{prob:capfp-gamma} the finite prize additionally needs the hierarchy to persist to $r\approx10^5$.  These values are calibration evidence only.
\end{remark}

\subsection{Slope elimination: reducing the aperiodic side to a support census}

The following results turn the first-form aperiodic input into an explicit support census.  They eliminate the slope quantifier unconditionally, remove the tangent stratum from the inequality outright, and subsume the counting content of every bounded- or growing-deficiency SPI chart.  The census is then refined further by the split-pencil and base-field-normalized formulations.

\begin{lemma}[interpolation functionals]\label{lem:capfp-functionals}
Let $T\subseteq D$, $|T|=m$, $w=m-K$, and for a word $f:D\to\F$ put
\[
        S_r(f,T)\ :=\ \sum_{x\in T}\frac{f(x)\,x^r}{\Lambda_T'(x)},\qquad 0\le r\le w-1 .
\]
Then $f|_T$ extends to a polynomial of degree $<K$ if and only if $S_r(f,T)=0$ for all $r<w$.  The functionals are $\F$-linear in $f$, and $\B$-valued whenever $f$ is $\B$-valued \textup(since $T\subseteq D\subseteq\B$\textup).
\end{lemma}

\begin{proof}
The interpolant is $I_f(X)=\sum_{x\in T}f(x)\Lambda_T(X)/((X-x)\Lambda_T'(x))$, of degree $\le m-1$, and extension of degree $<K$ means its top $w$ coefficients vanish.  From $\prod_{y\in T\setminus x}(1+uy)=\bigl(\sum_ie_i(T)u^i\bigr)\bigl(\sum_r(-x)^ru^r\bigr)$ one gets $e_j(T\setminus x)=\sum_{r\le j}(-1)^re_{j-r}(T)\,x^r$, so the coefficient of $X^{m-1-j}$ in $I_f$ is $(-1)^j\sum_{r\le j}(-1)^re_{j-r}(T)\,S_r(f,T)$: a unitriangular combination of $S_0,\ldots,S_j$ with $T$-symmetric coefficients.  Vanishing for $j=0,\ldots,w-1$ is therefore equivalent to $S_0=\cdots=S_{w-1}=0$.
\end{proof}

\begin{theorem}[slope elimination]\label{thm:capfp-slope-elim}
Fix a line $(u,v)$ over the challenge field, an agreement $a=m>n/2$, and $w=m-K$.  Call $T$ \emph{rank-one} if the vectors $\alpha(T):=(S_r(u,T))_{r<w}$ and $\beta(T):=(S_r(v,T))_{r<w}$ are proportional and not both zero, with $\beta(T)\ne0$; each rank-one support determines the single slope $z(T)$ with $\alpha(T)=-z(T)\beta(T)$.  Then:
\begin{enumerate}[label=\textup{(\alph*)}]
\item a slope $z$ is close at support $T$ \textup(i.e.\ $u+zv$ interpolable on $T$\textup) if and only if $\alpha(T)=-z\,\beta(T)$; supports with $\beta=0\ne\alpha$ carry no slope, and supports with $\alpha=\beta=0$ \textup(\emph{common} supports\textup) carry every slope;
\item if some $m$-support is common, the pair $(u,v)$ has correlated agreement at density $m/n$ on it, and \emph{no} slope is MCA-bad at radius $1-m/n$;
\item in every case,
\[
        N_{\mathrm{MCA\text{-}bad}}(u,v;\,m)\ \le\ \#\mathrm{R1}(u,v;m)\ :=\ \#\{\text{rank-one supports}\};
\]
\item \textup(pencil dictionary\textup)\ for any exact-agreement syndrome pencil $M(Z)=M_0+ZM_1$ of \emph{any} deficiency, parametrizing split kernel candidates by their root sets $A$: per candidate, $M_0\Lambda_A+zM_1\Lambda_A=0$ has at most one solution $z$ unless $M_0\Lambda_A=M_1\Lambda_A=0$ \textup(an identically-split residual branch as in \cref{thm:capf-spi}\textup); so the chart's slope count is bounded by the same rank-one census with columns $(M_0\Lambda_A,M_1\Lambda_A)$, at every deficiency.
\end{enumerate}
\end{theorem}

\begin{proof}
\textup{(a)}  By \cref{lem:capfp-functionals} applied to $u+zv$ and $\F$-linearity, closeness at $T$ is $\alpha(T)+z\beta(T)=0$, and the trichotomy is immediate.
\textup{(b)}  A common support gives polynomials $c_u,c_v$ of degree $<K$ with $u|_T=c_u|_T$, $v|_T=c_v|_T$; then $(u+zv)|_T=(c_u+zc_v)|_T$ for every $z$: correlated agreement of the pair at density $m/n$, which by definition excludes MCA-badness at that radius.
\textup{(c)}  If a common support exists, the left side is zero by \textup{(b)}.  Otherwise every MCA-bad slope is close, every close slope has an $m$-support inside its agreement set, no support is common, supports with $\beta=0$ carry no slope by \textup{(a)}, and distinct slopes have distinct rank-one supports since $z(T)$ is a function of $T$.
\textup{(d)}  The pencil is affine-linear in $Z$, so per split candidate the condition is $\mu_0+z\mu_1=0$ with $\mu_i=M_i\Lambda_A$: one slope, or none, or all \textup(the residual branch\textup).
\end{proof}

\begin{remark}[Rank-one support census after slope elimination]\label{prob:capfp-R1}
Prove, for every line at band agreements,
\[
        \#\mathrm{R1}(u,v;m)\ \le\ e^{o(n)}\cdot\max\Bigl(1,\ \binom nm\,q^{-(w-1)}\Bigr),
\]
with the finite form fitting under the deployed budget at $a_0+1$.  By \cref{thm:capfp-slope-elim} and the paid compilers, this rank-one census implies the first-form aperiodic input.  It is an intermediate algebraic reduction; the final active split-pencil input is \cref{prob:capg-active-BC}.  The model values at the conjectured first safe agreements are $2^{-10453222}$, $2^{-10453966}$, $2^{-6272178}$, $2^{-6272426}$ per line against per-line budgets of $2^{+58}$, $2^{+58}$, $2^{-4}$, $2^{-4}$: adversarial exclusion windows of $10.45$, $10.45$, $6.27$, $6.27$ million bits \textup(script \texttt{r1\_model.py}\textup).  Structure available to attacks: the columns are $\B$-valued on $\B$-valued coordinate words \textup(\cref{lem:capfp-functionals}\textup), so \textup{(R1)} is a base-field point count on an explicit determinantal incidence variety --- the Conjecture-F setting of \cref{sec:capf-m1} --- and the twisted vectors $(u(x)/\Lambda_T'(x))\mathbf1_T$, $(v(x)/\Lambda_T'(x))\mathbf1_T$ live in the order-$w$ moment kernel of \cref{prop:capfp-kernel}, tying \textup{(R1)} and \textup{(Q)} to the same object.  \Cref{lem:capfp-pair-descent} bounds the clique structure of the image of $z(\cdot)$; \textup{(R1)} bounds its domain.  Known unconditional strata: $\B$-rational lines \textup(confinement: at most $p+1$ slopes\textup), and deficiency-one pencils, where \cref{thm:capf-spi} gives a polynomial eliminant with no census needed.
\end{remark}

\subsection{The lattice resolution of the census: near-rational dichotomy and the balanced core}

The rank-one census, and with it input \textup{(A)}, now admits a second unconditional reduction: the twisted census of \emph{every} word is the split locus of an explicit rank-two $\F[X]$-lattice, its near-rational stratum resolves exactly, and per line at most one slope survives outside the balanced core.  Throughout, $\omega:=n-m$, $\widehat U$ denotes the degree-$\le n-1$ interpolant of the word $U$ on $D$, and we assume the band condition $2w\le n-K$.


\begin{proposition}[the census as a shifted lattice locus]\label{prop:capfp-lattice}
Let
\[
        M_U:=\{(W,N)\in\F[X]^2:\ W(x)U(x)=N(x)\ \textup{for all }x\in D\},
\]
and put
\[
        \operatorname{wdeg}(W,N):=\max\bigl(\deg W,\deg N-(K-1)\bigr).
\]
Then:
\begin{enumerate}[label=\textup{(\alph*)}]
\item $M_U$ is a free rank-two $\F[X]$-module with basis
$(1,\widehat U)$ and $(0,\Lambda_D)$.  A basis in shifted weak Popov
form has predictable degrees, and its row degrees $d_1\le d_2$ satisfy
$d_1+d_2=n-K+1$.
\item If $\omega=n-m$, then the agreement-support census is exactly
\[
\begin{aligned}
\mathrm{cen}(U;m)
&:=\#\{T\subseteq D:\ |T|=m,\ U|_T\in \RS[\F,T,K]\}\\
&=\#\Bigl\{(W,N)\in M_U:\
        W\textup{ monic},\ \deg W=\omega,\
        \operatorname{wdeg}(W,N)\le\omega,\
        W\mid\Lambda_D,\ W\mid N
        \Bigr\}.
\end{aligned}
\]
\item Equivalently,
\[
        \mathrm{cen}(U;m)=\sum_{c\in \RS[\F,D,K]}\binom{\operatorname{agr}(U,c)}m.
\]
\end{enumerate}
The shifted-degree cap in \textup{(b)} is part of the census: without it one
counts high-degree congruent representatives rather than degree-$<K$ codeword
explanations.
\end{proposition}

\begin{proof}
For \textup{(a)}, any $(W,N)\in M_U$ satisfies
$N\equiv W\widehat U\pmod{\Lambda_D}$, hence
$N=W\widehat U-C\Lambda_D$ for some $C\in\F[X]$.  Thus
$(W,N)=W(1,\widehat U)-C(0,\Lambda_D)$, and the displayed pair is a basis.
Its determinant is $\Lambda_D$, of shifted degree $n-(K-1)$, so a shifted weak
Popov basis has predictable degrees and row-degree sum $n-K+1$.

If $T$ is a valid support and $c\in\RS[\F,D,K]$ is the unique codeword agreeing
with $U$ on $T$ \textup{(uniqueness follows from $m\ge K$)}, put
$W=\Lambda_{D\setminus T}$ and $N=Wc$.  Then $(W,N)\in M_U$, $W$ is monic of
degree $\omega$, $W\mid\Lambda_D$, $W\mid N$, and
$\operatorname{wdeg}(W,N)\le\omega$ because $\deg c<K$.

Conversely, suppose $(W,N)$ satisfies the conditions in \textup{(b)}.  Since
$W\mid N$, write $N=Wc$.  The shifted-degree cap gives
$\deg N\le\omega+K-1$, hence $\deg c<K$.  The identity defining $M_U$ gives
$W(U-c)=0$ on $D$, so $U=c$ on $T=D\setminus Z(W)$, and $|T|=m$.  This is a
valid agreement support.  The two constructions are inverse, and the
binomial-moment identity follows by grouping the valid supports by the
codeword that explains them.
\end{proof}


\begin{theorem}[near-rational dichotomy]\label{thm:capfp-dichotomy}
Let $d_1=d_1(U)$ be the minimal shifted degree of $M_U\setminus\{0\}$.
\begin{enumerate}[label=\textup{(\roman*)}]
\item If $d_1\le w$, then either $U$ lies within Hamming distance exactly $d_1$ of a \emph{unique} codeword $c$, the minimal vector is the error-locator pair $(\Lambda_E,\Lambda_E\hat c)$ with $E=\mathrm{Disagr}(U,c)$, and
\[
        \mathrm{cen}(U;m)\ =\ \binom{\,n-d_1\,}{m}\qquad\text{exactly};
\]
or $\mathrm{cen}(U;m)=0$.
\item If $d_1\ge w+1$, every census element lies in the two-parameter family $\{Ag_1+Bg_2:\ \deg A\le\omega-d_1,\ \deg B\le\omega-d_2\}$, of $\F$-dimension exactly $\omega-w+1$.
\end{enumerate}
\end{theorem}

\begin{proof}
If $d_1\le w$ then $d_2=n-K+1-d_1\ge n-K+1-w=\omega+1$, so by predictable degrees every census element \textup(shifted degree $\omega$\textup) is $A\,g_1$ with $g_1=(W_1,N_1)$ of shifted degree $d_1$.  Then $W_1\mid W\mid\Lambda_D$, so $\mathrm{cen}=0$ unless $W_1$ is \textup(up to scalar\textup) monic squarefree and $D$-split; and $W\mid N\iff W_1\mid N_1$ by cancellation, so $\mathrm{cen}=0$ unless $c:=N_1/W_1$ is a polynomial.  If $\deg N_1-(K-1)>\deg W_1$, predictable degrees give $\operatorname{wdeg}(Ag_1)=\deg A+d_1$ with $\deg W=\deg A+\deg W_1<\omega$ at the cap, so no element reaches $\deg W=\omega$: $\mathrm{cen}=0$.  Otherwise $\deg W_1=d_1$, $\deg c\le K-1$: a codeword, agreeing with $U$ off the $d_1$ roots of $W_1$.  Minimality forces $\operatorname{wt}(U-c)=d_1$ with disagreement set exactly $\mathrm{roots}(W_1)$: a smaller disagreement set would make the true error-locator pair a nonzero lattice element of smaller shifted degree.  Uniqueness of $c$: two codewords within distance $w$ of $U$ differ in at most $2w\le n-K$ positions, below the code distance $n-K+1$.  Finally, valid supports are exactly the $m$-subsets of the agreement set: containment follows from $W_1\mid W$, and every $m$-subset $T\subseteq D\setminus E$ arises from $A=\Lambda_{(D\setminus E)\setminus T}$, which satisfies all constraints.  Part \textup{(ii)} is predictable degrees with both caps active, of total dimension $2\omega+2-(d_1+d_2)=\omega-w+1$.
\end{proof}

\begin{corollary}[per-line reduction to the balanced core]\label{cor:capfp-line}
For every line $(u,v)$ at band agreement $m$: if two finite slopes have $d_1(u+zv)\le w$ with nonzero census, then $v$ is within distance $2w$ and $u$ within distance $3w$ of codewords, the pair has correlated agreement at density $\ge1-3w/n>m/n$, and \emph{no} slope is MCA-bad.  Hence in every case
\[
        N_{\mathrm{MCA\text{-}bad}}(u,v;m)\ \le\ 1+\#\bigl\{z:\ d_1(u+zv)\ge w+1,\ \mathrm{cen}(u+zv;m)>0\bigr\}:
\]
input \textup{(A)} reduces, unconditionally, to the balanced core along pencils.
\end{corollary}

\begin{proof}
By \cref{thm:capfp-dichotomy}\textup{(i)}, a slope with $d_1\le w$ and nonzero census gives $u+z_iv=c_i+\eta_i$ with $\operatorname{wt}\eta_i\le w$.  Two such: $v=(c_1-c_2)/(z_1-z_2)+\eta_v$ with $\operatorname{wt}\eta_v\le2w$, then $u=c_1-z_1(c_1-c_2)/(z_1-z_2)+\eta_u$ with $\operatorname{wt}\eta_u\le3w$: the pair agrees with a codeword pair jointly off at most $3w$ points, giving correlated agreement at density $1-3w/n\ (\approx0.9035$ at the deployed pairs$)$, far above $m/n\approx0.5322$, so no slope is MCA-bad.  Otherwise at most one such slope exists, slopes with zero census are not close at agreement $m$, and the rest is \cref{thm:capfp-dichotomy}\textup{(ii)}.
\end{proof}

\begin{remark}[balanced core, primitive challenge-field form]\label{prob:capfp-balanced}
Prove, for every word $U$ with $d_1(U)\ge w+1$ at band agreements,
\[
        \#\bigl\{(A,B):\ \deg A\le\omega-d_1,\ \deg B\le\omega-d_2,\ W_{A,B}\mid\Lambda_D,\ W_{A,B}\mid N_{A,B}\bigr\}
        \ \le\ e^{o(n)}\max\Bigl(1,\ \binom nm q^{-w}\Bigr),
\]
where $(W_{A,B},N_{A,B})=Ag_1+Bg_2$.  This single statement implies the twisted census bound for every word, hence \cref{prob:capfp-R1}, hence \textup{(A)}; with \textup{(Q)} it closes the asymptotic frontier.  The family is explicit, of dimension $\omega-w+1$ \textup($913642$ at the KoalaBear adjacent pair, $913686$ at Mersenne-31\textup); the split condition $W\mid\Lambda_D$ is the same divisor variety as in \cref{thm:capf-spi}, now with a canonical global origin: the basis $(g_1,g_2)$ arises from one shifted-Popov reduction of $\bigl((1,\widehat U),(0,\Lambda_D)\bigr)$, and along a pencil $u+zv$ it varies algebraically in $z$, giving every SPI chart a common parent.  The common-GCD and sunflower paid strata appear inside the coordinates as $\gcd(A,B)$ and common-factor loci.  The dichotomy and the exact unbalanced census were verified on $200$ random small-field instances \textup(script \texttt{lattice\_check.py}\textup).
\end{remark}


\begin{remark}[status of the intermediate balanced-core form]
\label{rem:capfp-balanced-active-reading}
\Cref{prob:capfp-balanced} is the challenge-field primitive model suggested by
the shifted-lattice dichotomy.  It is not the all-strata final census.  The
boundary profile belongs to the quotient/prefix problem, and the interior
profiles must include the base-field floor terms isolated in
\cref{prop:capg-census-floor}.  The final active version is therefore
\cref{prob:capg-active-BC}, equivalently the base-field-normalized
split-pencil census of \cref{prob:capg-split-pencil-B}.
\end{remark}

\subsection{Split-pencil form of the balanced core}

Three further unconditional statements put the balanced core into its final shape.  All were verified exhaustively on $120$ random small-field instances \textup(script \texttt{lattice\_check2.py}\textup).

\begin{lemma}[divisibility is automatic]\label{lem:capfp-autodiv}
If $(W,N)\in M_U$ and $W$ is squarefree with all roots in $D$, then $W\mid N$.
\end{lemma}

\begin{proof}
Every root $x$ of $W$ lies in $D$, where $N(x)=W(x)U(x)=0$; as $W$ is squarefree, $W\mid N$.
\end{proof}

\begin{lemma}[unimodularity]\label{lem:capfp-unimodular}
For any basis $(W_1,N_1),(W_2,N_2)$ of $M_U$: $\gcd(W_1,W_2)=1$ and $W_1N_2-W_2N_1=\gamma\Lambda_D$ with $\gamma\in\F^\times$.
\end{lemma}

\begin{proof}
Expanding $(1,\widehat U)=A_0g_1+B_0g_2$ in the first coordinate gives $1=A_0W_1+B_0W_2$.  The cross determinant vanishes at every $x\in D$ \textup(both products equal $W_1(x)W_2(x)U(x)$\textup), so $\Lambda_D$ divides it; its degree is at most $n$ and it is the basis determinant, a unit multiple of $\Lambda_D$.
\end{proof}

\begin{proposition}[the adversary is a determinantal representation]\label{prop:capfp-detrep}
Quadruples $(W_1,N_1,W_2,N_2)$ arising as shifted-reduced bases of the lattices $M_U$ are exactly the solutions of
\[
        W_1N_2-W_2N_1=\gamma\,\Lambda_D,\qquad\gamma\in\F^\times,\quad\gcd(W_1,W_2)=1,
\]
with the shifted-degree profile $(d_1,d_2)$, $d_1+d_2=n-K+1$: the word is recovered pointwise by $U(x)=N_i(x)/W_i(x)$ for either $i$ with $W_i(x)\ne0$ \textup(some $i$ works at every $x$ by coprimality, and the two ratios agree by the identity\textup).  Consequently the entire aperiodic input is a statement about degree-capped determinantal representations of $\Lambda_D$: for the KoalaBear rows, of $X^n-\beta$.
\end{proposition}

\begin{proof}
The forward direction is \cref{lem:capfp-unimodular}.  Conversely, given such a quadruple, define $U$ pointwise as stated; at a root $x$ of $W_1$, the identity gives $W_2(x)N_1(x)=0$ with $W_2(x)\ne0$, so $N_1(x)=0$ and $(W_1,N_1)\in M_U$, likewise $(W_2,N_2)$.  The submodule they generate has full rank and determinant $\gamma\Lambda_D$, equal to that of $M_U$, hence equals $M_U$, and the degree profile makes the basis reduced.
\end{proof}

\begin{remark}[Primitive split-pencil formulation]\label{prob:capfp-split}
Let $(W_1,N_1,W_2,N_2)$ be a determinantal representation of $\Lambda_D=X^n-\beta$ as in \cref{prop:capfp-detrep} with balanced profile $d_1\ge w+1$.  Prove
\[
        \#\bigl\{(A,B):\ \deg A\le\omega-d_1,\ \deg B\le\omega-d_2,\ AW_1+BW_2\ \bigl|\ X^n-\beta\bigr\}
        \ \le\ e^{o(n)}\max\Bigl(1,\ \binom n\omega\,q^{1-w}\Bigr),
\]
counting monic normalizations of exact degree $\omega$.  By \cref{lem:capfp-autodiv} this is the full census condition --- the divisibility $W\mid N$ has been eliminated --- and the split condition is the single congruence $X^n\equiv\beta\pmod{AW_1+BW_2}$.  This intermediate problem implies \cref{prob:capfp-balanced}, hence \cref{prob:capfp-R1}, hence the first-form aperiodic input; the active final problem is the base-field-normalized version of \cref{prob:capg-active-BC}.  The circle rows replace $X^n-\beta$ by the twin-coset domain polynomial.  Internal structure, fully proved: the $\det$-degenerate directions $AB'-A'B=0$ organize the solutions into \emph{rays} $r\cdot(A_0,B_0)$, each ray corresponding to a single list codeword $c$ with census contribution exactly $\binom{\mathrm{agr}(U,c)}{m}$, and distinct primitive members obey $\gcd(f_0,f_0')\mid A_0B_0'-A_0'B_0$, which is precisely the classical bound of $K-1$ on pairwise codeword agreement in disguise.  The split-pencil problem is therefore exactly the $m$-th binomial moment of the agreement profile, carried by an explicit two-generator pencil --- the form in which partitioning, stability, or monodromy methods can now be aimed at a single algebraic family.
\end{remark}

\subsection{The \textup{(Q)} stratum: divisor form and unification of the package}

Applying the lattice machinery to the witness words themselves identifies \textup{(Q)} as the boundary stratum of the split-pencil formulation: the package rests on \emph{one} problem.  Assume throughout the band profile condition $2m\le n+K-1$, i.e.\ $w+1\le\omega$ \textup(deployed: $2m\approx2.23\times10^6\le3.15\times10^6$\textup).

\begin{theorem}[\textup{(Q)} is the extreme profile of the split pencil]\label{thm:capfp-Q-unify}
Fix $z\in\B^w$ and the witness word $U_z$, so $\widehat{U_z}=P_z:=X^m+\sum_{j\le w}z_jX^{m-j}$, and write $X^n=P_zQ_z+R_z$ with $\deg R_z<m$.  Then $(1,P_z)\in M_{U_z}$ has shifted degree $w+1$, and $g_2:=(Q_z,\ \beta-R_z)\in M_{U_z}$ has shifted degree $\omega$.  Exactly one of:
\begin{enumerate}[label=\textup{(\roman*)}]
\item $d_1(U_z)\le w$, and the fiber is the exactly classified branch of \cref{thm:capfp-dichotomy}: $\binom{n-d_1}m$ for a unique near codeword, or empty;
\item $d_1(U_z)=w+1$, the pair $\bigl((1,P_z),g_2\bigr)$ is a reduced basis of profile $(w+1,\omega)$, and
\[
        \bigl|\Phi_{m,w}^{-1}(z)\bigr|\ =\ \#\bigl\{A:\ \deg A\le\omega-w-1,\ \ Q_z+A\ \bigm|\ X^n-\beta\bigr\}.
\]
\end{enumerate}
In particular \textup{(Q)} is exactly the prescribed-top-coefficient divisor problem for $X^n-\beta$ --- count monic degree-$\omega$ divisors whose top $w+1$ coefficients equal those of the truncated quotient $Q_z$ --- and coincides with the split-pencil formulation \cref{prob:capfp-split} at the boundary profile $d_1=w+1$.  Inputs \textup{(Q)} and \textup{(A)} are therefore strata of one problem: \textup{(Q)} the extreme profile, the aperiodic balanced core the interior profiles.
\end{theorem}

\begin{proof}
$\deg P_z=m<n$ gives $\widehat{U_z}=P_z$ and $\operatorname{wdeg}(1,P_z)=m-(K-1)=w+1$.  For $g_2$: at $x\in D$, $Q_z(x)P_z(x)=x^n-R_z(x)=\beta-R_z(x)$, so $g_2\in M_{U_z}$; its shifted degree is $\max(\omega,\deg R_z-K+1)\le\max(\omega,w)=\omega$.  The determinant of the pair is $(\beta-R_z)-Q_zP_z=\beta-X^n=-\Lambda_D$, a unit multiple, so the pair is a basis; the shifted degrees sum to $(w+1)+\omega=n-K+1$, the first row is $N$-led \textup(as $w+1>0$\textup) and the second $W$-led \textup(as $\deg R_z-K+1\le w<\omega$\textup), so the basis is reduced with profile $(w+1,\omega)$, provided $d_1=w+1$; since $(1,P_z)$ shows $d_1\le w+1$, the only alternative is $d_1\le w$, which is branch \textup{(i)} by \cref{thm:capfp-dichotomy}.  In branch \textup{(ii)}, census elements have shifted degree $\omega$ and $W$-degree exactly $\omega$: writing them as $A(1,P_z)+Bg_2$, predictable degrees force $B$ scalar and nonzero, normalized to $1$ by monicity, and $\deg A\le\omega-(w+1)$; so $W=Q_z+A$, whose top $w+1$ coefficients are those of $Q_z$, and by \cref{lem:capfp-autodiv} the split condition $W\mid X^n-\beta$ is the entire census condition.  Conversely every monic degree-$\omega$ divisor with those top coefficients is $Q_z+A$.  Finally, the prefix of $\Lambda_T$ and the prefix of $W=\Lambda_{D\setminus T}$ determine each other unitriangularly through $\Lambda_TW=X^n-\beta$, so prescribing $z$ is prescribing the top of $W$.
\end{proof}

\begin{lemma}[orbit constancy]\label{lem:capfp-orbit}
For $\zeta\in\mu_n$, the twisted action $z_j\mapsto\zeta^jz_j$ preserves fiber sizes: $T\mapsto\zeta^{-1}T$ is a bijection of fibers, so all witness lists are constant along the $p^w/\!\sim$ orbits.
\end{lemma}

\begin{proof}
$\Lambda_{\zeta^{-1}T}(X)=\zeta^{-m}\Lambda_T(\zeta X)$ has $j$-th prefix coordinate $\zeta^{j}$ times that of $\Lambda_T$, and $\zeta^{-1}D=D$.
\end{proof}

\begin{remark}[equivalent forms, and a hardness calibration]\label{rem:capfp-Q-forms}
Three further faces of the same object, the first two proved.  \textup{(1)}  \emph{Additive form:} for $w<p$ the truncated logarithm is a group isomorphism $(1+u\F_p[u]/u^{w+1})^\times\cong(\F_p^w,+)$, under which the prefix of a divisor becomes the sum $\sum_{\xi\in S}v_\xi$ of the moment-curve vectors $v_\xi=(\xi,\xi^2/2,\ldots,\xi^w/w)$ over the root set: \textup{(Q)} is balanced subset-sum equidistribution on a moment curve, at near-full entropy \textup($\binom nm\approx p^w\cdot2^{192}$\textup).  \textup{(2)}  \emph{Fourier accounting:} Parseval pins typical character sums at $\binom Nm p^{-w/2}$, so fiber control requires square-root cancellation \emph{again} in the signed sum over the $p^w$ frequencies --- two independent layers, which is what near-full-entropy equidistribution means quantitatively.  \textup{(3)}  \emph{Calibration:} \textup{(Q)} is the function-field analogue of equidistribution of the divisors of a fixed integer in residue classes; the proved head depths $w\le w_0$ correspond to the classical regime of polylogarithmic moduli, while the band depth $w\approx67466$ corresponds to moduli of size $(\#\text{divisors})^{1-o(1)}$, open even over $\mathbf Z$.  The identity of \cref{thm:capfp-Q-unify} was verified on $140$ random prefixes across two parameter sets, and orbit constancy on all $84$ orbits \textup(script \texttt{q\_unify\_check.py}\textup).
\end{remark}

\subsection{Fixed moments and the monodromy wall}

Two final results delimit \textup{(Q)} exactly: the asymptotic form reduces to \emph{fixed}-moment collision bounds, and the strongest existing analytic framework is shown, by a proved expansion and exact arithmetic, to stop precisely one cohomological layer short.

\begin{proposition}[moment sandwich]\label{prop:capfp-Q-sandwich}
For all $r\ge2$, with $R:=|\B|^w\max_z\mu_{m,w}(z)$ the max-to-mean fiber ratio and $\beta=\log_2|\B|$,
\[
        \frac{\log_2\Gamma_r}{\,r-1\,}\ \le\ \log_2R\ \le\ \frac{\log_2\Gamma_r+w\beta}{r}.
\]
Consequently, since $w\beta=O(n)$ at the deployed band depths: the asymptotic census bound $R\le e^{o(n)}$ holds \emph{if and only if} $\Gamma_r(m,w)\le e^{o(n)}$ for every fixed $r$; while deciding a printed adjacent pair with margin $\Delta$ bits requires a row-dependent finite hierarchy of order about $w\beta/\Delta$.
\end{proposition}

\begin{proof}
Left: $\Gamma_r\le R^{r-1}$ as in \cref{prop:capfp-anchors}.  Right: $\max_z\mu\le(\sum_z\mu^r)^{1/r}=(\Gamma_r|\B|^{-w(r-1)})^{1/r}$, so $R\le(\Gamma_r|\B|^{w})^{1/r}$.  For the equivalence, take $r=\lceil1/\varepsilon\rceil$ in the right inequality and use the left for the converse.
\end{proof}

\begin{remark}[finite moment orders at the deployed rows]\label{rem:capfp-finite-orders}
At the active adjacent rows of \cref{cor:capg-adjacent-pairs}, the bit scale
$w\beta$ is about $2.09\cdot10^6$ bits.  Thus the finite moment depth needed
for a direct moment-to-max conversion depends on the row margin:
\[
\begin{array}{c|c|c|c}
\text{row} & w & \Delta\text{ bits} & w\beta/\Delta \\
\hline
\text{KoalaBear MCA} & 67470 & 22.2 & 9.4\cdot10^4 \\
\text{KoalaBear list} & 67470 & 22.0 & 9.5\cdot10^4 \\
\text{Mersenne-31 MCA} & 67446 & 3.3 & 6.3\cdot10^5 \\
\text{Mersenne-31 list} & 67446 & 3.1 & 6.7\cdot10^5 .
\end{array}
\]
Thus the finite adjacent rows require either direct constant-factor fiber
control or a finite hierarchy reaching the displayed order; a uniform slogan
such as ``order $10^5$'' is not canonical for the Mersenne-31 rows.
\end{remark}

Thus the asymptotic Frontier Conjecture now rests on \emph{fixed-order} collision counts --- for each constant $r$, the number of $r$-tuples of degree-$\omega$ divisors of $X^n-\beta$ with pairwise equal top-$(w{+}1)$ coefficients must not exceed $e^{o(n)}\binom nm^{\,r}|\B|^{-w(r-1)}$ --- together with the split-pencil bound at $e^{o(n)}$ for the interior profiles.  Fixed moments are exact algebraic objects \textup(via \cref{prop:capfp-kernel}, weighted counts of $\{0,\pm1\}$-configurations in the moment kernel\textup), the most attackable known form of the input.

\begin{proposition}[twisted-smooth expansion]\label{prop:capfp-Q-smooth}
Let $\Phi:\F_p^\times\to\{0,1\}$ be the indicator of $D$, a combination of the $(p-1)/n$ multiplicative characters of the quotient $\F_p^\times/(\text{$n$-th powers})\cdot$--- for the deployed KoalaBear rows, $1016$ characters.  Then for every monic $W$,
\[
        \mathbf1\{W\mid X^n-\beta\}\ =\ \mathbf1\{W\ \text{splits squarefree over }\F_p\}\cdot\prod_{W(\xi)=0}\Phi(\xi),
\]
so by \cref{thm:capfp-Q-unify} each fiber of \textup{(Q)} is a short-interval count of character-twisted completely-split polynomials: expanding the product over the characters writes it as a sum, over tuples from a $1016$-element set, of twisted factorization functions on the interval $\{Q_z+A:\deg A\le\omega-w-1\}$ --- term by term, exactly the objects of the Katz-monodromy short-interval theory of factorization statistics \textup(Keating--Rudnick; Sawin; Sawin--Shusterman\textup).  Verified on $4000$ random instances \textup(script \texttt{q\_smooth\_check.py}\textup).
\end{proposition}

\begin{remark}[the exact wall, located]\label{rem:capfp-Q-wall}
The expansion has $\approx1016^{\,\omega}$ tuples, $\log_2\approx9{,}799{,}978$ bits, and the maximal conceivable per-tuple estimate --- full square-root cancellation over the interval --- saves $\tfrac{h+1}2\log_2p\approx14{,}156{,}266$ bits against an interval of $28.3$ million bits.  The triangle inequality over tuples therefore lands $\approx23{,}956{,}053$ bits above the $192$-bit target: no per-term theorem, however deep, can finish \textup{(Q)}.  What remains is cancellation \emph{among} the twisted sums --- a statement about the cohomology of the whole character-tuple family over the prefix space, not of its fibers.  The theory is otherwise tight at both ends: at the deep end $\omega\lesssim50$ the trivial interval bound already meets the target \textup(the regime of \cref{thm:deep-mca}\textup), and at the head $w\le w_0$ the character route closes it \textup(\cref{cor:capff1-collision-head}\textup); the band is exactly the stretch where neither end reaches.
\end{remark}

\subsection{Conditional closing statement}

\begin{proposition}[what closes what]\label{prop:capfp-closing}
The asymptotic forms of \cref{prob:capfp-Q,prob:capfp-A} together imply the asymptotic statement of \cref{prob:capfp-F}, by the prefix ledger and the entropy--subfield separation of \cref{prop:capff1-closing}.  The finite forms, with all constants displayed and fitting inside the printed adjacent margins of $22.2$, $22.0$, $3.3$, $3.1$ bits, imply the four adjacent pairs via the one-step staircase compiler \textup(\cref{prop:onestep,thm:capf-staircase}\textup); a polynomial-loss version of either input does not decide any printed pair.  Conversely, failure of the adjacent-pair prediction locates either an exponentially heavy primitive fiber \textup(refuting \cref{prob:capfp-Q}, and by \cref{prop:capfp-kernel} producing exponentially many balanced ternary codewords of the moment kernel\textup) or an unpaid aperiodic configuration \textup(refuting \cref{prob:capfp-A}\textup).
\end{proposition}

\begin{proof}
The first two claims are \cref{prop:capff1-closing} applied per row.  For the converse, an excess of listed codewords at $a_0+1$ beyond the budget must, after the paid compilers, be carried either by a quotient/prefix stratum --- which by \cref{prop:capff1-witness-exact} and the ledger is a fiber excess, converted to kernel codewords by \cref{prop:capfp-kernel}\textup{(d)} --- or by an aperiodic residual configuration.  By \cref{thm:capfp-slope-elim}, input \textup{(A)} may be replaced throughout by the rank-one census \cref{prob:capfp-R1}, by \cref{thm:capfp-dichotomy,cor:capfp-line} further by the balanced core \cref{prob:capfp-balanced}, and by \cref{lem:capfp-autodiv,prop:capfp-detrep} finally by the split-pencil formulation \cref{prob:capfp-split}; and by \cref{thm:capfp-Q-unify}, input \textup{(Q)} is the boundary profile $d_1=w+1$ of the same problem.  At this intermediate stage the package rests on a split-pencil problem over the profile range $w+1\le d_1\le\floor{(n-K+1)/2}$, with the proved strata: head depths $w\le w_0$ at the boundary profile, the near-rational dichotomy, subfield confinement, and deficiency-one eliminants.  The final formulation separates the boundary profile into quotient/prefix flatness and inserts the base-field floor terms in the interior split-pencil model.  By \cref{prop:capfp-Q-sandwich}, the asymptotic form of the whole package further reduces to fixed-order collision bounds $\Gamma_r=e^{o(n)}$ together with the interior split-pencil bound at $e^{o(n)}$; the printed adjacent pairs additionally require the row-dependent hierarchy of \cref{rem:capfp-finite-orders} or direct constant-factor census control.
\end{proof}

\providecommand{\F}{\mathbb F}
\providecommand{\B}{\mathbb B}
\providecommand{\RS}{\mathrm{RS}}
\providecommand{\Hb}{H_{\mathrm b}}
\providecommand{\Lst}{\mathrm{Lst}}
\providecommand{\Cen}{\operatorname{cen}}
\providecommand{\Rone}{\mathrm{R1}}
\providecommand{\wdeg}{\operatorname{wdeg}}
\providecommand{\emca}{\varepsilon_{\mathrm{mca}}}
\providecommand{\eca}{\varepsilon_{\mathrm{ca}}}
\providecommand{\ceil}[1]{\left\lceil #1\right\rceil}
\providecommand{\floor}[1]{\left\lfloor #1\right\rfloor}

\section{Flexible-budget conversion and the MCA frontier}
\label{sec:capg-frontier}

The deep-point conversion has so far been used only at its threshold
operating point: \cref{thm:A} converts a list of size exceeding $q/k+1$
into $\eca>\tfrac1{2k}(1-n/q)$, and every deployed MCA certificate
(\cref{cor:deployed,cor:widened-deployed,prop:capff1-identity-frontier})
paid the full threshold $q/k+1\approx2^{166}$ even though the challenge
budgets are only $\floor{2^{-128}q}\approx2^{58}$ (KoalaBear) and
$\floor{2^{-100}q}\approx2^{24}$ (Mersenne-31).  The counting inside the
proof of \cref{thm:A} is in fact operating-point-free; extracting it
recovers the difference.

\begin{corollary}[flexible-budget deep-point conversion]
\label{cor:capg-budget-conversion}
Let $C=\RS[\F,D,k]$, $C^{+}=\RS[\F,D,k+1]$, $q=|\F|>n$, and let $m$ be an
integer agreement with $k+1\le m\le n$.  Suppose some received word
$U:D\to\F$ has
\[
        \bigl|\Lst\bigl(C^{+},\,1-m/n,\,U\bigr)\bigr|\ \ge\ L_0
        \qquad\text{with}\qquad
        \binom{L_0}{2}\,k\ <\ q-n .
\]
Then some simple-pole line $(f_\alpha,g_\alpha)$,
$f_\alpha=U/(x-\alpha)$, $g_\alpha=-1/(x-\alpha)$, $\alpha\in\F\setminus
D$, has at least $L_0$ distinct CA-bad --- hence, by \cref{fact:chain},
support-wise MCA-bad --- finite slopes at agreement threshold $m$.
Consequently
\[
        \eca\bigl(C,1-m/n\bigr)\ \ge\ \frac{L_0}{q},
        \qquad
        \emca\bigl(C,\delta\bigr)\ \ge\ \frac{L_0}{q}
        \quad\text{for every }\delta\in\bigl[1-\tfrac mn,\,1\bigr)
\]
by \cref{lem:mca-monotone}.  If moreover $\alpha$ is restricted to
$\F\setminus(D\cup\B)$, the same conclusion holds with $q-n$ replaced by
$q-|\B|$ in the hypothesis, and the certifying line is not $\B$-rational.
\end{corollary}

\begin{proof}
This is the count inside the proof of \cref{thm:A}, stopped before the
threshold specialization.  Let $P_1,\dots,P_{L_0}\in\F[X]_{<k+1}$ be
distinct listed polynomials.  For $i\ne j$ the difference $P_i-P_j$ is a
nonzero polynomial of degree at most $k$, so
$\sum_{\alpha\in\Omega}\#\{\,\{i,j\}:P_i(\alpha)=P_j(\alpha)\,\}
\le\binom{L_0}2k<|\Omega|$ for $\Omega=\F\setminus D$; hence some
$\alpha\in\Omega$ has \emph{no} collisions, and the $L_0$ values
$P_i(\alpha)$ are pairwise distinct.  As in the proof of \cref{thm:A},
each value $z_i=P_i(\alpha)$ is a slope at which
$f_\alpha+z_ig_\alpha=(P_i(x)-P_i(\alpha))/(x-\alpha)$ is explained by
$C$ on the agreement set of $P_i$ with $U$ (size $\ge m\ge k+1$), while no
polynomial of degree $<k$ agrees with $g_\alpha$ on more than $k$ points;
so all $L_0$ slopes are CA-bad at radius $1-m/n$ and support-wise MCA-bad
with the same witnesses.  The probability statements follow from
\cref{def:mca} with $\gamma$ uniform in $\F$.  For the last sentence run
the average over $\Omega'=\F\setminus(D\cup\B)$; the pole lines with
$\alpha\notin\B$ have $g_\alpha$ not $\B$-valued.
\end{proof}

Combining \cref{cor:capg-budget-conversion} with the identity-prefix floor
\cref{lem:capff1-identity-prefix-floor} at $K=k+1$, the MCA-row unsafe
certificate becomes the pure fiber-versus-budget comparison
\[
        \binom nm\ >\ p^{\,m-k-1}\,\Bigl\lfloor\frac{q}{2^{t}}\Bigr\rfloor ,
        \qquad t=\text{row target exponent},
\]
of exactly the same shape as the list-row certificates --- the threshold
$q/k+1$ has disappeared.  Optimizing $m$ under exact integer verification
moves the MCA frontier.

\begin{proposition}[moved MCA frontier, exact]\label{prop:capg-moved-frontier}
With $n=2^{21}$, $k=2^{20}$, and the deployed rows of
\cref{prop:capff1-identity-frontier}:
\begin{center}\footnotesize
\begin{tabular}{@{}llllll@{}}
\toprule
row / challenge & $m_1$ & $w=m_1-k-1$ & certified unsafe range &
margins (bits) \\
\midrule
KoalaBear, MCA, $\varepsilon^{*}=2^{-128}$ & $1116047$ & $67470$ &
$\delta\in[\,\tfrac{981105}{2097152},\tfrac12)=[0.4678273\ldots,\tfrac12)$
& $+9.0\ /\ -22.2$ \\
Mersenne-31 line round, MCA, $\varepsilon^{*}=2^{-100}$ & $1116023$ &
$67446$ &
$\delta\in[\,\tfrac{981129}{2097152},\tfrac12)=[0.4678387\ldots,\tfrac12)$
& $+27.9\ /\ -3.3$ \\
\bottomrule
\end{tabular}
\end{center}
Specifically, the exact integer comparisons
\[
        \binom{n}{1116047}\ >\ p^{67470}\Bigl\lfloor\frac{p^{6}}{2^{128}}\Bigr\rfloor,
        \qquad
        \binom{n}{1116023}\ >\ (p')^{67446}\Bigl\lfloor\frac{(p')^{4}}{2^{100}}\Bigr\rfloor
\]
hold, and both fail with $m_1$ replaced by $m_1+1$ and $w$ by $w+1$; the
auxiliary admissibility
$\binom{\floor{\varepsilon^{*}q}+1}{2}k<q-n$ holds exactly on both rows
(indeed with $q-n$ replaced by $q-|\B|$), and
$\varepsilon^{*}q\notin\mathbb Z$, so unsafety is strict.  Hence
$\emca>\varepsilon^{*}$ on the printed ranges: the MCA rows are certified at these stronger agreements, and in \cref{cor:capff1-sandwich-edges} the imported sandwich
upper edges improve to
\[
        \delta^{*}_{C}(2^{-128})\le\frac{981105}{2097152}
        \quad\text{(KoalaBear)},\qquad
        \delta^{*}_{C}(2^{-100})\le\frac{981129}{2097152}
        \quad\text{(Mersenne-31)}.
\]
The two list rows are unchanged.  Both new edges remain strictly below the
entropy--subfield envelope of \cref{def:capff1-gstar}: the deployed MCA
gaps now sit within $7.2\times10^{-7}$ (KoalaBear) and
$4.9\times10^{-7}$ (Mersenne-31) of the respective ceilings
$1-\rho-g^{*}$.
\end{proposition}

\begin{proof}
Apply \cref{cor:capg-budget-conversion} with
$L_0=\floor{\varepsilon^{*}q}+1$ and $U=U_{z^{*}}$ the identity witness
word of \cref{lem:capff1-identity-prefix-floor} at $K=k+1$ and $m=m_1$;
by \cref{prop:capff1-witness-exact} its $C^{+}$-list at radius $1-m_1/n$
has size $|\mathrm{Fib}|\ge\ceil{\binom n{m_1}p^{-w}}\ge L_0$, which is
the first displayed comparison ($\ceil{x}\ge B+1\iff x>B$).  The
admissibility hypothesis is the second displayed comparison.  Then
$\emca(C,\delta)\ge L_0/q>\varepsilon^{*}$ on the printed range.  For the
circle row the floor and the witness classification run verbatim on the
twin-coset $x$-domain ($D\subseteq\B=\F_{p'}$ is all that is used), and
the conversion is field-agnostic.  Failure at $m_1+1$ is the printed
failing comparison.  All six integer comparisons (two passes, two
failures, two admissibilities) were verified by exact arithmetic; the bit
margins are orientation values.  The envelope positions are
$g=67471/2097152=0.03217268<g^{*}(\tfrac12,\log_2p)=0.03217339$ and
$g=67447/2097152=0.03216124<g^{*}(\tfrac12,\log_2p')=0.03216172$.
\end{proof}

\begin{corollary}[updated adjacent pairs]\label{cor:capg-adjacent-pairs}
In \cref{prob:capff1-frontier} and \cref{prob:capfp-F}, the predicted
adjacent pairs become
\[
\begin{array}{c|cc|c}
\text{row} & \text{proved unsafe }a_0 & \text{conjectured first safe }a_0+1
& \text{fail margin (bits)}\\ \hline
\text{KoalaBear MCA} & 1116047 & 1116048 & 22.2\\
\text{KoalaBear list} & 1116046 & 1116047 & 22.0\\
\text{Mersenne-31 MCA} & 1116023 & 1116024 & 3.3\\
\text{Mersenne-31 list} & 1116022 & 1116023 & 3.1
\end{array}
\]
with the endpoint convention of \cref{cor:capf-endpoint}.  Two structural
remarks.  First, at equal $m$ the MCA comparison is exactly one factor of
$p$ easier than the list comparison (depth $m-k-1$ against $m-k$), while
one agreement step costs $\log_2p-\log_2\frac{n-m}{m+1}>\log_2p$ bits
for $m>n/2$; hence
$a_0^{\mathrm{list}}\le a_0^{\mathrm{MCA}}\le a_0^{\mathrm{list}}+1$
always in this regime, with the right equality realized on both deployed
rows: the unsafe edges interleave as
$\mathrm{list}<\mathrm{MCA}$ one step apart.  Second,
the two Mersenne-31 rows are now the tight rows ($3.1$ and $3.3$ bits);
any finite safe-side input must be exact-constant there.
\end{corollary}

\begin{proof}
Immediate from \cref{prop:capg-moved-frontier} and the unchanged list rows
of \cref{prop:capff1-identity-frontier}; the one-factor-of-$p$ relation is
the identity $w_{\mathrm{MCA}}(m)=w_{\mathrm{list}}(m)-1$.
\end{proof}

\section{Dual-weighted syndrome calculus}\label{sec:capg-calculus}

The exact-agreement machinery of
\cref{sec:aperiodic-hankel-certificates} manipulates window sums
$\sum_x y(x)L(x)x^m$.  Over a multiplicative coset the $m=0$ row of the
unweighted sum retains a constant-term defect, which the main text absorbs
by the shifted-window convention.  We adopt instead the dual-weighted
normalization, which is exact in the unshifted window over \emph{every}
domain $D$ and makes the dictionary below uniform.  The two conventions
differ by the fixed diagonal change $y\mapsto v\cdot y$ (a code equivalence
onto the dual-twisted generalized Reed--Solomon partner) together with an
invertible triangular change of window rows; all rank, kernel and cofactor
statements used from the main text are invariant under such changes, and
\cref{thm:capf-spi} is stated for arbitrary affine pencils, so it applies to
the realization below verbatim.

\begin{lemma}[residue pairing]\label{lem:capg-residue}
Let $D\subseteq\F$ with $|D|=n\geq 1$ and set $v_x:=1/P_D'(x)$ for $x\in D$
(well defined: $P_D'(x)=\prod_{y\in D\setminus\{x\}}(x-y)\neq 0$).  Then for
every $h\in\F[X]$ with $\deg h\leq n-1$,
\[
\sum_{x\in D} v_x\,h(x) \;=\; [X^{\,n-1}]\,h .
\]
In particular the sum vanishes whenever $\deg h\leq n-2$.
\end{lemma}

\begin{proof}
The Lagrange interpolant
$\sum_{x\in D}h(x)\,\ell_{D\setminus\{x\}}(X)/P_D'(x)$ has degree
$\leq n-1$ and agrees with $h$ on the $n$ points of $D$; since
$\deg h\leq n-1$, the difference has $n$ roots and degree $\leq n-1$, hence
is zero.  Each $\ell_{D\setminus\{x\}}$ is monic of degree $n-1$, so
comparing coefficients of $X^{n-1}$ gives the identity.
\end{proof}

\begin{definition}[weighted syndromes, annihilator spaces]\label{def:capg-syndrome}
For a word $y\colon D\to\F$ put $s_m(y):=\sum_{x\in D}v_x\,y(x)\,x^m$
($m\geq 0$).  For $L=\sum_i \ell_i X^i\in\F[X]$ and a window length
$\tau\geq 1$ we say $L$ \emph{annihilates $y$ in window $\tau$}, and write
$L\in\mathcal{A}_\tau(y)$, if
\[
\sum_{x\in D} v_x\,y(x)\,L(x)\,x^m \;=\;\sum_i \ell_i\, s_{m+i}(y)\;=\;0
\qquad (0\leq m\leq \tau-1).
\]
$\mathcal{A}_\tau(y)$ is an $\F$-subspace of $\F[X]$.  Restricted to
$\deg L\leq j$ the condition reads $H_{\tau,j+1}\big(s(y)\big)\,\tilde\ell=0$
with the Hankel matrix built from $(s_0(y),\dots,s_{\tau+j-1}(y))$; with
$u:=\big(s_m(f)\big)_m$ and $v:=\big(s_m(g)\big)_m$ this realizes the
exact-agreement pencils of \cref{sec:aperiodic-hankel-certificates} for the
line $(f,g)$.
\end{definition}

\begin{lemma}[annihilator--tangent dictionary]\label{lem:capg-dictionary}
Let $L$ be split over $D$ with root set $T$, $\gamma:=\deg L=|T|$, and let
$\tau\geq 1$ satisfy $\tau+\gamma\leq n$.  Then for every word $y$,
\[
L\in\mathcal{A}_\tau(y)
\quad\Longleftrightarrow\quad
y\in \mathrm{RS}[\F,D,\,n-\tau-\gamma]\;\oplus\;\F^{T},
\]
where $\F^{T}$ denotes the words supported on $T$ and
$\mathrm{RS}[\F,D,\kappa]:=\{p|_D:\deg p<\kappa\}$ (the zero space if
$\kappa\leq 0$).  The sum on the right is direct.
\end{lemma}

\begin{proof}
($\supseteq$) If $y$ is supported on $T$ then $y(x)L(x)=0$ for all $x\in D$.
If $y=p|_D$ with $\deg p\leq n-\tau-\gamma-1$, then for $0\leq m\leq\tau-1$
the polynomial $h:=p\,L\,X^{m}$ has degree
$\leq (n-\tau-\gamma-1)+\gamma+(\tau-1)=n-2$, so
$\sum_x v_x y(x)L(x)x^m=[X^{n-1}]h=0$ by \cref{lem:capg-residue}.

(Rank.)  The linear map $y\mapsto\big(\sum_x v_xy(x)L(x)x^m\big)_{m<\tau}$
has, in the standard basis of $\F^D$, the column at $x$ equal to
$v_xL(x)\cdot(x^m)_{m<\tau}$.  The columns over $T$ vanish; over
$E:=D\setminus T$ the scalars $v_xL(x)$ are nonzero, so the matrix is a
column-scaled $\tau\times(n-\gamma)$ Vandermonde on distinct nodes and has
rank $\min(\tau,n-\gamma)=\tau$.

(Dimensions.)  The kernel has dimension $n-\tau$.  The right side has
dimension $(n-\tau-\gamma)+\gamma=n-\tau$ once the sum is direct: a word in
the intersection is $p|_D$ vanishing on $E$, i.e.\ a polynomial of degree
$\leq n-\tau-\gamma-1<n-\gamma=|E|$ with $|E|$ roots, hence zero.  The
containment of a subspace of equal dimension forces equality.
\end{proof}

At the critical window $\tau=n-k-\gamma$ the dictionary reads: \emph{a split
annihilator of degree $\gamma$ certifies agreement with a codeword of $C$
off its root set}.  Shorter windows certify agreement with higher-degree
polynomials; this graded form is used in \cref{prop:capg-class0}.

\begin{lemma}[window cost of multiplication]\label{lem:capg-window}
If $L\in\mathcal{A}_\tau(y)$ and $A\in\F[X]$ with $d:=\deg A\leq\tau-1$,
then $AL\in\mathcal{A}_{\tau-d}(y)$.
\end{lemma}

\begin{proof}
$\sum_x v_xy(AL)(x)x^m=\sum_{c}A_c\sum_x v_xyL(x)x^{m+c}=0$ whenever
$m+d\leq\tau-1$.
\end{proof}

\begin{lemma}[gcd window]\label{lem:capg-gcdwindow}
Let $L_1,L_2\in\mathcal{A}_\tau(y)$ be nonzero with $\deg L_i\leq j$, let
$G:=\gcd(L_1,L_2)$ and $\gamma:=\deg G$, and suppose $\tau\geq j-\gamma$.
Then $G\in\mathcal{A}_{\tau-(j-\gamma-1)}(y)$ (with the convention that for
$L_2\mid L_1$ this is the trivial statement $G\in\mathcal{A}_\tau(y)$).
\end{lemma}

\begin{proof}
Write $L_i=GC_i$ with $C_1,C_2$ coprime.  If $C_2$ is constant then
$G$ is a scalar multiple of $L_2$ and the claim is trivial.  Otherwise
B\'ezout gives $A,B$ with $AC_1+BC_2=1$, $\deg A<\deg C_2\leq j-\gamma$ and
$\deg B<\deg C_1\leq j-\gamma$; hence $G=AL_1+BL_2$ with
$\deg A,\deg B\leq j-\gamma-1$, and \cref{lem:capg-window} plus linearity
give $G\in\mathcal{A}_{\tau-(j-\gamma-1)}(y)$.
\end{proof}

\begin{lemma}[small annihilator]\label{lem:capg-small}
If $\dim\big(\mathcal{A}_\tau(y)\cap\F[X]_{\leq j}\big)\geq d+1$ for some
$0\leq d\leq j$, then $\mathcal{A}_\tau(y)$ contains a nonzero element of
degree $\leq j-d$.
\end{lemma}

\begin{proof}
$\F[X]_{\leq j-d}$ has codimension $d$ in $\F[X]_{\leq j}$, so a subspace of
dimension $\geq d+1$ meets it nontrivially.
\end{proof}

\begin{remark}[support-side companions]\label{rem:capg-supportside}
The functionals $S_r(h,T)=\sum_{x\in T}h(x)x^r/\Lambda_T'(x)$ of
\cref{lem:capfr1-interpolation-test,lem:capfp-functionals} are the
\emph{support-side} instance of the same residue identity as
\cref{lem:capg-residue}: there the domain is $T$ itself, the weights are
$1/\Lambda_T'$, and window annihilation tests interpolability on $T$.
The dictionary \cref{lem:capg-dictionary} is the domain-side dual: the
domain is all of $D$, the weights are $1/P_D'$, and window annihilation
against a split locator $L$ tests membership in
$\RS[\F,D,n-\tau-\gamma]\oplus\F^{Z(L)}$.  The two calculi meet in the
rank-one census: the twisted vectors
$(u(x)/\Lambda_T'(x))\mathbf 1_T$ of \cref{prob:capfp-R1} are supported
window data in the support-side calculus, while the pencil realization
used in \cref{sec:capg-rigidity} is domain-side.  Statements proved in
either calculus transfer along the diagonal code equivalence recorded
above.
\end{remark}

\section{Split rigidity for kernel sections}\label{sec:capg-rigidity}

The residual branches of the exact-agreement atlas are populated by kernel
sections --- one-parameter families $W(Z,X)$ annihilating the line's
syndromes identically in the slope variable --- whose normalized locators
pass the split test identically.  We show these branches collapse: the
locator family is constant, and the line is a tangent pair.  The engine is
an elementary rigidity statement for divisors with constant coefficients.

\begin{lemma}[relative algebraic closure]\label{lem:capg-relclosed}
$\F$ is relatively algebraically closed in $\F(Z)$: if $c\in\F(Z)$ satisfies
$m(c)=0$ for some nonzero $m\in\F[Y]$, then $c\in\F$.
\end{lemma}

\begin{proof}
Write $c=a/b$ with $a,b\in\F[Z]$ coprime, $b$ monic, and
$m=\sum_{i\leq d}m_iY^i$ with $m_d\neq 0$, $d\geq 1$ (if $d=0$ then $m$ is a
nonzero constant and has no root, so $d\geq 1$).  Multiplying $m(c)=0$ by
$b^d$ gives $m_d a^d=-b\sum_{i<d}m_ia^ib^{d-1-i}$, so $b\mid m_d a^d$; by
coprimality $b\mid m_d$, hence $b=1$ and $c=a\in\F[Z]$.  If $\deg_Z a\geq 1$
then, letting $i_0$ be the largest index with $m_{i_0}\neq 0$ and
$i_0\geq 1$, the polynomial $m(a(Z))=\sum_i m_ia(Z)^i$ has a nonzero leading
term of degree $i_0\deg_Z a$ (the degrees $i\deg_Z a$ are distinct), so
$m(a)\not\equiv 0$, a contradiction; if $i_0=0$ then $m(a)=m_0\ne0$,
likewise a contradiction.  Hence $\deg_Z a=0$ and $c\in\F$.
\end{proof}

\begin{lemma}[constant-divisor rigidity]\label{lem:capg-constdiv}
Let $P\in\F[X]$ be nonzero and let $L\in\F(Z)[X]$ be monic with
$L\mid P$ in $\F(Z)[X]$.  Then $L\in\F[X]$ and $L\mid P$ in $\F[X]$.
\end{lemma}

\begin{proof}
Over $\overline{\F}(Z)$ the polynomial $P$ factors as
$\mathrm{lc}(P)\prod_r(X-\xi_r)$ with all $\xi_r\in\overline{\F}$ (they are
roots of $P$, hence algebraic over $\F$, hence constant).  Since
$\overline{\F}(Z)[X]$ is a unique factorization domain and $L$ is a monic
divisor of $P$ there, $L=\prod_{r\in I}(X-\xi_r)$ for some index set $I$;
its coefficients lie in $\overline{\F}$.  They also lie in $\F(Z)$, hence in
$\F$ by \cref{lem:capg-relclosed}.  Finally, division with remainder of $P$
by $L$ in $\F[X]$ must reproduce the (unique) division in $\F(Z)[X]$, whose
remainder is zero; so $L\mid P$ in $\F[X]$.
\end{proof}

\begin{theorem}[split kernel sections are tangent-borne]\label{thm:capg-tangentborne}
Fix an exact agreement $a$ for $C=\mathrm{RS}[\F,D,k]$, and set $j=n-a$,
$t=a-k\geq 1$.  Let $(f,g)$ be a received line and let
$W(Z,X)=\sum_{i=0}^{j}W_i(Z)X^i\in\F[Z][X]$ be nonzero with
\begin{equation}\label{eq:capg-pairing}
\sum_{x\in D} v_x\,(f+Zg)(x)\,W(Z,x)\,x^{m}\;\equiv\;0 \quad\text{in }\F[Z]
\qquad (0\leq m\leq t-1),
\end{equation}
equivalently $\big(H_t(u)+ZH_t(v)\big)\,\big(W_0(Z),\dots,W_j(Z)\big)^{\!\top}\equiv 0$
with $u,v$ the weighted syndrome sequences of $f,g$.  Let
$\gamma:=\deg_X W$ over $\F(Z)$ (so $W_\gamma\not\equiv 0$) and let
$\widehat{L}:=W/W_\gamma\in\F(Z)[X]$ be the monic normalization.  Assume
\[
\widehat{L}\ \text{divides}\ P_D\ \text{in}\ \F(Z)[X].
\]
Then:
\begin{enumerate}[label=(\roman*)]
\item $\widehat{L}=L^{*}\in\F[X]$ is a (constant) monic divisor of $P_D$,
hence split over $D$; write $T^{*}:=Z(L^{*})$, $|T^{*}|=\gamma$.
\item $f,g\in \mathrm{RS}[\F,D,\,n-t-\gamma]\oplus\F^{T^{*}}$; write
$f=c_f+e_f$, $g=c_g+e_g$ with $\deg c_f,\deg c_g<n-t-\gamma$ and
$e_f,e_g$ supported on $T^{*}$.
\item At every agreement threshold $a'$ with $n-a'\leq t$, the line $(f,g)$
has at most $\max(\gamma,1)$ MCA-bad finite slopes.  If $c_f,c_g\in C$,
every such slope is a tangent ratio $-e_f(x)/e_g(x)$ with $x\in T^{*}$,
$e_g(x)\neq 0$, and there are at most $\gamma$ of them; otherwise there is at
most one such slope, namely the unique $z$ (if any) with $c_f+zc_g\in C$.
\item If $c_f,c_g\in C$, the pair $(f,g)$ is column-close to $C$ at radius
$\gamma/n$, hence has no CA-bad slopes at any radius pair
$(\delta,\delta_{\mathrm{int}})$ with $\delta_{\mathrm{int}}\geq\gamma/n$.
\end{enumerate}
\end{theorem}

\begin{proof}
(i) is \cref{lem:capg-constdiv} applied to $P:=P_D$ (squarefree with root set
$D$, so any monic divisor is split over $D$).

(ii)  Dividing \eqref{eq:capg-pairing} by the unit $W_\gamma\in\F(Z)^\times$
and using $\widehat{L}=L^{*}$ constant,
\[
0=\sum_x v_x (f+Zg)(x)L^{*}(x)x^m
=\Big(\sum_x v_x f(x)L^{*}(x)x^m\Big)
+Z\Big(\sum_x v_x g(x)L^{*}(x)x^m\Big)
\quad\text{in }\F(Z)
\]
for $0\leq m\leq t-1$; both coefficients of this degree-$\leq1$ polynomial
in $Z$ vanish, i.e.\ $L^{*}\in\mathcal{A}_t(f)\cap\mathcal{A}_t(g)$.  Since
$t+\gamma\leq t+j=n-k\leq n$, \cref{lem:capg-dictionary} with $\tau=t$ gives
(ii); note $\kappa:=n-t-\gamma\geq n-t-j=k\geq 1$.

(iii)  Fix $a'$ with $r':=n-a'\leq t$ and let $z$ be an MCA-bad finite slope
at threshold $a'$, with witness support $S$, $|S|\geq a'$, and explaining
codeword $c\in C$ for the point $f+zg$ on $S$.  On
$S\setminus T^{*}$, which has size $\geq(n-r')-\gamma\geq n-t-\gamma=\kappa$,
both $c$ and $c_f+zc_g$ are restrictions of polynomials of degree
$<\kappa$ (using $C\subseteq \mathrm{RS}[\F,D,\kappa]$) agreeing at
$\geq\kappa$ points, hence equal as polynomials:
\begin{equation}\label{eq:capg-forced}
c \;=\; c_f+zc_g .
\end{equation}
Suppose first $c_f,c_g\in C$.  Explanation of the \emph{pair} on $S$ by
codewords $p_1,p_2\in C$ forces, by the same $\kappa$-point argument on
$S\setminus T^{*}$, $p_1=c_f$ and $p_2=c_g$; hence the pair is explained on
$S$ if and only if $e_f,e_g$ both vanish on $S\cap T^{*}$.  Since $z$ is
MCA-bad there is a witness $S$ on which the pair is \emph{not} explained, so
some $x\in S\cap T^{*}$ has $(e_f(x),e_g(x))\neq(0,0)$; and
\eqref{eq:capg-forced} restricted to $S$ gives $e_f(x)+ze_g(x)=0$ at that
$x$.  If $e_g(x)=0$ this is impossible; if $e_g(x)\neq0$ then
$z=-e_f(x)/e_g(x)$.  Thus every bad slope is one of at most
$|T^{*}|=\gamma$ tangent ratios.  (If $\gamma=0$ then $e_f=e_g=0$,
$f,g\in C$, and no slope is ever MCA-bad.)

Suppose now that $c_f\notin C$ or $c_g\notin C$.  By
\eqref{eq:capg-forced}, badness of $z$ forces $c_f+zc_g\in C$.  If
$c_g\in C$ this forces $c_f\in C$, a contradiction, so there is no bad
slope.  If $c_g\notin C$, then two distinct bad slopes $z\neq z'$ would give
$(z-z')c_g\in C$, again a contradiction; so there is at most one bad slope.
This proves (iii).

(iv)  If $c_f,c_g\in C$ then $(f,g)$ differs from the codeword pair
$(c_f,c_g)$ only on $T^{*}$, so its column distance to $C\times C$ is at
most $\gamma/n$; a pair that is $\delta_{\mathrm{int}}$-close for
$\delta_{\mathrm{int}}\geq\gamma/n$ is not
$\delta_{\mathrm{int}}$-far, hence supports no CA-bad slope at
$(\delta,\delta_{\mathrm{int}})$ by definition.
\end{proof}

\begin{corollary}[collapse of the identically split top chart at deficiency
one]\label{cor:capg-spib}
In the setting of \cref{thm:capf-spi} at deficiency one ($t=j$, realized by
the weighted pencil $M(Z)=H_t(u)+ZH_t(v)$ of \cref{def:capg-syndrome}),
assume the top chart is nonempty ($c_j\not\equiv 0$) and that case (b)
holds, i.e.\ the split test passes identically.  Then there is a fixed
$j$-set $T^{*}\subseteq D$ such that $f$ and $g$ each agree with a codeword
of $C$ off $T^{*}$; the line has at most $j$ MCA-bad finite slopes at every
agreement threshold $a'\geq a$, each a tangent ratio; and it has no CA-bad
slopes at any radius pair $(\delta,\delta_{\mathrm{int}})$ with
$\delta_{\mathrm{int}}\geq j/n$.  In particular the residual named in case
(b) of \cref{thm:capf-spi} is closed.
\end{corollary}

\begin{proof}
For the $t\times(j+1)$ pencil $M(Z)$ let
$c(Z)=\big(c_0(Z),\dots,c_j(Z)\big)$ be the cofactor vector,
$c_i=(-1)^i\det M^{(i)}(Z)$ with $M^{(i)}$ the maximal minor omitting column
$i$.  Then $M(Z)c(Z)\equiv 0$ identically: appending any row of $M(Z)$ on
top of $M(Z)$ produces a $(j+1)\times(j+1)$ matrix with a repeated row,
whose Laplace expansion along the top row is the corresponding entry of
$M(Z)c(Z)$.  Set $W(Z,X):=\sum_i c_i(Z)X^i$; on the top chart
$\deg_XW=j=\gamma$ and $\widehat{L}=W/c_j$.  Case (b) states that the
pseudo-division remainders of $c_j^{\Delta}P_D$ by $W$ vanish identically,
i.e.\ $c_j^{\Delta}P_D=QW$ in $\F[Z][X]$; since $c_j$ is a unit in $\F(Z)$
this is exactly $\widehat{L}\mid P_D$ in $\F(Z)[X]$.
\cref{thm:capg-tangentborne} applies with $\gamma=j$; here
$n-t-\gamma=n-2j=k$, so $c_f,c_g\in C$ automatically and parts (iii)--(iv)
give the stated conclusions for all $a'$ with $n-a'\leq t=j$, i.e.\ all
$a'\geq a$.
\end{proof}

\begin{corollary}[identically singular chains with split
sections]\label{cor:capg-singular}
In the setting of \cref{lem:singular-pencil} (any $t\geq j$), suppose some
polynomial kernel section $W(Z,X)$ of the pencil, of $X$-degree
$\gamma\leq j$ over $\F(Z)$, has identically split monic normalization
(equivalently, $\widehat{L}\mid P_D$ in $\F(Z)[X]$).  Then the conclusions
of \cref{thm:capg-tangentborne} hold; in particular the line has at most
$\max(\gamma,1)$ MCA-bad finite slopes at every threshold $a'$ with
$n-a'\leq t$.
\end{corollary}

\begin{proof}
A kernel section satisfies \eqref{eq:capg-pairing} by definition of the
pencil realization; apply \cref{thm:capg-tangentborne}.
\end{proof}

\begin{remark}[quantitative identically-split criterion]\label{rem:capg-quantsplit}
Let $W(Z,X)$ be as in \cref{thm:capg-tangentborne} with
$e_Z:=\deg_Z W$ and top coefficient $W_\gamma$.  Pseudo-dividing $P_D$ by
$W$ in $X$ takes $n-\gamma+1$ steps, each multiplying the running remainder
by $W_\gamma$ and subtracting a coefficient multiple of $W$, so every
remainder coefficient lies in $\F[Z]$ with degree at most
$(n-\gamma+1)e_Z$.  Consequently, if the specialized monic normalization
$W(z,\cdot)/W_\gamma(z)$ is split over $D$ for more than
$(n-\gamma+2)\,e_Z$ values of $z$ (the extra $e_Z$ absorbing the roots of
$W_\gamma$), then it is split identically and
\cref{thm:capg-tangentborne} applies.  In words: a kernel section cannot pass
the split test at superlinearly many slopes without collapsing to a tangent
pair.
\end{remark}

We now reduce the remaining deficiency-one residual, class (0)
(identically rank-deficient pencils), using the calculus of
\cref{sec:capg-calculus}.

\begin{lemma}[small kernel sections in class (0)]\label{lem:capg-class0-section}
Let $t=j$ and suppose the pencil $M(Z)$ has rank $\leq j-1$ over $\F(Z)$.
Then there is a nonzero $W(Z,X)\in\F[Z][X]$ with $\deg_XW\leq j-1$,
$\deg_ZW\leq j-1$, satisfying \eqref{eq:capg-pairing}.
\end{lemma}

\begin{proof}
The $\F(Z)$-kernel of $M(Z)$ has dimension $\geq 2$ inside
$\F(Z)^{j+1}$, hence meets the hyperplane $\{w_j=0\}$ nontrivially;
equivalently, the $t\times j$ subpencil $M'(Z)$ of the first $j$ columns has
rank $\rho\leq j-1$ over $\F(Z)$ and a nontrivial kernel.  Choose a
$\rho\times\rho$ minor of $M'$ that is not identically zero and solve the
pivot coordinates by Cramer's rule with one free coordinate set to that
minor; the resulting kernel vector has polynomial entries, each a signed
$\rho\times\rho$ or bordered minor of an affine pencil, hence of $Z$-degree
$\leq\rho\leq j-1$.
\end{proof}

\begin{proposition}[class-(0) dichotomy at deficiency one]\label{prop:capg-class0}
Let $t=j$, let the pencil of the line $(f,g)$ be identically rank-deficient
(class (0)), and let $W$ be as in \cref{lem:capg-class0-section}, with
$\gamma_W:=\deg_XW\leq j-1$ and $e_Z:=\deg_ZW\leq j-1$.  Consider the
MCA-bad finite slopes of $(f,g)$ at exact agreement $a$, i.e.\ those with a
witness support of size exactly $a$, and stratify each bad slope $z$ by
$\gamma_z:=\deg\gcd\big(\ell_{T},\,W(z,\cdot)\big)$, where $T$ is the
co-support of a witness (so $\ell_T\in\mathcal{A}_t(f+zg)$ by
\cref{lem:capg-dictionary}) and $z$ avoids the $\leq e_Z$ roots of the top
coefficient and the $\leq e_Z$ zeros of $W(z,\cdot)$.  Then:
\begin{enumerate}[label=(\alph*)]
\item (\emph{coprime stratum, $\gamma_z=0$}) if
$\big(s_0(f),s_0(g)\big)\neq(0,0)$, at most one bad slope lies in this
stratum: the unique $z$ with $s_0(f)+z\,s_0(g)=0$.
\item (\emph{section-split stratum, $W(z,\cdot)\mid\ell_T$}) these slopes
number at most $(n-\gamma_W+2)\,e_Z\leq (n+1)(j-1)$, unless the monic
normalization of $W$ is identically split, in which case
\cref{cor:capg-singular} bounds the line's bad slopes by
$\max(\gamma_W,1)\leq j-1$ outright.
\item (\emph{divisor-chain residual, $1\leq\gamma_z<\gamma_W$}) each such
slope carries a split $G_z:=\gcd(\ell_T,W(z,\cdot))$ of degree $\gamma_z$
with $G_z\in\mathcal{A}_{\gamma_z+1}(f+zg)$ and $G_z\mid W(z,\cdot)$; by the
dictionary, $f+zg$ agrees with a polynomial of degree $<n-2\gamma_z-1$ off
$Z(G_z)$.  We make no polynomial claim for this stratum; it is the named
residual.
\end{enumerate}
Lines with $s_0(f)=s_0(g)=0$ (equivalently, both interpolants of degree
$\leq n-2$) form the second named residual.
\end{proposition}

\begin{proof}
For a bad slope $z$ at exact agreement $a$ with witness support $S$ and
$T=D\setminus S$, $|T|=j$: the dictionary (\cref{lem:capg-dictionary},
$\tau=t$, $\gamma=j$, $n-t-j=k$) converts the explanation into
$\ell_T\in\mathcal{A}_t(f+zg)$.  Specializing \eqref{eq:capg-pairing} at $z$
gives $W(z,\cdot)\in\mathcal{A}_t(f+zg)$, nonzero for the admitted $z$, of
degree $\leq j-1$.  \cref{lem:capg-gcdwindow} with $\tau=t=j$ and degree cap
$j$ places $G_z$ in $\mathcal{A}_{\gamma_z+1}(f+zg)$.  If $\gamma_z=0$ then
$1\in\mathcal{A}_1(f+zg)$, i.e.\ $s_0(f+zg)=s_0(f)+zs_0(g)=0$, giving (a).
If $\gamma_z=\deg W(z,\cdot)$, i.e.\ $W(z,\cdot)\mid\ell_T$, then the monic
normalization of $W(z,\cdot)$ is split over $D$; by
\cref{rem:capg-quantsplit} either this happens for at most
$(n-\gamma_W+2)e_Z$ slopes or the section is identically split and
\cref{cor:capg-singular} applies --- giving (b).  In the remaining stratum,
$G_z$ divides the split $\ell_T$, hence is split, and the window
$\gamma_z+1$ instance of the dictionary ($\kappa=n-(\gamma_z+1)-\gamma_z$)
gives the stated agreement statement --- (c).
\end{proof}

\begin{remark}[closure of the identically-split residual of slope
elimination]\label{rem:capg-slopeelim-closure}
\Cref{thm:capfp-slope-elim}\textup{(d)} bounds the slope count of any
exact-agreement syndrome pencil $M(Z)=M_0+ZM_1$ by a rank-one census over
split candidates $\Lambda_A$, \emph{except} on the residual branch
$M_0\Lambda_A=M_1\Lambda_A=0$ ``as in \cref{thm:capf-spi}''.  For a fixed
split $\Lambda_A$ this branch is a common support in the sense of
\cref{prop:capfr1-slope-elimination} and is already harmless by
\cref{thm:capfp-slope-elim}\textup{(b)}.  The genuinely open case was the
\emph{family} version --- kernel sections $W(Z,X)$ that are identically
split without any single $Z$-independent split annihilator being given
--- and that case is closed by \cref{thm:capg-tangentborne}: any rational
identically split kernel section is constant, so the branch collapses to
the fixed-candidate case, at every deficiency admitting a rational kernel
section, with the tangent count $\le\max(\gamma,1)$.  In particular the
``known unconditional strata'' list of \cref{prob:capfp-R1} may now cite
deficiency one \emph{without} the residual caveat: by
\cref{cor:capg-spib}, deficiency-one charts are eliminant-or-tangent
outright.
\end{remark}

\section{Identity-scale aperiodic slope floors and the refuted band
conjecture}\label{sec:capg-floors}

\cref{lem:capff1-identity-prefix-floor} produces list-size floors at the
identity periodicity scale, and \cref{prop:capff1-witness-exact} shows the
witness lists are exactly the prefix fibers.  We now transport these
statements to \emph{per-line, per-slope} form: for the pole lines of
\cref{thm:A}, the threshold-$m$ witness supports of every MCA-bad slope
are exactly fiber members, and for $\gcd(m,n)=1$ every witness is
aperiodic.  The consequence is that \cref{prob:band}, which asks for a
$q$-independent polynomial bound on aperiodically witnessed slopes
throughout the band $k+2n/N<a<n-\floor{(n-k)/3}$, \emph{fails} below the
entropy--subfield envelope.  This upgrades two prior gestures to
theorems: the informal band-edge move recorded in
\cref{rem:entropy-frontier}, and the genericity heuristic in the preamble
of \cref{subsec:capfr1-band-flatness} (``generic supports are aperiodic
except for a negligible quotient-periodic subfamily'') --- at odd $m$ the
correct statement is that \emph{all} witnesses are aperiodic, with no
exceptional subfamily, so no periodic bucket can absorb any part of the
floor.

\begin{definition}[prefixes, fibers, prefix words, pole lines]\label{def:capg-prefix}
Let $B\subseteq\F$ be a subfield with $D\subseteq B$, and fix
$k+1\leq m\leq n$; put $K:=k+1$ and $w:=m-K\geq 0$.  For an $m$-subset
$M\subseteq D$ let $\Lambda_M:=\ell_M$ and define the depth-$w$ prefix
$\Phi_w(M):=\big([X^{m-h}]\Lambda_M\big)_{h=1,\dots,w}\in B^{w}$.  For
$z\in B^{w}$ set $\mathrm{Fib}_w(z):=\Phi_w^{-1}(z)$ and
\[
U_z(X):=X^{m}+\textstyle\sum_{h=1}^{w}z_hX^{m-h}.
\]
For $\alpha\in\F\setminus D$ the \emph{pole line} at $(z,\alpha)$ is the
line spanned by the words
\[
f_\alpha(x):=\frac{U_z(x)}{x-\alpha},\qquad
g_\alpha(x):=\frac{-1}{x-\alpha}\qquad(x\in D),
\]
i.e.\ the pole lines of \cref{thm:A} instantiated at the prefix word $U_z$.
\end{definition}

\begin{lemma}[pole-line witness classification]\label{lem:capg-witness}
Fix $z\in B^{w}$, $\alpha\in\F\setminus D$, and write
$(f,g):=(f_\alpha,g_\alpha)$.
\begin{enumerate}[label=(\roman*)]
\item For every $\zeta\in\F$ and every $S\subseteq D$ with $|S|\geq m$: the
word $f+\zeta g$ is explained by $C$ on $S$ if and only if
\[
|S|=m,\qquad S\in\mathrm{Fib}_w(z),\qquad
\Lambda_S(\alpha)=U_z(\alpha)-\zeta .
\]
\item No $S$ with $|S|\geq k+1$ admits a simultaneous explanation of the
pair $(f,g)$ by $C$; in particular the pair is column-far at radius
$1-m/n$.
\end{enumerate}
Consequently the MCA-bad finite slopes of $(f,g)$ at agreement threshold $m$
are exactly
\[
\big\{\,U_z(\alpha)-\Lambda_S(\alpha)\;:\;S\in\mathrm{Fib}_w(z)\,\big\},
\]
each of them CA-bad at radii $(1-m/n,\,1-m/n)$ as well; and \emph{every}
threshold-$m$ witness support of every such slope is a member of
$\mathrm{Fib}_w(z)$.
\end{lemma}

\begin{proof}
(i, $\Rightarrow$)  Let $c=P_c|_D$ with $\deg P_c\leq k-1$ explain
$f+\zeta g$ on $S$; on $S$ this reads $U_z(x)-\zeta=(x-\alpha)P_c(x)$.  Set
$\Psi:=U_z-\zeta-(X-\alpha)P_c$.  Since
$\deg\big((X-\alpha)P_c\big)\leq k\leq m-1$, $\Psi$ is monic of degree $m$,
and it vanishes on $S$; hence $|S|\leq m$, so $|S|=m$ and
$\Psi=\Lambda_S$.  The coefficients of $\Psi$ in degrees
$m-1,\dots,m-w=k+1$ are unaffected by the degree-$\leq k$ correction, so
they equal those of $U_z$, i.e.\ $\Phi_w(S)=z$.  Evaluating at $\alpha$
kills the correction term and gives $\Lambda_S(\alpha)=U_z(\alpha)-\zeta$.

(i, $\Leftarrow$)  Let $S\in\mathrm{Fib}_w(z)$ and
$\zeta:=U_z(\alpha)-\Lambda_S(\alpha)$.  Then
$\Theta:=U_z-\Lambda_S$ has degree $\leq m-w-1=k$ (the monic terms and the
depth-$w$ prefix cancel), and $\zeta=\Theta(\alpha)$.  The polynomial
$P_c:=\big(\Theta-\Theta(\alpha)\big)/(X-\alpha)$ has degree $\leq k-1$, and
for $x\in S$ (where $\Lambda_S(x)=0$):
$U_z(x)-\zeta=\Theta(x)-\Theta(\alpha)=(x-\alpha)P_c(x)$, i.e.\
$(f+\zeta g)(x)=P_c(x)$; so $P_c|_D$ explains $f+\zeta g$ on $S$, $|S|=m$.

(ii)  If $G\in\F[X]$, $\deg G\leq k-1$, has $G|_D=g$ on some $S$ with
$|S|\geq k+1$, then $(X-\alpha)G(X)+1$ has degree $\leq k$, vanishes on $S$
($\geq k+1$ points), hence is identically zero; but its value at $\alpha$ is
$1$ --- a contradiction.  Since a simultaneous explanation of the pair
requires an explanation of $g$, none exists on any $S$ with $|S|\geq k+1$;
in particular no pair of codewords agrees with $(f,g)$ on $\geq m>k$
columns, so the pair is column-far at radius $1-m/n$.

(Consequences.)  A slope $\zeta$ is MCA-bad at threshold $m$ iff some
$S$, $|S|\geq m$, carries an explanation of the point but not of the pair;
by (ii) the pair is never explained at these sizes, so badness is
equivalent to the existence of an explanation of the point, which by (i) is
equivalent to $\zeta\in U_z(\alpha)-\Lambda_\bullet(\alpha)
\big(\mathrm{Fib}_w(z)\big)$, with witnesses exactly the fiber members
mapping to $\zeta$.  CA-badness at $(1-m/n,1-m/n)$ follows since the point
is $(1-m/n)$-close by (i) and the pair is $(1-m/n)$-far by (ii)
(cf.\ \cref{fact:chain}).
\end{proof}

\begin{remark}[relation to the exact witness lists]\label{rem:capg-witnessrel}
\Cref{lem:capg-witness} is the pole-twisted form of
\cref{prop:capff1-witness-exact}: the map $c\mapsto(X-\alpha)c+\zeta$ is a
bijection from degree-$<k$ explanations of $f_\alpha+\zeta g_\alpha$ on
$S$ to $\RS[\F,D,k+1]$-explanations $c'$ of $U_z$ on $S$ with
$c'(\alpha)=\zeta$, and under this bijection the classification above is
the classification of \cref{prop:capff1-witness-exact} refined by the
evaluation-at-$\alpha$ fibration $c_M\mapsto c_M(\alpha)
=U_z(\alpha)-\Lambda_M(\alpha)$.  Part \textup{(ii)} (the pair is never
explained) is what the fibration adds: it is exactly the far condition of
\cref{def:mca}, and it is what makes the transported floor land in
$\emca$ rather than in list size.
\end{remark}

\begin{theorem}[identity-scale aperiodic slope floor]\label{thm:capg-aperiodic-floor}
Let $D\subseteq\B$ be either a smooth multiplicative coset
$D=\theta H$, $H=\mu_n(\B)$, or a twin-coset Chebyshev $x$-domain of
size $n$ as in \cref{thm:stabilizer-partition}, and let
$k+1\leq m\leq n$ satisfy $\gcd(m,n)=1$.  Let
$z^{*}\in B^{w}$ maximize $|\mathrm{Fib}_w(\cdot)|$, so that
\[
N_{\mathrm{fib}}:=|\mathrm{Fib}_w(z^{*})|\;\geq\;
\Big\lceil \binom{n}{m}\,|B|^{-w}\Big\rceil ,
\]
and set $N_0:=\min\big(N_{\mathrm{fib}},\,\lceil (q-n)/k\rceil\big)$, where
$q=|\F|>n$.  Then there exists $\alpha\in\F\setminus D$ such that the pole
line $(f_\alpha,g_\alpha)$ at $(z^{*},\alpha)$ has at least
\[
\tfrac12\,N_0\Big(1-\tfrac{k}{q-n}\Big)
\]
distinct MCA-bad finite slopes at agreement threshold $m$, each also CA-bad
at radii $(1-m/n,1-m/n)$, and such that \emph{every} threshold-$m$ witness
support of every one of these slopes is aperiodic ($c(S)=1$ in the sense of
\cref{def:periodicity-scale}).  In particular all these slopes lie in the
aperiodic bucket of the stabilizer partition
(\cref{thm:stabilizer-partition}).
\end{theorem}

\begin{proof}
By \cref{lem:capg-witness} every threshold-$m$ witness of every bad slope of
such a line is a member of $\mathrm{Fib}_w(z^{*})$.  For any $m$-subset
$S\subseteq D$, \cref{thm:stabilizer-partition}\textup{(i)} gives
$c(S)\mid n$ and exhibits $S$ as a union of complete scale-$c(S)$ fibers;
every scale-$c$ fiber of $D$ has exactly $c$ points ($H_c$-orbits in the
multiplicative case, uniform $T_c$-fibers on the twin-coset $x$-domain by
\cref{lem:cheb-fibers}), so $c(S)$ divides $|S|=m$.  Hence
$c(S)\mid\gcd(m,n)=1$: every witness is aperiodic.

For the count, fix any subfamily
$\mathcal{G}\subseteq\mathrm{Fib}_w(z^{*})$ with $|\mathcal{G}|=N_0$.  For
distinct $S,S'\in\mathcal{G}$ the polynomial $\Lambda_S-\Lambda_{S'}$ is
nonzero of degree $\leq m-w-1=k$ (monic terms and prefixes cancel), so
$\Lambda_S(\alpha)=\Lambda_{S'}(\alpha)$ for at most $k$ values
$\alpha\in\F$.  Summing over $\alpha\in\Omega:=\F\setminus D$,
\[
\sum_{\alpha\in\Omega}\#\{\text{colliding pairs at }\alpha\}
\;\leq\;\binom{N_0}{2}k,
\]
so some $\alpha\in\Omega$ has at most $\binom{N_0}{2}k/(q-n)$ colliding
pairs, whence the map $S\mapsto U_{z^*}(\alpha)-\Lambda_S(\alpha)$ takes at
least
\[
N_0-\frac{N_0^{2}k}{2(q-n)}\;\geq\;
N_0\Big(1-\tfrac12\big(1+\tfrac{k}{q-n}\big)\Big)
\;=\;\tfrac12 N_0\Big(1-\tfrac{k}{q-n}\Big)
\]
distinct values on $\mathcal{G}$, using
$N_0\leq (q-n)/k+1$.  By \cref{lem:capg-witness} each distinct value is an
MCA-bad finite slope at threshold $m$ (and CA-bad at the stated radii), with
all witnesses in the fiber, hence aperiodic; membership in the aperiodic
bucket follows from the bucket convention of
\cref{thm:stabilizer-partition}.
\end{proof}


\begin{corollary}[deployed MCA frontier floors]
\label{cor:capg-deployed-floor}
Let the rows be as in \cref{prop:capg-moved-frontier}.  At
$m=1116047$ for KoalaBear and $m=1116023$ for Mersenne-31, the fiber
pigeonhole gives
\[
        N_1:=\ceil{\binom nm p^{-w}}
\]
with $N_1\ge2^{66}$ and $N_1\ge2^{51}$ respectively.  The exact comparisons
$N_1(N_1-1)k<2(q-n)$ hold, so the zero-collision average in the proof of
\cref{cor:capg-budget-conversion} yields a pole $\alpha$ at which all $N_1$
fiber slopes are distinct.  Thus, at the two MCA frontier agreements, a
single received line carries at least $N_1$ aperiodically witnessed MCA-bad
finite slopes.  These agreements lie strictly inside the unnormalized band of
\cref{prob:band}; this is the deployed obstruction that forces the normalized
band problem \cref{prob:capfr1-normalized-band}.
\end{corollary}

\begin{proof}
The two comparisons $\binom nm>p^{w}\floor{\varepsilon^*q}$ of
\cref{prop:capg-moved-frontier} give the stated $N_1$ values.  The two
zero-collision inequalities were included in the exact certificate of
\cref{prop:capg-moved-frontier}.  Hence some pole has no collisions among the
fiber evaluations, and \cref{lem:capg-witness} makes each value an MCA-bad
slope with witnesses inside the fiber.  Since both frontier agreements are
odd while $n$ is a power of two, the witnesses are aperiodic.
\end{proof}


\begin{corollary}[no uniform polynomial bound below the
envelope]\label{cor:capg-asymptotic}
For $\nu\to\infty$ set $n=2^{\nu}$, $k=n/2$, let $p_\nu$ be the least prime
$\equiv 1\pmod n$ (so $\beta_\nu:=\log_2 p_\nu=O(\log n)$ by Linnik's
theorem), $B_\nu:=\F_{p_\nu}$, $D_\nu:=\mu_n\subseteq B_\nu^{\times}$; let
$m_\nu$ be an odd integer with
$m_\nu-k\in\big[n/(8\beta_\nu),\,n/(4\beta_\nu)\big]$, and let
$\F_\nu:=\F_{p_\nu^{d_\nu}}$ with $d_\nu$ minimal such that
$q_\nu\geq k\cdot 2^{n/8}$.  Then for all large $\nu$:
\begin{enumerate}[label=(\alph*)]
\item $N_{\mathrm{fib}}\geq 2^{n/2}\geq\lceil(q_\nu-n)/k\rceil$, so
$N_0=\lceil(q_\nu-n)/k\rceil$;
\item some received line of the rate-$1/2$ multiplicative row has at least
$2^{n/8-2}$ MCA-bad finite slopes at agreement threshold $m_\nu$, all of
whose witness supports are aperiodic;
\item $a=m_\nu$ lies in the band of \cref{prob:band} for every smoothness
convention with $2n/N\leq n/(16\beta_\nu)$, in particular for the finest
scale $N=n$.
\end{enumerate}
Consequently no polynomial $n\mapsto P(n)$, uniform over this family (and a
fortiori none independent of $q$), bounds the number of aperiodically
witnessed MCA-bad finite slopes of a received line on the band of
\cref{prob:band}: the conjecture is false as stated below the
entropy--subfield envelope.
\end{corollary}

\begin{proof}
Write $g:=(m_\nu-k)/n\leq 1/(4\beta_\nu)\leq 1/4$.  The elementary bound
$H_2(\tfrac12+x)\geq 1-4x^{2}$ on $[0,\tfrac14]$ (the difference vanishes at
$x=0$ and has positive derivative $8x-\log_2\frac{1+2x}{1-2x}$ there) gives
\[
\log_2 N_{\mathrm{fib}}\;\geq\;nH_2(m_\nu/n)-w\beta_\nu
\;\geq\;n\big(1-4g^{2}\big)-ng\beta_\nu
\;\geq\;n\Big(1-\tfrac{1}{4\beta_\nu^{2}}-\tfrac14\Big)\;\geq\;\tfrac n2 ,
\]
while $\lceil(q_\nu-n)/k\rceil\leq q_\nu/k\leq p_\nu 2^{n/8}=
2^{n/8+O(\log n)}<2^{n/2}$ for large $\nu$ --- giving (a).  Then
\cref{thm:capg-aperiodic-floor} (with $\gcd(m_\nu,2^\nu)=1$ since $m_\nu$ is
odd) yields at least
$\tfrac12\lceil(q_\nu-n)/k\rceil(1-k/(q_\nu-n))\geq
(q_\nu-n-k)/(2k)\geq 2^{n/8-2}$ slopes for large $\nu$ --- (b).  For (c),
$m_\nu-k\geq n/(8\beta_\nu)>2n/N$ under the stated convention, and
$m_\nu<n-\lfloor(n-k)/3\rfloor$ since $m_\nu-k\leq n/(4\beta_\nu)=o(n)$
while $n-\lfloor(n-k)/3\rfloor-k=n/2-\lfloor n/6\rfloor=\Theta(n)$.  The
count in (b) is $2^{\Omega(n)}$, which exceeds every polynomial in $n$.
\end{proof}

\begin{remark}[the band conjecture after the floors]\label{rem:capg-band}
\Cref{cor:capg-asymptotic} refutes \cref{prob:band} as stated: no
polynomial $P$, uniform over the row family and independent of $q$, can
bound the aperiodically witnessed slope counts on the sub-envelope part
of the band.  At the deployed size the refutation is quantitative rather
than formal --- any $P$ certifying the band at $n=2^{21}$ must satisfy
$P(2^{21})>2^{164}$, i.e.\ degree $\ge8$ --- but that already destroys
the intended consumption: \cref{thm:conditional-mca,thm:conditional-list}
would inherit an aperiodic numerator of at least $2^{164}$ at every band
agreement, far above both challenge budgets, so the conditional theorems
would certify nothing.  Accordingly:
\begin{enumerate}[label=\textup{(\roman*)}]
\item the operative band input is \cref{prob:capfr1-normalized-band}, now
by necessity rather than by convention; its left-edge reserve $R(n)$
absorbs exactly the identity-prefix floors proved here, and the sentence
of \cref{subsec:capfr1-band-flatness} assigning the sub-band
$A\lesssim k+g^{*}n$ to the unsafe side is a theorem
(\cref{cor:capg-deployed-floor,cor:capg-asymptotic});
\item every citation of \cref{prob:band} inside
\cref{thm:conditional-mca,thm:conditional-list} and the closing
statements \cref{prop:capfpr-AQ-closing,prop:capff1-closing,%
cor:capfr1-Q-R1-closing,cor:capfr1-terminal-package,prop:capfp-closing} must
be read as \cref{prob:capfr1-normalized-band}, with the finite deployed
target at the \emph{new} adjacent values of
\cref{cor:capg-adjacent-pairs};
\item the flatness form of the band input is
\cref{prob:capfr1-master-flatness}, whose base-field normalization
$|\B|^{-s}$ is confirmed --- rather than refuted --- by the floors: the
first term of its model crosses one exactly at the envelope, which is
where the floors stop.
\end{enumerate}
\end{remark}

\begin{remark}[even $m$: the two list rows]\label{rem:capg-evenm}
At the deployed list rows $m$ is even with $\gcd(m,n)=2$, so
\cref{thm:capg-aperiodic-floor} does not apply verbatim.  The witness
classification (\cref{lem:capg-witness}) still confines all threshold-$m$
witnesses to a single fiber, and a $2$-periodic member $S$ (stabilizer
$\{\pm1\}$, $S=-S$) has $\Lambda_S(X)=G(X^{2})$ up to the coset twist, so
its prefix vanishes at every odd index $h\leq w$.  Hence either the
pigeonholed $z^{*}$ has some nonzero odd-index coordinate --- and then the
fiber contains \emph{no} periodic support and the floor is purely
aperiodic exactly as before --- or $z^{*}$ is even-supported, in which case
the periodic members of the fiber are precisely the scale-$2$ pullbacks of
an image fiber at half parameters and are charged to the periodic image
ledger (\cref{thm:fiber-descent,cor:periodic-support-count}); the aperiodic
residue of the fiber remains subject to the same collision count.  We do not
claim a floor for the list rows here; only that periodic contamination of
identity-scale fibers is confined to even-supported prefixes and is
quotient-borne.
\end{remark}

\begin{remark}[relation to the conditional theorems]
\cref{cor:capg-deployed-floor,cor:capg-asymptotic} do not touch the
unconditional results (\cref{thm:deep-mca,thm:mca-from-ca,thm:A}); they
show that the hypothesis consumed by
\cref{thm:conditional-mca,thm:conditional-list} must be the
normalized \cref{prob:capfr1-normalized-band} rather than \cref{prob:band}
(see \cref{rem:capg-band}).
The deployed conditional rows already sit at the envelope frontier with the
recorded margins, so the substitution is a restatement of what those rows
in fact require, not a weakening; see the final package in
\cref{sec:frontier-next-steps}.
\end{remark}

\section{Prefix-collision rigidity and exact second moments}
\label{sec:capg-ledger}

The floors of \cref{sec:capg-floors} and the closing accounting of
\cref{prop:capff1-closing,prop:capfp-closing} both turn on the sizes of the
prefix fibers $\mathrm{Fib}_w(z)$.  \Cref{rem:capfr1-packing-gap} names,
as a low-risk \textup{Q1} deliverable, the sharpening of the kernel
side-distance $\delta_w=\ceil{(w+1)/2}$ of
\cref{lem:capfr1-kernel-distance,prop:capfp-kernel} to $e\ge w+1$ by
locator subtraction, together with an exact collision ledger built on it.
This section pays that deliverable: the side-distance $w+1$ is proved
characteristic-free, the packing cap of
\cref{prop:capfr1-prefix-packing,cor:capfp-packing} is replaced with the
packing exponent doubled from $\delta_w-1\approx w/2$ to $w$, and the
second moment is stratified exactly with strata starting at $e=w+1$
instead of $\delta_w$ and with the top stratum identified.

\begin{lemma}[prefix-collision rigidity]\label{lem:capg-prefix-rigidity}
Let $D\subseteq\F$, $|D|=n$, and $1\leq w<m\leq n$.  If $M\neq M'$ are
$m$-subsets of $D$ with $\Phi_w(M)=\Phi_w(M')$, then
$|M\setminus M'|=|M'\setminus M|\geq w+1$.  Moreover if
$|M\setminus M'|=w+1$ then
$\ell_{M\setminus M'}-\ell_{M'\setminus M}$ is a nonzero constant
(a \emph{constant-shift split pair}).
\end{lemma}

\begin{proof}
Let $R:=M\cap M'$, $A:=\ell_{M\setminus M'}$, $B:=\ell_{M'\setminus M}$,
$e:=|M\setminus M'|=|M'\setminus M|$, and $G:=\ell_R$ of degree $m-e$.  Then
$\Lambda_M-\Lambda_{M'}=G\,(A-B)$.  Prefix equality (together with
monicity) means $\deg(\Lambda_M-\Lambda_{M'})\leq m-w-1$, hence
$\deg(A-B)\leq e-w-1$.  If $e\leq w$ this forces $A=B$, hence $M=M'$
(disjoint nonempty root sets cannot coincide unless $e=0$), a
contradiction; so $e\geq w+1$.  If $e=w+1$ then $\deg(A-B)\leq 0$ with
$A\neq B$ (their root sets are disjoint and nonempty), so $A-B$ is a
nonzero constant.
\end{proof}

\begin{remark}[relation to the moment-kernel bounds]
\label{rem:capg-rigidcompare}
\Cref{lem:capg-prefix-rigidity} strengthens the fiber statements of
\cref{lem:capfr1-kernel-distance} and \cref{prop:capfp-kernel}\textup{(b)}
in two directions at once.  First, the per-side distance doubles: the
moment-kernel argument bounds only the \emph{support} of
$\mathbf 1_M-\mathbf 1_{M'}\in K_w(D)$, giving $|M\triangle M'|\ge w+1$,
i.e.\ $e\ge\delta_w$ per side, whereas locator subtraction uses the
multiplicative structure of the balanced word and gives $e\ge w+1$ per
side, i.e.\ $|M\triangle M'|\ge2(w+1)$.  Equivalently: the
\emph{balanced-ternary} minimum weight of the moment kernel $K_w(D)$ is
at least $2(w+1)$, twice its BCH-style minimum weight $w+1$ --- a
statement about $K_w(D)$ itself, relevant to
\cref{prop:capfp-kernel}\textup{(d)}.  Second, the hypothesis weakens: the
locator argument never invokes Newton identities, so it is valid in every
characteristic and at every depth $w$, whereas the power-sum transfer
requires $w<\operatorname{char}\F$ (harmless at the deployed rows,
$w\approx6.7\times10^4<2^{31}$, but not free in general).
\end{remark}

\begin{corollary}[unconditional anticode fiber cap; replaces the
packing bounds]\label{cor:capg-anticode-cap}
For every $z\in B^{w}$,
\[
|\mathrm{Fib}_w(z)|\;\le\;\binom{n}{m-w}\Big/\binom{m}{w}.
\]
This replaces \cref{prop:capfr1-prefix-packing} and
\cref{cor:capfp-packing}, whose packing parameter is
$\delta_w-1$: at the deployed frontier rows the saving over the trivial
bound $\binom nm$, in bits, improves as follows (orientation values;
first line the four rows of \cref{cor:capfp-packing}, second line the
moved MCA rows of \cref{prop:capg-moved-frontier}):
\begin{center}\footnotesize
\begin{tabular}{@{}lllll@{}}
\toprule
row & $(m,w)$ & saved, \cref{cor:capfp-packing} & saved, this cap &
gap to conjectured average \\
\midrule
KoalaBear list & $(1116046,67470)$ & $213510$ & $361177$ & $1729629$ \\
Mersenne-31 list & $(1116022,67446)$ & $213453$ & $361085$ & $1729741$ \\
\addlinespace
KoalaBear MCA & $(1116047,67470)$ & $213510$ & $361177$ & $1729629$ \\
Mersenne-31 MCA & $(1116023,67446)$ & $213453$ & $361085$ & $1729741$ \\
\bottomrule
\end{tabular}
\end{center}
The cap thus removes about $17\%$ of the trivial-to-average gap, roughly
$1.48\times10^5$ bits more per row than the $\delta_w$-packing; the
remaining $\approx1.73\times10^6$ bits per row are the flatness targets
of \cref{prob:capfp-Q,prob:capfr1-mode-null}.
\end{corollary}

\begin{proof}
By \cref{lem:capg-prefix-rigidity} the fiber is a family of $m$-subsets
with pairwise $|M\setminus M'|>w$; apply \cref{prop:capf-anticode} with
$s=w$.  The comparison numbers are
$\log_2\binom nm-\log_2\bigl[\binom n{m-s}/\binom ms\bigr]$ at $s=w$
versus $s=\delta_w-1$, evaluated by the entropy sandwich as in
\cref{cor:capfp-packing}.
\end{proof}

\begin{remark}[calibration at the deployed rows]\label{rem:capg-calibration}
Average fibers (bits): $191.61$, $67.10$, $114.30$, $52.11$ at the four
rows of \cref{prop:capff1-identity-frontier}, and $66.91$, $51.93$ at the
moved MCA edges; at the new safe-side rungs $a_0+1$ the averages are
$35.7$ (KoalaBear MCA) and $20.7$ (Mersenne-31 MCA).  Neither the
$\delta_w$-packing nor this cap decides any adjacent pair --- the caps
sit $\approx1.73\times10^6$ bits above average --- but the ledger they
anchor is exact: \cref{thm:capg-second-moment} expresses the entire
$L^2$ mass in closed form, so the flatness inputs of
\cref{prob:capfr1-mode-null,prob:capfp-gamma} can now be posed about a
single explicitly named quantity (the shift-pair counts below) rather
than about an unstructured second moment.
\end{remark}

\begin{theorem}[exact second-moment stratification]\label{thm:capg-second-moment}
Let $N_w(z):=|\mathrm{Fib}_w(z)|$.  Then
\[
\sum_{z\in B^{w}} N_w(z)^{2}
\;=\;\binom{n}{m}\;+\;\sum_{e=w+1}^{\min(m,\,n-m)} P_e ,
\qquad
P_e\;=\;\sum_{\substack{R\subseteq D\\ |R|=m-e}}
\mathrm{sp}_w\big(e;\,D\setminus R\big),
\]
where for $D'\subseteq D$
\[
\mathrm{sp}_w(e;D')\;:=\;\#\Big\{(A,B):\ \text{$A,B$ monic split over $D'$,
degree $e$, disjoint root sets, }\deg(A-B)\leq e-w-1\Big\}
\]
counts ordered \emph{depth-$w$ shift pairs}.  In particular $P_e=0$ for
$1\leq e\leq w$, and the top stratum is
\[
\mathrm{sp}_w(w+1;D')=\#\big\{(A,c):\ \text{$A$ monic split over $D'$ of
degree $w+1$},\ c\in\F^{\times},\ A-c\ \text{split over }D'\big\}.
\]
\end{theorem}

\begin{proof}
The left side counts ordered pairs $(M,M')$ with equal prefixes.  The
diagonal contributes $\binom{n}{m}$.  For $M\neq M'$ the map
$(M,M')\mapsto (R,A,B):=(M\cap M',\,\ell_{M\setminus M'},\,
\ell_{M'\setminus M})$ is a bijection onto triples with $R$ an
$(m-e)$-subset and $(A,B)$ monic split of degree $e$ over $D\setminus R$
with disjoint root sets, for the unique $e=|M\setminus M'|$; and as in the
proof of \cref{lem:capg-prefix-rigidity}, prefix equality is equivalent to
$\deg(A-B)\leq e-w-1$.  Rigidity kills $1\leq e\leq w$; $e$ cannot exceed
$m$ nor $n-m$.  At $e=w+1$ the condition reads $A-B=c\in\F^{\times}$
(nonzero since disjoint nonempty root sets force $A\neq B$); conversely for
any split $A$ and $c\neq 0$ with $A-c$ split, the root sets are
automatically disjoint (a common root would force $c=0$).
\end{proof}

\begin{remark}[relation to the balanced-ternary collision ledger, and
the exact $\Gamma_2$]\label{rem:capg-secondmom-rel}
\Cref{thm:capg-second-moment} refines
\cref{prop:capfr1-collision-ledger} and
\cref{prop:capfp-kernel}\textup{(c)} in three ways.  \textup{(i)}  The
collision strata are proved \emph{empty} for
$\delta_w\le e\le w$: in the notation there, $T_w(e)=0$ for all
$e\le w$ (for $w<\operatorname{char}\F$ via Newton; the
elementary-prefix statement is characteristic-free), so the ledger
starts at $e=w+1$.  \textup{(ii)}  The top stratum $e=w+1$ is
identified: its pairs are exactly the constant-shift split pairs
$(A,\,A-c)$, $c\in\B^{\times}$, both split over $D$ --- an explicitly
algebraic family --- so the leading correction to fiber flatness is a
single named census.  \textup{(iii)}  In the notation of
\cref{prob:capfp-gamma}, the exact collision-gap value is
\[
        \Gamma_2(m,w)
        \;=\;
        \frac{\binom nm+\sum_{e\ge w+1}P_e}
             {\binom nm^{2}\,|\B|^{-w}}
        \;=\;
        \frac1{\bar F}
        +\frac{\sum_{e\ge w+1}P_e}{\binom nm\,\bar F},
        \qquad
        \bar F:=\binom nm|\B|^{-w},
\]
so in the dense bulk ($\bar F\gg1$) the collision gap \emph{is} the
shift-pair ledger normalized by $\binom nm\bar F$: the exact $r=2$ input
to
\cref{prob:capfp-gamma} and extends the anchor
\cref{prop:capfp-anchors} beyond the head depths as an identity (though
not yet as a bound).  A proof of \cref{prob:capg-shiftpairs} below
therefore bounds $\Gamma_2$ exactly, and, by
\cref{prop:capfp-Q-sandwich}, fixed-$r$ analogues of the same
stratification are the asymptotic form of \textup{(Q)}.
\end{remark}

\begin{remark}[quotient prototype and the primitive part]\label{rem:capg-prototype}
Let $D'\supseteq\mu_n$ contain the full group and suppose $(w+1)\mid n$.
Then for every $\beta\in\mu_{n/(w+1)}$ the binomial $X^{w+1}-\beta$ is split
over $\mu_n$ (its root set is a full fiber of the $(w+1)$-power map), and
any two distinct $\beta,\beta'\in\mu_{n/(w+1)}$ give a shift pair with
$c=\beta-\beta'$; hence
$\mathrm{sp}_w(w+1;\mu_n)\geq \tfrac{n}{w+1}\big(\tfrac{n}{w+1}-1\big)$.
These pairs are exactly scale-$(w+1)$ pullbacks: both root sets are unions
of complete fibers, i.e.\ periodic supports in the sense of
\cref{def:periodicity-scale}, so this mass belongs to the quotient ledger
(\cref{thm:fiber-descent}).  When $(w+1)\nmid n$ the binomial prototype is
unavailable and no pullback pair of scale $w+1$ exists; shift pairs may
still exist at composite scales dividing $\gcd$-data of $(e,n)$, and are
pullbacks whenever both root sets are periodic at a common scale
$c\geq 2$.  We call a shift pair \emph{primitive} if it is not a pullback at
any scale $c\geq 2$.
\end{remark}

\begin{remark}[Primitive shift-pair ledger]\label{prob:capg-shiftpairs}
For multiplicative cosets $D$ and the ranges of
\cref{prop:capff1-identity-frontier}, bound the primitive part of
$\mathrm{sp}_w(e;D')$ for $D'=D\setminus R$: show it is dominated, up to a
factor $\mathrm{poly}(n)$, by the pullback part plus the diagonal-scale term
implicit in $\binom{n}{m}$.  Via \cref{thm:capg-second-moment} and
Cauchy--Schwarz this is the exact combinatorial content of second-moment
flatness; it refines \cref{lem:capf-gcd} and
\cref{prob:capf-primitive-image-fiber} to the support side, and is the
support-level counterpart of the alignment moment calculus
(\cref{thm:capf-first-moment,thm:capf-second-moment,cor:capf-exact-variance};
scope as in \cref{rem:capf-moment-scope}).
\end{remark}

\begin{remark}[flatness wiring after capg]\label{rem:capg-flatness-wiring}
The first-pass capg flatness problem is withdrawn as a separate item: its
content is covered by \cref{prob:capfr1-master-flatness} (asymptotic,
base-field-normalized form) and \cref{prob:capfr1-mode-null} (finite
extremality form), and this part's contribution to them is structural:
\begin{enumerate}[label=\textup{(\roman*)}]
\item the collision side of \cref{prob:capfr1-mode-null} now has the
exact ledger of \cref{thm:capg-second-moment}, whose sole unbounded
ingredient is \cref{prob:capg-shiftpairs}; a proof of the latter with
constant $\kappa'$ bounds every fiber's excess over average via
$\Gamma_2$ and Cauchy--Schwarz, though, per
\cref{rem:capfpr-low-moments,prop:capfp-Q-sandwich}, second-moment control
alone cannot decide the printed pairs;
\item the finite targets sit at the \emph{new} rungs
$a_0+1=1116048,\,1116024$ with fail margins $22.2$ and $3.3$ bits
(\cref{cor:capg-adjacent-pairs}), and the rung audit
\cref{prob:capfr1-rung-audit} is a prerequisite there exactly as before;
\item nothing in this part bounds $\mathrm{sp}_w(e)$ for the primitive
part; that residual is named in \cref{prob:capg-shiftpairs} and is a
strict sub-problem of the mode-at-null input.
\end{enumerate}
\end{remark}

\section{Subfield floors for split-pencil censuses}\label{sec:capg-subfield}

The aperiodic side of the safe programme has been reduced, in
\cref{subsec:capfr1-lattice-split} and the final-form subsections of
\cref{sec:capfp-prize}, to census problems whose conjectured models are
normalized at the challenge-field scale:
\cref{prob:capfr1-rank-one-census,prob:capfp-R1} with
$\max\bigl(1,\binom nmq^{-(w'-1)}\bigr)$,
\cref{prob:capfr1-balanced-core,prob:capfp-balanced} with
$\max\bigl(1,\binom nmq^{-w'}\bigr)$, and
\cref{prob:capfr1-split-pencil,prob:capfp-split} with
$\max\bigl(1,\binom n\omega q^{1-w'}\bigr)$, where in this section
$K$ is the code dimension of the census ($K=k$ for the deployed rows),
$w':=m-K$, and $\omega=n-m$.  These normalizations are calibrated to
challenge-field-generic words and lines.  The identity-prefix mechanism,
however, produces \emph{base-field} instances at every balanced profile,
and their censuses sit at the $|\B|$-scale, which is larger than the
$q$-scale by a factor exponential in $n$ whenever $[\F:\B]\ge2$.  As
stated, all three problems are therefore refuted; the corrections are
recorded at the end of the section, and they align the census problems
with the base-field normalization that
\cref{prob:capfr1-master-flatness,prob:capfp-Q} already carry.

\begin{proposition}[subfield census floors at every balanced profile]
\label{prop:capg-census-floor}
Let $C=\RS[\F,D,K]$, $D\subseteq\B$, $|D|=n$, $q=|\F|$, $p=|\B|$, and
fix an agreement $m$ with $K\le m$ and $2(m-K)\le n-K$; write $w'=m-K$.
For a profile parameter $d_1$ with $w'+1\le d_1\le\floor{(n-K+1)/2}$ put
$m':=K-1+d_1$ and assume $m'+d_1\le n$.  Then:
\begin{enumerate}[label=\textup{(\alph*)}]
\item \textup{(Boundary profile, $d_1=w'+1$.)}  For every $z\in\B^{w'}$
the witness word $U_z$ of \cref{lem:capff1-identity-prefix-floor} has
minimal shifted lattice degree $d_1(U_z)=w'+1$
(\cref{thm:capfp-Q-unify}), and
$\Cen(U_z;m)=|\Phi_{m,w'}^{-1}(z)|$; at the heaviest prefix,
\[
        \Cen(U_{z^{*}};m)\ \ge\
        \Bigl\lceil\binom nm\,p^{-w'}\Bigr\rceil .
\]
\item \textup{(Interior profiles, $d_1\ge w'+2$.)}  For every
$z'\in\B^{\,d_1-1}$ the level-$m'$ witness word $U^{(m')}_{z'}$ has
$d_1\bigl(U^{(m')}_{z'}\bigr)=d_1$, is a balanced determinantal
representation datum in the sense of
\cref{prop:capfr1-detrep,prop:capfp-detrep}, and at the heaviest level-$m'$
prefix its size-$m$ census satisfies
\[
        \Cen\bigl(U^{(m')}_{z'};\,m\bigr)\ \ge\
        \binom{m'}{m}\,
        \Bigl\lceil\binom n{m'}\,p^{-(d_1-1)}\Bigr\rceil
        \;=:\;M_{\B}(d_1).
\]
At $d_1=w'+1$ the same formula reproduces \textup{(a)}.
\item \textup{(Per-line $\Rone$ floors.)}  Let $w'\ge1$, $m\ge K+1$, and
let $(f_\alpha,g_\alpha)$ be the pole line of $U_{z^{*}}$ at any
$\alpha\in\F\setminus(D\cup\B)$.  Then every member $T$ of the heaviest
fiber is a non-common rank-one support of $(f_\alpha,g_\alpha)$ in the
sense of \cref{prop:capfr1-slope-elimination}, so
\[
        \#\Rone(f_\alpha,g_\alpha;m)\ \ge\
        \Bigl\lceil\binom nm\,p^{-w'}\Bigr\rceil ,
\]
and the line is not $\B$-rational, hence outside the subfield-confined
stratum of \cref{lem:confine} cited in \cref{prob:capfp-R1}.
\end{enumerate}
Consequently, along the asymptotic family of \cref{cor:capg-asymptotic}
(where $\binom nm p^{-w'}\ge2^{n/2-1}$ while
$\binom nm q^{-(w'-1)}\le1$), any final census model omitting the base-field
floor is too small.  At the deployed list rows the missing floor contributes,
in bits, $67.1$ (KoalaBear) and $52.1$ (Mersenne-31) at the boundary profile,
and $56.0$, $43.9$, $31.3$ (KoalaBear, $d_1=w'+2,w'+3,w'+4$), resp.\
$41.0$, $28.9$, $16.2$ (Mersenne-31), at the first interior profiles, against
challenge-field models below one.
\end{proposition}

\begin{proof}
\textup{(a)}  The lattice statement and the census identity are
\cref{thm:capfp-Q-unify} (the hypothesis $m+w'<n$ needed for
$d_1=w'+1$ follows from $2w'\le n-K$); the fiber lower bound is the
pigeonhole of \cref{lem:capff1-identity-prefix-floor}.

\textup{(b)}  $U^{(m')}_{z'}$ is defined by the same recipe at level
$m'$: its interpolant is the monic polynomial
$P^{(m')}_{z'}=X^{m'}+\sum_{j\le d_1-1}z'_jX^{m'-j}$ of degree $m'<n$.
The minimal-degree computation of \cref{thm:capfp-Q-unify} applies
verbatim at level $m'$ (using $m'+d_1\le n$, so the reduction modulo
$\Lambda_D$ is trivial in the relevant degree range), giving
$d_1(U^{(m')}_{z'})=m'-K+1=d_1$, together with the reduced basis
$\bigl((1,P^{(m')}_{z'}),\,g_2\bigr)$, which is a coprime unimodular
quadruple by \cref{lem:capfp-unimodular}; this is the claimed
determinantal datum.  For the census: by
\cref{prop:capff1-witness-exact} at level $m'$, each member $M'$ of the
level-$m'$ fiber of $z'$ yields a codeword $c_{M'}$ of $\RS[\F,D,K]$
agreeing with $U^{(m')}_{z'}$ on exactly the $m'$ points of $M'$, and
distinct $M'$ give distinct codewords.  Each such codeword contributes
$\binom{m'}m$ valid size-$m$ supports to the binomial-moment formula of
\cref{prop:capfr1-lattice-census}\textup{(c)}, all distinct across
distinct $M'$ since the codeword attached to a valid support is unique
($m\ge K$).  The heaviest level-$m'$ fiber has size at least
$\ceil{\binom n{m'}p^{-(d_1-1)}}$ by pigeonhole.

\textup{(c)}  By \cref{lem:capg-witness}\textup{(i)} each fiber member
$T$ carries the single close slope
$\zeta_T=U_{z^{*}}(\alpha)-\Lambda_T(\alpha)$, so by
\cref{prop:capfr1-slope-elimination} (or directly
\cref{lem:capfr1-interpolation-test}) $\alpha(T)+\zeta_T\beta(T)=0$.  The
support is non-common: $\beta(T)=0$ would mean $g_\alpha$ extends on $T$
to a polynomial $G$ of degree $<K$, whence $(X-\alpha)G+1$, of degree at
most $K\le m-1<|T|$, vanishes on $T$, forcing the zero polynomial, which
contradicts its value $1$ at $\alpha$.  Thus $\alpha(T),\beta(T)$ are
proportional with $\beta(T)\ne0$: a non-common rank-one support, and
distinct $T$ are distinct supports.  Non-$\B$-rationality:
$g_\alpha=-1/(x-\alpha)$ takes a value outside $\B$ at any $x\in D$ with
$x-\alpha\notin\B$, which holds since $\alpha\notin\B$.

The deployed violation values are orientation computations of
$\log_2M_{\B}(d_1)$ at $(K,m)=(k,1116046)$ resp.\ $(k,1116022)$; the
asymptotic claim substitutes the family parameters of
\cref{cor:capg-asymptotic}, for which
$q\ge k\cdot2^{n/8}$ makes every challenge-field model term at most one
while the base-field floor is at least $2^{n/2-1}$ at the boundary and at
least $2^{n/4}$ at the bounded-offset interior profiles.
\end{proof}

\begin{remark}[boundary profile and active split-pencil range]
\label{rem:capg-boundary-offbyone}
\Cref{thm:capfp-Q-unify} states that \textup{(Q)} is the boundary
profile $d_1=w'+1$ of the split pencil and the aperiodic balanced core
consists of \emph{interior} profiles; but
\cref{prob:capfp-split,prob:capfr1-split-pencil} quantify over
``balanced profile $d_1\ge w'+1$'', which includes the boundary and
hence includes \textup{(Q)} itself with the wrong ($q$-scale) model.
Part \textup{(a)} above identifies the boundary contribution.  The intended
reading --- boundary delegated to \textup{(Q)}, split-pencil census over
$d_1\ge w'+2$ --- is adopted in the restatement below; part \textup{(b)}
shows that the interior model must also be base-field-normalized.
\end{remark}

\begin{remark}[Base-field-normalized split-pencil model]
\label{prob:capg-split-pencil-B}
Let $(W_1,N_1,W_2,N_2)$ be a determinantal representation of
$\Lambda_D$ as in \cref{prop:capfr1-detrep} with interior balanced
profile $w'+2\le d_1\le\floor{(n-K+1)/2}$.  Prove
\[
        \#\bigl\{(A,B):\ \deg A\le\omega-d_1,\ \deg B\le\omega-d_2,\
        AW_1+BW_2\mid\Lambda_D,\ \deg(AW_1+BW_2)=\omega\bigr\}
        \ \le\ e^{o(n)}\,
        \max\Bigl(1,\ M_{\B}(d_1),\ \binom n\omega\,q^{1-w'}\Bigr),
\]
counting monic normalizations, with
$M_{\B}(d_1)=\binom{K-1+d_1}{m}\binom n{K-1+d_1}p^{-(d_1-1)}$ the
subfield floor of \cref{prop:capg-census-floor}, achieved there; the
finite deployed form replaces $e^{o(n)}$ by exact constants at the new
adjacent values of \cref{cor:capg-adjacent-pairs}.  The same
correction applies \emph{mutatis mutandis} to
\cref{prob:capfr1-balanced-core,prob:capfp-balanced} (model
$\max(1,M_{\B}(d_1))$ at profile $d_1$, plus the $q$-generic term) and to
\cref{prob:capfr1-rank-one-census,prob:capfp-R1} (model
$\max\bigl(1,\binom nm p^{-w'},\binom nm q^{-(w'-1)}\bigr)$ per line,
the middle term achieved by the non-$\B$-rational pole lines of
\cref{prop:capg-census-floor}\textup{(c)}).  With these models the
implication chain of
\cref{cor:capfr1-terminal-package,prop:capfp-closing} is unchanged: the
$|\B|$-scale strata are exactly the strata that the prefix ledger and
\textup{(Q)} already price, so the corrected problems still imply the
normalized band input \cref{prob:capfr1-normalized-band}, and they are
now consistent with \cref{prob:capfr1-master-flatness}, whose affine
subspaces live over $\B$ with model $\binom nj|\B|^{-s}+1$.
\end{remark}

\begin{remark}[what the subfield floors do not say]
\label{rem:capg-subfield-scope}
The floors identify the correct model; they do not weaken the programme.
In particular: \textup{(i)} the reductions
\cref{prop:capfr1-slope-elimination,thm:capfp-slope-elim,%
thm:capfr1-near-rational-dichotomy,thm:capfp-dichotomy,%
prop:capfr1-detrep,prop:capfp-detrep,lem:capfr1-autodiv,lem:capfp-autodiv}
are unconditional and untouched --- only the conjectured \emph{sizes}
change; \textup{(ii)} the per-line slope counts remain bounded by the
per-line census, and on the certifying pole lines the slope count and
the support count agree to within the collision correction of
\cref{thm:capg-aperiodic-floor}, so the corrected $\Rone$ model is tight
at the floors; \textup{(iii)} for genuinely challenge-field-primitive
representations --- no member of the pencil proportional to a
$\B$-rational datum in the sense of \cref{lem:confine} --- the original
$q$-scale models remain the plausible targets, and separating that
stratum cleanly is part of \cref{prob:capg-split-pencil-B}.
\end{remark}

\section{Frontier conjectures and next steps}\label{sec:frontier-next-steps}

We now state the remaining package in the form used by the closing argument.
The unsafe side is theorem-level: it is supplied by the identity-prefix floor,
the list certificates of \cref{prop:capff1-identity-frontier}, and the
flexible-budget MCA certificates of \cref{prop:capg-moved-frontier}.  The
safe side is conditional on three precise inputs.

\begin{residualinput}[quotient and prefix flatness]\label{prob:capg-active-Q}
For the identity prefix map
\[
        \Phi_w(M)=\bigl((-1)^j e_j(M)\bigr)_{1\le j\le w},
        \qquad |M|=m,
\]
prove the max-fiber bound
\[
        \max_z |\Phi_w^{-1}(z)|
        \le \kappa\binom n m |\mathbb B|^{-w}
\]
with \(\kappa=n^{O(1)}\) in the reserve/asymptotic regime and with an
explicit row-dependent constant in the finite adjacent regime.  The finite
constant must fit inside the margins in \cref{cor:capg-adjacent-pairs}.
Equivalently, one may prove the mode-at-null or exchange-compression
extremality statement of \cref{prob:capfr1-mode-null} together with an exact
null-fiber ledger.
\end{residualinput}

\begin{residualinput}[base-field-normalized split-pencil census]\label{prob:capg-active-BC}
For the interior balanced profiles of the split-pencil reduction, prove the
base-field-normalized census of \cref{prob:capg-split-pencil-B}.  The boundary
profile belongs to \cref{prob:capg-active-Q}; the interior model must include
the \(|\mathbb B|\)-scale floor terms of \cref{prop:capg-census-floor}.
\end{residualinput}

\begin{residualinput}[primitive shift-pair control]\label{prob:capg-active-shiftpairs}
For the shift-pair strata \(\mathrm{sp}_w(e;D')\) in
\cref{thm:capg-second-moment}, bound the primitive part after all
quotient-pullback pairs have been removed.  The asymptotic form may lose a
polynomial factor with reserve; the finite form requires explicit constants
compatible with \cref{cor:capg-adjacent-pairs}.
\end{residualinput}

\begin{proposition}[conditional frontier closure]\label{prop:capg-final-active-package}
Assume the asymptotic forms of
\cref{prob:capg-active-Q,prob:capg-active-BC,prob:capg-active-shiftpairs} with
polynomial or \(e^{o(n)}\) loss.  Then the asymptotic frontier conjecture of
\cref{prob:capff1-frontier} holds with any reserve \(R\gg\log n\).  If the
finite forms hold with explicit constants below the margins in
\cref{cor:capg-adjacent-pairs}, then the finite adjacent pairs in that
corollary follow from the staircase compiler \cref{thm:capf-staircase} and the
one-step certificate mechanism \cref{prop:onestep}.
\end{proposition}

\begin{proof}
The lower side is supplied by the exact unsafe certificates.  For the upper
side, quotient and prefix fibers are paid by \cref{prob:capg-active-Q}; the
interior aperiodic support census is paid by
\cref{prob:capg-active-BC}; and the primitive part of the exact second-moment
ledger is paid by \cref{prob:capg-active-shiftpairs}.  The normalized band
input begins at the entropy-subfield envelope, where the identity-floor
obstruction no longer supplies exponential aperiodic mass.  With reserve, a
polynomial factor costs only \(O(\log n)\) agreement bits and is absorbed.  At
the printed finite adjacent pairs, no such absorption is available, so the
explicit constants must lie below the row margins.  The integer staircase then
localizes the first safe agreement exactly.
\end{proof}

The next mathematical tasks are therefore sharply delimited: prove finite
prefix extremality or exchange compression for \cref{prob:capg-active-Q}; prove
the base-field-normalized split-pencil census for the interior balanced-core
profiles; bound the primitive shift-pair terms in the second-moment ledger;
and keep all quotient-rung accounting exact at the finite adjacent pairs.



\section{Discussion}\label{sec:discussion}

\paragraph{Values versus fibers.}
Above the per-rate value-set reach of \cite{Cho26b} (roughly $2^{150}$--$2^{162}$), the value route --- showing that the slope multiset $\eone(\ell^\wedge Q)$ covers a $2^{-128}$ fraction of $\F$ at a fixed large prime --- runs into the per-fiber collision problem (\cite[Rem.\ ``The remaining band'']{Cho26b}): characteristic-zero locator values can collide arbitrarily badly under reduction at a worst-case prime. \cite{Cho26b} therefore closed the band at the $\delta^*$ grade only conditionally and with weaker constants --- error $2^{-42}$ at gap $2^{-11}$, assuming $|\F|\ge2n$, via quotient-core lists and the same conversion \cref{thm:A} --- leaving the \emph{error-one} grade outside the present method. The present route also never counts values, and improves each parameter of that cap: pigeonhole guarantees one heavy fiber regardless of how the values $\eone(A)$ distribute, \cref{thm:A} converts heavy fibers into correlated-agreement error directly, and the subfield pigeonhole carries the construction to extension fields with no condition relating $n$ and $|\F|$. The cost relative to the error-one results of \cite{Cho26a,Cho26b} below $2^{150}$ is quantitative: error $\approx\frac1{2k}$ instead of $1-o(1)$, and gap $2^{-9}$ (uniformly $2^{-10}$ across all four rates) instead of $\frac1{64}$. For the prize question --- is $\emca\le2^{-128}$ achievable above the stated gaps? --- this cost is immaterial.

\paragraph{Limits of the method.}
Hypothesis \eqref{eq:hyp} needs $N\,\Hb(\rho)\gtrsim2\log_2q-\log_2k$, so the smallest certifiable gap is
\[
\frac2N\;\approx\;\frac{2\,\Hb(\rho)}{2\log_2q-\log_2k}\;\approx\;\frac{\Hb(\rho)}{\log_2 q},
\]
and the certified error is capped at $\approx\frac1k$ by the shape of \cref{thm:A} (any $\eta<1$). Since $q$ enters only as the size of the field from which the line challenge $\gamma$ is drawn, this is also a protocol-facing floor: enlarging the challenge field cannot push a deep-point-conversion failure below gap $\approx\Hb(\rho)/\log_2 q_{\mathrm{line}}$, which quantitatively bounds the ``larger challenge field'' fallback in certificate frameworks. The construction also needs $N\le n$, i.e.\ $\log_2q\lesssim\Hb(\rho)\,n/2$; this covers the entire cryptographic regime $\mathrm{poly}(n)\le q\le2^{o(n)}$, but for $q\ge2^{\Hb(\rho)n/2}$ the method is vacuous and the value-route caps of \cite{Cho26a,Cho26b} are the only ones known. Within the band, the \emph{error-one} question at gap $\frac1{64}$ (the error-one question \cref{prob:errorone}) remains exactly as open as before: the present result bounds $\delta^*$, not the failure profile near capacity.

\paragraph{Comparison with the cyclotomic sieve.}
The sieve of \cite{Cho26a} produces, for infinitely many primes, failures of error $(\log p)^{-6}$ at gap $\Theta(1/\log p)$, but its prime set is given by Siegel--Walfisz and is ineffective. \Cref{cor:grand} trades error strength for universality: every field below $2^{256}$, fully effective, and elementary given \cref{thm:A}. The two results triangulate the failure landscape from different directions.

\paragraph{Isogeny-smooth reach and the lacunarity dichotomy.}
The counting side of the method extends beyond multiplicative smoothness. The fiber lemma needs only a polynomial folding map with complete uniform fibers inside the domain (\cref{def:map-smooth,lem:phi-fiber}), and the divisibility hypothesis $a\mid k$ is removed by rounding, so the universal cap now covers mixed-radix multiplicative rows without integral $\rho N$ (\cref{cor:mixed-radix}), the $x$-coordinate line rounds of circle FRI over $p\equiv3\pmod4$, and --- through the torus uniformization, over any challenge field containing $i$ --- the bivariate circle codes themselves, at the deployed Mersenne-$31$/QM31 parameters with error $>2^{-23}$ at gap $2^{-7}$ (\cref{cor:circle-grand,cor:circle-deployed}). The boundary of the method is now a clean dichotomy in the degree type $(a,e)$ of the folding map: polynomial towers ($e=0$: powers, Chebyshev) give gap $\approx2/N$, while elliptic isogeny towers ($e=a-1$: ECFFT) collapse the slack-two window to gap $2/n$ (\cref{rem:ecfft-obstruction}). Subfield confinement and the extension-pole mechanism also carry over to circle rows (\cref{cor:circle-Fvalued}), so the deployed circle failures are again fiber-borne, genuinely extension-valued, and constructively reachable from any supplied list via extension poles.

\paragraph{Pressure on the packing conjecture.}
By the exact normal form of \cite{Cho26b}, $\emca(C,\delta)=\max_t\Lambda^{\mathrm{NC}}_t(\delta)/q$, where $\Lambda^{\mathrm{NC}}$ is a noncontained residue-line packing number. \Cref{cor:deployed} therefore forces $\max_t\Lambda^{\mathrm{NC}}\ge(q-n)/2k\approx2^{164}$ at gap $2^{-7}$ on the deployed sextic --- superpolynomially many noncontained residue lines at $n=2^{21}$. This puts no direct pressure on the final MCA conjecture of \cite[Conj.\ ``Final floor- and quotient-corrected MCA asymptotic'']{Cho26b}: at gap $2^{-7}$ both of its corrected-reserve hypotheses fail ($\sigma=2^{14}$ lies below $Cn/\log_2n\approx2^{16.6}$, and $\sigma\log_2 q_{\mathrm{gen}}=2^{14}\cdot31\approx2^{19}$ is far below $\log_2\binom n{k+\sigma}\approx2^{21}$), so the conjectured ceiling of the form $\bigl(n^{1+o(1)}+2^{(\beta/\Hb)\,Q_{\Hb}(\eta)}\bigr)/q$ is simply not asserted there --- the packing mass extracted through \cref{thm:A} lives in the region the framework already marks as below-floor. The genuine structural question this paper raises for the framework of \cite{Cho26b} is whether conversion-extracted fiber mass persists \emph{above} the corrected reserve: whether, at reserves where the conjecture does apply, the fiber route still produces packings beyond the $n^{1+o(1)}$ aperiodic term plus the quotient-periodic term, which would force an additional fiber-borne term into the conjectured ceiling. This is not settled here.

\paragraph{Verification caveat.}
The main cap no longer depends on importing the Crites--Stewart theorem: \cref{thm:A} proves the needed deep-point list-to-CA conversion directly, including the integer-radius condition $f\le n-k-1$.  \Cref{thm:B} remains a separate slacked fallback imported through the BCHKS/ABF route, and readers using that fallback should consult the underlying source directly.  The main composition is deliberately shallow: it uses the self-contained conversion, the locator-fiber pigeonhole, and support-wise MCA monotonicity.

\paragraph{What remains after the frontier package.}
The structural ledgers, extension-line accounting, scanner, and certificate grammar are closed at the level needed for the printed unsafe certificates and for conditional safe-side compilation.  After the entropy--subfield normalization, the remaining safe-side package is more precise than the band statement by itself.  The active residue is the three-part package
\[
\begin{gathered}
        \text{quotient/prefix max-fiber flatness},\\
        \text{base-field-normalized interior split-pencil census},\\
        \text{primitive shift-pair control in the exact second-moment ledger}.
\end{gathered}
\]
They are formulated precisely in \cref{sec:frontier-next-steps}.  Polynomial or \(e^{o(n)}\) losses suffice only with reserve in the asymptotic frontier.  The finite adjacent rows require explicit constants below the printed margins of \cref{cor:capg-adjacent-pairs}.  Separate auxiliary questions remain: the error-one profile \cref{prob:errorone}, the explicit deployed-witness strengthening \cref{prob:explicit}, and the genus-one/circle transport problem \cref{prob:isogeny}.

\paragraph{Questions resolved by the present methods.}
The genus-one macroscopic-cap and circle field-of-definition questions are resolved by \cref{cor:ecfft-macroscopic,cor:circle-anyfield}.  The explicit-witness question is resolved at the first staircase step and for simple-pole certifying pairs at gap $2^{-8}$ \textup(\cref{cor:explicit-deployed,thm:explicit-pairs}\textup); its full-gap strengthening is restated below.

\paragraph{Further frontier questions.}

\begin{researchquestion}[explicit deployed witnesses at the full deployed gap]\label{prob:explicit}
For the KoalaBear-sextic deployed row, exhibit explicit pairs $(f,g)$ with MCA-bad-slope density $>2^{-22}$ at the gap $2^{-7}$ certified by \cref{cor:deployed}; on the integer-admissible subinterval, exhibit the corresponding CA-bad-slope density as well.  By \cref{cor:Fvalued} such pairs exist and are not $\B$-rational, and the natural candidates are residue-line normal forms with denominator $E\in\F[X]$ not defined over $\B$.  The present paper gives explicit received words and simple-pole certifying pairs of exactly this normal form, $E=X-\alpha$, at the first staircase step and at gap $2^{-8}$ \textup(\cref{cor:explicit-deployed,thm:explicit-pairs}\textup); the remaining task is to push the explicit pair construction to the full deployed gap, derandomize the choice of the pole $\alpha$, and print the deployed-size bucket tables.
\end{researchquestion}

\begin{researchquestion}\label{prob:errorone}
Error one in the band: does $\emca(C,1-\rho-\frac1{64})=1-o(1)$ hold for all primes $2^{150}\le p\le2^{256}$? This is the per-fiber collision problem proper; \cref{cor:grand} caps $\delta^*$ there but says nothing above error $\approx\frac1{2k}$.  \Cref{cor:profile-deployed} raises the certified profile to $(1-2^{-54})\frac1k$ at gap $2^{-7}$, and \cref{rem:profile-ceiling} shows the question lies strictly beyond both mechanisms of this paper; it remains outside the present methods.
\end{researchquestion}

\paragraph{Further direction: improving the conversion threshold.}\label{ques:conversion-threshold}
Improve the threshold of \cref{thm:A}: any strengthening of the hypothesis range $\eca\le\eta/k$ toward $\eca\le\eta\cdot\mathrm{polylog}(q)/k$, or of the list conclusion, immediately tightens the certified error of \cref{thm:main} toward $1/\mathrm{poly}\log$; the gap, by contrast, is governed by the binomial in \eqref{eq:hyp} and would need a denser fiber construction to shrink.  Cf.\ \cref{prop:profile}: optimizing the dichotomy already attains $(1-o(1))/k$, so passing $1/k$ is exactly what a polylog strengthening would buy.


\begin{researchquestion}[aperiodic band on genus-one and circle tangent domains]\label{prob:isogeny}
Extend the active safe-side package---quotient/prefix flatness, base-field-normalized split-pencil control, and primitive shift-pair control---to the tangent-domain Reed--Solomon models arising from genus-one isogeny towers and to the stereographic circle models over challenge fields not containing $i$.  The macroscopic unsafe floors on these rows are supplied by \cref{cor:ecfft-macroscopic,cor:circle-anyfield}; the remaining task is to transport the normalized safe-side input through the rational scale map.
\end{researchquestion}

\paragraph{Acknowledgements.} This paper builds on the Crites--Stewart conversion perspective \cite{CS25}, the BCHKS theorem \cite{BCHKS25}, and the locator framework of \cite{Cho26a,Cho26b}, following the synthesis of Arnon--Boneh--Fenzi \cite{ABF26}; it proves the main conversion needed here directly, sharpens the conditional universal-cap theorem, and refines the subfield-confinement theorem of \cite{Cho26b}.


\begin{thebibliography}{BCHKS25}

\bibitem[PP26]{ProximityPrize}
The Proximity Prize.
\newblock The Proximity Prize: \$1M Reed--Solomon Challenge.
\newblock Preliminary release, 2026.
\newblock \url{https://proximityprize.org/}.


\bibitem[ABF26]{ABF26}
G.~Arnon, D.~Boneh, and G.~Fenzi.
\newblock Remaining questions in list decoding and correlated agreement.
\newblock IACR Cryptology ePrint Archive, Report 2026/680, 2026.
\newblock \url{https://eprint.iacr.org/2026/680}.

\bibitem[BCHKS25]{BCHKS25}
E.~Ben-Sasson, D.~Carmon, U.~Hab{\"o}ck, S.~Kopparty, and S.~Saraf.
\newblock On proximity gaps for Reed--Solomon codes.
\newblock IACR Cryptology ePrint Archive, Report 2025/2055, 2025.
\newblock \url{https://eprint.iacr.org/2025/2055}.


\bibitem[BCIKS20]{BCIKS20}
E.~Ben-Sasson, D.~Carmon, Y.~Ishai, S.~Kopparty, and S.~Saraf.
\newblock Proximity gaps for Reed--Solomon codes.
\newblock In \emph{Proc.\ 61st IEEE Symposium on Foundations of Computer Science (FOCS)}, 2020.
\newblock Extended manuscript: IACR Cryptology ePrint Archive, Report 2020/654.

\bibitem[BCKL21]{BCKL21}
E.~Ben-Sasson, D.~Carmon, S.~Kopparty, and D.~Levit.
\newblock Elliptic curve fast Fourier transform (ECFFT) part~I: Fast polynomial algorithms over all finite fields.
\newblock arXiv:2107.08473, 2021.

\bibitem[Cho26a]{Cho26a}
P.~Chojecki.
\newblock Capacity-edge obstructions to Reed--Solomon mutual correlated agreement over smooth multiplicative domains.
\newblock Manuscript, 2026.

\bibitem[Cho26b]{Cho26b}
P.~Chojecki.
\newblock Slack, quotient cores, and the entropy gap for smooth-domain Reed--Solomon codes: list fibers, cyclotomic rigidity, Galois-amplified collisions, and the corrected slack mutual-correlated-agreement theory.
\newblock Manuscript, merged corrected edition, 2026.

\bibitem[CS25]{CS25}
E.~Crites and A.~Stewart.
\newblock On Reed--Solomon proximity gaps conjectures.
\newblock IACR Cryptology ePrint Archive, Report 2025/2046, 2025.
\newblock \url{https://eprint.iacr.org/2025/2046}.

\bibitem[HLP24]{HLP24}
U.~Hab{\"o}ck, D.~Levit, and S.~Papini.
\newblock Circle STARKs.
\newblock IACR Cryptology ePrint Archive, Report 2024/278, 2024.
\newblock \url{https://eprint.iacr.org/2024/278}.

\end{thebibliography}

\end{document}
